%% This is a LaTeX document.  Hey, Emacs, -*- latex -*- , get it?
\documentclass[12pt,a4paper]{article}
\usepackage[a4paper,twoside,inner=2.75cm,outer=2.25cm,vmargin=3cm]{geometry}
\usepackage[francais]{babel}
\usepackage[utf8]{inputenc}
\usepackage[T1]{fontenc}
%\usepackage{ucs}
\usepackage{times}
% A tribute to the worthy AMS:
\usepackage{amsmath}
\usepackage{amsfonts}
\usepackage{amssymb}
\usepackage{amsthm}
%
\usepackage{mathrsfs}
\usepackage{wasysym}
\usepackage{url}
%
\usepackage{makeidx}
%% Self-note: compile index with:
%% xindy -M texindy -C utf8 -L french notes-inf105.idx
%
\usepackage{graphics}
\usepackage[usenames,dvipsnames]{xcolor}
\usepackage{tikz}
\usetikzlibrary{arrows,automata,positioning,calc}
\usepackage[hyperindex=false]{hyperref}
%
\theoremstyle{definition}
\newtheorem{comcnt}{Tout}[subsection]
\newcommand\thingy{%
\refstepcounter{comcnt}\medskip\noindent\textbf{\thecomcnt.} }
\newcommand\exercice{%
\refstepcounter{comcnt}\bigskip\noindent\textbf{Exercice~\thecomcnt.}}
\newtheorem{defn}[comcnt]{Définition}
\newtheorem{prop}[comcnt]{Proposition}
\newtheorem{lem}[comcnt]{Lemme}
\newtheorem{thm}[comcnt]{Théorème}
\newtheorem{cor}[comcnt]{Corollaire}
\newtheorem{scho}[comcnt]{Scholie}
\newtheorem{algo}[comcnt]{Algorithme}
\renewcommand{\qedsymbol}{\smiley}
%
\newcommand{\limp}{\mathrel{\Longrightarrow}\relax}
\newcommand{\liff}{\mathrel{\Longleftrightarrow}\relax}
%
\DeclareUnicodeCharacter{00A0}{~}
\DeclareUnicodeCharacter{03B5}{$\varepsilon$}
\DeclareUnicodeCharacter{0101}{\=a}
\DeclareUnicodeCharacter{1E47}{\d{n}}
%
\DeclareMathSymbol{\tiret}{\mathord}{operators}{"7C}
\DeclareMathSymbol{\traitdunion}{\mathord}{operators}{"2D}
%
%
\DeclareFontFamily{U}{manual}{} 
\DeclareFontShape{U}{manual}{m}{n}{ <->  manfnt }{}
\newcommand{\manfntsymbol}[1]{%
    {\fontencoding{U}\fontfamily{manual}\selectfont\symbol{#1}}}
\newcommand{\dbend}{\manfntsymbol{127}}% Z-shaped
\newcommand{\danger}{\noindent\hangindent\parindent\hangafter=-2%
  \hbox to0pt{\hskip-\hangindent\dbend\hfill}}
%
\newcommand{\defin}[2][]{\def\latexsucks{#1}\ifx\latexsucks\empty\index{#2}\else\index{\latexsucks}\fi\textbf{#2}}
\newcommand{\spaceout}{\hskip1emplus2emminus.5em}
%
\newif\ifcorrige
\corrigetrue
\newenvironment{corrige}%
{\ifcorrige\relax\else\setbox0=\vbox\bgroup\fi%
\smallbreak\footnotesize\noindent{\underbar{\textit{Corrigé.}}\quad}}
{{\hbox{}\nobreak\hfill\checkmark}%
\ifcorrige\relax\else\egroup\fi\par}
%
%
% NOTE: compile dot files with
% dot2tex --figonly -f tikz --tikzedgelabels --graphstyle=automaton file.dot  > file.tex
\tikzstyle{automaton}=[>=stealth',initial text={},thick,every loop/.style={min distance=7mm,looseness=5}]
\tikzstyle{state}=[]
\tikzstyle{final}=[accepting by arrow]
%
%
%
\makeindex
\begin{document}
\title{THL (Théorie des langages)\\Notes de cours}
\author{David A. Madore}
\maketitle

\centerline{\textbf{INF105}}

\vskip2cm

{\footnotesize
\immediate\write18{sh ./vc > vcline.tex}
\begin{center}
Git: \input{vcline.tex}
\\
(Recopier la ligne ci-dessus dans tout commentaire sur ce document)
\end{center}
\immediate\write18{echo ' (stale)' >> vcline.tex}
\par}

\pretolerance=8000
\tolerance=50000
\linepenalty=5  % Default is supposedly 10

\vskip3cm


%
%
%

{\footnotesize
\tableofcontents
\par}

\bigbreak

\setcounter{comcnt}{0}

\thingy \textbf{Présentation générale et plan :} On se propose de
développer ici des techniques mathématiques et algorithmiques ayant à
la fois un intérêt théorique comme théorie des langages formels, et
une application informatique pratique à l'analyse de textes ou de
chaînes de caractères (qu'on appellera par la suite « mots »,
cf. §\ref{subsection-introduction-and-words}).

\smallbreak

Après des définitions générales (sections
\ref{subsection-introduction-and-words} à \ref{subjection-languages}),
ces notes sont divisées en grandes parties (inégales) de la manière
suivante :
\begin{itemize}
\item l'étude des expressions rationnelles et des automates finis, et
  des langages qu'ils définissent, dans les sections
  \ref{section-languages-and-rational-languages} à \ref{section-recognizable-languages},
  qui sert notamment dans la recherche de texte et de motifs dans une
  chaîne de caractères ou un document ;
\item l'étude de certaines grammaires formelles dans la
  section \ref{section-context-free-grammars}, qui sert notamment à
  définir et analyser la syntaxe de langages de programmation ;
\item une introduction à la calculabilité dans la
  section \ref{section-computability}, qui place des limites
  théoriques sur ce qu'un algorithme peut faire.
\end{itemize}

\smallbreak

À ces parties seront associées la définition de différentes classes de
« langages » (voir §\ref{subjection-languages} pour la définition de
ce terme) qu'il s'agit d'étudier :
\begin{itemize}
\item les langages « rationnels » (définis par des expressions
  rationnelles, §\ref{subsection-rational-languages}) et
  « reconnaissables » (définis par des automates finis,
  §\ref{section-finite-automata}), qui s'avèrent coïncider
  (\ref{kleenes-theorem}) ;
\item les langages « algébriques » (définis par des grammaires hors
  contexte, §\ref{subsection-context-free-grammars}) ;
\item les langages « décidables » (définis par un algorithme
  quelconque, \ref{definition-computable-function-or-set}) et
  « semi-décidables » (définis par un algorithme qui pourrait ne pas
  terminer).
\end{itemize}
La progression logique vient de l'inclusion qui donne une hiérarchie
entre ces classes : tout langage rationnel/reconnaissable est
algébrique (\ref{rational-languages-are-algebraic}), tout langage
algébrique est décidable (\ref{algebraic-languages-are-decidable}), et
tout langage décidable est semi-décidable
(\ref{decidable-iff-semidecidable-and-complement}).  On procède donc
du plus particulier au plus général.

\thingy \textbf{Notations :} On écrit $A:=B$, ou parfois $B=:A$, pour
dire que $A$ est défini comme égal à $B$ (cela signifie simplement
$A=B$ mais en insistant sur le fait que c'est la définition du membre
du côté duquel se situent les deux points).  On note $\mathbb{N}$
l'ensemble $\{0,1,2,3\ldots\}$ des entiers naturels.

Le symbole \qedsymbol{} marque la fin d'une démonstration.

Les passages en petits caractères (ainsi que les notes en bas de page)
sont destinés à apporter des éclaircissements, des précisions ou des
compléments mais peuvent être ignorés sans nuire à la compréhension
globale du texte.

\section{Alphabets, mots et langages ; langages rationnels}\label{section-languages-and-rational-languages}

\subsection{Introduction, alphabets, mots et longueur}\label{subsection-introduction-and-words}

\thingy L'objet de ces notes, au confluent des mathématiques et de
l'informatique, est l'étude des \textbf{langages} : un langage étant
un ensemble de \textbf{mots}, eux-mêmes suites finies de
\textbf{lettres} choisies dans un \textbf{alphabet}, on va commencer
par définir ces différents termes avant de décrire plus précisément
l'objet de l'étude.

\thingy Le point de départ est donc ce que les mathématiciens
appelleront un \defin{alphabet}, et qui correspond pour les
informaticiens à un \textbf{jeu de caractères}.  Il s'agit d'un
ensemble \emph{fini}, sans structure particulière, dont les éléments
s'appellent \defin[lettre]{lettres}, ou encore
\index{caractère|see{lettre}}\textbf{caractères} dans une terminologie
plus informatique, ou parfois aussi \defin[symbole]{symboles}.

Les exemples mathématiques seront souvent donnés sur un alphabet tel
que l'alphabet à deux lettres $\{a,b\}$ ou à trois lettres
$\{a,b,c\}$.  On pourra aussi considérer l'alphabet $\{0,1\}$ appelé
\defin{binaire} (puisque l'alphabet n'a pas de structure
particulière, cela ne fait guère de différence par rapport à n'importe
quel autre alphabet à deux lettres).  Dans un contexte informatique,
des jeux de caractères (=alphabets) souvent importants sont ASCII,
Latin-1 ou Unicode : en plus de former un ensemble, ces jeux de
caractère attribuent un numéro à chacun de leurs éléments (par
exemple, la lettre A majuscule porte le numéro 65 dans ces trois jeux
de caractères), mais cette structure supplémentaire ne nous
intéressera pas ici.  Dans tous les cas, il est important pour la
théorie que l'alphabet soit \emph{fini}.

\smallskip

L'alphabet sera généralement fixé une fois pour toutes dans la
discussion, et désigné par la lettre $\Sigma$ (sigma majuscule).

\thingy Un \defin{mot} sur l'alphabet $\Sigma$ est une suite finie de
lettres (éléments de $\Sigma$) ; dans la terminologie informatique, on
parle plutôt de \index{caractères (chaîne de)|see{chaîne de caractères}}\defin{chaîne de caractères}, qui est une suite finie
(=liste) de caractères.  Le mot est désigné en écrivant les lettres
les unes à la suite des autres : autrement dit, si $x_1,\ldots,x_n \in
\Sigma$ sont des lettres, le mot formé par la suite finie
$x_1,\ldots,x_n$ est simplement écrit $x_1\cdots x_n$.

À titre d'exemple, $abbcab$ est un mot sur l'alphabet $\Sigma =
\{a,b,c,d\}$, et \texttt{foobar} est un mot (= chaîne de caractères)
sur l'alphabet ASCII.  (Dans un contexte informatique, il est fréquent
d'utiliser une sorte de guillemet pour délimiter les chaînes de
caractères : on écrira donc \texttt{\char`\"foobar\char`\"} pour
parler du mot en question.  Dans ces notes, nous utiliserons peu cette
convention.)

\medskip

L'ensemble de tous les mots sur un alphabet $\Sigma$ sera
désigné $\Sigma^*$ (on verra en \ref{kleene-star} ci-dessous cette
notation comme un cas particulier d'une construction « étoile » plus
générale).  Par exemple, si $\Sigma = \{0,1\}$, alors $\Sigma^*$ est
l'ensemble (infini !) dont les éléments sont toutes les suites finies
binaires (= suites finies de $0$ et de $1$).  Ainsi, écrire « $w \in
\Sigma^*$ » signifie « $w$ est un mot sur l'alphabet $\Sigma$ ».

\medskip

{\footnotesize\thingy Typographiquement, on essaiera autant que
  possible de désigner des mots par des variables mathématiques telles
  que $u,v,w$, tandis que $x,y,z$ désigneront plutôt des lettres
  quelconques dans un alphabet (quant à $a,b,c$, ils serviront de
  lettres dans les exemples).  Il n'est malheureusement pas possible
  d'être complètement systématique (il arrivera que $x$ désigne un
  mot) : on cherchera donc à toujours rappeler le type de toute
  variable en écrivant, par exemple, $t \in\Sigma$ ou $t
  \in\Sigma^*$.\par}

\thingy\label{length-of-word} Le nombre $n$ de lettres dans un mot $w
\in \Sigma^*$ est appelé la \defin{longueur} du mot, et généralement
notée $|w|$ ou bien $\ell(w)$ : autrement dit, si $x_1,\ldots,x_n \in
\Sigma$, alors la longueur $|x_1\cdots x_n|$ du mot $x_1\cdots x_n$,
vaut $n$.  Ceci coïncide bien avec la notion usuelle de longueur d'une
chaîne de caractères en informatique.  À titre d'exemple, sur
l'alphabet $\Sigma=\{a,b,c,d\}$, la longueur du mot $abbcab$ vaut $6$
(on écrira $|abbcab|=6$ ou bien $\ell(abbcab)=6$).

\thingy Quel que soit l'alphabet, il existe un unique mot de
longueur $0$, c'est-à-dire un unique mot n'ayant aucune lettre, appelé
le \defin{mot vide} (ou la \textbf{chaîne [de caractères] vide}).
Étant donné qu'il n'est pas commode de désigner un objet par une
absence de symbole, on introduit un symbole spécial, généralement
$\varepsilon$, pour désigner ce mot vide : on a donc
$|\varepsilon|=0$.  On souligne que le symbole $\varepsilon$
\underline{ne fait pas partie} de l'alphabet $\Sigma$, c'est un
symbole \emph{spécial} qui a été introduit pour désigner le mot vide.
(Lorsque les mots sont délimités par des guillemets, comme il est
usage pour les chaînes de caractères en informatique, le mot vide n'a
pas besoin d'un symbole spécial : il s'écrit juste
\texttt{\char`\"\char`\"} — sans aucun caractère entre les
guillemets.)

{\footnotesize Lorsque l'alphabet $\Sigma$ est \emph{vide},
  c'est-à-dire $\Sigma=\varnothing$, alors le mot vide est le seul mot
  qui existe : on a $\Sigma^*=\{\varepsilon\}$ dans ce cas.  C'est la
  seule situation où l'ensemble $\Sigma^*$ des mots est un ensemble
  fini.  Dans la suite, nous négligerons parfois ce cas particulier,
  qu'on pourra oublier : c'est-à-dire que nous ferons parfois
  l'hypothèse tacite que $\Sigma \neq \varnothing$.\par}

La notation $\Sigma^+$ est parfois utilisée pour désigner l'ensemble
des mots \emph{non vides} sur l'alphabet $\Sigma$ (par opposition à
$\Sigma^*$ qui désigne l'ensemble de tous les mots, y compris le mot
vide) ; on verra en \ref{kleene-plus} ci-dessous que c'est un cas
particulier d'une construction plus générale.

\thingy\label{convention-on-words-of-length-one} Les mots d'une seule
lettre sont naturellement en correspondance avec les lettres
elles-mêmes : on identifiera souvent tacitement, quoique un peu
abusivement, une lettre $x\in\Sigma$ et le mot de longueur $1$ formé
de la seule lettre $x$.  (En informatique, cette identification entre
\emph{caractères} et \emph{chaînes de caractères de longueur $1$} est
faite par certains langages de programmation, mais pas par tous :
\textit{caveat programmator}.)  Cette convention permet d'écrire par
exemple $\Sigma \subseteq \Sigma^*$ ou bien $|x|=1 \liff x\in\Sigma$.

\thingy\label{number-of-words-of-length-n} Si le cardinal de
l'alphabet $\Sigma$ vaut $\#\Sigma = N$, alors, pour chaque $n$, le
nombre de mots de longueur exactement $n$ est égal à $N^n$
(combinatoire classique).  Le nombre de mots de longueur $\leq n$ vaut
donc $1 + N + \cdots + N^n = \frac{N^{n+1}-1}{N-1}$ (somme d'une série
géométrique).


\subsection{Concaténation de mots, préfixes, suffixes, facteurs, sous-mots}

\thingy Si $u := x_1\cdots x_m$ et $v := y_1\cdots y_n$ sont deux
mots, de longueurs respectives $m$ et $n$, sur un même
alphabet $\Sigma$, alors on définit un mot $uv := x_1\cdots x_m
y_1\cdots y_n$ de longueur $m+n$, dont les lettres sont obtenues en
mettant bout à bout celles de $u$ puis celles de $v$ (dans cet ordre),
et on l'appelle \defin{concaténation} (ou, si cela ne prête pas à
confusion, simplement \index{produit (de mots)|see{concaténation}}\textbf{produit}) des mots $u$ et $v$.  (Dans un
contexte informatique, on parle de concaténation de chaînes de
caractères.)

\thingy Parmi les propriétés de la concaténation, signalons les faits
suivants :
\begin{itemize}
\item le mot vide $\varepsilon$ est « \textbf{neutre} » pour la
  concaténation, ce qui signifie par définition : $\varepsilon w = w
  \varepsilon = w$ quel que soit le mot $w \in \Sigma^*$ ;
\item la concaténation est « \textbf{associative} », ce qui signifie
  par définition : $u(vw) = (uv)w$ quels que soient les mots $u,v,w
  \in \Sigma^*$ (on peut donc noter $uvw$ sans parenthèse).
\end{itemize}

On peut traduire de façon savante ces deux propriétés en une phrase :
l'ensemble $\Sigma^*$ est un \defin{monoïde}, d'élément
neutre $\varepsilon$, pour la concaténation (cela signifie exactement
ce qui vient d'être dit).

\thingy On a par ailleurs $|uv| = |u| + |v|$ (la longueur de la
concaténation de deux mots est la somme des concaténations), et on
rappelle par ailleurs que $|\varepsilon| = 0$ ; on peut traduire cela
de manière savante : la longueur est un \textbf{morphisme de monoïdes}
entre le monoïde $\Sigma^*$ des mots (pour la concaténation) et le
monoïde $\mathbb{N}$ des entiers naturels (pour l'addition) (cela
signifie exactement ce qui vient d'être dit).

{\footnotesize\thingy\label{universal-property-remark}
  \textbf{Complément :} Le monoïde $\Sigma^*$
  possède la propriété suivante par rapport à l'ensemble $\Sigma$ : si
  $M$ est un monoïde quelconque (c'est-à-dire un ensemble muni d'une
  opération binaire associative $\cdot$ et d'un élément $e$ neutre
  pour cette opération), et si $\psi\colon \Sigma\to M$ est une
  application quelconque, alors il existe un unique morphisme de
  monoïdes $\hat\psi\colon \Sigma^* \to M$ (c'est-à-dire une
  application préservant le neutre et l'opération binaire) tel que
  $\hat\psi(x) = \psi(x)$ si $x\in\Sigma$.  (Démonstration : on a
  nécessairement $\hat\psi(x_1\cdots x_n) = \psi(x_1)\cdots
  \psi(x_n)$, or ceci définit bien un morphisme comme annoncé.)  On
  dit qu'il s'agit là d'une propriété « universelle », et plus
  précisément que $\Sigma^*$ est le \textbf{monoïde libre} sur
  l'ensemble $\Sigma$.  Par exemple, le morphisme « longueur »
  $\ell\colon\Sigma^*\to\mathbb{N}$ est le $\ell = \hat\psi$ obtenu en
  appliquant cette propriété à la fonction (constante) $\psi(x) = 1$
  pour tout $x\in\Sigma$.\par}

\thingy\label{powers-of-a-word} Lorsque $w \in \Sigma^*$ et $r \in
\mathbb{N}$, on définit un mot $w^r$ comme la concaténation de $r$
mots tous égaux à $w$, autrement dit, comme la répétition $r$
fois du mot $w$.  Formellement, on définit par récurrence :
\begin{itemize}
\item $w^0 = \varepsilon$ (le mot vide),
\item $w^{r+1} = w^r w$.
\end{itemize}
(Ces définitions valent, d'ailleurs, dans n'importe quel monoïde.  On
peut constater que $w^r w^s = w^{r+s}$ quels que soient
$r,s\in\mathbb{N}$.)  On a bien sûr $|w^r| = r|w|$.

\smallskip

Cette définition sert notamment à désigner de façon concise les mots
comportant des répétitions d'une même lettre : par exemple, le mot
$aaaaa$ peut s'écrire tout simplement $a^5$, et le mot $aaabb$ peut
s'écrire $a^3 b^2$.  (De même que pour le mot vide, il faut souligner
que ces exposants \emph{ne font pas partie} de l'alphabet.)

\thingy Lorsque $u,v,w \in \Sigma^*$ vérifient $w = uv$, autrement dit
lorsque le mot $w$ est la concaténation des deux mots $u$ et $v$, on
dira également :
\begin{itemize}
\item que $u$ est un \defin{préfixe} de $w$, ou
\item que $v$ est un \defin{suffixe} de $w$.
\end{itemize}

De façon équivalente, si $w = x_1\cdots x_n$ (où $x_1,\ldots,x_n \in
\Sigma$) est un mot de longueur $n$, et si $0\leq k\leq n$ est un
entier quelconque compris entre $0$ et $n$, on dira que $u :=
x_1\cdots x_k$ (c'est-à-dire, le mot formé des $k$ premières lettres
de $w$, dans le même ordre) est le \textbf{préfixe de longueur $k$}
de $w$, et que $v := x_{k+1}\cdots x_n$ (mot formé des $n-k$ dernières
lettres de $w$, dans le même ordre) est le \textbf{suffixe de
  longueur $n-k$} de $w$.  Il est clair qu'il s'agit bien là de
l'unique façon d'écrire $w = uv$ avec $|u|=k$ et $|v|=n-k$, ce qui
fait le lien avec la définition donnée au paragraphe précédent ;
parfois on dira que $v$ est le suffixe \defin[correspondant (préfixe
  ou suffixe)]{correspondant} à $u$ ou que $u$ est le préfixe
correspondant à $v$ (dans le mot $w$).

Le mot vide est préfixe et suffixe de n'importe quel mot.  Le mot $w$
lui-même est aussi un préfixe et un suffixe de lui-même.  Entre les
deux, pour n'importe quelle longueur $k$ donnée, il existe un unique
préfixe et un unique suffixe de longueur $k$.  (Il peut tout à fait se
produire que le préfixe et le suffixe de longueur $k$ soient égaux
pour d'autres $k$ que $0$ et $|w|$, comme le montre l'exemple qui
suit.)

\smallskip

À titre d'exemple, le mot $abbcab$ sur l'alphabet $\Sigma=\{a,b,c,d\}$
a les sept préfixes suivants, rangés par ordre croissant de longueur :
$\varepsilon$ (le mot vide), $a$, $ab$, $abb$, $abbc$, $abbca$ et
$abbcab$ lui-même ; il a les sept suffixes suivants, rangés par ordre
croissant de longueur : $\varepsilon$ (le mot vide), $b$, $ab$, $cab$,
$bcab$, $bbcab$ et $abbcab$ lui-même.  Le suffixe correspondant au
préfixe $abb$ est $cab$ puisque $abbcab = (abb)(cab)$.

\thingy Comme généralisation à la fois de la notion de préfixe et de
celle de suffixe, on a la notion de facteur : si $u_0,v,u_1 \in
\Sigma^*$ sont trois mots quelconques sur un même alphabet $\Sigma$,
et si $w = u_0 v u_1$ est leur concaténation, on dira que $v$ est un
\defin{facteur} de $w$.  Alors qu'un préfixe ou suffixe du mot $w$
est déterminé simplement par sa longueur, un facteur est déterminé par
sa longueur et l'emplacement à partir duquel il commence.

\smallskip

À titre d'exemple, les facteurs du mot $abbcab$ sont : $\varepsilon$
(le mot vide), $a$, $b$, $c$, $ab$, $bb$, $bc$, $ca$, $abb$, $bbc$,
$bca$, $cab$, $abbc$, $bbca$, $bcab$, $abbca$, $bbcab$ et $abbcab$
lui-même.

\smallskip

Dans un contexte informatique, ce que nous appelons ici « facteur »
est souvent appelé « sous-chaîne [de caractères] ».  Il ne faut
cependant pas confondre ce concept avec celui de sous-mot défini
ci-dessous.

\thingy\label{definition-subword} Si $w = x_1\cdots x_n$ est un mot de
longueur $n$, on appelle \defin{sous-mot} de $w$ un mot de la forme
$x_{i_1} \cdots x_{i_k}$ où $1\leq i_1 < \cdots < i_k \leq n$.  En
plus clair, cela signifie que $v$ est obtenu en ne gardant que
certaines lettres du mot $w$, dans le même ordre, mais en en effaçant
d'autres ; à la différence du concept de facteur, celui de sous-mot
n'exige pas que les lettres gardées soient consécutives.

\smallskip

À titre d'exemple, le mot $acb$ est un sous-mot du mot $abbcab$
(obtenu en gardant les lettres soulignées ici :
$\underline{a}bb\underline{c}a\underline{b}$ ; pour se rattacher à la
définition ci-dessus, on pourra prendre $(i_1,i_2,i_3) = (1,4,6)$).

{\footnotesize\thingy Les remarques de dénombrement suivantes peuvent
  aider à mieux comprendre les notations de préfixe, suffixe, facteur
  et sous-mot : si $w$ est un mot de longueur $n$, alors il a
\begin{itemize}
\item exactement $n+1$ préfixes (car un préfixe est déterminé par sa
  longueur $k$ entre $0$ et $n$),
\item exactement $n+1$ suffixes (raison analogue),
\item au plus $\sum_{k=1}^n (n+1-k) + 1 = \frac{1}{2}(n^2+n+2)$
  facteurs (car un facteur est déterminé par sa longueur $k$ et son
  point de départ qui peut être choisi parmi $n+1-k$ possibilités, le
  $+1$ final étant mis pour le facteur vide),
\item au plus $2^n$ sous-mots (car un sous-mot est déterminé en
  choisissant, pour chacune des $n$ lettres, si on l'efface ou la
  conserve).
\end{itemize}
Le nombre exact peut être plus petit en cas de coïncidences entre
certains choix (par exemple, $aaa$ n'a que $4$ facteurs, $\varepsilon,
a, aa, aaa$ alors que $abc$ en a bien $\frac{1}{2}(3^2+3+2) = 7$) ;
mais les bornes ci-dessus sont effectivement atteintes pour certains
mots.\par}

\thingy\label{definition-mirror-word} Si $w = x_1\cdots x_n$, où
$x_1,\ldots,x_n \in \Sigma$, est un mot de longueur $n$ sur un
alphabet $\Sigma$, alors on définit son mot \defin{miroir} ou
\index{transposé (mot)|see{miroir}}\textbf{transposé}, parfois noté
$w^{\textsf{R}}$ ou $w^{\textsf{T}}$ (parfois les exposants sont
écrits à gauche), comme le mot $x_n\cdots x_1$ dont les lettres sont
les mêmes que celles de $w$ mais dans l'ordre inverse.  À titre
d'exemple, $(ababb)^{\textsf{R}} = bbaba$.  On remarquera que $(w_1
w_2)^{\textsf{R}} = w_2^{\textsf{R}} w_1^{\textsf{R}}$ si $w_1,w_2$
sont deux mots quelconques.

\smallskip

Un mot $w$ est dit \defin{palindrome} lorsque $w = w^{\textsf{R}}$.
Par exemple, $abba$ est un palindrome sur $\{a,b,c,d\}$ (ou bien le
mot « ressasser » sur l'alphabet du français).

{\footnotesize\thingy\label{number-of-occurrences-of-letter} La
  notation suivante est souvent utile : si $w$ est un mot sur un
  alphabet $\Sigma$ et si $z$ est une lettre (= élément de $\Sigma$),
  on note $|w|_z$ le \index{occurrences (nombre d')}nombre total
  d'occurrences de la lettre $z$ dans $w$.  À titre d'exemple, sur
  l'alphabet $\Sigma=\{a,b,c,d\}$, le nombre d'occurrences des
  différentes lettres dans le mot $abbcab$ sont : $|abbcab|_a=2$,
  $|abbcab|_b=3$, $|abbcab|_c=1$ et $|abbcab|_d=0$.

Formellement, on peut définir $|w|_z$ de la façon suivante : si $w =
x_1 \cdots x_n$ où $x_1,\ldots,x_n \in \Sigma$, alors $|w|_z$ est le
cardinal de l'ensemble $\{i : x_i = z\}$.  On peut remarquer qu'on a :
$|w| = \sum_{z\in\Sigma} |w|_z$ (i.e., la longueur de $w$ est la somme
des nombres d'occurrences dans celui-ci des différentes lettres de
l'alphabet).\par}


\subsection{Langages et opérations sur les langages}\label{subjection-languages}

\thingy Un \defin{langage} $L$ sur l'alphabet $\Sigma$ est simplement
un ensemble de mots sur $\Sigma$.  Autrement dit, il s'agit d'un
sous-ensemble (= une partie) de l'ensemble $\Sigma^*$ (de tous les mots
sur $\Sigma$) : en symboles, $L \subseteq \Sigma^*$.

On souligne qu'on ne demande pas que $L$ soit fini (mais il peut
l'être, auquel cas on parlera fort logiquement de « langage fini »).

\thingy À titre d'exemple, l'ensemble $\{d,dc,dcc,dccc,dcccc,\ldots\}
= \{dc^r \colon r\in\mathbb{N}\}$ des mots formés d'un $d$ suivi d'un
nombre quelconque (éventuellement nul) de $c$ est un langage sur
l'alphabet $\Sigma = \{a,b,c,d\}$.  On verra plus loin que ce langage
est « rationnel » (et pourra être désigné par l'expression rationnelle
$dc{*}$).

Voici quelques autres exemples de langages :
\begin{itemize}
\item Le langage (fini) $\{foo,bar,baz\}$ constitué des seuls trois
  mots $foo$, $bar$, $baz$ sur l'alphabet $\Sigma = \{a,b,f,o,r,z\}$.
\item Le langage (fini) constitué des mots de longueur exactement $42$
  sur l'alphabet $\Sigma = \{p,q,r\}$.  Comme on l'a vu
  en \ref{number-of-words-of-length-n}, cet ensemble a pour cardinal
  exactement $3^{42}$.
\item Le langage constitué des mots de longueur exactement $1$ sur un
  alphabet $\Sigma$ (= mots de une seule lettre), qu'on peut identifier
  à $\Sigma$ lui-même (en identifiant un mot de une lettre à la lettre
  en question, cf. \ref{convention-on-words-of-length-one}).
\item Le langage (fini) constitué du seul mot vide (= mot de longueur
  exactement $0$) sur l'alphabet, disons, $\Sigma = \{p,q,r\}$.  Ce
  langage $\{\varepsilon\}$ a pour cardinal $1$ (ou $3^0$ si on veut).
  Il ne faut pas le confondre avec le suivant :
\item Le langage vide, qui ne contient aucun mot (sur un alphabet
  quelconque).  Ce langage a pour cardinal $0$.
\item Le langage sur l'alphabet $\Sigma=\{a,b\}$ constitué des mots
  qui commencent par trois $a$ consécutifs : ou, si on préfère, qui
  ont le mot $aaa$ comme préfixe.
\item Le langage sur l'alphabet $\Sigma=\{a,b\}$ constitué des mots
  qui contiennent trois $a$ consécutifs ; ou, si on préfère, qui ont
  $aaa$ comme facteur.
\item Le langage sur l'alphabet $\Sigma=\{a,b\}$ constitué des mots
  qui contiennent au moins trois $a$, non nécessairement consécutifs ;
  ou, si on préfère, qui ont $aaa$ comme sous-mot.
\item Le langage sur l'alphabet $\Sigma=\{a\}$ constitué de tous les
  mots dont la longueur est un nombre premier ($L = \{aa, aaa, a^5,
  a^7, a^{11},\ldots\}$).  Ce langage est infini.
\item Le langage sur l'alphabet $\Sigma=\{0,1\}$ constitué de tous les
  mots commençant par un $1$ et qui, interprétés comme un nombre écrit
  en binaire, désignent un nombre premier ($L = \{10, 11, 101, 111,
  1011, \ldots\}$).
\item Le langage sur l'alphabet Unicode constitué de tous les mots qui
  représentent un document XML bien-formé d'après la spécification
  XML 1.0.
\end{itemize}

\thingy On pourrait aussi considérer un langage (sur
l'alphabet $\Sigma$) comme une \emph{propriété} des mots (sur
l'alphabet en question).  Précisément, si $P$ est une propriété qu'un
mot $w \in \Sigma^*$ peut ou ne pas avoir, on considère le langage
$L_P = \{w \in \Sigma^* : w \text{~a la propriété~} P\}$, et
inversement, si $L \subseteq \Sigma^*$ est un langage, on considère la
propriété « appartenir à $L$ » : en identifiant la propriété et le
langage qu'on vient d'associer l'un à l'autre (par exemple, le langage
des mots commençant par $a$ et la propriété « commencer par $a$ »), un
langage pourrait être considéré comme une propriété des mots.

\smallskip

{\footnotesize(Ce qui précède n'a rien de spécifique aux langages :
  une partie d'un ensemble $E$ quelconque peut être identifiée à une
  propriété que les éléments de $E$ peuvent ou ne pas avoir, à savoir,
  appartenir à la partie en question.)\par}

\smallskip

On évitera de faire cette identification pour ne pas introduire de
complication, mais il est utile de la garder à l'esprit : par exemple,
dans un langage de programmation fonctionnel, un « langage » au sens
de ces notes peut être considéré comme une fonction (pure,
c'est-à-dire, déterministe et sans effet de bord) prenant en entrée
une chaîne de caractères et renvoyant un booléen.

\thingy\label{union-and-intersection-of-languages} Si $L_1$ et $L_2$
sont deux langages sur un même alphabet $\Sigma$ (autrement dit,
$L_1,L_2 \subseteq \Sigma^*$), on peut former les langages
\index{réunion (de langages)|see{union}}\defin[union (de langages)]{union} $L_1\cup L_2$ et
\defin[intersection (de langages)]{intersection} $L_1\cap L_2$ qui
sont simplement les opérations ensemblistes usuelles (entre parties
de $\Sigma^*$).

\smallskip

Les opérations correspondantes sur les propriétés de mots sont
respectivement le « ou logique » (=disjonction) et le « et logique »
(=conjonction) : à titre d'exemple, sur $\Sigma = \{a,b\}$ si $L_1$
est le langage des mots commençant par $a$ et $L_2$ le langage des
mots finissant par $b$, alors $L_1 \cup L_2$ est le langage des mots
commençant par $a$ \emph{ou bien} finissant par $b$, tandis que $L_1
\cap L_2$ est le langage des mots commençant par $a$ \emph{et}
finissant par $b$.

\thingy Si $L$ est un langage sur l'alphabet $\Sigma$, autrement dit
$L \subseteq \Sigma^*$, on peut former le langage $\Sigma^*\setminus
L$, parfois noté simplement $\overline L$ si ce n'est pas ambigu, dit
\defin[complémentaire (langage)]{complémentaire} de $L$, et qui est
simplement l'ensemble des mots sur $\Sigma$ \emph{n'appartenant pas}
à $L$.  L'opération correspondante sur les propriétés de mots est la
négation logique.

\smallskip

À titre d'exemple, sur $\Sigma=\{a,b\}$, si $L$ est le langage des
mots commençant par $a$, alors $\overline{L}$ est le langage des mots
ne commençant pas par $a$, c'est-à-dire, la réunion de
$\{\varepsilon\}$ et du langage des mots commençant par $b$ (car sur
$\Sigma=\{a,b\}$, un mot ne commençant pas par $a$ est vide ou bien
commence par $b$).

\thingy\label{concatenation-of-languages} Si $L_1$ et $L_2$ sont deux
langages sur un même alphabet $\Sigma$ (autrement dit, $L_1,L_2
\subseteq \Sigma^*$), on peut former le langage \defin[concaténation
  (de langages)]{concaténation} $L_1 L_2$ : il est défini comme
l'ensemble des mots $w$ qui peuvent s'écrire comme concaténation d'un
mot $w_1$ de $L_1$ et d'un mot $w_2$ de $L_2$, soit
\[
\begin{aligned}
L_1 L_2 &:= \{w_1 w_2 : w_1 \in L_1,\, w_2 \in L_2\}\\
 &= \{w \in \Sigma^* : \exists w_1 \in L_1\, \exists w_2 \in L_2\,(w = w_1 w_2)\}\\
\end{aligned}
\]

\smallskip

À titre d'exemple, sur l'alphabet $\Sigma = \{a,b,c,d\}$, si on a $L_1
= \{a,bb\}$ et $L_2 = \{bc, cd\}$ alors $L_1 L_2 = \{abc, acd, bbbc,
bbcd\}$.

\thingy Si $L$ est un langage sur l'alphabet $\Sigma$, autrement dit
$L \subseteq \Sigma^*$, et si $r \in \mathbb{N}$, on peut définir un
langage $L^r$, par analogie avec \ref{powers-of-a-word}, comme le
langage $L^r = \{w_1\cdots w_r : w_1,\ldots,w_r \in L\}$ constitué des
concaténation de $r$ mots appartenant à $L$, ou si on préfère, par
récurrence :
\begin{itemize}
\item $L^0 = \{\varepsilon\}$,
\item $L^{r+1} = L^r L$.
\end{itemize}

\smallskip

À titre d'exemple, sur l'alphabet $\Sigma = \{a,b,c,d\}$, si on a $L =
\{a,bb\}$, alors $L^2 = \{aa, abb, bba, bbbb\}$ et $L^3 = \{aaa, aabb,
abba, abbbb, \penalty-100 bbaa, bbabb, bbbba, bbbbbb\}$.

\medskip

\emph{Attention}, $L^r$ n'est pas le langage $\{w^r : w\in L\}$
constitué des répétitions $r$ fois ($w^r$) des mots $w$ de $L$ : c'est
le langage des concaténations de $r$ mots appartenant à $L$ \emph{mais
  ces mots peuvent être différents}.  À titre d'exemple, si $L =
\{a,b\}$ alors $L^r$ est le langage constitué des $2^r$ mots de
longueur exactement $r$ sur $\{a,b\}$, ce n'est pas l'ensemble à deux
éléments $\{a^r, b^r\}$ constitué des seuls deux mots $a^r = aaa\cdots
a$ et $b^r = bbb\cdots b$.

\thingy\label{kleene-star} Si $L$ est un langage sur
l'alphabet $\Sigma$, on définit enfin l'\index{Kleene (étoile
  de)|see{étoile de Kleene}}\defin{étoile de Kleene} $L^*$ de $L$
comme le langage constitué des concaténations d'un nombre
\emph{quelconque} de mots appartenant à $L$, c'est-à-dire la réunion
de tous les langages $L^r$ mentionnés ci-dessus :
\[
\begin{aligned}
L^* &:= \bigcup_{r=0}^{+\infty} L^r = \bigcup_{r\in\mathbb{N}} L^r\\
&= \{w_1\cdots w_r : r\in\mathbb{N},\, w_1,\ldots,w_r\in L\}\\
\end{aligned}
\]

\smallskip

À titre d'exemple, sur l'alphabet $\Sigma = \{a,b,c,d\}$, si on a $L =
\{a,bb\}$, alors on a $L^* = \{\varepsilon, a, bb, \penalty-200 aa,
abb, bba, bbbb, \penalty-200 aaa, aabb, abba, abbbb, \penalty-100
bbaa, bbabb, bbbba, bbbbbb, \ldots\}$.

\smallskip

Comme ci-dessus, il faut souligner que les mots $w_1,\ldots,w_r$
concaténés n'ont pas à être égaux : notamment, $\{a,b\}^*$ est le
langage constitué de tous les mots sur l'alphabet $\{a,b\}$, pas le
langage $\{a\}^* \cup \{b\}^*$ constitué des mots obtenus en répétant
la lettre $a$ ou en répétant la lettre $b$.

\smallskip

On remarquera que la définition de $L^*$ ci-dessus redonne bien,
lorsqu'on l'applique à l'alphabet $\Sigma$ lui-même (considéré comme
langage des mots de longueur $1$), l'ensemble $\Sigma^*$ de tous les
mots : la notation $\Sigma^*$ est donc justifiée \textit{a
  posteriori}.

\smallskip

Le mot vide appartient toujours à $L^*$ (quel que soit $L$) puisque
$L^0 = \{\varepsilon\}$ et qu'on peut prendre $r=0$ ci-dessus
(autrement dit, le mot vide est la concaténation de zéro mots de $L$).

\smallskip

{\footnotesize On peut par ailleurs montrer que $L^*$ est le plus
  petit (pour l'inclusion) langage $M$ tel que $M = \{\varepsilon\}
  \cup LM$.  Cette observation sera implicite par exemple
  dans \ref{cfg-union-concatenation-and-star} ci-dessous, et motive
  aussi la construction $L^+$ qui suit.\par}

\thingy\label{kleene-plus} On introduit parfois la notation $L^+ :=
\bigcup_{r=1}^{+\infty} L^r = \{w_1\cdots w_r : r>0,\penalty-100\,
w_1,\ldots,w_r\in L\}$ pour l'ensemble des mots formés par
concaténation d'un nombre \emph{non nul} de mots de $L$.  Lorsque le
mot vide $\varepsilon$ n'appartient pas déjà à $L$, ce langage $L^+$
diffère de $L^*$ seulement en ce qu'il ne contient pas $\varepsilon$ ;
tandis que si $\varepsilon$ appartient déjà à $L$, alors $L^+$ est
égal à $L^*$.  En toute généralité, on a $L^+ = LL^*$.

\thingy\label{definition-mirror-language} En rappelant la définition
du mot miroir faite en \ref{definition-mirror-word}, si $L$ est un
langage sur l'alphabet $\Sigma$, on définit le langage miroir
$L^{\mathsf{R}}$ comme l'ensemble des mots miroirs des mots de $L$,
c'est-à-dire $L^{\mathsf{R}} = \{w^{\mathsf{R}} : w \in L\}$.

\bigbreak

\thingy\label{set-of-languages} De même que l'ensemble des mots sur un
alphabet $\Sigma$ admet une notation, à savoir $\Sigma^*$, on peut
introduire une notation pour l'ensemble de tous les \emph{langages}
(= ensembles de mots) sur $\Sigma$ : ce sera $\mathscr{P}(\Sigma^*)$.
Il s'agit d'un cas particulier de la construction « ensemble des
parties » (si $E$ est un ensemble, $\mathscr{P}(E) = \{A : A\subseteq
E\}$ est l'ensemble de tous les sous-ensembles de $A$ ; c'est un
axiome de la théorie des ensembles qu'un tel ensemble existe bien).
On pourrait donc écrire « $L \in \mathscr{P}(\Sigma^*)$ » comme
synonyme de « $L\subseteq \Sigma^*$ » ou de « $L$ est un langage
sur $\Sigma$ » ; on évitera cependant de le faire, car cette notation
est plus lourde qu'utile.

\smallskip

{\footnotesize Il sera marginalement question dans ces notes de
  « classes de langages » : une classe de langages est un ensemble de
  langages (c'est-à-dire une partie de $\mathscr{P}(\Sigma^*)$, ou si
  on préfère un élément de $\mathscr{P}(\mathscr{P}(\Sigma^*))$).  On
  ne développera pas de théorie générale des classes de langages et on
  n'en parlera pas de façon systématique, mais on parlera de certaines
  classes de langages importantes : la classe des langages rationnels
  ou reconnaissables (§\ref{subsection-rational-languages} à
  §\ref{section-recognizable-languages}), la classe des langages
  algébriques (§\ref{section-context-free-grammars}), la classe des
  langages décidables (§\ref{section-computability}) et la classe des
  langages semi-décidable (\textit{ibid.}).\par}


\subsection{Langages rationnels et expressions rationnelles}\label{subsection-rational-languages}

\thingy \textbf{Introduction :} Les langages qui vont jouer le rôle le
plus important dans ces notes sont les langages dits « rationnels »
(qui seront définis de \ref{rational-languages} à
\ref{regular-expressions} ci-dessous, et qui s'avéreront être les
mêmes que les langages « reconnaissables » par automates finis définis
en \ref{definition-recognizable-language}).

Pour donner un premier aperçu informel de ce dont il s'agit avant de
passer à une définition précise, commençons par donner quelques
exemples, disons, sur l'alphabet $\Sigma = \{a,b,c,d\}$, de langages
rationnels :
\begin{itemize}
\item le langage constitué des mots contenant au moins un $d$ ;
\item le langage constitué des mots ne contenant aucun $d$ ;
\item le langage constitué des mots commençant par un $d$ ;
\item le langage constitué des mots commençant par un $d$ et
  finissant par un $c$ ;
\item le langage constitué des mots commençant par un $d$ et finissant
  par un $c$, et dont toutes les lettres intermédiaires sont des $a$
  ou des $b$, c'est-à-dire, le langage constitué des mots de la forme
  suivante : un $d$, puis un nombre quelconque de $a$ et de $b$, et
  enfin un $c$ ;
\item le langage constitué des mots de la forme suivante : un $d$,
  puis un nombre quelconque de mots qui peuvent être soit $a$ soit
  $bb$, et enfin un $c$.
\end{itemize}

Ce dernier exemple est assez typique : on dira plus loin qu'il s'agit
du langage « dénoté » par l'expression rationnelle $d(a|bb){*}c$, à
lire comme « un $d$, puis autant qu'on veut (“$*$”) de mots qui
peuvent être soit (“$|$”) un $a$ soit $bb$, et enfin un $c$ ».  Les
constructions essentielles qui permettront de fabriquer les langages
rationnels sont : la réunion (j'autorise ceci \emph{ou} cela, par
exemple « un $a$ \emph{ou bien} $bb$ »,
cf. \ref{union-and-intersection-of-languages}), la concaténation (je
demande ceci \emph{puis} cela, par exemple « un $d$, \emph{puis} une
suite quelconque de $a$ ou de $b$ »,
cf. \ref{concatenation-of-languages}), et l'étoile de Kleene
(représentant une répétition quelconque d'un certain motif,
cf. \ref{kleene-star}).

\medskip

L'importance des langages rationnels, et des expressions rationnelles
(=régulières) qui les décrivent, vient :
\begin{itemize}
\item du point de vue théorique : de ce qu'ils forment une classe à la
  fois suffisamment simple pour pouvoir être étudiée, suffisamment
  riche pour contenir toutes sortes de langages intéressants,
  suffisamment naturelle pour être définie de plusieurs manières
  différentes, et suffisamment flexible pour être stable un certain
  nombre d'opérations ;
\item du point de vue pratique et informatique : de ce qu'ils sont
  algorithmiquement commodes à manier (grâce à la théorie) et
  suffisent à représenter beaucoup de « recherches » qu'on a
  naturellement envie de faire dans un texte, de « motifs » qu'on peut
  vouloir trouver, ou de « contraintes » qu'on peut vouloir imposer à
  une chaîne de caractères.
\end{itemize}

\smallskip

Passons maintenant à une définition plus précise.

\bigbreak

\thingy\label{rational-languages} Soit $\Sigma$ un alphabet.  On va
considérer les langages de base triviaux suivants :
\begin{itemize}
\item le langage vide $\varnothing$,
\item le langage constitué du seul mot vide, $\{\varepsilon\}$, et
\item les langages constitués d'un seul mot lui-même formé d'une seule
  lettre, $\{x\}$ pour chaque $x\in\Sigma$,
\end{itemize}
et on va les combiner par les opérations dites « rationnelles »
suivantes :
\begin{itemize}
\item la réunion $(L_1,L_2) \mapsto L_1 \cup L_2$
  (discutée en \ref{union-and-intersection-of-languages}),
\item la concaténation $(L_1,L_2) \mapsto L_1 L_2$
  (définie en \ref{concatenation-of-languages}), et
\item l'étoile de Kleene $L \mapsto L^*$
  (définie en \ref{kleene-star}).
\end{itemize}

On obtient ainsi une certaine famille de langages (cf. ci-dessous pour
une définition plus précise), qu'on appelle \defin[rationnel
  (langage)]{langages rationnels} : les langages rationnels sont par
définition exactement ceux qui peuvent s'obtenir à partir des langages
de base énumérés ci-dessus par application (un nombre \emph{fini} de
fois) des opérations qu'on vient de dire.

Autrement dit, la réunion de deux langages rationnels, la
concaténation de deux langages rationnels, et l'étoile de Kleene d'un
langage rationnel, sont rationnels ; et les langages rationnels sont
exactement ceux qu'on obtient ainsi à partir des langages de base.

\smallskip

À titre d'exemple, sur l'alphabet $\{a,b,c,d\}$, comme le langage
$\{c\}$ (constitué du seul mot $c$) est rationnel, son étoile de
Kleene, c'est-à-dire $\{c\}^* = \{\varepsilon, c, cc, ccc,
cccc,\ldots\}$, est rationnel, et comme $\{d\}$ l'est aussi, la
concaténation $\{d\}(\{c\}^*) = \{d, dc, dcc, dccc, \ldots\}$ est
encore un langage rationnel.

\thingy\label{stable-under-rational-operations} Formellement, la
définition des langages rationnelles est la suivante : un ensemble
$\mathscr{C} \subseteq \mathscr{P}(\Sigma^*)$ de langages (où
$\mathscr{P}(\Sigma^*)$ est l'ensemble des parties de $\Sigma^*$,
i.e., l'ensemble de tous les langages sur $\Sigma$,
cf. \ref{set-of-languages}) est dit \emph{stable par opérations
  rationnelles} lorsqu'il est stable par les opérations de réunion,
concaténation et étoile de Kleene, i.e., si $L_1,L_2 \in \mathscr{C}$
alors $L_1\cup L_2 \in \mathscr{C}$ et $L_1 L_2 \in \mathscr{C}$, et
si $L \in \mathscr{C}$ alors $L^* \in \mathscr{C}$ ; le \emph{plus
  petit}\footnote{La classe des langages rationnelles (qu'on cherche à
  définir) n'est pas le seul ensemble de langages stable par
  opérations rationnelles : l'ensemble $\mathscr{P}(\Sigma^*)$ de tous
  les langages est aussi évidemment stable par opérations
  rationnelles ; on s'intéresse au plus petit $\mathscr{C}$ possible
  pour n'avoir que ce qu'on peut construire à partir des langages de
  base triviaux par un nombre fini d'opérations rationnelles.}  (pour
l'inclusion) ensemble de langages stable par opérations rationnelles
et contenant les langages $\varnothing$, $\{\varepsilon\}$ et $\{x\}$
pour $x \in \Sigma$ (i.e. $\varnothing\in\mathscr{C}$,
$\{\varepsilon\} \in \mathscr{C}$ et si $x\in\Sigma$ alors
$\{x\}\in\mathscr{C}$), ou plus exactement, l'intersection de tous les
ensembles $\mathscr{C}$ vérifiant tous ces propriétés, est la classe
$\mathscr{R}$ des langages rationnels (et un langage rationnel est
simplement un élément de $\mathscr{R}$).

\medskip

\emph{Attention !}, le fait que la classe $\mathscr{R}$ des langages
rationnels soit stable par concaténation signifie que si $L_1$ et
$L_2$ sont rationnels alors le langage $L_1 L_2$ (constitué de tous
les mots concaténés d'un mot de $L_1$ et d'un mot de $L_2$) est
rationnel ; \emph{cela ne signifie pas} qu'un langage rationnel donné
soit stable par concaténation (un langage stable $L$ par concaténation
est un langage tel que si $w_1,w_2\in L$ alors $w_1 w_2 \in L$).

\thingy\label{regular-expressions} Pour décrire la manière dont un
langage rationnel est fabriqué (à partir des langages de base par les
opérations rationnelles), comme il est malcommode d'écrire quelque
chose comme $\{d\}(\{c\}^*)$, on introduit un nouvel objet, les
\index{expression rationnelle|see{rationnelle
    (expression)}}\defin[rationnelle (expression)]{expressions
  rationnelles} (certains préfèrent le terme d'\index{régulière
  (expression)|see{rationnelle}}\textbf{expressions régulières}), qui
sont des expressions servant à dénoter un langage rationnel.  Par
exemple, plutôt que d'écrire « $\{d\}(\{c\}^*)$ », on parlera du
langage « dénoté par l'expression rationnelle $dc{*}$ ».  Ceci fournit
du même coup une nouvelle définition des langages rationnels : ce sont
les langages dénotés par une expression rationnelle.

\smallskip

Plus exactement, une expression rationnelle (sur un alphabet $\Sigma$)
est un mot sur l'alphabet $\Sigma \cup \{\bot,
\underline{\varepsilon}, {(}, {)}, {|}, {*}\}$, où $\bot,
\underline{\varepsilon}, {(}, {)}, {|}, {*}$ sont de nouveaux
caractères \emph{n'appartenant pas} à l'alphabet $\Sigma$, appelés
\defin[métacaractère]{métacaractères}, et qui servent à marquer la
manière dont est formée l'expression rationnelle.  On définit
simultanément la notion d'expression rationnelle $r$ et de
\defin[dénoté (langage)]{langage dénoté} (ou \textbf{désigné} ou
simplement \textbf{défini}) par l'expression $r$, noté $L(r)$
(ou $L_r$), de la manière suivante\footnote{Si on veut être tout à
  fait rigoureux, il faudrait démontrer, pour que cette définition ait
  un sens, qu'une expression rationnelle $r$ ne désigne qu'un seul
  langage $L(r)$, c'est-à-dire, qu'il n'y a qu'une seule façon de la
  lire : pour cela, il faudrait faire appel aux techniques qui seront
  introduites en §\ref{section-context-free-grammars} et dire que la
  grammaire que nous définissons ici pour les expressions rationnelles
  est \emph{inambiguë} (cf. \ref{ambiguous-grammar}).} :
\begin{itemize}
\item $\bot$ est une expression rationnelle et son langage dénoté
  est $L(\bot) := \varnothing$,
\item $\underline{\varepsilon}$ est une expression rationnelle et son
  langage dénoté est $L(\underline{\varepsilon}) :=
  \{\varepsilon\}$,
\item si $x\in\Sigma$ est une lettre de l'alphabet $\Sigma$, alors le
  mot $x$ est une expression rationnelle et son langage dénoté
  est $L(x) := \{x\}$,
\item si $r_1,r_2$ sont deux expressions rationnelles et $L_1 =
  L(r_1)$ et $L_2 = L(r_2)$ les langages dénotés correspondants,
  alors $r_1 r_2$ est une expression rationnelle et son langage
  dénoté est $L(r_1 r_2) := L_1 L_2$,
\item si $r_1,r_2$ sont deux expressions rationnelles et $L_1 =
  L(r_1)$ et $L_2 = L(r_2)$ les langages dénotés correspondants,
  alors $(r_1|r_2)$ est une expression rationnelle\footnote{Certains
    préfèrent la notation $r_1+r_2$ ici,
    cf. \ref{regexp-notation-variations} pour une discussion.} et son
  langage dénoté est $L((r_1|r_2)) := L_1\cup L_2$,
\item si $r$ est une expression rationnelle et $L = L(r)$ les langage
  dénoté correspondant, alors $(r){*}$ est une expression
  rationnelle\footnote{Notons que certains auteurs préfèrent mettre le
    $*$ en exposant ici, c'est-à-dire noter $(r)^*$ (comme nous
    l'avons fait pour l'étoile de Kleene $L^*$ d'un langage $L$) ;
    comme nous avons choisi de définir une expression rationnelle
    comme un mot sur l'alphabet $\Sigma \cup \{\bot,
    \underline{\varepsilon}, {(}, {)}, {|}, {*}\}$, il semble plus
    raisonnable de ne pas rentrer dans la considération byzantine de
    savoir si un mot peut contenir des symboles en exposant.  En tout
    état de cause, il n'apparaîtra jamais aucune circonstance où les
    deux notations $X^*$ et $X{*}$ auraient tous les deux un sens et
    que ces sens diffèrent !  On peut donc ignorer purement et
    simplement la question de savoir si oui ou non l'étoile est en
    exposant.} et son langage dénoté est $L((r){*}) := L^*$.
\end{itemize}

Un langage rationnel est par construction la même chose qu'un langage
pour lequel il existe une expression rationnelle qui le dénote.

\thingy À titre d'exemple, sur l'alphabet $\Sigma = \{a,b,c,d\}$, $c$ est une
expression rationnelle qui dénote le langage $\{c\}$, donc $(c){*}$ en
est une qui dénote le langage $\{c\}^* = \{\varepsilon, c, cc,
ccc,\ldots\}$, et enfin $d(c){*}$ en est une qui dénote le langage $\{d,
dc, dcc, \ldots\}$ des mots formés d'un $d$ et d'une succession
quelconques de $c$.  Voici quelques autres exemples, toujours sur
$\Sigma = \{a,b,c,d\}$ :
\begin{itemize}
\item l'expression rationnelle $(a|b)$ dénote le langage $\{a\} \cup
  \{b\} = \{a,b\}$ constitué des deux mots d'une seule lettre $a$ et
  $b$ ;
\item l'expression rationnelle $(a|b)c$ dénote le langage
  $\{a,b\}\{c\} = \{ac,bc\}$, de même que $(ac|bc)$ ;
\item l'expression rationnelle $(bc){*}$ dénote le langage $\{bc\}^* =
  \{\varepsilon, bc, bcbc, bcbcbc, \ldots\}$ ;
\item l'expression rationnelle $(a|(bc){*})$ dénote le langage $\{a\}
  \cup \{bc\}^* = \{a,\penalty0 \varepsilon,\penalty1000
  bc,\penalty1000 bcbc,\penalty1000 bcbcbc, \ldots\}$ ;
\item l'expression rationnelle $(a|(bc){*})d$ dénote le langage $\{ad, d,
  bcd, bcbcd, bcbcbcd, \ldots\}$ ;
\item l'expression rationnelle $\bot d$ dénote le langage
  vide $\varnothing$ (car il n'y a pas de mot dans le langage vide,
  donc pas non plus de mot dans sa concaténation avec le
  langage $\{d\}$) ;
\item l'expression rationnelle $\underline{\varepsilon} d$ dénote le
  langage $\{d\}$ ;
\item l'expression rationnelle $(\bot|c)$ dénote le langage $\{c\}$ ;
\item l'expression rationnelle $(\underline{\varepsilon}|c)$ dénote le
  langage $\{\varepsilon, c\}$.
\end{itemize}

\thingy On dira qu'un mot $w$ \defin[vérifier (une expression
  rationnelle)]{vérifie} une expression rationnelle $r$ lorsque ce mot
appartient au langage qu'elle dénote (i.e., $w \in L(r)$).  Par
exemple, $dccc$ vérifie l'expression rationnelle $d(c){*}$.

\thingy Deux expressions rationnelles $r_1,r_2$ sont dites
\defin[équivalentes (expressions rationnelles)]{équivalentes}
lorsqu'elles dénotent le même langage.  À titre d'exemple, sur
l'alphabet $\{a,b\}$, les deux expressions rationnelles $(ab){*}a$ et
$a(ba){*}$ sont équivalentes (toutes les deux dénotent le langage
$\{a, aba, ababa, \ldots\}$ constitué des mots commençant et finissant
par un $a$ et dans lesquels chaque paire de $a$ est séparée par un
unique $b$).

On verra plus loin (en \ref{equivalence-of-regexps-is-decidable})
qu'on dispose d'un algorithme permettant de décider si deux
expressions rationnelles sont équivalentes.

{\footnotesize\thingy\label{regexp-identities} Voici quelques exemples
  d'équivalences d'expressions rationnelles (valables en remplaçant
  les différents $r$ qui interviennent dedans par des expressions
  rationnelles quelconques) :

\begin{itemize}
\item identités « triviales » :\quad $(r|\bot)\equiv r$,\quad
  $(\bot|r)\equiv r$,\quad $r\bot\equiv \bot$,\quad $\bot r\equiv
  \bot$, \quad $r\underline{\varepsilon} \equiv r$,\quad
  $\underline{\varepsilon}r \equiv r$,\quad $(\bot){*}\equiv
  \underline{\varepsilon}$ ;
\item identité d'associativité :\quad $((r_1|r_2)|r_3)\equiv
  (r_1|(r_2|r_3))$ ;
\item identités de distributivité :\quad $r(r_1|r_2)\equiv
  (rr_1|rr_2)$,\quad $(r_1|r_2)r\equiv (r_1r|r_2r)$ ;
\item identité de commutativité :\quad $(r_1|r_2)\equiv (r_2|r_1)$ ;
\item identités « apériodiques » :\quad $((r_1|r_2)){*} \equiv
  (r_1){*}(r_2(r_1){*}){*}$,\quad $((r_1|r_2)){*} \equiv
  ((r_1){*}r_2){*}(r_1){*}$\quad et\quad $(r_1 r_2){*} \equiv
  (\underline{\varepsilon}|r_1(r_2 r_1){*}r_2)$ ;
\item identités « cycliques » :\quad $(r){*} \equiv
  (\underline{\varepsilon}|r)(rr){*} \equiv
  (\underline{\varepsilon}|r|rr)(rrr){*} \equiv
  (\underline{\varepsilon}|r|rr|rrr)(rrrr){*} \equiv \cdots$ ;
\item identités d'« idempotence » :\quad $(r|r)\equiv r$,\quad
  $((r){*}){*} \equiv (r){*}$ ;
\end{itemize}

Ces différentes identités présentent un intérêt théorique dont
l'explication dépasserait le cadre de ces notes ; cependant, il peut
être intéressant de réfléchir à ce que chacune signifie.

Signalons au passage que toutes ces identités \emph{ne suffisent pas}
à produire toutes les équivalences entre expressions rationnelles.
Par exemple, elles ne permettent pas de démontrer l'équivalence
suivante : $(a|b){*} \equiv
((a|b)(b|ab{*}ab{*}ab{*}a)){*}(\underline{\varepsilon}|(a|b)(|ab{*}|ab{*}ab{*}|ab{*}ab{*}ab{*}))$
(écrite avec les conventions de \ref{convention-on-writing-regexps}) ;
la question d'arriver à trouver un système d'axiomes qui permet de
déduire toutes les équivalences entre expressions rationnelles est un
problème délicat (il n'y en a notamment pas qui soit fini).

\par}

\thingy\label{convention-on-writing-regexps}
La convention de parenthésage introduite ci-dessus est
inambiguë mais parfois inutilement lourde : on se permettra parfois de
l'alléger, par exemple d'écrire $(r_1|r_2|r_3)$ pour $((r_1|r_2)|r_3)$
(ou pour $(r_1|(r_2|r_3))$, ce qui n'a guère d'importance vu qu'elles
dénotent le même langage), ou encore $x{*}$ pour $(x){*}$ lorsque $x$
est formé d'un seul caractère.

Pour retirer des parenthèses, la convention sur la priorité des
opérations est la suivante : l'opération d'étoile ${*}$ est la plus
prioritaire (c'est-à-dire que $ab{*}$ se lit comme $a(b){*}$ et non
pas comme $(ab){*}$), la concaténation est de priorité intermédiaire,
et la barre de disjonction $|$ est la moins prioritaire (c'est-à-dire
que $ab|cd$ se lit comme $(ab|cd)$ et pas comme $a(b|c)d$).

\thingy\label{regexp-notation-variations} La disjonction que nous
avons notée « ${|}$ » est plus souvent notée « $+$ » par les
mathématiciens\footnote{Dans le même contexte mathématique, il est
  alors fréquent de noter « $0$ » pour ce que nous avons
  noté « $\bot$ » (c'est un élément neutre pour la disjonction), et on
  en profite souvent pour noter « $1$ » pour « $\varepsilon$ » et/ou
  « $\underline{\varepsilon}$ » (c'est un élément neutre pour la
  concaténation).  Par exemple, beaucoup des identités énoncées
  en \ref{regexp-identities} sembleront plus naturelles si on les
  écrit avec ces notations.}.  Cette notation est plus élégante pour
faire le lien avec l'algèbre, et plus logique dans des cadres plus
généraux (expressions rationnelles et automates « avec
multiplicités ») ; nous avons cependant préféré l'éviter parce qu'elle
prête à confusion : en effet, on a introduit en \ref{kleene-plus} une
opération $L \mapsto L^+$ qui n'a rien à voir avec la disjonction, et
de nombreux moteurs d'expressions régulières en informatique (par
exemple, \texttt{egrep},
cf. §\ref{subsection-remarks-on-regexp-syntax}) utilisent
effectivement le symbole \texttt{+} dans ce sens, c'est-à-dire pour
désigner « au moins une répétition »
(cf. \ref{remarks-egrep-plus-etc}).  Même si nous avons cherché à ne
pas utiliser le symbole $+$ pour éviter l'ambiguïté, il faut savoir
que ce double sens existe (soit comme disjonction, soit comme
l'opération de \ref{kleene-plus}).

{\footnotesize Les métacaractères $\bot$ et $\underline{\varepsilon}$
  sont introduits ici par souci de complétude mais sont rarement
  utilisés dans les expressions rationnelles (le métacaractère
  $\underline{\varepsilon}$ a été souligné parce qu'il s'agit d'une
  vraie lettre et non pas du mot vide ; on peut ignorer cette
  subtilité qui n'a que très peu d'importance).\par}


\subsection{Remarques sur les moteurs d'expressions régulières en informatique}\label{subsection-remarks-on-regexp-syntax}

\thingy Dans le monde informatique, il existe de nombreux
\emph{moteurs d'expressions régulières}, c'est-à-dire outils (qu'il
s'agisse de primitives d'un langage, de bibliothèques externes, de
programmes en ligne de commande, ou autres) permettant de savoir si un
mot est reconnu par une expression régulière (=rationnelle), autrement
dit, s'il appartient au langage dénoté par elle.  L'un de ces moteurs
est le programme \texttt{egrep} standard sous Unix/POSIX ; un autre
est le moteur \texttt{java.util.regex} de Java.

\thingy Les expressions régulières au sens de ces différents moteurs
sont généralement plus puissantes que les expressions rationnelles au
sens mathématique défini en \ref{regular-expressions} ci-dessus :
différentes extensions permettent de désigner des langages qui ne sont
pas rationnels au sens mathématique.

L'extension la plus fréquente est celle des \emph{références arrière}
(ou \emph{backreferences} en anglais) qui permettent de demander qu'un
facteur du mot se retrouve aussi à un autre emplacement.  Par exemple,
pour beaucoup de moteurs (notamment \texttt{egrep}), l'expression
régulière « \texttt{(a*)b\char"5C\relax 1} » désigne le langage
$\{a^nba^n : n\in\mathbb{N}\} = \{b,aba, aabaa,\ldots\}$ des mots
formés d'un nombre quelconque de $a$ puis d'un $b$ puis de la
\emph{même suite de $a$} (le « \texttt{\char"5C\relax 1} » désigne
« une copie de la chaîne de caractères qui a été capturée par le
premier jeu de parenthèses ») ; or ce langage \emph{n'est pas
  rationnel} au sens mathématique (ce sera une conséquence du « lemme
de pompage » \ref{pumping-lemma}), c'est donc bien le signe qu'on sort
du cadre théorique décrit plus haut.

\thingy\label{remarks-egrep-plus-etc} Il existe aussi un certain
nombre de constructions qui, sans dépasser la puissance des
expressions rationnelles au sens mathématique, apportent des
commodités d'écriture dans un contexte informatique.

Parmi les caractères les plus courants de cette sorte, mentionnons le
\texttt{+} qui signifie « au moins une répétition » (alors que
\texttt{*} signifie « au moins zéro répétitions ») : ainsi,
« $r\texttt{+}$ » est équivalent à $rr{*}$ ; on trouve aussi
\texttt{?} pour « au plus une répétition » : c'est-à-dire que
« $r\texttt{?}$ » est équivalent à $(\underline{\varepsilon}|r)$ (qui
sera d'ailleurs simplement noté $(|r)$ car il n'y a jamais de symbole
spécial pour $\underline{\varepsilon}$).

Parmi d'autres constructions qui ne sortent pas du cadre théorique des
langages rationnels, mentionnons encore le point « \texttt{.} » qui
désigne un caractère quelconque de l'alphabet (on peut le voir comme
une abréviation pour $(x_1|x_2|\ldots|x_N)$ où $x_1,x_2,\ldots,x_N$
sont tous les éléments de $\Sigma$ — ce qui serait très fastidieux à
écrire si on devait énumérer tout Unicode).  Les crochets tels que
« \texttt{[$xyz$]} » désignent un caractère quelconque parmi ceux
listés (c'est donc la même chose que $(x|y|z)$ mais cela ne fonctionne
qu'avec des caractères individuels ; en contrepartie, on peut écrire
des intervalles comme « \texttt{[a-z]} » qui désigne un caractère
quelconque entre \texttt{a} et \texttt{z} dans l'ordre
ASCII/Unicode\footnote{...ou parfois l'ordre de tri défini par la
  « locale » en cours : c'est malheureusement une source supplémentaire de
  confusion et de bugs.}, ou bien des négations, comme
« \texttt{[\char"5Ea-z]} » qui désigne un caractère qui \emph{n'est
  pas} entre \texttt{a} et \texttt{z}) ou
« \texttt{[\char"5Eaeio-z]} » qui désigne un caractère qui \emph{n'est
  ni} \texttt{a} ni \texttt{e} ni\texttt{i} ni entre \texttt{o}
et \texttt{z}).

Toutes sortes d'autres raccourcis ou commodités de notation peuvent
exister, par exemple \texttt{\char"5C<} et \texttt{\char"5C>} pour
désigner un début et une fin de mot (la définition précise de « mot »
pouvant varier), ou encore \texttt{$r$\{$n_1$,$n_2$\}} qui cherche
entre $n_1$ et $n_2$ répétitions de $r$ : ainsi,
« $r\texttt{\{2,5\}}$ » est équivalent à $(rr|rrr|rrrr|rrrrr)$.

\thingy Une autre subtilité est que la plupart des moteurs
d'expressions régulières en informatique vont, par défaut,
\emph{rechercher un facteur} (appelé « sous-chaîne » en informatique)
vérifiant l'expression \emph{à l'intérieur} de la chaîne donnée,
plutôt que tester si la chaîne tout entière vérifie l'expression.
Pour éviter ce comportement, on peut utiliser des \index{ancre
  (métacaractère)}\emph{ancres}, typiquement commandées par les
(méta)caractères « \texttt{\char"5E} » et « \texttt{\char"24} » qui
servent à ancrer l'expression au début et à la fin de la chaîne
respectivement : c'est-à-dire que rechercher « \texttt{a} » recherche
un \texttt{a} quelconque à l'intérieur de la chaîne donnée, rechercher
« \texttt{\char"5E\relax a} » demande que le \texttt{a} soit au début
de la chaîne donnée, rechercher « \texttt{a\char"24} » demande que le
\texttt{a} soit à la fin de la chaîne donnée, et rechercher
« \texttt{\char"5E\relax a\char"24} » demande que la chaîne donnée
soit exactement \texttt{a} (cet exemple n'est donc pas très utile,
mais de façon générale, trouver si une chaîne vérifie une expression
rationnelle $r$, revient à y chercher « \texttt{\char"5E\relax
  $r$\char"24} »).

\thingy Comme les expressions régulières en informatique sont
représentées par des chaînes de caractères qui appartiennent au même
alphabet (ASCII ou Unicode) que les chaînes sur lesquelles on effectue
la recherche, le problème se pose de distinguer les métacaractères
(l'étoile de Kleene \texttt{*}, par exemple) des caractères eux-mêmes
(comment rechercher les chaînes contenant le caractère « \texttt{*} »
si \texttt{*} est utilisé par l'étoile de Kleene ?).  La solution est
d'introduire un mécanisme d'\emph{échappement}, normalement par le
caractère \emph{backslash} : ainsi, « \texttt{x\char"5C*} » recherche
un \texttt{x} suivi d'un astérisque \texttt{*}, tandis que
« \texttt{x*} » recherche un nombre quelconque de répétitions de la
lettre \texttt{x}.

\thingy Il existe malheureusement de nombreuses différences, parfois
très subtiles, entre moteurs, ne serait-ce que dans les notations : un
moteur pourra par exemple noter « \texttt{(?)} » ce qu'un autre note
« \texttt{\char"5C(\char"5C?\char"5C)} » et vice versa.  La seule
solution est de consulter attentivement la documentation de chaque
moteur d'expressions régulières pour connaître la syntaxe utilisée.

Signalons tout de même qu'il existe deux principales familles de
syntaxes d'expressions régulières en informatique : les expressions
régulières « POSIX étendues », utilisée notamment par le programme
Unix \texttt{egrep}, et les expressions régulières de Perl, qui ont
été réadaptées dans beaucoup de langages, notamment Java, JavaScript,
Python et d'autres.

\thingy Signalons comme complication supplémentaire que dans de
nombreux langages, les expressions régulières sont saisies comme des
chaînes de caractères plutôt que d'avoir une syntaxe spéciale, et ceci
a pour effet d'introduire un niveau supplémentaire d'échappement : par
exemple, en Java, pour rechercher si une chaîne de caractères $s$
contient un astérisque, on utilisera
\texttt{$s$.matches("\char"5C\char"5C*")} puisque l'expression
régulière à utiliser est \texttt{\char"5C*} et que cette chaîne de
caractères s'écrit \texttt{"\char"5C\char"5C*"} en Java.

De même, dans le shell Unix, pour rechercher une ligne d'un fichier
contenant un \texttt{x} suivi d'un nombre quelconque de \texttt{y}
suivi d'un \texttt{z}, on utilisera \texttt{egrep \char"27xy*z\char"27
  \ monfichier}, les guillemets (simples) ayant pour fonction d'éviter
que le shell Unix interprète lui-même l'astérisque comme un caractère
spécial (dont le rôle est différent de l'étoile de Kleene, quoique
suffisamment proche pour pouvoir causer une confusion).


\section{Automates finis}\label{section-finite-automata}

\setcounter{comcnt}{0}

\thingy Les automates finis sont un modèle de calcul particulièrement
simple et particulièrement appropriés à l'étude et à l'analyse des
langages rationnels.  Il faut imaginer un automate fini comme une
machine disposant d'une quantité finie (et sous-entendu, très limitée)
de mémoire : la configuration complète de cette mémoire est décrite
par un \emph{état}, qui appartient à un ensemble fini d'états
possibles.  On fournit un mot à l'automate en le lui faisant consommer
lettre par lettre (de la première à la dernière), et à chaque lettre
qu'il reçoit, l'automate prend une décision (à savoir, passer dans un
nouvel état) en fonction de son état actuel et de la lettre qu'on lui
donne à consommer ; l'automate commence par un état dit
\emph{initial} ; lorsque le mot a été entièrement consommé, l'état
dans lequel l'automate se trouve a pour conséquence soit que
l'automate \emph{accepte} le mot (s'il se trouve dans un état
\emph{final}), soit qu'il le \emph{rejette}.  L'ensemble des mots
acceptés par l'automate est un langage dit \emph{reconnaissable}.  On
va voir en \ref{kleenes-theorem} que ces langages reconnaissables sont
en fait les mêmes que les langages rationnels définis
en \ref{rational-languages}.

\thingy Une subtilité est qu'il existe plusieurs sortes d'automates
finis.  Nous allons en définir quatre ou cinq, de plus en plus
généraux, mais nous allons voir qu'ils sont, finalement, tous
équivalents, c'est-à-dire que leur puissance de calcul, ou les
langages qu'ils sont capables de reconnaître, sont toujours les mêmes
(en revanche, cette équivalence peut avoir un coût algorithmique, car
la conversion d'une sorte d'automate en une autre n'est pas gratuite).
Du plus particulier au plus général, les automates que nous allons
définir sont informellement décrits ainsi :
\begin{itemize}
\item les automates finis déterministes (complets)
  (cf. \ref{definition-dfa}) : ceux-ci partent d'un état initial bien
  défini, suivent à chaque lettre consommée une transition bien
  définie vers un nouvel état, et leur comportement est donc
  entièrement déterministe (d'où le nom) ;
\item les automates finis déterministes à spécification incomplète
  (cf. \ref{definition-dfai}) : la différence est qu'ici l'automate
  n'a pas forcément d'instruction sur quoi faire quand il est dans un
  certain état et reçoit un certain symbole, il peut manquer des
  transitions, auquel cas l'automate cesse de fonctionner et rejette
  le mot ;
\item les automates non-déterministes (cf. \ref{definition-nfa}) :
  ceux-ci peuvent admettre \emph{plusieurs} possibilités de nouvel
  état dans lequel transitionner lorsqu'on leur donne une lettre
  donnée à consommer depuis un état donné, et l'automate accepte le
  mot s'il y a une manière de faire les transitions qui résulte en ce
  que le mot soit accepté ;
\item les automates non-déterministes à transitions spontanées
  (cf. \ref{definition-enfa}) : ceux-ci ont, en plus des précédents,
  le pouvoir d'évoluer spontanément, sans consommer de lettre ;
\item on introduira aussi des automates encore plus généraux, à
  transitions étiquetées par des expressions rationnelles
  (cf. \ref{definition-rnfa}) pour montrer que tout langage
  reconnaissable est rationnel.
\end{itemize}

\thingy Dans tous les cas, les automates se représenteront comme des
graphes orientés (admettant des boucles et arêtes multiples ;
cf. \ref{discussion-example1} ci-dessous pour un exemple simple), dont
les sommets sont les états de l'automate (on leur donne généralement
des noms pour les reconnaître, mais ces noms sont arbitraires), dont
les arêtes sont étiquetées par des symboles (des lettres de
l'alphabet, ou éventuellement le mot vide $\varepsilon$ ou une
expression rationnelle sur l'alphabet), et dont certains sommets
(=états) sont distingués comme « initiaux » et/ou « finaux » (par une
flèche pointant depuis nulle part vers l'état en question ou depuis
l'état en question vers nulle part).

Un mot sera accepté par l'automate toujours selon la même logique :
s'il y a moyen de relier un état initial à un état final par un chemin
orienté dans le graphe tel que le mot corresponde à la lecture des
étiquettes des arêtes du chemin dans l'ordre.

\subsection{Automates finis déterministes complets (=DFA)}

\thingy\label{definition-dfa} Un \defin{automate fini déterministe}
(complet), ou en abrégé \index{DFA|see{automate fini
    déterministe}}\textbf{DFA} (pour \textit{deterministic finite
  automaton}), sur un alphabet $\Sigma$ est la donnée
\begin{itemize}
\item d'un ensemble fini $Q$ dont les éléments sont appelés
  \defin[état]{états},
\item d'un état $q_0 \in Q$ appelé \defin[initial (état)]{état
  initial} (un DFA a un \emph{unique} état initial),
\item d'un ensemble $F \subseteq Q$ d'états appelés états \defin[final
  (état)]{finaux}\footnote{Le pluriel de « final » est indifféremment
  « finaux » ou « finals ».} ou \index{acceptant
  (état)|see{final}}\textbf{acceptants},
\item d'une fonction $\delta \colon Q\times\Sigma \to Q$ appelée
  \defin[transition (fonction de)]{fonction de transition}.
\end{itemize}

\thingy\label{graphical-representation-of-dfa} Graphiquement, on
représente un DFA comme un graphe orienté aux arêtes étiquetées par
des éléments de $\Sigma$ : plus exactement, on trace un nœud pour
chaque élément $q \in Q$, et lorsque $\delta(q,x) = q'$ on introduit
une flèche $q \to q'$ étiquetée par la lettre $x$.

La condition cruciale (pour être un DFA) est que, pour chaque état $q
\in Q$ et chaque lettre $x \in \Sigma$, \emph{il existe une unique}
arête partant de $q$ et étiquetée par $x$ (c'est essentiellement une
reformulation du fait que $\delta$ est une fonction).

En outre, on introduit une flèche pointant de nulle part vers $q_0$
(pour marquer celui-ci comme l'état initial), et pour chaque $q\in F$
une flèche pointant de $q$ vers nulle part\footnote{Certains auteurs
  préfèrent d'autres conventions, par exemple celle consistant à
  entourer deux fois les états finaux.} (pour marquer ceux-ci comme
états finaux).

Pour abréger la représentation graphique, lorsque plusieurs arêtes
étiquetées par des lettres $x,y$ différentes relient les mêmes sommets
$q,q'$ (i.e., lorsqu'on a à la fois $\delta(q,x) = q'$ et $\delta(q,y)
= q'$), on pourra écrire « $x,y$ », voire « $x|y$ », sur l'arête en
question et ne la tracer qu'une seule fois.
(Voir \ref{discussion-example2} ci-dessous pour un exemple.)

\thingy\label{discussion-example1} Pour donner un exemple simple,
l'automate sur $\Sigma = \{a,b\}$ représenté ci-dessous a $Q =
\{0,1\}$ et $q_0 = 0$ et $F = \{1\}$ et la fonction de transition
$\delta$ donnée par $\delta(0,a) = 0$ et $\delta(0,b) = 1$ et
$\delta(1,a) = 1$ et $\delta(1,b) = 0$.  On pourra se convaincre (une
fois lues les définitions plus loin) que cet automate accepte les mots
dont le nombre de $b$ est impair.

\begin{center}
%%% begin example1 %%%

\begin{tikzpicture}[>=latex,line join=bevel,automaton]
%%
\node (q1) at (98bp,20.306bp) [draw,circle,state,final] {$1$};
  \node (q0) at (18bp,20.306bp) [draw,circle,state,initial] {$0$};
  \draw [->] (q1) ..controls (75.212bp,3.6347bp) and (64.284bp,-1.3057bp)  .. (54bp,1.3057bp) .. controls (50.042bp,2.3107bp) and (46.047bp,3.8633bp)  .. node[auto] {$b$} (q0);
  \draw [->] (q0) ..controls (46.106bp,20.306bp) and (58.578bp,20.306bp)  .. node[auto] {$b$} (q1);
  \draw [->] (q1) to[loop above] node[auto] {$a$} (q1);
  \draw [->] (q0) to[loop above] node[auto] {$a$} (q0);
%
\end{tikzpicture}

%%% end example1 %%%
\end{center}

\thingy Il faut comprendre le fonctionnement d'un DFA de la manière
suivante : initialement, l'automate est dans l'état initial $q_0$.  On
va lui présenter un mot $w \in \Sigma^*$, lettre par lettre, de la
gauche vers la droite : i.e., si $w = x_1\cdots x_n$ on va faire
consommer à l'automate les lettres $x_1,x_2,\ldots,x_n$ dans cet
ordre.  Le fait de consommer une lettre $x$ fait passer l'automate de
l'état $q$ à l'état $\delta(q,x)$ ; autrement dit, l'automate passe
successivement dans les états $q_0$ puis $q_1 := \delta(q_0,x_1)$ puis
$q_2 := \delta(q_1,x_2)$, et ainsi de suite jusqu'à $q_n :=
\delta(q_{n-1},x_n)$ : on dit que l'automate effectue les transitions
$q_0\to q_1$ (en consommant $x_1$) puis $q_1\to q_2$ (en
consommant $x_2$) et ainsi de suite.  Si $q_n$ est l'état dans lequel
se trouve l'automate une fois qu'il a consommé le mot $w$, on dira que
l'automate \emph{accepte} ou \emph{rejette} le mot selon que $q_n \in
F$ ou que $q_n \not\in F$.

\smallskip

Graphiquement, on peut présenter la procédure de la manière suivante :
on part de l'état $q_0$ (sommet du graphe représentant l'automate)
indiqué par la flèche entrante (pointant de nulle part), et pour
chaque lettre du mot $w = x_1\cdots x_n$ considéré, on suit l'arête
portant cette lettre pour étiquette (et partant de l'état où on se trouve
actuellement).  Si à la fin l'état $q_n$ est acceptant (représenté par
une flèche pointant vers nulle part), le mot $w$ est accepté, sinon il
est rejeté.

Cela revient encore à dire que le mot $w$ est accepté lorsqu'il existe
un chemin orienté dans l'automate, reliant l'état $q_0$ initial à un
état $q_n$ final, et tel que le mot $w = x_1 \cdots x_n$ soit obtenu
en lisant dans l'ordre les étiquettes $x_i$ des différentes arêtes
$q_{i-1} \to q_i$ de ce chemin.  Cette définition pourra resservir
presque à l'identique pour les autres sortes d'automates, la
spécificité des DFA est que l'état initial $q_0$ est unique et que
connaissant l'étiquette $x_i$ il n'y a qu'une seule transition
$q_{i-1} \to q_i$ possible (il n'y a pas de choix à faire, pas de
non-déterminisme).

\thingy\label{definition-multiple-transition-function} Pour donner de
façon plus formelle la définition du langage accepté par un automate,
il sera utile d'introduire une fonction $\delta^*$ de transition
multiple (une généralisation de $\delta$), qui définit l'état
$\delta^*(q,w)$ dans lequel aboutit un DFA si à partir de l'état $q$
on lui fait consommer successivement les lettres d'un mot $w$.

Formellement : si $A = (Q,q_0,F,\delta)$ est un DFA sur
l'alphabet $\Sigma$, on définit une fonction $\delta^* \colon
Q\times\Sigma^* \to Q$ par $\delta^*(q,x_1\cdots x_n) =
\delta(\cdots\delta(\delta(q,x_1),x_2)\cdots,x_n)$ ou, ce qui revient
au même (par récurrence sur la longueur du second argument) :
\begin{itemize}
\item $\delta^*(q,\varepsilon) = q$ quel que soit $q\in Q$ (où
  $\varepsilon$ désigne le mot vide),
\item $\delta^*(q,wx) = \delta(\delta^*(q,w),x)$ quels que soient
  $q\in Q$, $w\in\Sigma^*$ et $x\in\Sigma$,
\end{itemize}
(en particulier, $\delta^*(q,x) = \delta(q,x)$ si $x\in\Sigma$, donc
avec la convention faite en \ref{convention-on-words-of-length-one},
on peut dire que $\delta^*$ prolonge $\delta$ ; il sera par ailleurs
utile de remarquer que $\delta^*(q,ww') =
\delta^*(\delta^*(q,w),w')$, ce qui se démontre par récurrence).

Cette fonction $\delta^*$ étant définie, on dira que l'automate $A$
\defin[accepter]{accepte} ou
\index{reconnaître|see{accepter}}\textbf{reconnaît} un mot $w$ lorsque
$\delta^*(q_0,w) \in F$ ; dans le cas contraire, on dira qu'il
\textbf{rejette} le mot $w$.

\thingy\label{definition-recognizable-language} L'ensemble $L(A)$
(ou $L_A$) des mots acceptés par l'automate $A$ s'appelle
\textbf{langage accepté}, ou \textbf{reconnu}, ou \textbf{défini}, par
l'automate $A$.

\smallskip

Un langage $L \subseteq \Sigma^*$ qui peut s'écrire sous la forme du
langage $L(A)$ accepté par un DFA $A$ s'appelle \defin[reconnaissable
  (langage)]{reconnaissable} (sous-entendu : par automate déterministe
fini).

\smallskip

On dit que deux DFA $A,A'$ sont \defin[équivalents
  (automates)]{équivalents} lorsqu'ils reconnaissent le même langage,
i.e., $L(A) = L(A')$.

\thingy\label{discussion-example2} Donnons un exemple : l'automate
fini ci-dessous sur $\Sigma := \{a,b,c\}$ a trois états, $Q =
\{0,1,2\}$.  On peut en faire la description informelle suivante :
l'automate commence dans l'état $0$, où il reste jusqu'à rencontrer
un $a$ qui le fait passer dans l'état $1$, où il reste ensuite jusqu'à
rencontrer un $b$ qui le fait passer dans l'état $2$, où il reste
définitivement et qui constitue le seul état acceptant.

\begin{center}
%%% begin example2 %%%

\begin{tikzpicture}[>=latex,line join=bevel,automaton]
%%
\node (q1) at (98bp,18bp) [draw,circle,state] {$1$};
  \node (q0) at (18bp,18bp) [draw,circle,state,initial] {$0$};
  \node (q2) at (178bp,18bp) [draw,circle,state,final] {$2$};
  \draw [->] (q0) ..controls (46.106bp,18bp) and (58.578bp,18bp)  .. node[auto] {$a$} (q1);
  \draw [->] (q1) to[loop above] node[auto] {$a,c$} (q1);
  \draw [->] (q1) ..controls (126.11bp,18bp) and (138.58bp,18bp)  .. node[auto] {$b$} (q2);
  \draw [->] (q0) to[loop above] node[auto] {$b,c$} (q0);
  \draw [->] (q2) to[loop above] node[auto] {$a,b,c$} (q2);
%
\end{tikzpicture}

%%% end example2 %%%
\end{center}

Cette description rend clair le fait que l'automate en question
accepte exactement les mots contenant un $a$ suivi, pas forcément
immédiatement, d'un $b$ ; autrement dit, les mots dont $ab$ est un
sous-mot (cf. \ref{definition-subword}).  Ce langage est donc
reconnaissable.  (Il est aussi rationnel puisque dénoté par
l'expression rationnelle $(b|c){*}a(a|c){*}b(a|b|c){*}$.)

\thingy\label{definition-dfa-accessible-state} Un état $q$ d'un DFA
est dit \defin[accessible (état)]{accessible} lorsqu'il existe un mot
$w \in \Sigma^*$ tel que $q = \delta(q_0,w)$, autrement dit,
graphiquement, lorsqu'il existe un chemin orienté
$q_0,q_1,\ldots,q_n=q$ reliant l'état initial $q_0$ à l'état $q$
considéré : bref, cela correspond à un état auquel il est possible que
l'automate arrive (en partant de l'état initial et en consommant un
certain mot).  Dans le cas contraire, l'état est dit
\textbf{inaccessible}.  Il est évident qu'ajouter ou supprimer (ou
modifier) les états inaccessibles dans un DFA ne change rien au
langage reconnu (on obtient des automates équivalents).

\smallskip

Par exemple, dans le DFA qui suit, l'état $2$ est inaccessible
(l'automate est donc équivalent à celui représenté
en \ref{discussion-example1}).  On remarquera qu'il ne change rien que
cet état soit final ou non.

\begin{center}
%%% begin example1b %%%

\begin{tikzpicture}[>=latex,line join=bevel,automaton]
%%
\node (q1) at (98bp,20.306bp) [draw,circle,state,final] {$1$};
  \node (q0) at (18bp,20.306bp) [draw,circle,state,initial] {$0$};
  \node (q2) at (172bp,20.306bp) [draw,circle,state,final] {$2$};
  \draw [->] (q1) ..controls (75.212bp,3.6347bp) and (64.284bp,-1.3057bp)  .. (54bp,1.3057bp) .. controls (50.042bp,2.3107bp) and (46.047bp,3.8633bp)  .. node[auto] {$b$} (q0);
  \draw [->] (q0) ..controls (46.106bp,20.306bp) and (58.578bp,20.306bp)  .. node[auto] {$b$} (q1);
  \draw [->] (q1) to[loop above] node[auto] {$a$} (q1);
  \draw [->] (q0) to[loop above] node[auto] {$a$} (q0);
%
\end{tikzpicture}

%%% end example1b %%%
\end{center}

\thingy On va maintenant introduire différentes variations sur le
thème des automates finis, c'est-à-dire différentes généralisations de
la définition faite en \ref{definition-dfa} correspondant à des types
d'automates finis plus puissants que les DFA mais dont on va montrer,
à chaque fois, qu'ils peuvent se ramener à des DFA au sens où pour
chacun de ces automates généralisés on pourra construire
algorithmiquement un DFA qui reconnaît le même langage (si bien que la
classe des langages reconnaissables par n'importe laquelle de ces
sortes d'automates sera toujours la même).  Les plus simples sont les
automates déterministes finis incomplets et on va donc commencer par
eux.


\subsection{Automates finis déterministes à spécification incomplète (=DFAi)}

\thingy\label{definition-dfai} Un \textbf{automate fini déterministe à
  spécification incomplète} ou \textbf{...partielle}, ou simplement
\defin{automate fini déterministe
  incomplet}\footnote{\label{footnote-on-word-incomplete}Le mot
  « incomplet » signifie en fait « non nécessairement complet », i.e.,
  l'automate a le droit de manquer certaines transitions, mais il peut
  très bien être complet (un DFA est un DFAi particulier, pas le
  contraire).}, en abrégé \index{DFAi|see{automate fini déterministe
    incomplet}}\textbf{DFAi}, sur un alphabet $\Sigma$ est la donnée
\begin{itemize}
\item d'un ensemble fini $Q$ d'états,
\item d'un état initial $q_0 \in Q$,
\item d'un ensemble $F \subseteq Q$ d'états finaux,
\item d'une fonction de transition \emph{partielle}\footnote{Une
  « fonction partielle » $f\colon X\dasharrow Y$, où $X, Y$ sont deux
  ensembles est, par définition, la même chose qu'une fonction
  $f\colon D\to Y$ où $D\subseteq X$ est un sous-ensemble de $X$
  appelé \textbf{ensemble de définition} de $f$.  (Lorsque en fait
  $D=X$, la fonction est dite « totale ».)} $\delta \colon
  Q\times\Sigma \dasharrow Q$,
\end{itemize}
autrement dit, la seule différence avec la définition faite
en \ref{definition-dfa} est que la fonction $\delta$ est partielle, ce
qui signifie qu'elle n'est pas obligatoirement définie sur tout couple
$(q,x) \in Q\times\Sigma$.

\smallskip

Un DFA est considéré comme un DFAi particulier où la fonction de
transition $\delta$ se trouve être définie partout.

\thingy Graphiquement, on représente un DFAi comme un DFA, à la
différence près que pour chaque $q\in Q$ et chaque $x\in \Sigma$, il y
a maintenant \emph{au plus une} (et non plus exactement une) arête
partant de $q$ et étiquetée par $x$.

\smallskip

L'intérêt informatique des DFAi est de ne pas s'obliger à stocker
inutilement des transitions et des états inutiles au sens où ils ne
permettront jamais d'accepter le mot (voir la notion d'automate
« émondé » en \ref{coaccessible-and-useful-states} ci-dessous).  C'est
la raison pour laquelle, même si les DFA complets sont théoriquement
plus satisfaisants à manier (pour la même raison qu'une fonction
totale est plus satisfaisante qu'une fonction partielle), il est
souvent algorithmiquement plus judicieux de travailler sur des DFAi.

\thingy Le fonctionnement d'un DFAi est le même que celui d'un DFA, à
la modification suivante près : si on donne à consommer à l'automate
une lettre pour laquelle la transition n'est pas définie, i.e., s'il
rencontre un $x$ pendant qu'il se trouve dans un état $q$ pour lequel
$\delta(q,x)$ n'est pas défini, alors l'automate cesse de
fonctionner : l'automate n'a plus d'état, n'effectue plus de
transition, et n'acceptera pas le mot quelles que soient les lettres
ultérieures.

\smallskip

Cela revient une fois de plus à dire que le mot $w$ est accepté
lorsqu'il existe un chemin orienté dans l'automate, reliant l'état
$q_0$ initial à un état $q_n$ final, et tel que le mot $w = x_1 \cdots
x_n$ soit obtenu en lisant dans l'ordre les étiquettes $x_i$ des
différentes arêtes $q_{i-1} \to q_i$ de ce chemin.  La différence des
DFAi par rapport aux DFA est que, cette fois, la tentative de chemin
pourrait s'arrêter prématurément (on ne peut plus consommer un symbole
et passer dans un nouvel état, et à plus forte raison, aboutir à un
état final).

\thingy Formellement : si $A = (Q,q_0,F,\delta)$ est un DFAi sur
l'alphabet $\Sigma$, on définit une fonction $\delta^* \colon
Q\times\Sigma^* \dasharrow Q$ par $\delta^*(q,x_1\cdots x_n) =
\delta(\cdots\delta(\delta(q,x_1),x_2)\cdots,x_n)$ avec la convention
que dès qu'une sous-expression n'est pas définie, toute l'expression
n'est pas définie, ou, ce qui revient au même (par récurrence sur la
longueur du second argument) :
\begin{itemize}
\item $\delta^*(q,\varepsilon) = q$ quel que soit $q\in Q$ (où
  $\varepsilon$ désigne le mot vide),
\item $\delta^*(q,wx) = \delta(\delta^*(q,w),x)$ à condition que $q'
  := \delta^*(q,w)$ soit défini et que $\delta(q',x)$ le soit (et si
  ces deux conditions ne sont pas satisfaites, $\delta^*(q,wx)$ n'est
  pas défini).
\end{itemize}

Enfin, l'automate $A$ accepte un mot $w$ lorsque $\delta^*(q_0,w)$
\emph{est défini} et appartient à $F$ ; dans le cas contraire (que ce
soit parce que $\delta^*(q_0,w)$ n'est pas défini ou parce qu'étant
défini il n'appartient pas à $F$), l'automate rejette le mot.

\smallskip

Le langage accepté $L(A)$ et l'équivalence de deux automates sont
définis de façon analogue aux DFA
(cf. \ref{definition-recognizable-language}).

\thingy\label{discussion-example2b} Voici un exemple de DFAi sur
l'alphabet $\Sigma = \{a,b,c\}$.  Cet automate reconnaît exactement
les mots formés d'un nombre quelconque de $c$, suivis d'un $a$, suivis
d'un nombre quelconque de $c$, suivis d'un $b$, suivis d'un nombre
quelconque de $c$.

\begin{center}
%%% begin example2b %%%

\begin{tikzpicture}[>=latex,line join=bevel,automaton]
%%
\node (q1) at (98bp,18bp) [draw,circle,state] {$1$};
  \node (q0) at (18bp,18bp) [draw,circle,state,initial] {$0$};
  \node (q2) at (178bp,18bp) [draw,circle,state,final] {$2$};
  \draw [->] (q0) ..controls (46.106bp,18bp) and (58.578bp,18bp)  .. node[auto] {$a$} (q1);
  \draw [->] (q1) to[loop above] node[auto] {$c$} (q1);
  \draw [->] (q1) ..controls (126.11bp,18bp) and (138.58bp,18bp)  .. node[auto] {$b$} (q2);
  \draw [->] (q0) to[loop above] node[auto] {$c$} (q0);
  \draw [->] (q2) to[loop above] node[auto] {$c$} (q2);
%
\end{tikzpicture}

%%% end example2b %%%
\end{center}

(Ce langage est aussi dénoté par l'expression rationnelle
$c{*}ac{*}bc{*}$.)

\begin{prop}\label{completion-of-dfai}
Soit $A = (Q,q_0,F,\delta)$ un DFAi sur un alphabet $\Sigma$.  Alors
il existe un DFA $A' = (Q',q'_0,F',\delta')$ (sur le même
alphabet $\Sigma$) qui soit équivalent à $A$ au sens où il reconnaît
le même langage $L(A') = L(A)$.  De plus, $A'$ se déduit
algorithmiquement de $A$ en ajoutant au plus un état « puits » à $A$ :
on a $\#Q' \leq \#Q + 1$.
\end{prop}
\begin{proof}
On définit $Q' = Q \cup \{q_\bot\}$ où $q_\bot$ est un nouvel état
(n'appartenant pas à $Q$), qu'on appellera « puits ».  On garde l'état
initial $q'_0 = q_0$.  On garde l'ensemble $F' = F$ d'états finaux,
c'est-à-dire notamment que le puits n'est pas acceptant.  Enfin, on
définit $\delta'(q,x)$ pour $q\in Q'$ et $x\in\Sigma$ par
\[
\begin{aligned}
\delta'(q,x) &= \delta(q,x)\text{ si $\delta(q,x)$ est défini}\\
\delta'(q,x) &= q_\bot\text{ sinon}\\
\end{aligned}
\]
(notamment, $\delta'(q_\bot,x) = q_\bot$ quel que soit $x$).

Il est alors facile de voir que $A'$ a le même comportement que $A$ au
sens où $\delta^{\prime*}(q,w) = \delta^*(q,w)$ lorsque le terme de
droite est défini et $\delta^{\prime*}(q,w) = q_\bot$ sinon (le DFA
$A'$ « tombe dans le puits » lorsque le DFAi $A$ cesse de
fonctionner).  En particulier, ils reconnaissent les mêmes langages.
\end{proof}

\thingy De façon générale, un état \defin[puits (état)]{puits} dans un
DFA est un état $q_\bot$ tel que $\delta(q_\bot,x)=q_\bot$ pour toute
lettre $x$.

On dit que le DFA $A'$ est obtenu en \defin[compléter (un
  DFAi)]{complétant} le DFAi $A$ lorsqu'il est obtenu par la procédure
décrite dans la démonstration de cette proposition, c'est-à-dire par
l'addition d'un état puits vers lequel on fait pointer toutes les
transitions manquantes (sauf si $A$ est déjà complet, auquel cas on
convient qu'il est son propre complété, i.e., on n'ajoute un puits que
quand c'est réellement nécessaire).

\thingy À titre d'exemple, le DFA suivant représente la complétion du
DFAi représenté en \ref{discussion-example2b} :

\begin{center}
%%% begin example2c %%%

\begin{tikzpicture}[>=latex,line join=bevel,automaton]
%%
\begin{scope}
  \pgfsetstrokecolor{black}
  \definecolor{strokecol}{rgb}{1.0,1.0,1.0};
  \pgfsetstrokecolor{strokecol}
  \definecolor{fillcol}{rgb}{1.0,1.0,1.0};
  \pgfsetfillcolor{fillcol}
  \filldraw (0bp,0bp) -- (0bp,182bp) -- (214bp,182bp) -- (214bp,0bp) -- cycle;
\end{scope}
  \node (q1) at (102bp,131bp) [draw,circle,state] {$1$};
  \node (q0) at (18bp,85bp) [draw,circle,state,initial] {$0$};
  \node (q2) at (196bp,85bp) [draw,circle,state,final] {$2$};
  \node (qbot) at (102bp,22bp) [draw,circle,state] {$\bot$};
  \draw [->] (q1) to[loop above] node[auto] {$c$} (q1);
  \draw [->] (q2) to[loop above] node[auto] {$c$} (q2);
  \draw [->] (qbot) to[loop below] node[auto] {$a,b,c$} (qbot);
  \draw [->] (q0) ..controls (44.565bp,65.359bp) and (61.506bp,52.343bp)  .. node[auto] {$b$} (qbot);
  \draw [->] (q0) to[loop above] node[auto] {$c$} (q0);
  \draw [->] (q1) ..controls (102bp,96.993bp) and (102bp,73.356bp)  .. node[auto] {$a$} (qbot);
  \draw [->] (q0) ..controls (46.061bp,100.18bp) and (63.141bp,109.76bp)  .. node[auto] {$a$} (q1);
  \draw [->] (q2) ..controls (166.87bp,65.735bp) and (145.76bp,51.281bp)  .. node[auto] {$a,b$} (qbot);
  \draw [->] (q1) ..controls (132.83bp,116.08bp) and (154.08bp,105.46bp)  .. node[auto] {$b$} (q2);
%
\end{tikzpicture}

%%% end example2c %%%
\end{center}

\thingy\label{coaccessible-and-useful-states} On définit un état
accessible d'un DFAi comme pour un DFA
(cf. \ref{definition-dfa-accessible-state}).

\smallskip

On dira en outre d'un état $q$ d'un DFAi qu'il est
\defin[co-accessible (état)]{co-accessible} lorsqu'il existe un mot $w
\in \Sigma^*$ tel que $\delta^*(q,w)$ soit défini et soit final,
autrement dit, graphiquement, lorsqu'il existe un chemin orienté
reliant l'état $q$ considéré à un état final (remarquer que les états
finaux eux-mêmes sont co-accessibles : prendre $w=\varepsilon$ dans ce
qu'on vient de dire).  Un état non co-accessible est donc un état à
partir duquel il est impossible de faire accepter le mot.  Cette
définition pourrait également être faite pour les DFA, mais pour les
DFAi elle présente l'intérêt qu'on peut supprimer les états
non co-accessibles (ainsi, bien sûr, que toutes les transitions qui y
conduisent) et obtenir de nouveau un DFAi.

\smallskip

Un DFAi dont tous les états sont à la fois accessibles et
co-accessibles (on les dit aussi \defin[utile (état)]{utiles}) est
parfois appelé \defin[émondé (automate)]{émondé}.  On peut émonder un
DFAi en ne conservant que ses états utiles\footnote{Si on aime
  pinailler, il y a un petit problème pour émonder un DFAi n'ayant
  aucun état final (donc aucun état co-accessible, donc aucun état
  utile) ; pour résoudre ce problème, on peut modifier légèrement la
  définition d'un DFAi et autoriser que l'état initial ne soit pas
  défini (auquel cas l'automate n'accepte évidemment aucun mot, i.e.,
  reconnaît le langage $\varnothing$), si bien que l'automate émondé
  d'un DFAi sans état final est l'automate vide (sans aucun état).} :
ainsi, tout DFAi est équivalent à un DFAi émondé.

\thingy Il faut prendre garde au fait que certains auteurs définissent
les « automates finis déterministes » (sans précision supplémentaire)
comme étant complets par défaut, d'autres comme étant incomplets par
défaut.  Le plus prudent est de préciser systématiquement « complet »
ou « incomplet » (en se rappelant qu'« incomplet » signifie « non
nécessairement complet », cf. la
note \ref{footnote-on-word-incomplete} sous \ref{definition-dfai}) dès
qu'il est important de ne pas confondre.


\subsection{Automates finis non-déterministes (=NFA)}

\thingy\textbf{Idée générale :} Le principe directeur général du
non-déterminisme en informatique est, grossièrement parlant, le
suivant : on a un mécanisme de calcul (ici, un automate fini) qui doit
accepter ou rejeter une certaine donnée (ici, un mot qui lui est
présenté), mais plutôt que d'évoluer de façon déterministe d'une
configuration à une autre (ici, d'un état à un autre), il se peut
qu'il existe plusieurs possibilités différentes, autrement dit, une
configuration peut évoluer de plusieurs manières différentes, ce qui
donne lieu à plusieurs chemins de calcul différents.  Selon le point
de vue qu'on adopte, et éventuellement la stratégie algorithmique pour
modéliser le non-déterminisme, on peut considérer que ces chemins sont
empruntés tous simultanément (et qu'on est dans tous les états
possibles à la fois\footnote{Certains sont parfois tentés de comparer
  avec une superposition quantique : il existe effectivement une
  analogie, mais le quantique implique une description bien précise
  des états (avec des amplitudes associées) qui n'est pas pertinente
  ici, et l'analogie mal appliquée conduirait à des erreurs en
  complexité algorithmique ; il vaut donc mieux ne pas y penser.},
c'est d'ailleurs l'approche qui servira
en \ref{determinization-of-nfa}), ou que tous les chemins possibles
sont testés successivement, ou encore que l'appareil non-déterministe
« devine » quel est le bon ; mais en tout cas la règle fondamentale du
non-déterminisme est toujours la suivante : \emph{la donnée est
  acceptée dès lors qu'il existe au moins un chemin de calcul possible
  qui conduit à l'accepter} (dans notre cas, cela signifiera au moins
un chemin acceptant le mot).

\thingy\label{definition-nfa} Un \defin{automate fini
  non-déterministe}, en abrégé \index{NFA|see{automate fini
    non-déterministe}}\textbf{NFA}, sur un alphabet $\Sigma$ est la
donnée
\begin{itemize}
\item d'un ensemble fini $Q$ d'états,
\item d'un ensemble $I \subseteq Q$ d'états dits initiaux,
\item d'un ensemble $F \subseteq Q$ d'états dits finaux,
\item d'une \emph{relation} de \index{transition (relation
  de)}transition $\delta \subseteq Q \times \Sigma \times Q$
  (c'est-à-dire une partie du produit cartésien $Q \times \Sigma
  \times Q$, i.e., un ensemble de triplets $(q,x,q') \in Q \times
  \Sigma \times Q$) ; lorsque $(q,x,q') \in \delta$, on dira qu'il
  existe une transition étiquetée par $x$ de $q$ vers $q'$.
\end{itemize}
Autrement dit, on autorise maintenant un ensemble quelconque d'états
initiaux, et de même, au lieu qu'un état $q$ et une lettre $x$
déterminent un unique état $q' = \delta(q,x)$, on a maintenant affaire
à un ensemble $\delta$ quelconque de triplets $(q,x,q')$ pour les
transitions possibles.

\thingy Un DFAi (ou \textit{a fortiori} un DFA) est considéré comme un
NFA particulier en définissant l'ensemble des états initiaux du NFA
comme un singleton, $I_{\mathrm{NFA}} = \{q_{0,\mathrm{DFAi}}\}$, et
en définissant la relation de transition du NFA comme le graphe de la
fonction de transition du DFAi, c'est-à-dire $(q,x,q') \in
\delta_{\mathrm{NFA}}$ lorsque $\delta_{\mathrm{DFAi}}(q,x)$ est
défini et vaut $q'$.

\thingy Graphiquement, on représente un NFA comme un DFA : comme un
graphe orienté dont les nœuds sont les éléments de $Q$, et où on place
une arête étiquetée $x$ de $q$ vers $q'$ pour chaque triplet $(q,x,q')
\in \delta$ ; comme précédemment, on marque les états initiaux par une
flèche entrante (i.e., pointant de nulle part) et les états finaux par
une flèche sortante (i.e., pointant vers nulle part).

\thingy Il faut comprendre le fonctionnement d'un NFA de la manière
suivante : un mot $w$ est accepté par l'automate lorsqu'\emph{il
  existe} un chemin orienté conduisant d'\emph{un} état initial $q_0$
à un état final $q_n$ et tel que le mot $w = x_1 \cdots x_n$ soit
obtenu en lisant dans l'ordre les étiquettes $x_i$ des différentes
arêtes $q_{i-1} \to q_i$ de ce chemin ; autrement dit, $w$ est accepté
lorsqu'\emph{il existe} $q_0,\ldots,q_n \in Q$ tels que $q_0 \in I$ et
$q_n\in F$ et $(q_{i-1},x_i,q_i) \in \delta$ pour chaque $1\leq i\leq
n$.

\smallskip

La différence cruciale avec les DFAi est donc que, maintenant, il
pourrait exister plusieurs chemins possibles partant d'un état initial
dont les transitions sont étiquetées par les lettres du même mot.
Insistons bien sur le fait que le mot est accepté dès lors que
\emph{l'un} au moins de ces chemins relie un état initial à un état
final.

\thingy De même que dans le cadre des DFA et DFAi on a étendu la
fonction de transition $\delta$ à une fonction $\delta^*$ qui décrit
le comportement de l'automate quand on consomme plusieurs caractères
(cf. \ref{definition-multiple-transition-function}), on va étendre la
relation de transition $\delta$ d'un NFA à la situation où on consomme
plusieurs caractères.  On a $(q,w,q') \in \delta^*$ lorsqu'il existe
un chemin orienté conduisant de $q$ à $q'$ et tel que le mot $w$ soit
obtenu en lisant dans l'ordre les étiquettes des différentes arêtes de
ce chemin.

\smallskip

Formellement : si $A = (Q,I,F,\delta)$ est un NFA sur
l'alphabet $\Sigma$, on définit une relation $\delta^* \subseteq Q
\times \Sigma^* \times Q$ par $(q,w,q') \in \delta^*$ lorsque $w =
x_1\cdots x_n$ et qu'il existe $(q_0,\ldots,q_n)$ tels que $q_0 = q$
et $q_n = q'$ et $(q_{i-1},x_i,q_i) \in \delta$ pour chaque $1\leq
i\leq n$ ; ou, ce qui revient au même (par récurrence sur la longueur
de $w$) :
\begin{itemize}
\item $(q,\varepsilon,q') \in \delta^*$ si et seulement si $q'=q$,
\item $(q,wx,q') \in \delta^*$ si et seulement si il existe $q^\sharp$
  tel que $(q,w,q^\sharp) \in \delta^*$ et $(q^\sharp,x,q') \in
  \delta$.
\end{itemize}

Enfin, l'automate $A$ accepte un mot $w$ lorsqu'il existe $q_0\in I$
et $q_\infty\in F$ tels que $(q_0,w,q_\infty) \in \delta^*$.

\smallskip

Le langage accepté $L(A)$ et l'équivalence de deux automates sont
définis de façon analogue aux DFA
(cf. \ref{definition-recognizable-language}).

\thingy\label{discussion-example4} Pour illustrer le fonctionnement
des NFA, considérons l'automate à trois états sur $\Sigma=\{a,b\}$
représenté par la figure suivante : on a $Q = \{0,1,2\}$ avec
$I=\{0\}$ et $F=\{2\}$ et $\delta = \{(0,a,0),\penalty0
(0,b,0),\penalty-100 (0,a,1),\penalty-50 (1,a,2),\penalty0 (1,b,2)\}$.

\begin{center}
%%% begin example4 %%%

\begin{tikzpicture}[>=latex,line join=bevel,automaton]
%%
\node (q1) at (98bp,18bp) [draw,circle,state] {$1$};
  \node (q0) at (18bp,18bp) [draw,circle,state,initial] {$0$};
  \node (q2) at (188bp,18bp) [draw,circle,state,final] {$2$};
  \draw [->] (q0) ..controls (46.106bp,18bp) and (58.578bp,18bp)  .. node[auto] {$a$} (q1);
  \draw [->] (q1) ..controls (128.76bp,18bp) and (145.63bp,18bp)  .. node[auto] {$a,b$} (q2);
  \draw [->] (q0) to[loop above] node[auto] {$a,b$} (q0);
%
\end{tikzpicture}

%%% end example4 %%%
\end{center}

Cet automate n'est pas déterministe car il existe deux transitions
étiquetées $a$ partant de l'état $0$.  En considérant les différents
chemins possibles entre $0$ et $2$ dans ce graphe, on comprend que le
langage qu'il reconnaît est le langage des mots sur $\{a,b\}$ dont
l'avant-dernière lettre est un $a$ (c'est aussi le langage dénoté par
l'expression rationnelle $(a|b){*}a(a|b)$).  Une façon de présenter le
non-déterminisme est que l'automate « devine », quand il est dans
l'état $0$ et qu'on lui fait consommer un $a$, si ce $a$ sera
l'avant-dernière lettre, et, dans ce cas, passe dans l'état $1$ pour
pouvoir accepter le mot.

\bigbreak

Comme les DFAi avant eux, les NFA sont en fait équivalents aux DFA ;
mais cette fois-ci, le coût algorithmique de la transformation peut
être bien plus important :

\begin{prop}\label{determinization-of-nfa}
Soit $A = (Q,I,F,\delta)$ un NFA sur un alphabet $\Sigma$.  Alors il
existe un DFA $A' = (Q',\textbf{q}'_0,F',\delta')$ (sur le même
alphabet $\Sigma$) qui soit équivalent à $A$ au sens où il reconnaît
le même langage $L(A') = L(A)$.  De plus, $A'$ se déduit
algorithmiquement de $A$ avec une augmentation au plus exponentielle
du nombre d'états : $\#Q' \leq 2^{\#Q}$.
\end{prop}
\begin{proof}
On définit $Q' = \mathscr{P}(Q) = \{\mathbf{q} \subseteq Q\}$
l'\emph{ensemble des parties} de $Q$ : c'est ce qui servira d'ensemble
d'états du DFA $A'$ qu'on construit (i.e., un état de $A'$ est un
ensemble d'états de $A$ — intuitivement, c'est l'ensemble des états
dans lesquels on se trouve simultanément).  On pose $\textbf{q}'_0 = I$ (l'état
initial de $A'$ sera l'ensemble de tous les états initiaux de $A$, vu
comme un élément de $Q'$) ; et on pose $F' = \{\mathbf{q}\subseteq Q
:\penalty-100 \mathbf{q}\cap F \neq\varnothing\}$ : les états finaux
de $A'$ sont les états $\mathbf{q}$ qui, vus comme des ensembles
d'états de $A$, contiennent \emph{au moins un} état final.  Enfin,
pour $\mathbf{q}\subseteq Q$ et $x \in \Sigma$, on définit
$\delta'(\mathbf{q},x) = \{q_1\in Q : \exists q_0\in\mathbf{q}\,
((q_0,x,q_1) \in \delta)\}$ comme l'ensemble de tous les états $q_1$
de $A$ auxquels on peut accéder depuis un état $q_0$ dans $\mathbf{q}$
par une transition $(q_0,x,q_1)$ (étiquetée par $x$) de $A$.

Il est alors facile de voir (par récurrence sur $|w|$) que
$\delta^{\prime*}(\mathbf{q},w)$ est l'ensemble de tous les les états
$q_1 \in Q$ tels que $(q_0,w,q_1)\in\delta^*$, i.e., auxquels on peut
accéder depuis un état $q_0$ dans $\mathbf{q}$ par une suite de
transitions de $A$ étiquetées par les lettres de $w$.  En particulier,
$\delta^{\prime*}(I,w)$ est l'ensemble de tous les états de $A$
auxquels on peut accéder depuis un état initial de $A$ par une suite
de transitions de $A$ étiquetées par les lettres de $w$ : le mot $w$
appartient à $L(A)$ si et seulement si cet ensemble contient un élément
de $F$, ce qui par définition de $F'$ signifie exactement
$\delta^{\prime*}(I,w) \in F'$.  On a bien prouvé $L(A') = L(A)$.

Enfin, $\#Q' = \#\mathscr{P}(Q) = 2^{\#Q}$ (car une partie de $Q$ peut
se voir à travers sa fonction indicatrice, qui est une fonction $Q \to
\{0,1\}$).
\end{proof}

\thingy On dit que le DFA $A'$ est obtenu en \defin[déterminiser (un
  NFA)]{déterminisant} le NFA $A$ lorsqu'il est obtenu par la
procédure décrite dans la démonstration de cette proposition en ne
gardant que les états accessibles.

\smallskip

Algorithmiquement, la déterminisation de $A$ s'obtient par la
procéduire suivante :
\begin{itemize}
\item créer une file (ou une pile) d'ensembles d'états de $A$ ;
  initialiser cette file avec le seul élément $I$ (vu comme un
  sous-ensemble de $Q$) ; et créer l'automate $A'$ avec initialement
  l'unique état $I$, marqué comme état initial, et aussi comme final
  s'il contient un état final de $A$ ;
\item tant que la file/pile n'est pas vide : en extraire un élément
  $\mathbf{q}$, et, pour chaque lettre $x\in\Sigma$,
\begin{itemize}
\item calculer l'ensemble $\mathbf{q}' = \{q_1\in Q : \exists
  q_0\in\mathbf{q}\, ((q_0,x,q_1) \in \delta)\}$ (en listant tous les
  triplets $(q_0,x,q_1) \in\delta$ dont le premier élément est
  dans $\mathbf{q}$ et le second élément est $x$),
\item si $\mathbf{q}'$ n'existe pas déjà comme état de $A'$, l'y
  ajouter, et dans ce cas, l'ajouter à la file/pile, et de plus, si
  $\mathbf{q}'$ contient un état final de $A$, marquer cet état comme
  final pour $A'$,
\item et ajouter à $A'$ la transition $\delta'(\mathbf{q},x) =
  \mathbf{q}'$.
\end{itemize}
\end{itemize}

La file/pile sert à stocker les états de $A'$ qui ont été créés mais
pour lesquels les transitions sortantes n'ont pas encore été
calculées.  L'algorithme se termine quand la file/pile se vide, ce qui
se produit toujours en au plus $2^{\#Q}$ étapes car chaque $\mathbf{q}
\subseteq Q$ ne peut apparaître qu'une seule fois dans la file/pile.

Il se peut que l'état $\varnothing$ soit créé : cet état servira
effectivement de puits, au sens où on aura $\delta'(\varnothing,x) =
\varnothing$ quel que soit $x$ (et l'état n'est pas acceptant).

Il arrive souvent que l'automate déterminisé soit plus petit que les
$2^{\#Q}$ états qu'il a dans le pire cas.

\thingy À titre d'exemple, déterminisons le NFA $A$ présenté
en \ref{discussion-example4} : on commence par construire un état
$\{0\}$ pour $A'$ car le NFA a $\{0\}$ pour ensemble d'états
initiaux ; on a $\delta'(\{0\},a) = \{0,1\}$ car $0$ et $1$ sont les
deux états auxquels on peut arriver dans $A$ par une transition
partant de $0$ et étiquetée par $a$, tandis que $\delta'(\{0\},b) =
\{0\}$ ; ensuite, $\delta'(\{0,1\},a) = \{0,1,2\}$ car $0,1,2$ sont
les trois états auxquels on peut arriver dans $A$ par une transition
partant de $0$ ou $1$ et étiquetée par $a$ ; et ainsi de suite.  En
procédant ainsi, on construit l'automate à $4$ états qui suit :

\begin{center}
\footnotesize
%%% begin example4det %%%

\begin{tikzpicture}[>=latex,line join=bevel,automaton]
%%
\begin{scope}
  \pgfsetstrokecolor{black}
  \definecolor{strokecol}{rgb}{1.0,1.0,1.0};
  \pgfsetstrokecolor{strokecol}
  \definecolor{fillcol}{rgb}{1.0,1.0,1.0};
  \pgfsetfillcolor{fillcol}
  \filldraw (0bp,0bp) -- (0bp,163bp) -- (212bp,163bp) -- (212bp,0bp) -- cycle;
\end{scope}
  \node (q02) at (188bp,27bp) [draw,circle,state,final] {$\{0,2\}$};
  \node (q0) at (18bp,18bp) [draw,circle,state,initial] {$\{0\}$};
  \node (q01) at (100bp,56bp) [draw,circle,state] {$\{0,1\}$};
  \node (q012) at (188bp,106bp) [draw,circle,state,final] {$\{0,1,2\}$};
  \draw [->] (q01) ..controls (127.39bp,49.663bp) and (137.27bp,46.959bp)  .. (146bp,44bp) .. controls (150.77bp,42.384bp) and (155.77bp,40.486bp)  .. node[auto] {$b$} (q02);
  \draw [->] (q02) ..controls (159.35bp,21.739bp) and (147.76bp,21.451bp)  .. (138bp,25bp) .. controls (132.05bp,27.164bp) and (126.38bp,30.754bp)  .. node[auto] {$a$} (q01);
  \draw [->] (q0) ..controls (45.691bp,30.678bp) and (60.407bp,37.668bp)  .. node[auto] {$a$} (q01);
  \draw [->] (q0) to[loop above] node[auto] {$b$} (q0);
  \draw [->] (q01) ..controls (128.64bp,72.07bp) and (144.34bp,81.198bp)  .. node[auto] {$a$} (q012);
  \draw [->] (q012) to[loop above] node[auto] {$a$} (q012);
  \draw [->] (q02) ..controls (155.56bp,18.853bp) and (136.78bp,14.735bp)  .. (120bp,13bp) .. controls (102.32bp,11.172bp) and (97.763bp,12.273bp)  .. (80bp,13bp) .. controls (68.968bp,13.451bp) and (56.843bp,14.36bp)  .. node[auto] {$b$} (q0);
  \draw [->] (q012) ..controls (188bp,73.926bp) and (188bp,64.965bp)  .. node[auto] {$b$} (q02);
%
\end{tikzpicture}

%%% end example4det %%%
\end{center}

On remarquera qu'on a construit moins que les $2^3 = 8$ états qu'on
pouvait craindre.

{\footnotesize Il est par ailleurs instructif de regarder comment
  fonctionne l'automate $A'$ ci-dessus pour déterminer si
  l'avant-dernière lettre d'un mot est un $a$ : essentiellement, il
  retient deux bits d'information, à savoir « est-ce que la dernière
  lettre consommée était un $a$ ? » et « est-ce que l'avant-dernière
  lettre consommée était un $a$ ? » ; chacune des quatre combinaisons
  de ces deux bits correspond à un état de $A'$ : non/non est l'état
  $\{0\}$, oui/non est l'état $\{0,1\}$, non/oui est l'état $\{0,2\}$,
  et oui/oui est l'état $\{0,1,2\}$.\par}

\thingy On vient donc de voir que les NFA sont équivalents aux DFA en
ce sens qu'ils permettent de définir les mêmes langages.  Il est
néanmoins intéressant de les définir et de les étudier pour plusieurs
raisons : d'une part, l'équivalence qu'on vient de voir, c'est-à-dire
la déterminisation d'un NFA en DFA, a un coût algorithmique important
(exponentiel en nombre d'états dans le pire cas), d'autre part,
certaines opérations sur les automates se définissent plus
naturellement sur les NFA que sur les DFA, et certaines constructions
conduisent naturellement à un NFA.  On verra en particulier
en \ref{subsection-recognizable-languages-under-rational-operations}
que pour transformer une expression rationnelle en automate, on
passera naturellement par un NFA, même si on peut souhaiter le
déterminiser ensuite.


\subsection{Automates finis non-déterministes à transitions spontanées (=εNFA)}

\thingy\label{definition-enfa} Un \defin{automate fini
  non-déterministe à transitions spontanées} ou \textbf{...à
  $\varepsilon$-transitions}, en abrégé
\index{epsilon-NFA@$\varepsilon$NFA|see{automate fini non-déterministe
    à transitions spontanées}}\textbf{εNFA}, sur un alphabet $\Sigma$
est la donnée
\begin{itemize}
\item d'un ensemble fini $Q$ d'états,
\item d'un ensemble $I \subseteq Q$ d'états dits initiaux,
\item d'un ensemble $F \subseteq Q$ d'états dits finaux,
\item d'une relation de transition $\delta \subseteq Q \times
  (\Sigma\cup\{\varepsilon\}) \times Q$.
\end{itemize}

Autrement dit, on autorise maintenant des transitions étiquetées par
le mot vide $\varepsilon$ plutôt que par une lettre $x \in\Sigma$ :
ces transitions sont dites \defin[spontanée (transition)]{spontanées},
ou
\index{epsilon-transition@$\varepsilon$-transition|see{spontanée}}\textbf{ε-transitions}.

\smallskip

Soulignons qu'on ne définit les ε-transitions \emph{que} pour les
automates non-déterministes : ou, pour dire les choses autrement,
\emph{un automate qui possède des ε-transitions est par nature même
  non-déterministe}.

\smallskip

La représentation graphique des εNFA est évidente (on utilisera le
symbole « $\varepsilon$ » pour étiqueter les transitions spontanées).
Un NFA est considéré comme un εNFA particulier pour lequel il n'y a
pas de ε-transition.

\thingy Intuitivement, les transitions spontanées signifient que
l'automate a la possibilité de passer spontanément, c'est-à-dire
\emph{sans consommer de lettre}, d'un état $q$ à un état $q'$, lorsque
ces états sont reliés par une ε-transition.  (On comprend, du coup,
pourquoi un automate à transition spontanée est forcément
non-déterministe : ces transitions spontanées ne sont qu'une
\emph{possibilité}, ce qui sous-entend le non-déterminisme.)

\smallskip

De façon plus précise, un εNFA accepte un mot $w$ lorsqu'\emph{il
  existe} un chemin orienté conduisant d'un état initial $q_0$ à un
état final $q_n$ et tel que $w$ coïncide avec le mot obtenu en lisant
dans l'ordre les étiquettes $t_i$ des différentes arêtes $q_{i-1} \to
q_i$ de ce chemin, quitte à ignorer les $\varepsilon$.

Autrement dit, $w$ est accepté lorsqu'\emph{il existe} $q_0,\ldots,q_n
\in Q$ et $t_1,\ldots,t_n \in (\Sigma\cup\{\varepsilon\})$ tels que
$q_0 \in I$ et $q_n\in F$ et $(q_{i-1},t_i,q_i) \in \delta$ pour
chaque $1\leq i\leq n$ et $w = t_1\cdots t_n$ (\emph{attention} : dans
cette écriture, $t_1,\ldots,t_n$ ne sont pas forcément les lettres
de $w$, certains des $t_i$ peuvent être le mot vide $\varepsilon$, les
lettres de $w$ sont obtenues en ignorant ces symboles).

\thingy On veut de nouveau définir une relation $\delta^*$ telle que
$(q,w,q') \in \delta^*$ lorsqu'il existe un chemin orienté conduisant
de $q$ à $q'$ et tel que le mot $w$ soit obtenu en lisant dans l'ordre
les étiquettes des différentes arêtes de ce chemin.

Voici une formalisation possible : si $A = (Q,I,F,\delta)$ est
un εNFA sur l'alphabet $\Sigma$, on définit une relation $\delta^*
\subseteq Q \times \Sigma^* \times Q$ par $(q,w,q') \in \delta^*$
lorsqu'il existe $q_0,\ldots,q_n \in Q$ et $t_1,\ldots,t_n \in
(\Sigma\cup\{\varepsilon\})$ tels que $q_0 = q$ et $q_n = q'$ et
$(q_{i-1},t_i,q_i) \in\delta$ pour chaque $1\leq i\leq n$, et enfin $w
= t_1\cdots t_n$.

Enfin, l'automate $A$ accepte un mot $w$ lorsqu'il existe $q_0\in I$
et $q_\infty\in F$ tels que $(q_0,w,q_\infty) \in \delta^*$.

\smallskip

Le langage accepté $L(A)$ et l'équivalence de deux automates sont
définis de façon analogue aux DFA
(cf. \ref{definition-recognizable-language}).

\thingy\label{discussion-example5} Voici un exemple de εNFA
particulièrement simple sur $\Sigma = \{a,b,c\}$:

\begin{center}
%%% begin example5 %%%

\begin{tikzpicture}[>=latex,line join=bevel,automaton]
%%
\node (q1) at (98bp,18bp) [draw,circle,state] {$1$};
  \node (q0) at (18bp,18bp) [draw,circle,state,initial] {$0$};
  \node (q2) at (178bp,18bp) [draw,circle,state,final] {$2$};
  \draw [->] (q0) ..controls (46.106bp,18bp) and (58.578bp,18bp)  .. node[auto] {$\varepsilon$} (q1);
  \draw [->] (q1) to[loop above] node[auto] {$b$} (q1);
  \draw [->] (q1) ..controls (126.11bp,18bp) and (138.58bp,18bp)  .. node[auto] {$\varepsilon$} (q2);
  \draw [->] (q0) to[loop above] node[auto] {$a$} (q0);
  \draw [->] (q2) to[loop above] node[auto] {$c$} (q2);
%
\end{tikzpicture}

%%% end example5 %%%
\end{center}

En considérant les différents chemins possibles entre $0$ et $2$ sur
ce graphe, on comprend que le langage qu'il reconnaît est celui des
mots sur $\{a,b,c\}$ formés d'un nombre quelconque de $a$ suivi d'un
nombre quelconque de $b$ suivi d'un nombre quelconque de $c$, ou, si
on préfère, le langage dénoté par l'expression rationnelle
$a{*}b{*}c{*}$.

\thingy Si $q$ est un état d'un εNFA, on appelle
\defin[epsilon-fermeture@$\varepsilon$-fermeture]{ε-fermeture} de $q$
l'ensemble des états $q'$ (y compris $q$ lui-même) accessibles depuis
$q$ par une succession quelconque de ε-transitions, c'est-à-dire, si
on veut, $\{q'\in Q :\penalty0 (q,\varepsilon,q') \in\delta^*\}$.  On
notera temporairement $C(q)$ cet ensemble.  (Par exemple, dans
l'exemple \ref{discussion-example5} ci-dessus, on a $C(0) = \{0,1,2\}$
et $C(1) = \{1,2\}$ et $C(2) = \{2\}$.  Dans tout NFA sans
ε-transitions, on a $C(q) = \{q\}$ pour tout état $q$.)

\smallskip

Il est clair qu'on peut calculer algorithmiquement $C(q)$ (par exemple
par un algorithme de Dijkstra / parcours en largeur, sur le graphe
orienté dont l'ensemble des sommets est $Q$ et l'ensemble des arêtes
est l'ensemble des ε-transitions de $A$ : la ε-fermeture $C(q)$ est
simplement l'ensemble des sommets accessibles depuis $q$ dans ce
graphe).

\begin{prop}\label{removal-of-epsilon-transitions}
Soit $A = (Q,I,F,\delta)$ un εNFA sur un alphabet $\Sigma$.  Alors il
existe un NFA $A^\S = (Q,I^\S,F^\S,\delta^\S)$ (sur le même
alphabet $\Sigma$) ayant le même ensemble d'états $Q$ que $A$ et qui
soit équivalent à $A$ au sens où il reconnaît le même langage
$L(A^\S) = L(A)$.  De plus, $A^\S$ se déduit algorithmiquement de $A$.
\end{prop}
\begin{proof}
On pose $I^\S = I$ (mêmes états initiaux).  L'idée est maintenant de
faire une transition $(q,x,q') \in \delta^\S$ à chaque fois qu'on peut
atteindre $q'$ à partir de $q$ dans $A$ par une suite quelconque de
ε-transitions suivie d'une unique transition étiquetée par $x$,
autrement dit, $(q,\varepsilon,q^\sharp) \in \delta^*$ (c'est-à-dire
$q^\sharp \in C(q)$) et $(q^\sharp,x,q') \in \delta$.

On définit donc $\delta^\S \subseteq Q\times\Sigma\times Q$ par
$(q,x,q') \in \delta^\S$ lorsqu'il existe $q^\sharp \in C(q)$ tel que
$(q^\sharp,x,q') \in \delta$ : autrement dit, pour créer les
transitions $q\to q'$ dans $A^\S$, on parcourt tous les $q^\sharp \in
C(q)$, et on crée une transition $q\to q'$ étiquetée par $x \in \Sigma$
dans $A^\S$ lorsqu'il existe une transition $q^\sharp\to q'$ étiquetée
par ce $x$ dans $A$.  De même, on définit $F^\S \subseteq Q$ comme
l'ensemble des $q\in Q$ tels que $C(q) \cap F \neq \varnothing$,
c'est-à-dire, depuis lesquels peut atteindre un état final par une
succession de ε-transitions.

Si on a un chemin $q_0 \to q_1 \to \cdots \to q_n$ dans $A$ menant
d'un état initial $q_0 \in I$ à un état final $q_n \in F$ et
étiquetées par $t_1,\ldots,t_n \in (\Sigma\cup\{\varepsilon\})$
(c'est-à-dire $(q_{i-1},t_i,q_i) \in \delta$), appelons
$j_1<\ldots<j_m$ les indices tels que $t_j \in\Sigma$, autrement dit,
tels que la transition $q_{j-1} \to q_j$ ne soit pas spontanée, et
posons $j_0 = 0$.  Alors on passe de $q_{j_{(i-1)}}$ à $q_{j_i}$ par une
succession de ε-transitions (de $q_{j_{(i-1)}}$ à $q_{(j_i)-1}$) suivie
par une unique transition non spontanée : on a $q_{(j_i)-1} \in
C(q_{j_{(i-1)}})$ et $(q_{(j_i)-1},t_{j_i},q_{j_i}) \in \delta$,
autrement dit $(q_{j_{(i-1)}},t_{j_i},q_{j_i}) \in \delta^\S$ ; et comme
le mot $w = t_1\cdots t_n$ s'écrit aussi $t_{j_1}\cdots t_{j_m}$, on a
un chemin reliant $q_{j_0} = q_0 \in I$ à $q_{j_m} \in F^\S$ (puisque
$q_n \in C(q_{j_m}) \cap F$).  Le mot $w$ supposé accepté par $A$ est
donc accepté par $A^\S$.  La réciproque est analogue.
\end{proof}

\thingy On dit que le NFA $A^\S$ est obtenu en \defin[éliminer (les
  transitions spontanées)]{éliminant les ε-transitions} dans le
εNFA $A$ lorsqu'il est obtenu par la procédure décrite dans la
démonstration de cette proposition, et en supprimant tous les états
non-initiaux de $A^\S$ auxquels n'aboutissent dans $A$ que des
ε-transitions (ces états sont devenus inaccessibles dans $A^\S$).
Algorithmiquement, il s'agit donc, pour chaque état $q\in Q$ et chaque
$q^\sharp$ dans la ε-fermeture $C(q)$ de $q$, de créer une transition
$q\to q'$ étiquetée par $x$ dans $A^\S$ pour chaque transition
$q^\sharp\to q'$ étiquetée par $x$ dans $A$.

\thingy\label{removal-of-epsilon-transitions-from-example5}
À titre d'exemple, éliminons les ε-transitions du εNFA $A$
présenté en \ref{discussion-example5} : comme $C(0) = \{0,1,2\}$, on
fait partir de $0$ toutes les transitions partant d'un des états
$0,1,2$ et étiquetées par une lettre, et de même, comme $C(1) =
\{1,2\}$, on fait partir de $1$ toutes les transitions partant d'un
des états $1,2$ et étiquetées par une lettre.  On obtient finalement
l'automate suivant :

\begin{center}
%%% begin example5ne %%%

\begin{tikzpicture}[>=latex,line join=bevel,automaton]
%%
\node (q1) at (97bp,48bp) [draw,circle,state,final,accepting below] {$1$};
  \node (q0) at (18bp,18bp) [draw,circle,state,initial,final,accepting below] {$0$};
  \node (q2) at (176bp,18bp) [draw,circle,state,final] {$2$};
  \draw [->] (q0) ..controls (47.643bp,10.917bp) and (64.2bp,7.4655bp)  .. (79bp,6bp) .. controls (94.922bp,4.4234bp) and (99.078bp,4.4234bp)  .. (115bp,6bp) .. controls (125.98bp,7.0877bp) and (137.94bp,9.2693bp)  .. node[auto] {$c$} (q2);
  \draw [->] (q2) to[loop above] node[auto] {$c$} (q2);
  \draw [->] (q0) to[loop above] node[auto] {$a$} (q0);
  \draw [->] (q1) to[loop above] node[auto] {$b$} (q1);
  \draw [->] (q1) ..controls (124.28bp,37.762bp) and (137.94bp,32.438bp)  .. node[auto] {$c$} (q2);
  \draw [->] (q0) ..controls (45.279bp,28.238bp) and (58.943bp,33.562bp)  .. node[auto] {$b$} (q1);
%
\end{tikzpicture}

%%% end example5ne %%%
\end{center}

(Sur cet exemple précis, on obtient un automate déterministe
incomplet, mais ce n'est pas un phénomène général : en général il faut
s'attendre à obtenir un NFA.)

{\footnotesize

\thingy\textbf{Remarque :} La manière dont on a éliminé les
ε-transitions ci-dessus consiste à remplacer \emph{une succession de
  ε-transitions suivie d'une unique transition étiquetée $x$} par une
transition étiquetée $x$ (et de même, on modifie les états finaux,
mais pas les états initiaux).  Il existait un moyen « dual »
d'éliminer les ε-transitions, à savoir remplacer \emph{une unique
  transition étiquetée $x$ suivie d'une succession de ε-transitions}
par une transition étiquetée $x$ (et de même, on modifie les états
initiaux, mais pas les états finaux).  Autrement dit, la construction
$A \mapsto A^\S$ décrite en \ref{removal-of-epsilon-transitions}
définit $(q,x,q') \in \delta^\S$ lorsqu'il existe $q^\sharp \in C(q)$
tel que $(q^\sharp,x,q') \in \delta$, et $I^\S=I$ et $F^\S$ comme
l'ensemble des états $q$ tels que $C(q) \cap F \neq \varnothing$ ; la
construction « duale » $A \mapsto A^\P$ consiste à poser $(q,x,q') \in
\delta^\P$ lorsqu'il existe $q^{\flat\prime}$ tel que
$(q,x,q^{\flat\prime}) \in \delta$ et $q' \in C(q^{\flat\prime})$, et
$F^\P=F$ et $I^\P$ comme la réunion des $C(q)$ pour tout $q\in I$.

Ces deux manières d'éliminer les ε-transitions donnent des NFA
équivalents.  En fait, on peut définir l'une en termes de l'autre en
utilisant la définition de l'automate transposé $A^{\mathsf{R}}$
présentée en \ref{nfa-mirror} plus bas (et qui consiste simplement à
inverser le sens de toutes les flèches de l'automate) : si on inverse
les flèches, qu'on élimine les ε-transitions à la manière décrite
en \ref{removal-of-epsilon-transitions}, et qu'on inverse de nouveau
les flèches, on obtient l'élimination « duale », autrement dit, $A^\P
= ((A^{\mathsf{R}})^\S)^{\mathsf{R}}$.

L'une ou l'autre manière d'éliminer les ε-transitions était possible,
mais il vaut mieux ne pas les mélanger.  C'est pour cette raison qu'on
a fait un choix en \ref{removal-of-epsilon-transitions} ; la présente
remarque a principalement pour objectif d'expliquer la raison d'une
perte de symétrie (notamment entre états initiaux et finaux) dans ce
choix.

\par}

\thingy L'intérêt des εNFA, même s'ils sont finalement équivalents aux
NFA, est que certaines constructions sur les automates sont plus
simples ou plus transparentes lorsqu'on s'autorise à fabriquer des
εNFA (à titre d'exemple, l'automate présenté
en \ref{discussion-example5} est beaucoup plus transparent à lire que
le NFA équivalent donné
en \ref{removal-of-epsilon-transitions-from-example5}).  On verra
notamment en \ref{subsection-thompson-construction} qu'il est facile
(quoique inefficace) de construire un εNFA reconnaissant le langage
dénoté par une expression rationnelle quelconque ; et
en \ref{subsection-recognizable-languages-under-rational-operations}
que, même si on fabrique directement un NFA, il peut être utile
d'introduire temporairement des transitions spontanées dans la
construction de ce NFA, quitte à les éliminer immédiatement (cela
simplifie la description).


\section{Langages reconnaissables et langages rationnels}\label{section-recognizable-languages}

\subsection{Stabilité des langages reconnaissables par opérations booléennes et miroir}

\thingy On rappelle qu'on a défini un langage \index{reconnaissable
  (langage)}reconnaissable comme un langage $L$ pour lequel il existe
un DFA $A$ tel que $L = L(A)$.  D'après \ref{completion-of-dfai},
\ref{determinization-of-nfa} et \ref{removal-of-epsilon-transitions},
on peut remplacer « DFA » dans cette définition par « DFAi », « NFA »
ou « εNFA » sans changer la définition.

\smallskip

Nous allons maintenant montrer que les langages reconnaissables sont
stables par différentes opérations.  Dans cette section, nous traitons
le cas des opérations booléennes (complémentaire, union, intersection)
et l'opération « miroir » ; la
section \ref{subsection-recognizable-languages-under-rational-operations}
traite des opérations rationnelles.

\begin{prop}\label{dfa-complement}
Si $L$ est un langage reconnaissable sur un alphabet $\Sigma$, alors
le \index{complémentaire (langage)}complémentaire $\Sigma^*\setminus L$ de $L$ est reconnaissable ; de
plus, un DFA reconnaissant l'un se déduit algorithmiquement d'un DFA
reconnaissant l'autre.
\end{prop}
\begin{proof}
Par hypothèse, il existe un DFA (complet !) $A = (Q,q_0,F,\delta)$ tel
que $L = L(A)$.  Considérons le DFA $A'$ défini par l'ensemble d'états
$Q' = Q$, l'état initial $q'_0 = q_0$, la fonction de transition
$\delta' = \delta$ et pour seul changement l'ensemble d'états finaux
$F' = Q \setminus F$ complémentaire de $F$.

Pour $w \in \Sigma^*$, on a $w \in L(A')$ si et seulement si
$\delta^{\prime*}(q_0,w) \in F'$, c'est-à-dire $\delta^*(q_0,w) \in
F'$ (puisque $\delta' = \delta$), c'est-à-dire $\delta^*(q_0,w)
\not\in F$ (par définition du complémentaire), c'est-à-dire $w \not\in
L(A)$.  Ceci montre bien que $L(A')$ est le complémentaire de $L(A)$.
\end{proof}

\thingy Cette démonstration a utilisé la caractérisation des langages
reconnaissables par les DFA : il était crucial de le faire, et les
autres sortes d'automates définis plus haut n'auraient pas permis
d'arriver (simplement) à la même conclusion.  Il est intéressant de
réfléchir à pourquoi.  (Essentiellement, dans un NFA, un mot est
accepté dès qu'\emph{il existe} un chemin qui l'accepte, or
l'existence d'un chemin aboutissant à un état non-final n'est pas la
même chose que l'inexistence d'un chemin aboutissant à un état final.)

\begin{prop}\label{dfa-union-and-intersection}
Si $L_1,L_2$ sont des langages reconnaissables (sur un même
alphabet $\Sigma$), alors la \index{union (de langages)}réunion
$L_1\cup L_2$ et \index{intersection (de langages)}l'intersection
$L_1\cap L_2$ sont reconnaissables ; de plus, un DFA reconnaissant
l'un comme l'autre se déduit algorithmiquement de DFA reconnaissant
$L_1$ et $L_2$.
\end{prop}
\begin{proof}
Traitons le cas de l'intersection.  Par hypothèse, il existe des DFA
(complets !) $A_1 = (Q_1,q_{0,1},F_1,\delta_1)$ et $A_2 =
(Q_2,q_{0,2},F_2,\delta_2)$ tels que $L_1 = L(A_1)$ et $L_2 =
L(A_2)$.  Considérons le DFA $A'$ défini par l'ensemble d'états $Q' =
Q_1 \times Q_2$ (c'est-à-dire l'ensemble des couples formés d'un état
de $A_1$ et d'un état de $A_2$), l'état initial $q'_0 =
(q_{0,1},q_{0,2})$, la fonction de transition $\delta' \colon
((p_1,p_2),x) \mapsto (\delta_1(p_1,x), \delta_2(p_2,x))$ et pour
ensemble d'états finaux $F' = F_1\times F_2$.  Remarquons que
$(p_1,p_2) \in Q'$ appartient à $F' = F_1\times F_2$ si et seulement
si $p_1 \in F_1$ \emph{et} $p_2 \in F_2$ (i.e., $F' \subseteq Q'$ est
l'ensemble des couples dont les deux composantes sont finales).  Par
ailleurs, si $w\in \Sigma^*$, on a $\delta^{\prime*}(q_0',w) =
(\delta_1^*(q_{0,1},w), \delta_2^*(q_{0,2},w))$, et par ce qui vient
d'être dit, ceci appartient à $F'$ si et seulement
$\delta_1^*(q_{0,1},w) \in F_1$ et $\delta_2^*(q_{0,2},w) \in F_2$.
On voit donc qu'un mot $w$ appartient à $L(A')$ si et seulement il
appartient à la fois à $L_1$ et à $L_2$, ce qu'il fallait démontrer.

Pour la réunion, on peut invoquer le fait que la réunion est le
complémentaire de l'intersection des complémentaires, et
utiliser \ref{dfa-complement} ; si on déroule cette démonstration, on
voit qu'on construit un DFA $A''$ exactement égal à $A'$ construit
ci-dessus, à la seule différence près que son ensemble d'états finaux
est $F'' = (F_1\times Q_2) \cup (Q_1\times F_2)$, qui est le
sous-ensemble de $Q'' = Q_1\times Q_2$ formé des couples dont
\emph{l'une au moins} des deux composantes est finale.
\end{proof}

\thingy La construction $A'$ ci-dessus est parfois appelée automate
\index{produit (d'automates)}\emph{produit} des DFA $A_1$ et $A_2$.
Intuitivement, il faut comprendre que faire fonctionner l'automate
$A'$ revient à faire fonctionner les automates $A_1$ et $A_2$
simultanément (en parallèle), en leur donnant à consommer les mêmes
symboles.

La construction de l'automate produit pour fabriquer le langage union
ou intersection utilise la caractérisation des langages
reconnaissables par les DFA ; une construction du même type pourrait
être définie pour les NFA, mais uniquement pour le langage
intersection.  (Pour la réunion effectuée sur les NFA, on renvoie à la
construction présentée en \ref{nfa-union}.)

\begin{prop}\label{nfa-mirror}
Si $L$ est un langage reconnaissable sur un alphabet $\Sigma$, alors
le langage miroir (cf. \ref{definition-mirror-language})
$L^{\mathsf{R}}$ de $L$ est reconnaissable ; de plus, un NFA ou εNFA
reconnaissant l'un se déduit algorithmiquement d'un NFA ou εNFA
reconnaissant l'autre : il s'agit simplement d'inverser le sens de
toutes les flèches (y compris celles qui marquent les états initiaux
et finaux).
\end{prop}

L'automate ainsi construit en inversant toutes les flèches d'un
automate $A$ (la définition précise est donnée dans la démonstration
qui suit) et qui reconnaît le langage miroir de celui reconnu par $A$
peut s'appeler automate \defin[transposé (automate)]{transposé}
$A^{\mathsf{R}}$ de $A$.

\begin{proof}
Par hypothèse, il existe un NFA ou un εNFA $A = (Q,I,F,\delta)$ tel
que $L = L(A)$.  Considérons l'automate $A^{\mathsf{R}}$ de même type
défini par l'ensemble d'états $Q^{\mathsf{R}} = Q$ et inversant toutes
les flèches de $A$, c'est-à-dire $I^{\mathsf{R}} = F$ et
$F^{\mathsf{R}} = I$ et $(q,t,q') \in \delta^{\mathsf{R}}$ si et
seulement si $(q',t,q) \in \delta$.  Un chemin existe dans
$A^{\mathsf{R}}$ si et seulement si le même chemin inversé existe
dans $A$, ce qui montre qu'un mot appartient à $L(A^{\mathsf{R}})$ si
et seulement si son miroir appartient à $L(A)$.  On a donc bien
$L(A^{\mathsf{R}})=L^{\mathsf{R}}$, langage miroir de $L$.
\end{proof}

\thingy Alors que les constructions du complémentaire et de
l'intersection s'effectuaient naturellement sur les DFA, celle du
langage miroir s'effectue naturellement sur les NFA.  (On peut, bien
sûr, considérer un DFA comme un NFA particulier, et effectuer dessus
l'opération d'inversion des flèches qu'on vient de décrire, mais en
général on n'obtiendra pas un DFA, seulement un NFA ; les NFA dont
l'automate transposé est déterministe — c'est-à-dire tels que, pour
chaque état $q$ et chaque lettre $x$, il existe une unique arête
aboutissant à $q$ et étiquetée par $x$ — sont parfois dits
« co-déterministes ».)

\subsection{Stabilité des langages reconnaissables par opérations rationnelles, automates standards, construction de Glushkov}\label{subsection-recognizable-languages-under-rational-operations}

\thingy Nous allons maintenant montrer que la classe des langages
reconnaissables est stable par les opérations rationnelles (union,
concaténation et étoile de Kleene,
cf. \ref{stable-under-rational-operations} ; on l'a déjà vu sur les
DFA en \ref{dfa-union-and-intersection} pour la réunion, mais on va en
donner une nouvelle démonstration, cette fois basée sur les NFA, qui a
un contenu algorithmique utile dans des circonstances différentes).
Cela permettra de conclure
en \ref{rational-languages-are-recognizable} que les langages
rationnels sont reconnaissables (la réciproque faisant l'objet de la
section \ref{subsection-rnfa-and-kleenes-algorithm}).

\smallskip

Pour établir ces stabilités, on va travailler sur les NFA et utiliser
la construction parfois appelée « de Glushkov » ou « automate
standard » ; ceci fournira un « automate de Glushkov » pour chaque
expression rationnelle $r$.  Cette construction travaille, en fait,
sur des NFA vérifiant une propriété supplémentaire facile à assurer,
et on va commencer par un lemme dans ce sens :

\begin{lem}\label{standard-automaton-lemma}
Soit $A$ un NFA.  Alors il existe un NFA $A'$ (sur le même
alphabet $\Sigma$) qui soit équivalent à $A$ et qui possède la
propriété supplémentaire d'avoir un \emph{unique} état initial $q_0$
et qu'aucune transition n'aboutit à $q_0$.  De plus, $A'$ se déduit
algorithmiquement de $A$ et a au plus un état de plus que $A$.

(On pourra appeler \defin[standard (automate)]{standard} un NFA
vérifiant cette propriété d'avoir un unique état initial qui n'est la
cible d'aucune transition.  L'affirmation est donc que tout NFA est
équivalent à un NFA \emph{standard} qui s'en déduit algorithmiquement
par l'ajout d'au plus un état.)
\end{lem}
\begin{proof}
On fabrique $A'$ en reprenant le même ensemble d'états $Q$ que
dans $A$ auquel on ajoute un unique nouvel état $q_0$ qui sera le seul
état initial de $A'$ ; pour chaque transition partant d'un état
initial de $A$, on ajoute dans $A'$ une transition identiquement
étiquetée, et de même cible, partant de $q_0$.

Formellement : soit $A = (Q,I,F,\delta)$.  On définit alors $A' =
(Q',\{q_0\},F',\delta')$ de la manière suivante : $Q' = Q \uplus
\{q_0\}$ (où $\uplus$ désigne une réunion
disjointe\footnote{C'est-à-dire qu'on exige que $q_0$ n'appartienne
  pas à $Q$ (c'est un nouvel état).}), $\delta'$ est la réunion de
l'ensemble des transitions $(q,x,q')$ qui étaient déjà dans $\delta$
et de l'ensemble des $(q_0,x,q')$ telles qu'il existe une transition
$(q,x,q') \in \delta$ avec $q\in I$, et enfin $F' = F$ ou $F' =
F\cup\{q_0\}$ selon que $I\cap F = \varnothing$ ou $I\cap F \neq
\varnothing$ respectivement (i.e., $q_0$ est final lorsqu'il existait
déjà un état à la fois initial et final).

La figure suivante illustre la transformation en question :
\begin{center}
\begin{tikzpicture}[>=latex,line join=bevel,automaton,baseline=(A.base)]
\node (A) at (0bp,0bp) [draw,dotted,circle,minimum size=75bp] {$A$};
\node (q) at (-15bp,15bp) [draw,circle,state,initial] {\footnotesize $q$};
\node (qp) at (15bp,15bp) [draw,circle,state] {\footnotesize\vbox to0pt{\vss\hbox to0pt{$q'$\hss}}\phantom{$q$}};
\draw [->] (A.east) -- ($(A.east)+(3ex,0)$);
\draw [->] (q) -- node[auto] {\footnotesize $x$} (qp);
\end{tikzpicture}
\quad devient\quad
\begin{tikzpicture}[>=latex,line join=bevel,automaton,baseline=(A.base)]
\node (A) at (0bp,0bp) [draw,dotted,circle,minimum size=75bp] {\phantom{$A$}};
\node (q) at (-15bp,15bp) [draw,circle,state] {\footnotesize $q$};
\node (qp) at (15bp,15bp) [draw,circle,state] {\footnotesize\vbox to0pt{\vss\hbox to0pt{$q'$\hss}}\phantom{$q$}};
\node (q0) at (-40bp,0bp) [draw,circle,state,initial,fill=white] {$q_0$};
\draw [->] (A.east) -- ($(A.east)+(3ex,0)$);
\draw [->] (q) -- node[auto] {\footnotesize $x$} (qp);
\draw [->] (q0) to[out=0,in=270] node[auto,swap] {\footnotesize $x$} (qp.south);
\end{tikzpicture}
\end{center}

(\emph{Remarque : } De façon équivalente, on peut fabriquer $A'$ en
ajoutant d'abord un unique état initial $q_0$ et des ε-transitions de
$q_0$ vers chacun des états qui étaient initiaux dans $A$, puis en
éliminant les ε-transitions qu'on vient d'ajouter
(cf. \ref{removal-of-epsilon-transitions}).  Cela donne le même
résultat que ce qui vient d'être dit.)
\end{proof}

\bigbreak

On a vu en \ref{dfa-union-and-intersection} une preuve, à base de DFA,
que $L_1 \cup L_2$ est reconnaissable lorsque $L_1$ et $L_2$ le sont.
Donnons maintenant une autre preuve de ce fait, à base de NFA :

\begin{prop}\label{nfa-union}
Si $L_1,L_2$ sont des langages reconnaissables (sur un même
alphabet $\Sigma$), alors la \index{union (de langages)}réunion $L_1
\cup L_2$ est reconnaissable ; de plus, un NFA la reconnaissant se
déduit algorithmiquement de NFA reconnaissant $L_1$ et $L_2$.
\end{prop}
\begin{proof}
Par hypothèse, il existe des NFA reconnaissant $L_1$ et $L_2$ :
d'après \ref{standard-automaton-lemma}, on peut supposer qu'ils sont
\emph{standards} en ce sens qu'ils ont un unique état initial qui
n'est la cible d'aucune transition.  Disons que $A_1 =
(Q_1,\{q_1\},F_1,\delta_1)$ et $A_2 = (Q_2,\{q_2\},F_2,\delta_2)$ sont
des NFA standards tels que $L_1 = L(A_1)$ et $L_2 = L(A_2)$.

L'automate $A'$ s'obtient réunissant $A_1$ et $A_2$ mais en
« fusionnant » les états initiaux $q_1$ et $q_2$ de $A_1$ et $A_2$ en
un unique état initial $q'_0$, d'où partent les mêmes transitions
(avec les mêmes étiquettes) que depuis l'un ou l'autre de
$q_1$ ou $q_2$.  Graphiquement :

\begin{center}
\begin{tikzpicture}[>=latex,line join=bevel,automaton,baseline=(A.base)]
\node (A) at (0bp,0bp) [draw,dotted,circle,minimum size=60bp] {$A_1$};
\node (q0) at (-35bp,0bp) [draw,circle,state,initial,fill=white] {$q_1$};
\node (i1) at (-10bp,20bp) [circle,inner sep=2bp,fill] {};
\node (i2) at (-10bp,-20bp) [circle,inner sep=2bp,fill] {};
\node (o1) at (15bp,20bp) [circle,inner sep=2bp,fill] {};
\node (o2) at (25bp,0bp) [circle,inner sep=2bp,fill] {};
\node (o3) at (15bp,-20bp) [circle,inner sep=2bp,fill] {};
\draw [->] (q0) -- (i1);
\draw [->] (q0) -- (i2);
\draw [->] (o1) -- ($(o1.east)+(3ex,0)$);
\draw [->] (o2) -- ($(o2.east)+(3ex,0)$);
\draw [->] (o3) -- ($(o3.east)+(3ex,0)$);
\end{tikzpicture}
et
\begin{tikzpicture}[>=latex,line join=bevel,automaton,baseline=(A.base)]
\node (A) at (0bp,0bp) [draw,dotted,circle,minimum size=60bp] {$A_2$};
\node (q0) at (-35bp,0bp) [draw,circle,state,initial,fill=white] {$q_2$};
\node (i1) at (-10bp,20bp) [circle,inner sep=2bp,fill] {};
\node (i2) at (-10bp,-20bp) [circle,inner sep=2bp,fill] {};
\node (o1) at (15bp,20bp) [circle,inner sep=2bp,fill] {};
\node (o2) at (25bp,0bp) [circle,inner sep=2bp,fill] {};
\node (o3) at (15bp,-20bp) [circle,inner sep=2bp,fill] {};
\draw [->] (q0) -- (i1);
\draw [->] (q0) -- (i2);
\draw [->] (o1) -- ($(o1.east)+(3ex,0)$);
\draw [->] (o2) -- ($(o2.east)+(3ex,0)$);
\draw [->] (o3) -- ($(o3.east)+(3ex,0)$);
\end{tikzpicture}
\\deviennent\\
\begin{tikzpicture}[>=latex,line join=bevel,automaton,baseline=(A.base)]
\node (A1) at (0bp,0bp) [draw,dotted,circle,minimum size=60bp] {\phantom{$A_1$}};
\node (q0) at (-35bp,0bp) [draw,circle,state,initial,fill=white] {$q'_0$};
\node (A1) at (90bp,0bp) [draw,dotted,circle,minimum size=60bp] {\phantom{$A_2$}};
\node (i1) at (-10bp,20bp) [circle,inner sep=2bp,fill] {};
\node (i2) at (-10bp,-20bp) [circle,inner sep=2bp,fill] {};
\node (l1) at (15bp,20bp) [circle,inner sep=2bp,fill] {};
\node (l2) at (25bp,0bp) [circle,inner sep=2bp,fill] {};
\node (l3) at (15bp,-20bp) [circle,inner sep=2bp,fill] {};
\node (r1) at (80bp,20bp) [circle,inner sep=2bp,fill] {};
\node (r2) at (80bp,-20bp) [circle,inner sep=2bp,fill] {};
\node (o1) at (105bp,20bp) [circle,inner sep=2bp,fill] {};
\node (o2) at (115bp,0bp) [circle,inner sep=2bp,fill] {};
\node (o3) at (105bp,-20bp) [circle,inner sep=2bp,fill] {};
\draw [->] (q0) -- (i1);
\draw [->] (q0) -- (i2);
\draw [->] (l1) -- ($(l1.east)+(3ex,0)$);
\draw [->] (l2) -- ($(l2.east)+(3ex,0)$);
\draw [->] (l3) -- ($(l3.east)+(3ex,0)$);
\draw [->] (o1) -- ($(o1.east)+(3ex,0)$);
\draw [->] (o2) -- ($(o2.east)+(3ex,0)$);
\draw [->] (o3) -- ($(o3.east)+(3ex,0)$);
\draw [->] (q0) to[out=10,in=180] (r1);
\draw [->] (q0) to[out=350,in=180] (r2);
\end{tikzpicture}
\end{center}

De façon plus formelle, considérons un nouvel ensemble d'états $Q' =
(Q_1 \uplus Q_2) \setminus \{q_1,q_2\} \uplus \{q'_0\}$ où $\uplus$
désigne la réunion disjointe (autrement dit, on prend la réunion
disjointe des états non-initiaux de $A_1$ et $A_2$ et on ajoute un
nouvel état $q'_0$), et la fonction $\varphi_1\colon Q_1 \to Q'$ qui
envoie $q_1$ sur $q'_0$ et tout autre état de $Q_1$ sur lui-même, et
$\varphi_2\colon Q_2 \to Q'$ de façon analogue.  On définit alors
l'automate $A'$ dont l'ensemble d'états est $Q'$, l'état initial
est $q'_0$, les états finaux $F' = \varphi_1(F_1) \cup
\varphi_2(F_2)$, et la relation de transition $\delta'$ est formée des
triplets $(\varphi_1(q),x,q')$ où $(q,x,q') \in \delta_1$ et des
$(\varphi_2(q),x,q')$ où $(q,x,q') \in \delta_2$.

(\emph{Remarque : } De façon équivalente, on peut fabriquer $A'$ en
ajoutant d'abord un unique état initial $q'_0$ à la réunion disjointe
de $A_1$ et $A_2$ et des ε-transitions de $q'_0$ vers $q_1$ et $q_2$
(qui cessent d'être initiaux), puis en éliminant les ε-transitions
qu'on vient d'ajouter (cf. \ref{removal-of-epsilon-transitions}) ainsi
que les états $q_1$ et $q_2$ devenus inutiles.  Cela donne le même
résultat que ce qui vient d'être dit.)

Il est alors clair qu'un chemin de l'état initial à un état final dans
cet automate $A'$ consiste soit en un chemin d'un état initial à un
état final dans $A_1$ soit en un tel chemin dans $A_2$.  On a donc
bien $L(A') = L_1 \cup L_2$.
\end{proof}

\begin{prop}\label{nfa-concatenation}
Si $L_1,L_2$ sont des langages reconnaissables (sur un même
alphabet $\Sigma$), alors la \index{concaténation (de
  langages)}concaténation $L_1 L_2$ est reconnaissable ; de plus, un
NFA la reconnaissant se déduit algorithmiquement de NFA reconnaissant
$L_1$ et $L_2$.
\end{prop}
\begin{proof}
Par hypothèse, il existe des NFA reconnaissant $L_1$ et $L_2$ :
d'après \ref{standard-automaton-lemma}, on peut supposer qu'ils sont
\emph{standards} en ce sens qu'ils ont un unique état initial qui
n'est la cible d'aucune transition.  Disons que $A_1 =
(Q_1,\{q_1\},F_1,\delta_1)$ et $A_2 = (Q_2,\{q_2\},F_2,\delta_2)$ sont
des NFA standards tels que $L_1 = L(A_1)$ et $L_2 = L(A_2)$.

L'automate $A'$ s'obtient en réunissant $A_1$ et $A_2$, en ne gardant
que les états finaux de $A_2$, en supprimant $q_2$ et en remplaçant
chaque transition sortant de $q_2$ par une transition identiquement
étiquetée, et de même cible, partant de \emph{chaque} état final
de $A_1$ (ces derniers seront marqués finaux si $q_2$ était final
dans $A_2$).  Graphiquement :

\begin{center}
\begin{tikzpicture}[>=latex,line join=bevel,automaton,baseline=(A.base)]
\node (A) at (0bp,0bp) [draw,dotted,circle,minimum size=60bp] {$A_1$};
\node (q0) at (-35bp,0bp) [draw,circle,state,initial,fill=white] {$q_1$};
\node (i1) at (-10bp,20bp) [circle,inner sep=2bp,fill] {};
\node (i2) at (-10bp,-20bp) [circle,inner sep=2bp,fill] {};
\node (o1) at (15bp,20bp) [circle,inner sep=2bp,fill] {};
\node (o2) at (25bp,0bp) [circle,inner sep=2bp,fill] {};
\node (o3) at (15bp,-20bp) [circle,inner sep=2bp,fill] {};
\draw [->] (q0) -- (i1);
\draw [->] (q0) -- (i2);
\draw [->] (o1) -- ($(o1.east)+(3ex,0)$);
\draw [->] (o2) -- ($(o2.east)+(3ex,0)$);
\draw [->] (o3) -- ($(o3.east)+(3ex,0)$);
\end{tikzpicture}
et
\begin{tikzpicture}[>=latex,line join=bevel,automaton,baseline=(A.base)]
\node (A) at (0bp,0bp) [draw,dotted,circle,minimum size=60bp] {$A_2$};
\node (q0) at (-35bp,0bp) [draw,circle,state,initial,fill=white] {$q_2$};
\node (i1) at (-10bp,20bp) [circle,inner sep=2bp,fill] {};
\node (i2) at (-10bp,-20bp) [circle,inner sep=2bp,fill] {};
\node (o1) at (15bp,20bp) [circle,inner sep=2bp,fill] {};
\node (o2) at (25bp,0bp) [circle,inner sep=2bp,fill] {};
\node (o3) at (15bp,-20bp) [circle,inner sep=2bp,fill] {};
\draw [->] (q0) -- (i1);
\draw [->] (q0) -- (i2);
\draw [->] (o1) -- ($(o1.east)+(3ex,0)$);
\draw [->] (o2) -- ($(o2.east)+(3ex,0)$);
\draw [->] (o3) -- ($(o3.east)+(3ex,0)$);
\end{tikzpicture}
\\deviennent\\
\begin{tikzpicture}[>=latex,line join=bevel,automaton,baseline=(A.base)]
\node (A1) at (0bp,0bp) [draw,dotted,circle,minimum size=60bp] {\phantom{$A_1$}};
\node (q0) at (-35bp,0bp) [draw,circle,state,initial,fill=white] {$q_1$};
\node (A1) at (90bp,0bp) [draw,dotted,circle,minimum size=60bp] {\phantom{$A_2$}};
\node (i1) at (-10bp,20bp) [circle,inner sep=2bp,fill] {};
\node (i2) at (-10bp,-20bp) [circle,inner sep=2bp,fill] {};
\node (l1) at (15bp,20bp) [circle,inner sep=2bp,fill] {};
\node (l2) at (25bp,0bp) [circle,inner sep=2bp,fill] {};
\node (l3) at (15bp,-20bp) [circle,inner sep=2bp,fill] {};
\node (r1) at (80bp,20bp) [circle,inner sep=2bp,fill] {};
\node (r2) at (80bp,-20bp) [circle,inner sep=2bp,fill] {};
\node (o1) at (105bp,20bp) [circle,inner sep=2bp,fill] {};
\node (o2) at (115bp,0bp) [circle,inner sep=2bp,fill] {};
\node (o3) at (105bp,-20bp) [circle,inner sep=2bp,fill] {};
\draw [->] (q0) -- (i1);
\draw [->] (q0) -- (i2);
\draw [->] (l1) to[out=0,in=180] (r1);
\draw [->] (l1) to[out=0,in=180] (r2);
\draw [->] (l2) to[out=0,in=180] (r1);
\draw [->] (l2) to[out=0,in=180] (r2);
\draw [->] (l3) to[out=0,in=180] (r1);
\draw [->] (l3) to[out=0,in=180] (r2);
\draw [->] (o1) -- ($(o1.east)+(3ex,0)$);
\draw [->] (o2) -- ($(o2.east)+(3ex,0)$);
\draw [->] (o3) -- ($(o3.east)+(3ex,0)$);
\end{tikzpicture}
\end{center}

De façon plus formelle, considérons un nouvel ensemble d'états $Q' =
(Q_1 \uplus Q_2) \setminus \{q_2\}$ où $\uplus$ désigne la réunion
disjointe.  On définit alors l'automate $A'$ dont l'ensemble d'états
est $Q'$, l'état initial est $q_1$, l'ensemble $F'$ des états finaux
est $F_2$ si $q_2$ n'était pas final dans $A_2$ et $F_1 \cup F_2$ si
$q_2$ était final dans $A_2$, la relation de transition $\delta'$
est la réunion de $\delta_1$, de l'ensemble des triplets $(q,x,q') \in
\delta_2$ tels que $q\neq q_2$, et enfin de l'ensemble des triplets
$(q,x,q')$ tels que $(q_2,x,q') \in \delta_2$ et que $q\in F_1$.

(\emph{Remarque : } De façon équivalente, on peut fabriquer $A'$ en
ajoutant d'abord à la réunion disjointe de $A_1$ et $A_2$ une
ε-transition de chaque état final de $A_1$ vers $q_2$ (qui cesse
d'être initial), puis en éliminant les ε-transitions qu'on vient
d'ajouter (cf. \ref{removal-of-epsilon-transitions}) ainsi que l'état
$q_2$ devenu inutile.  Cela donne le même résultat que ce qui vient
d'être dit.)

Il est alors clair qu'un chemin de l'état initial $q_1$ à un état
final dans cet automate $A'$ consiste en un chemin de $q_1$ à un état
final dans $A_1$ suivi d'un chemin de $q_2$ à un état final dans $A_2$
(moins $q_2$ lui-même).  On a donc bien $L(A') = L_1 L_2$.
\end{proof}

\begin{prop}\label{nfa-star}
Si $L$ est un langage reconnaissable (sur un alphabet $\Sigma$), alors
\index{étoile de Kleene}l'étoile de Kleene $L^*$ est reconnaissable ;
de plus, un NFA la reconnaissant se déduit algorithmiquement de NFA
reconnaissant $L$.
\end{prop}
\begin{proof}
Par hypothèse, il existe un NFA reconnaissant $L$ :
d'après \ref{standard-automaton-lemma}, on peut supposer qu'ils est
\emph{standard} en ce sens qu'il a un unique état initial qui n'est la
cible d'aucune transition.  Disons que $A = (Q,\{q_0\},F,\delta)$ est
un NFA standard tel que $L = L(A)$.

L'automate $A'$ s'obtient en ajoutant à $A$, pour chaque transition
sortant de $q_0$, une transition identiquement étiquetée, et de même
cible, partant de chaque état final de $A$, et en rendant $q_0$ final
s'il ne l'était pas déjà :

\begin{center}
\begin{tikzpicture}[>=latex,line join=bevel,automaton,baseline=(A.base)]
\node (A) at (0bp,0bp) [draw,dotted,circle,minimum size=60bp] {$A$};
\node (q0) at (-35bp,0bp) [draw,circle,state,initial,fill=white] {$q_0$};
\node (i1) at (-10bp,20bp) [circle,inner sep=2bp,fill] {};
\node (i2) at (-10bp,-20bp) [circle,inner sep=2bp,fill] {};
\node (o1) at (15bp,20bp) [circle,inner sep=2bp,fill] {};
\node (o2) at (25bp,0bp) [circle,inner sep=2bp,fill] {};
\node (o3) at (15bp,-20bp) [circle,inner sep=2bp,fill] {};
\draw [->] (q0) -- (i1);
\draw [->] (q0) -- (i2);
\draw [->] (o1) -- ($(o1.east)+(3ex,0)$);
\draw [->] (o2) -- ($(o2.east)+(3ex,0)$);
\draw [->] (o3) -- ($(o3.east)+(3ex,0)$);
\end{tikzpicture}
devient
\begin{tikzpicture}[>=latex,line join=bevel,automaton,baseline=(A.base)]
\node (A) at (0bp,0bp) [draw,dotted,circle,minimum size=60bp] {\phantom{$A$}};
\node (q0) at (-35bp,0bp) [draw,circle,state,initial,accepting below,fill=white] {$q_0$};
\node (i1) at (-10bp,20bp) [circle,inner sep=2bp,fill] {};
\node (i2) at (-10bp,-20bp) [circle,inner sep=2bp,fill] {};
\node (o1) at (15bp,20bp) [circle,inner sep=2bp,fill] {};
\node (o2) at (25bp,0bp) [circle,inner sep=2bp,fill] {};
\node (o3) at (15bp,-20bp) [circle,inner sep=2bp,fill] {};
\draw [->] (q0) -- (i1);
\draw [->] (q0) -- (i2);
\draw [->] (o1) -- ($(o1.east)+(3ex,0)$);
\draw [->] (o2) -- ($(o2.east)+(3ex,0)$);
\draw [->] (o3) -- ($(o3.east)+(3ex,0)$);
\draw [->] (o1) -- (i1);
\draw [->] (o1) -- (i2);
\draw [->] (o2) -- (i1);
\draw [->] (o2) -- (i2);
\draw [->] (o3) -- (i1);
\draw [->] (o3) -- (i2);
\end{tikzpicture}
\end{center}

De façon plus formelle, on considère l'automate $A'$ dont l'ensemble
d'états est $Q' := Q$, l'état initial est $q_0$, les états finaux $F'
:= F \cup \{q_0\}$, et la relation de transition $\delta'$ est la
réunion de $\delta$ et de l'ensemble des triplets $(q,x,q')$ tels que
$(q_0,x,q') \in \delta$ et que $q \in F$.

(\emph{Remarque : } De façon équivalente, on peut fabriquer $A'$ en
ajoutant d'abord à $A$ une ε-transition de chaque état final de $A$
vers $q_0$, puis en éliminant les ε-transitions qu'on vient d'ajouter
(cf. \ref{removal-of-epsilon-transitions}), et enfin en marquant $q_0$
comme final.  Cela donne le même résultat que ce qui vient d'être
dit.)

Il est alors clair qu'un chemin de l'état initial $q_0$ à un état
final dans cet automate $A'$ consiste en un nombre quelconque
(éventuellement nul) de chemins de l'état initial à un état final
dans $A$.  On a donc bien $L(A') = L^*$.
\end{proof}

\thingy\label{trivial-standard-automata} Il sera utile de fixer
également des NFA (« standards » au sens de
\ref{standard-automaton-lemma}) reconnaissant les langages de base
triviaux $\varnothing$, $\{\varepsilon\}$ et $\{x\}$ (pour
chaque $x\in\Sigma$), c'est-à-dire ceux dénotés par les expressions
rationnelles $\bot$, $\underline{\varepsilon}$ et $x$ respectivement.
On prendra les suivants :

\begin{center}
\begin{tabular}{ll}
\begin{tikzpicture}[>=latex,line join=bevel,automaton,baseline=(q0.base)]
\node (q0) at (0bp,0bp) [draw,circle,state,initial] {$q_0$};
\end{tikzpicture}
&pour le langage $\varnothing$ (i.e., pour l'expression rationnelle $\bot$),\\[1.75ex]
\begin{tikzpicture}[>=latex,line join=bevel,automaton,baseline=(q0.base)]
\node (q0) at (0bp,0bp) [draw,circle,state,initial,final] {$q_0$};
\end{tikzpicture}
&pour le langage $\{\varepsilon\}$ (i.e., pour l'expression
rationnelle $\underline{\varepsilon}$), et\\[1.75ex]
\begin{tikzpicture}[>=latex,line join=bevel,automaton,baseline=(q0.base)]
\node (q0) at (0bp,0bp) [draw,circle,state,initial] {$q_0$};
\node (q1) at (60bp,0bp) [draw,circle,state,final] {$q_1$};
\draw [->] (q0) to node[auto,swap] {$x$} (q1);
\end{tikzpicture}
&pour le langage $\{x\}$ (i.e., pour l'expression rationnelle $x$).
\end{tabular}
\end{center}

\bigbreak

Récapitulons le contenu essentiel de ce que nous avons montré :

\begin{cor}\label{rational-languages-are-recognizable}
Tout langage rationnel est reconnaissable ; de plus, un NFA le
reconnaissant se déduit algorithmiquement d'une expression rationnelle
le dénotant.

En particulier, il existe un algorithme qui, donnée une expression
rationnelle $r$ (sur un alphabet $\Sigma$) et un mot $w \in \Sigma^*$,
décide si $w\in L(r)$, c'est-à-dire si $w$ vérifie $r$.  (Autrement
dit, les langages rationnels sont \emph{décidables} au sens
de \ref{definition-computable-function-or-set}.)
\end{cor}
\begin{proof}
L'affirmation du premier paragraphe résulte de façon évidente de la
définition des langages rationnels
(cf. §\ref{subsection-rational-languages}), du fait que les langages
$\varnothing$, $\{\varepsilon\}$ et $\{x\}$ (pour chaque $x\in\Sigma$)
sont reconnaissables par automates finis
(cf. \ref{trivial-standard-automata}), et grâce aux propositions
\ref{nfa-union}, \ref{nfa-concatenation} et \ref{nfa-star}.  On
renvoie à \ref{glushkov-construction} ci-dessous (ou à
§\ref{subsection-thompson-construction}) pour une description
algorithmique plus précise.

Pour décider si un mot $w$ vérifie une expression rationnelle $r$, on
commence par transformer cette expression rationnelle en NFA comme on
vient de l'expliquer (c'est-à-dire construire un NFA $A(r)$
reconnaissant le langage $L(r)$ dénoté par $r$), puis on déterminise
cet automate (cf. \ref{determinization-of-nfa}), après quoi il est
facile de tester si le DFA résultant de la déterminisation accepte le
mot $w$ considéré (il suffit de suivre l'unique chemin dans l'automate
partant de l'état initial et étiqueté par $w$, et de voir si l'état
auquel il aboutit est final).
\end{proof}

\thingy\label{glushkov-construction} Les constructions que nous avons
décrites dans cette section associent naturellement un NFA standard à
chaque expression rationnelle : il s'obtient en partant des automates
de base décrits en \ref{trivial-standard-automata} et en appliquant
les constructions décrites dans les démonstrations de \ref{nfa-union},
\ref{nfa-concatenation} et \ref{nfa-star}.

\medskip

Plus exactement, on associe à chaque expression rationnelle $r$ (sur
un alphabet $\Sigma$ fixé) un automate $A(r)$ (ou
$A_{\textrm{Glushkov}}(r)$) standard, appelé \defin[Glushkov
  (construction d'automate de)]{automate de Glushkov}, qui reconnaît
le langage $L(r)$ dénoté par $r$, de la manière suivante (par
induction sur la complexité de l'expression rationnelle telle que
définie en \ref{regular-expressions}) :
\begin{itemize}
\item les automates de Glushkov de $\bot$, $\underline{\varnothing}$
  et $x$ (pour $x\in\Sigma$) sont définis comme ceux décrits et
  illustrés en \ref{trivial-standard-automata},
\item connaissant les automates de Glushkov $A_1$ et $A_2$ de $r_1$ et
  $r_2$ respectivement, celui de $(r_1|r_2)$ est celui décrit et
  illustré dans la démonstration de \ref{nfa-union},
\item connaissant les automates de Glushkov $A_1$ et $A_2$ de $r_1$ et
  $r_2$ respectivement, celui de $r_1 r_2$ est celui décrit et
  illustré dans la
  démonstration de \ref{nfa-concatenation},
\item connaissant l'automate de Glushkov $A$ de $r$, celui de $(r){*}$
  est celui décrit et illustré dans la démonstration
  de \ref{nfa-star}.
\end{itemize}

\smallskip

Cet automate de Glushkov $A_{\textrm{Glushkov}}(r)$ possède les
propriétés suivantes :
\begin{itemize}
\item c'est un NFA reconnaissant le langage $L(r)$ dénoté par
  l'expression rationnelle $r$ dont on est parti,
\item il est standard au sens de \ref{standard-automaton-lemma},
  c'est-à-dire qu'il possède un unique état initial auquel n'aboutit
  aucune transition,
\item les transitions aboutissant à n'importe quel état donné sont
  toutes étiquetées par la même lettre,
\item son nombre d'états est égal à $1$ plus le nombre de lettres (à
  l'exclusion des métacaractères) contenues dans l'expression $r$.
\end{itemize}

Les deux dernières propriétés se vérifient inductivement, c'est-à-dire
qu'on observe qu'elles sont satisfaites sur les automates de base
décrits en \ref{trivial-standard-automata} qu'elles sont préservées
par les constructions de \ref{nfa-union}, \ref{nfa-concatenation}
et \ref{nfa-star}.  En fait, on peut être un peu plus précis : chaque
état, autre que l'état initial, de l'automate de Glushkov associé à
l'expression rationnelle $r$ correspond à une lettre $x$ (à
l'exclusion des métacaractères) de $r$, l'état en question provient de
l'état $q_1$ de l'automate décrit en \ref{trivial-standard-automata}
pour le langage $\{x\}$, et toutes les transitions menant à cet état
sont étiquetées par $x$.

Ces observations sont utiles pour détecter des erreurs lors de la
construction de l'automate.

\thingy À titre d'exemple pour illustrer la construction de Glushkov,
construisons l'automate qu'elle associe à l'expression rationnelle
$((a{*}|b)c){*}$ sur l'alphabet $\Sigma = \{a,b,c\}$.  On doit obtenir
un automate ayant $4$ états.

L'automate de Glushkov de $a{*}$ est le suivant, obtenu en appliquant
la construction de \ref{nfa-star} à l'automate trivial pour le
langage $\{a\}$ (donné en \ref{trivial-standard-automata}) :

\begin{center}
%%% begin example8a %%%

\begin{tikzpicture}[>=latex,line join=bevel,automaton]
%%
\node (q1) at (97bp,18bp) [draw,circle,state,final] {$1$};
  \node (q0) at (18bp,18bp) [draw,circle,state,initial,final,accepting below] {$0$};
  \draw [->] (q0) ..controls (45.659bp,18bp) and (57.817bp,18bp)  .. node[auto] {$a$} (q1);
  \draw [->] (q1) to[loop above] node[auto] {$a$} (q1);
%
\end{tikzpicture}

%%% end example8a %%%
\end{center}

On en déduit au moyen de \ref{nfa-union} l'automate suivant
pour $a{*}|b$ :

\begin{center}
%%% begin example8b %%%

\begin{tikzpicture}[>=latex,line join=bevel,automaton]
%%
\node (q1) at (97bp,72bp) [draw,circle,state,final] {$1$};
  \node (q0) at (18bp,41bp) [draw,circle,state,initial,final,accepting below] {$0$};
  \node (q2) at (97bp,18bp) [draw,circle,state,final] {$2$};
  \draw [->] (q0) ..controls (45.279bp,51.58bp) and (58.943bp,57.081bp)  .. node[auto] {$a$} (q1);
  \draw [->] (q0) ..controls (40.761bp,31.564bp) and (47.615bp,28.946bp)  .. (54bp,27bp) .. controls (58.828bp,25.529bp) and (64.04bp,24.206bp)  .. node[auto] {$b$} (q2);
  \draw [->] (q1) to[loop above] node[auto] {$a$} (q1);
%
\end{tikzpicture}

%%% end example8b %%%
\end{center}

On en déduit au moyen de \ref{nfa-concatenation} l'automate suivant
pour $(a{*}|b)c$ :

\begin{center}
%%% begin example8c %%%

\begin{tikzpicture}[>=latex,line join=bevel,automaton]
%%
\node (q1) at (97bp,99.608bp) [draw,circle,state] {$1$};
  \node (q0) at (18bp,45.608bp) [draw,circle,state,initial] {$0$};
  \node (q3) at (176bp,45.608bp) [draw,circle,state,final] {$3$};
  \node (q2) at (97bp,45.608bp) [draw,circle,state] {$2$};
  \draw [->] (q0) ..controls (42.871bp,23.135bp) and (60.567bp,9.4181bp)  .. (79bp,3.6077bp) .. controls (94.26bp,-1.2026bp) and (99.74bp,-1.2026bp)  .. (115bp,3.6077bp) .. controls (129.69bp,8.2379bp) and (143.91bp,17.889bp)  .. node[auto] {$c$} (q3);
  \draw [->] (q2) ..controls (124.66bp,45.608bp) and (136.82bp,45.608bp)  .. node[auto] {$c$} (q3);
  \draw [->] (q0) ..controls (45.659bp,45.608bp) and (57.817bp,45.608bp)  .. node[auto] {$b$} (q2);
  \draw [->] (q1) to[loop above] node[auto] {$a$} (q1);
  \draw [->] (q0) ..controls (44.341bp,63.379bp) and (60.249bp,74.535bp)  .. node[auto] {$a$} (q1);
  \draw [->] (q1) ..controls (123.34bp,81.836bp) and (139.25bp,70.68bp)  .. node[auto] {$c$} (q3);
%
\end{tikzpicture}

%%% end example8c %%%
\end{center}

Et enfin, de nouveau par \ref{nfa-star} l'automate suivant
pour $((a{*}|b)c)*$ :

\begin{center}
%%% begin example8d %%%

\begin{tikzpicture}[>=latex,line join=bevel,automaton]
%%
\begin{scope}
  \pgfsetstrokecolor{black}
  \definecolor{strokecol}{rgb}{1.0,1.0,1.0};
  \pgfsetstrokecolor{strokecol}
  \definecolor{fillcol}{rgb}{1.0,1.0,1.0};
  \pgfsetfillcolor{fillcol}
\end{scope}
  \node (q1) at (97bp,101.45bp) [draw,circle,state] {$1$};
  \node (q0) at (18bp,45.453bp) [draw,circle,state,initial,final,accepting below] {$0$};
  \node (q3) at (176bp,63.453bp) [draw,circle,state,final] {$3$};
  \node (q2) at (97bp,45.453bp) [draw,circle,state] {$2$};
  \draw [->] (q0) ..controls (49.576bp,16.777bp) and (84.747bp,-9.0778bp)  .. (115bp,3.4533bp) .. controls (133.15bp,10.97bp) and (148.67bp,26.932bp)  .. node[auto,below,near start] {$c$} (q3);
  \draw [->] (q2) ..controls (124.6bp,51.671bp) and (137.34bp,54.651bp)  .. node[auto] {$c$} (q3);
  \draw [->] (q0) ..controls (45.659bp,45.453bp) and (57.817bp,45.453bp)  .. node[auto] {$b$} (q2);
  \draw [->] (q3) ..controls (148.91bp,76.331bp) and (134.76bp,83.311bp)  .. node[auto] {$a$} (q1);
  \draw [->] (q1) to[loop above] node[auto] {$a$} (q1);
  \draw [->] (q3) ..controls (156.78bp,44.679bp) and (148.66bp,38.493bp)  .. (140bp,35.453bp) .. controls (134.61bp,33.562bp) and (128.7bp,33.684bp)  .. node[auto] {$b$} (q2);
  \draw [->] (q3) to[loop above] node[auto] {$c$} (q3);
  \draw [->] (q0) ..controls (44.52bp,64.013bp) and (60.758bp,75.823bp)  .. node[auto] {$a$} (q1);
  \draw [->] (q1) ..controls (120.07bp,117.96bp) and (130.95bp,122.37bp)  .. (140bp,117.45bp) .. controls (151.16bp,111.39bp) and (159.31bp,100.06bp)  .. node[auto,above] {$c$} (q3);
%
\end{tikzpicture}

%%% end example8d %%%
\end{center}

Toutes les transitions aboutissant à l'état $1$ sont étiquetées $a$,
toutes celles aboutissant à l'état $2$ sont étiquetées $b$, et toutes
celles aboutissant à $3$ sont étiquetées $c$.


\subsection{L'automate de Thompson (alternative à l'automate de Glushkov)}\label{subsection-thompson-construction}

\thingy La construction de Glushkov (exposée
en \ref{glushkov-construction}) d'un automate reconnaissant le langage
dénoté par expression rationnelle $r$ fabrique un NFA.  Cette
construction produit un automate raisonnablement compact (en nombre
d'états), mais il peut être intéressant de disposer d'une autre
construction, plus transparente mais moins efficace : la
\defin[Thompson (construction d'automate de)]{construction de
  Thompson} fournit un autre moyen d'associer à une expression
rationnelle $r$ un automate $A_{\textrm{Thompson}}(r)$ reconnaissant
le langage qu'elle dénote.  Elle possède pour sa part les propriétés
suivantes :
\begin{itemize}
\item c'est un εNFA reconnaissant le langage $L(r)$ dénoté par
  l'expression rationnelle $r$ dont on est parti,
\item il possède un unique état initial auquel n'aboutit aucune
  transition, et un unique état final duquel ne part aucune
  transition,
\item son nombre d'états est égal au double du nombre de symboles
  autres que les parenthèses constituant l'expression $r$ (en comptant
  aussi bien les lettres de $\Sigma$ que les métacaractères $\bot$,
  $\underline{\varepsilon}$, $|$ et $*$ ; mais sans compter la
  concaténation implicite).
\end{itemize}

\smallskip

Dans les dessins qui suivent, on symbolisera de la manière suivante un
automate de Thompson $A$ quelconque :
\begin{center}
\begin{tikzpicture}[>=latex,line join=bevel,automaton,baseline=(A.base)]
\node (A) at (30bp,0bp) [draw,dotted,circle,minimum size=50bp] {$A$};
\node (qi) at (0bp,0bp) [draw,circle,state,initial,fill=white] {$q_0$};
\node (qf) at (60bp,0bp) [draw,circle,state,final,fill=white] {\vbox to0pt{\vss\hbox to0pt{$q_\infty$\hss}}\phantom{$q_0$}};
\end{tikzpicture}
\end{center}

\thingy\label{trivial-thompson-automata} Les automates de Thompson des
expressions rationnelles $\bot$, $\underline{\varepsilon}$ et $x$
(pour $x\in\Sigma$) seront les suivants :

\begin{center}
\begin{tabular}{ll}
\begin{tikzpicture}[>=latex,line join=bevel,automaton,baseline=(qi.base)]
\node (qi) at (0bp,0bp) [draw,circle,state,initial] {$q_0$};
\node (qf) at (60bp,0bp) [draw,circle,state,final] {\vbox to0pt{\vss\hbox to0pt{$q_\infty$\hss}}\phantom{$q_0$}};
\end{tikzpicture}
&pour le langage $\varnothing$ (i.e., pour l'expression rationnelle $\bot$),\\
\begin{tikzpicture}[>=latex,line join=bevel,automaton,baseline=(qi.base)]
\node (qi) at (0bp,0bp) [draw,circle,state,initial] {$q_0$};
\node (qf) at (60bp,0bp) [draw,circle,state,final] {\vbox to0pt{\vss\hbox to0pt{$q_\infty$\hss}}\phantom{$q_0$}};
\draw [->] (qi) to node[auto] {$\varepsilon$} (qf);
\end{tikzpicture}
&pour le langage $\{\varepsilon\}$ (i.e., pour l'expression
rationnelle $\underline{\varepsilon}$), et\\
\begin{tikzpicture}[>=latex,line join=bevel,automaton,baseline=(qi.base)]
\node (qi) at (0bp,0bp) [draw,circle,state,initial] {$q_0$};
\node (qf) at (60bp,0bp) [draw,circle,state,final] {\vbox to0pt{\vss\hbox to0pt{$q_\infty$\hss}}\phantom{$q_0$}};
\draw [->] (qi) to node[auto,swap] {$x$} (qf);
\end{tikzpicture}
&pour le langage $\{x\}$ (i.e., pour l'expression rationnelle $x$).
\end{tabular}
\end{center}

\thingy\label{thompson-union} Si $A_1$ et $A_2$ sont les automates de
Thompson pour les expressions rationnelles $r_1$ et $r_2$, celui de
$(r_1|r_2)$ sera construit de la manière suivante :
\begin{center}
\begin{tikzpicture}[>=latex,line join=bevel,automaton,baseline=(A.base)]
\node (A) at (30bp,0bp) [draw,dotted,circle,minimum size=50bp] {$A_1$};
\node (qi) at (0bp,0bp) [draw,circle,state,initial,fill=white] {$q_0$};
\node (qf) at (60bp,0bp) [draw,circle,state,final,fill=white] {\vbox to0pt{\vss\hbox to0pt{$q_\infty$\hss}}\phantom{$q_0$}};
\end{tikzpicture}
et
\begin{tikzpicture}[>=latex,line join=bevel,automaton,baseline=(A.base)]
\node (A) at (30bp,0bp) [draw,dotted,circle,minimum size=50bp] {$A_2$};
\node (qi) at (0bp,0bp) [draw,circle,state,initial,fill=white] {$q_0$};
\node (qf) at (60bp,0bp) [draw,circle,state,final,fill=white] {\vbox to0pt{\vss\hbox to0pt{$q_\infty$\hss}}\phantom{$q_0$}};
\end{tikzpicture}
\\deviennent\\
\begin{tikzpicture}[>=latex,line join=bevel,automaton,baseline=(qi.base)]
\node (qi) at (-35bp,0bp) [draw,circle,state,initial,fill=white] {$q_0$};
\node (qf) at (95bp,0bp) [draw,circle,state,final,fill=white] {\vbox to0pt{\vss\hbox to0pt{$q_\infty$\hss}}\phantom{$q_0$}};
\node (A1) at (30bp,35bp) [draw,dotted,circle,minimum size=50bp] {$A_1$};
\node (qi1) at (0bp,35bp) [draw,circle,state,fill=white] {\phantom{$q_0$}};
\node (qf1) at (60bp,35bp) [draw,circle,state,fill=white] {\phantom{$q_0$}};
\node (A2) at (30bp,-35bp) [draw,dotted,circle,minimum size=50bp] {$A_2$};
\node (qi2) at (0bp,-35bp) [draw,circle,state,fill=white] {\phantom{$q_0$}};
\node (qf2) at (60bp,-35bp) [draw,circle,state,fill=white] {\phantom{$q_0$}};
\draw[->] (qi) to node[auto] {$\varepsilon$} (qi1);  \draw[->] (qi) to node[auto] {$\varepsilon$} (qi2);
\draw[->] (qf1) to node[auto] {$\varepsilon$} (qf);  \draw[->] (qf2) to node[auto] {$\varepsilon$} (qf);
\end{tikzpicture}
\end{center}

\thingy\label{thompson-concatenation} Si $A_1$ et $A_2$ sont les
automates de Thompson pour les expressions rationnelles $r_1$ et
$r_2$, celui de $r_1 r_2$ sera construit de la manière suivante :
\begin{center}
\begin{tikzpicture}[>=latex,line join=bevel,automaton,baseline=(A.base)]
\node (A) at (30bp,0bp) [draw,dotted,circle,minimum size=50bp] {$A_1$};
\node (qi) at (0bp,0bp) [draw,circle,state,initial,fill=white] {$q_0$};
\node (qf) at (60bp,0bp) [draw,circle,state,final,fill=white] {\vbox to0pt{\vss\hbox to0pt{$q_\infty$\hss}}\phantom{$q_0$}};
\end{tikzpicture}
et
\begin{tikzpicture}[>=latex,line join=bevel,automaton,baseline=(A.base)]
\node (A) at (30bp,0bp) [draw,dotted,circle,minimum size=50bp] {$A_2$};
\node (qi) at (0bp,0bp) [draw,circle,state,initial,fill=white] {$q_0$};
\node (qf) at (60bp,0bp) [draw,circle,state,final,fill=white] {\vbox to0pt{\vss\hbox to0pt{$q_\infty$\hss}}\phantom{$q_0$}};
\end{tikzpicture}
\\deviennent\\
\begin{tikzpicture}[>=latex,line join=bevel,automaton,baseline=(qi.base)]
\node (A1) at (30bp,0bp) [draw,dotted,circle,minimum size=50bp] {$A_1$};
\node (qi) at (0bp,0bp) [draw,circle,state,initial,fill=white] {$q_0$};
\node (ql) at (60bp,0bp) [draw,circle,state,fill=white] {\phantom{$q_0$}};
\node (A2) at (150bp,0bp) [draw,dotted,circle,minimum size=50bp] {$A_2$};
\node (qr) at (120bp,0bp) [draw,circle,state,fill=white] {\phantom{$q_0$}};
\node (qf) at (180bp,0bp) [draw,circle,state,final,fill=white] {\vbox to0pt{\vss\hbox to0pt{$q_\infty$\hss}}\phantom{$q_0$}};
\draw[->] (ql) to node[auto] {$\varepsilon$} (qr);
\end{tikzpicture}
\end{center}

\thingy\label{thompson-star} Si $A$ est l'automate de Thompson pour
l'expression rationnelle $r$, celui de $(r){*}$ sera construit de la
manière suivante :
\begin{center}
\begin{tikzpicture}[>=latex,line join=bevel,automaton,baseline=(A.base)]
\node (A) at (30bp,0bp) [draw,dotted,circle,minimum size=50bp] {$A$};
\node (qi) at (0bp,0bp) [draw,circle,state,initial,fill=white] {$q_0$};
\node (qf) at (60bp,0bp) [draw,circle,state,final,fill=white] {\vbox to0pt{\vss\hbox to0pt{$q_\infty$\hss}}\phantom{$q_0$}};
\end{tikzpicture}
devient
\begin{tikzpicture}[>=latex,line join=bevel,automaton,baseline=(A.base)]
\node (qi) at (0bp,0bp) [draw,circle,state,initial,fill=white] {$q_0$};
\node (qf) at (180bp,0bp) [draw,circle,state,final,fill=white] {\vbox to0pt{\vss\hbox to0pt{$q_\infty$\hss}}\phantom{$q_0$}};
\node (A) at (90bp,0bp) [draw,dotted,circle,minimum size=50bp] {$A$};
\node (qai) at (60bp,0bp) [draw,circle,state,fill=white] {\phantom{$q_0$}};
\node (qaf) at (120bp,0bp) [draw,circle,state,fill=white] {\phantom{$q_0$}};
\draw[->] (qi) to node[auto] {$\varepsilon$} (qai);
\draw[->] (qaf) to node[auto] {$\varepsilon$} (qf);
\draw[->] (qi) ..controls (60bp,-60bp) and (120bp,-60bp)  .. node[auto] {$\varepsilon$} (qf);
\draw[->] (qaf) ..controls (120bp,60bp) and (60bp,60bp)  .. node[auto,above] {$\varepsilon$} (qai);
\end{tikzpicture}
\end{center}

\thingy Comme on le voit ci-dessus, la construction de Thompson est
très simple à appliquer ; mais elle conduit à des automates rapidement
énormes, comportant un nombre considérable d'états et de transitions
spontanées « stupides ».

\medskip

À titre d'exemple, voici l'automate de Thompson, déjà gros, de
l'expression rationnelle $(a|b){*}b$ :

\begin{center}
\scalebox{0.70}{%
%%% begin example9 %%%

\begin{tikzpicture}[>=latex,line join=bevel,automaton]
%%
\node (q1) at (97bp,61bp) [draw,circle,state] {$1$};
  \node (q0) at (18bp,23bp) [draw,circle,state,initial] {$0$};
  \node (q3) at (255bp,138bp) [draw,circle,state] {$3$};
  \node (q2) at (176bp,138bp) [draw,circle,state] {$2$};
  \node (q5) at (255bp,84bp) [draw,circle,state] {$5$};
  \node (q4) at (176bp,84bp) [draw,circle,state] {$4$};
  \node (q7) at (413bp,23bp) [draw,circle,state] {$7$};
  \node (q6) at (334bp,61bp) [draw,circle,state] {$6$};
  \node (q9) at (571bp,23bp) [draw,circle,state,final] {$9$};
  \node (q8) at (492bp,23bp) [draw,circle,state] {$8$};
  \draw [->] (q3) ..controls (280.5bp,113.5bp) and (299.16bp,94.836bp)  .. node[auto] {$\varepsilon$} (q6);
  \draw [->] (q2) ..controls (203.66bp,138bp) and (215.82bp,138bp)  .. node[auto] {$a$} (q3);
  \draw [->] (q6) ..controls (303.95bp,58.621bp) and (287.5bp,57.483bp)  .. (273bp,57bp) .. controls (221.92bp,55.299bp) and (209.08bp,55.299bp)  .. (158bp,57bp) .. controls (147.24bp,57.358bp) and (135.4bp,58.078bp)  .. node[auto] {$\varepsilon$} (q1);
  \draw [->] (q8) ..controls (519.66bp,23bp) and (531.82bp,23bp)  .. node[auto] {$b$} (q9);
  \draw [->] (q7) ..controls (440.66bp,23bp) and (452.82bp,23bp)  .. node[auto] {$\varepsilon$} (q8);
  \draw [->] (q5) ..controls (282.27bp,76.154bp) and (295.19bp,72.293bp)  .. node[auto] {$\varepsilon$} (q6);
  \draw [->] (q6) ..controls (361.12bp,48.108bp) and (375.27bp,41.127bp)  .. node[auto] {$\varepsilon$} (q7);
  \draw [->] (q4) ..controls (203.66bp,84bp) and (215.82bp,84bp)  .. node[auto] {$b$} (q5);
  \draw [->] (q1) ..controls (124.27bp,68.846bp) and (137.19bp,72.707bp)  .. node[auto] {$\varepsilon$} (q4);
  \draw [->] (q0) ..controls (49.792bp,8.7612bp) and (73.956bp,0bp)  .. (96bp,0bp) .. controls (96bp,0bp) and (96bp,0bp)  .. (335bp,0bp) .. controls (352.91bp,0bp) and (372.22bp,5.7837bp)  .. node[auto] {$\varepsilon$} (q7);
  \draw [->] (q0) ..controls (45.123bp,35.892bp) and (59.268bp,42.873bp)  .. node[auto] {$\varepsilon$} (q1);
  \draw [->] (q1) ..controls (122.5bp,85.495bp) and (141.16bp,104.16bp)  .. node[auto] {$\varepsilon$} (q2);
%
\end{tikzpicture}

%%% end example9 %%%
}
\end{center}

(Il a $10$ états puisqu'il y a $5$ symboles autres que les parenthèses
dans $(a|b){*}b$.)

Pour comparaison, voici son automate de Glushkov :

\begin{center}
\scalebox{0.85}{%
%%% begin example9b %%%

\begin{tikzpicture}[>=latex,line join=bevel,automaton]
%%
\begin{scope}
  \pgfsetstrokecolor{black}
  \definecolor{strokecol}{rgb}{1.0,1.0,1.0};
  \pgfsetstrokecolor{strokecol}
  \definecolor{fillcol}{rgb}{1.0,1.0,1.0};
  \pgfsetfillcolor{fillcol}
\end{scope}
  \node (q9) at (176bp,45.608bp) [draw,circle,state,final] {$9$};
  \node (q0) at (18bp,45.608bp) [draw,circle,state,initial] {$0$};
  \node (q3) at (97bp,150.61bp) [draw,circle,state] {$3$};
  \node (q5) at (97bp,45.608bp) [draw,circle,state] {$5$};
  \draw [->] (q0) ..controls (42.244bp,77.303bp) and (64.303bp,107.38bp)  .. node[auto] {$a$} (q3);
  \draw [->] (q5) ..controls (124.66bp,45.608bp) and (136.82bp,45.608bp)  .. node[auto] {$b$} (q9);
  \draw [->] (q3) ..controls (121.24bp,118.91bp) and (143.3bp,88.83bp)  .. node[auto] {$b$} (q9);
  \draw [->] (q5) to[loop below] node[auto] {$b$} (q5);
  \draw [->] (q0) ..controls (45.659bp,45.608bp) and (57.817bp,45.608bp)  .. node[auto] {$b$} (q5);
  \draw [->] (q0) ..controls (42.871bp,23.135bp) and (60.567bp,9.4181bp)  .. (79bp,3.6077bp) .. controls (94.26bp,-1.2026bp) and (99.74bp,-1.2026bp)  .. (115bp,3.6077bp) .. controls (129.69bp,8.2379bp) and (143.91bp,17.889bp)  .. node[auto] {$b$} (q9);
  \draw [->] (q3) to[loop above] node[auto] {$a$} (q3);
  \draw [->] (q5) ..controls (112.04bp,72.544bp) and (117.47bp,87.733bp)  .. (115bp,101.61bp) .. controls (113.63bp,109.34bp) and (111.12bp,117.45bp)  .. node[auto,swap,pos=-0.2] {$a$} (q3);
  \draw [->] (q3) ..controls (97bp,116.41bp) and (97bp,92.55bp)  .. node[auto,swap,near end] {$b$} (q5);
%
\end{tikzpicture}

%%% end example9b %%%
}
\end{center}

Il a $4$ états puisqu'il y a $3$ lettres dans $(a|b){*}b$.  Ces états
ont été étiquetés de manière à illustrer la proposition suivante, qui
fait le lien entre les deux constructions :

\begin{prop}
L'élimination des transitions spontanées (au sens
de \ref{removal-of-epsilon-transitions}, suivie de la suppression des
états devenus inutiles) dans l'automate de Thompson d'une expression
rationnelle conduit à l'automate de Glushkov de cette même expression.
\end{prop}

Sans que cela constitue une démonstration, on peut comprendre l'idée
en considérant les « remarques » qui ont été faites dans les
démonstrations de \ref{nfa-union}, \ref{nfa-concatenation}
et \ref{nfa-star}.




\subsection{Automates à transitions étiquetées par des expressions rationnelles (=RNFA), algorithme d'élimination des états}\label{subsection-rnfa-and-kleenes-algorithm}

\thingy On cherche dans cette section à montrer la réciproque
de \ref{rational-languages-are-recognizable}, c'est-à-dire, que les
langages rationnels sont reconnaissables.  On va pour cela donner
un algorithme (très coûteux !) qui transforme un automate en
expression rationnelle (dénotant le langage qu'il reconnaît).  Cette
algorithme « d'élimination des états » fonctionne naturellement sur
une sorte d'automate encore plus générale que tous ceux que nous avons
définis jusqu'à présent : on va donc commencer par définir ces
automates, même si leur intérêt réside presque uniquement en la preuve
de \ref{recognizable-languages-are-rational} ci-dessous.

\thingy\label{definition-rnfa} Un \defin[automate fini à transitions
  étiquetées par des expressions rationnelles]{automate fini
  (non-déterministe) à transitions étiquetées par des expressions
  rationnelles}, en abrégé \index{RNFA|see{automate fini à transitions
    étiquetées par des expressions rationnelles}}\textbf{RNFA}, sur un
alphabet $\Sigma$ est la donnée
\begin{itemize}
\item d'un ensemble fini $Q$ d'états,
\item d'un ensemble $I \subseteq Q$ d'états dits initiaux,
\item d'un ensemble $F \subseteq Q$ d'états dits finaux,
\item d'un ensemble \emph{fini} de transitions $\delta \subseteq Q
  \times (\mathrm{regexp}(\Sigma)) \times Q$ où
  $(\mathrm{regexp}(\Sigma))$ désigne l'ensemble des expressions
  rationnelles sur $\Sigma$.
\end{itemize}

Autrement dit, on autorise maintenant des transitions étiquetées par
des expressions rationnelles quelconques sur $\Sigma$.  Remarquons
qu'on doit maintenant demander \emph{explicitement} que l'ensemble
$\delta$ des transitions permises soit fini car l'ensemble $Q \times
(\mathrm{regexp}(\Sigma)) \times Q$, lui, ne l'est pas.

\thingy\label{rnfa-multiple-transition-relation} Pour un tel automate, on définit une relation $\delta^*
\subseteq Q \times \Sigma^* \times Q$ par $(q,w,q') \in \delta^*$
lorsqu'il existe $q_0,\ldots,q_n \in Q$ et $r_1,\ldots,r_n \in
\mathrm{regexp}(\Sigma)$ tels que $q_0 = q$ et $q_n = q'$ et
$(q_{i-1},r_i,q_i) \in\delta$ pour chaque $1\leq i\leq n$, et enfin $w
\in L(r_1\cdots r_n)$.

Concrètement, $(q,w,q') \in \delta^*$ signifie que le RNFA peut
passer de l'état $q$ à l'état $q'$ en effectuant des transitions
($q_0\to q_1 \to \cdots \to q_n$ étiquetées par $r_1,\ldots,r_n \in
\mathrm{regexp}(\Sigma)$) et en consommant le mot $w$ au sens où ce
dernier se décompose comme concaténation d'autant de facteurs que de
transitions ($w = v_1\cdots v_n$), chacun vérifiant l'expression
rationnelle qui étiquette la transition (soit $v_i \in L(r_i)$).

Enfin, l'automate $A$ accepte un mot $w$ lorsqu'il existe $q_0\in I$
et $q_\infty\in F$ tels que $(q_0,w,q_\infty) \in \delta^*$.

\smallskip

Le langage accepté $L(A)$ et l'équivalence de deux automates sont
définis de façon analogue aux DFA
(cf. \ref{definition-recognizable-language}).

\thingy Un εNFA (ou \textit{a fortiori} un NFA, DFAi ou DFA) est
considéré comme un RNFA particulier dont les transitions sont
étiquetées soit par une unique lettre (considérée comme expression
rationnelle) soit par le symbole $\underline{\varepsilon}$ (dénotant
le langage $\{\varepsilon\}$) dans le cas des transitions spontanées.

\smallskip

Une expression rationnelle $r$ peut aussi être considérée comme un
RNFA particulier comportant un unique état initial, un unique état
final, et une unique transition de l'un vers l'autre, étiquetée par
l'expression $r$ elle-même.  Il est évident que ce RNFA reconnaît
exactement le langage dénoté par $r$.

\smallskip

La représentation graphique des RNFA ne pose pas de problème
particulier (voir en \ref{example-of-state-elimination} pour
différents exemples).

\thingy\label{give-rnfa-single-transitions} On peut toujours modifier
un RNFA de manière à ce qu'il y ait au plus une, ou même si on le
souhaite, exactement une, transition entre deux états $q$ et $q'$
donnés.  En effet, s'il existe plusieurs transitions
$(q,r_1,q'),\ldots, \penalty0 (q,r_k,q') \in \delta$ possibles entre
$q$ et $q'$, on peut les remplacer par une unique transition
$(q,(r_1|\cdots|r_k),q')$, cela ne change visiblement rien au
fonctionnement de l'automate (et notamment pas le langage reconnu).
S'il n'y a \emph{aucune} transition de $q$ vers $q'$, on peut toujours
choisir d'en ajouter une $(q,\bot,q')$ (qui ne peut pas être
empruntée !) si c'est commode.

\medskip

Comme les εNFA, les NFA et les DFAi avant eux, les RNFA peuvent se
ramener aux automates précédemment définis :

\begin{prop}
Soit $A = (Q,I,F,\delta)$ un RNFA sur un alphabet $\Sigma$.  Alors il
existe un εNFA $A' = (Q',I',F',\delta')$ (sur le même
alphabet $\Sigma$) et qui soit équivalent à $A$ au sens où il
reconnaît le même langage $L(A') = L(A)$.  De plus, $A'$ se déduit
algorithmiquement de $A$.
\end{prop}
\begin{proof}
On a vu en \ref{rational-languages-are-recognizable} que pour chaque
expression rationnelle $r$ on peut trouver (algorithmiquement) un εNFA
$A(r)$ (par exemple l'automate de Glushkov ou l'automate de Thompson)
qui reconnaît le langage dénoté par $r$.  On peut donc construire $A'$
en remplaçant chaque transition $(q,r,q')$ de $A$ par une copie de
l'automate $A(r)$ placée entre les états $q$ et $q'$.  Symboliquement :

\begin{center}
\begin{tikzpicture}[>=latex,line join=bevel,automaton,baseline=(q1.base)]
\node (q1) at (0bp,0bp) [draw,circle,state] {$q$};
\node (q2) at (70bp,0bp) [draw,circle,state] {$q'$};
  \draw [->] (q1) to node[auto] {$r$} (q2);
\end{tikzpicture}
\quad devient\quad
\begin{tikzpicture}[>=latex,line join=bevel,automaton,baseline=(q1.base)]
\node (q1) at (0bp,0bp) [draw,circle,state] {$q$};
\node (q2) at (100bp,0bp) [draw,circle,state] {$q'$};
\node (A) at (50bp,0bp) [draw,dotted,circle] {$A(r)$};
  \draw [->] (q1) to node[auto] {$\varepsilon$} (A);
  \draw [->] (A) to node[auto] {$\varepsilon$} (q2);
\end{tikzpicture}
\end{center}

Plus précisément, si $\{(q_j,r_j,q'_j) : 1\leq j\leq M\}$ est une
énumération de $\delta$, on construit $A'$ en lui donnant pour
ensemble d'états $Q \uplus \biguplus_{j=1}^M Q_{r_j}$ où $Q_{r_j}$ est
l'ensemble d'états de l'automate $A(r_j)$ construit pour
reconnaître $r_j$, les ensembles d'états initiaux et finaux sont
$I'=I$ et $F'=F$ comme dans $A$, et la relation de transition
$\delta'$ est la réunion de chacune $\delta_{r_j}$ de celle des
εNFA $A(r_j)$ à quoi on ajoute encore des transitions spontanées
$(q_j,\varepsilon,q^\sharp)$ pour tout état initial $q^\sharp \in
I_{r_j}$ de $A(r_j)$ et des transitions spontanées
$(q^\flat,\varepsilon,q'_j)$ pour tout état final $q^\flat \in
F_{r_j}$ de $A(r_j)$.  Il est clair que faire un chemin dans $A'$
revient à un faire un chemin dans $A$ où, à chaque fois qu'on fait la
transition $q_j\to q'_j$ étiquetée par $r_j$, on la remplace par un
chemin $q_j \to q^\sharp \to \cdots \to q^\flat \to q'_j$ formé d'une
transition spontanée vers un état initial de $A(r_j)$ suivi d'un
chemin dans ce dernier, suivi d'une transition spontanée depuis un
état final de $A(r_j)$.
\end{proof}

\medskip

Mais la surprise des RNFA est qu'ils peuvent aussi se ramener à des
expressions rationnelles !

\begin{prop}\label{recognizable-languages-are-rational}
Soit $A = (Q,I,F,\delta)$ un RNFA (ou, en particulier, un NFA ou
DFA(i)) sur un alphabet $\Sigma$.  Alors il existe une expression
rationnelle $r$ sur $\Sigma$ qui dénote le langage reconnu par $A$,
soit $L(r) = L(A)$.  De plus, $r$ se déduit algorithmiquement de $A$.
\end{prop}
\begin{proof}
On a expliqué qu'on pouvait considérer une expression rationnelle
comme un RNFA ayant un unique état initial, un unique état final, et
une unique transition de l'un vers l'autre (étiquetée par l'expression
rationnelle en question).  On va montrer que $A$ est équivalent à un
RNFA de cette nature, ce qui montrera bien qu'il est équivalent à une
expression rationnelle.

Remarquons tout d'abord qu'on peut supposer que $A$ a un unique état
initial $q_0$, qui ne soit pas final, et qui n'ait aucune transition
qui y aboutisse (si ce n'est pas le cas, il suffit de créer un nouvel
état $q_0$, d'en faire le seul état initial, et de le munir de
transitions spontanées — c'est-à-dire étiquetées par
$\underline{\varepsilon}$ — vers tous les états précédemment
initiaux).  De même (symétriquement), on peut supposer que $A$ a un
unique état final $q_\infty$, qui ne soit pas initial, et sans aucune
transition qui en part.  On fera l'hypothèse que $A$ a ces propriétés,
et on s'arrangera pour les préserver dans ce qui suit.

Soient maintenant $q$ un état de $A$ qui n'est ni l'état initial $q_0$
ni l'état final $q_\infty$.  On va montrer qu'on peut \emph{éliminer}
$q$, c'est-à-dire, quitte à ajouter des transitions, remplacer $A$ par
un automate équivalent $A'$ qui n'a pas cet état.  Pour cela, soient
$q_1,q_2$ deux états quelconques de $A$, autres que $q$ mais
possiblement égaux entre eux, où $q_1$ peut être l'état initial (mais
pas l'état final) et $q_2$ peut être l'état final (mais pas l'état
initial).  On a vu en \ref{give-rnfa-single-transitions} qu'on pouvait
supposer qu'il existait une unique transition $(q_1,r_{12},q_2)$ et de
même $(q_1,r_1,q)$ et $(q,r_2,q_2)$ et $(q,s,q)$ (transition de $q$
vers lui-même).  En même temps qu'on élimine $q$, on met dans $A'$ la
transition $(r_{12}|r_1(s){*}r_2)$ entre $q_1$ et $q_2$.
Symboliquement :

\begin{center}
\begin{tikzpicture}[>=latex,line join=bevel,automaton,baseline=(q1.base)]
\node (q1) at (0bp,0bp) [draw,circle,state] {$q_1$};
\node (q2) at (70bp,0bp) [draw,circle,state] {$q_2$};
\node (q) at (35bp,50bp) [draw,circle,state] {$q$};
  \draw [->] (q1) to node[auto] {$r_{12}$} (q2);
  \draw [->] (q1) to node[auto] {$r_1$} (q);
  \draw [->] (q) to node[auto] {$r_2$} (q2);
  \draw [->] (q) to[loop above] node[auto] {$s$} (q);
\end{tikzpicture}
\quad devient\quad
\begin{tikzpicture}[>=latex,line join=bevel,automaton,baseline=(q1.base)]
\node (q1) at (0bp,0bp) [draw,circle,state] {$q_1$};
\node (q2) at (100bp,0bp) [draw,circle,state] {$q_2$};
  \draw [->] (q1) to node[auto] {$\scriptstyle (r_{12}|r_1(s){*}r_2)$} (q2);
\end{tikzpicture}
\end{center}

Pour éliminer l'état $q$, cette transformation doit être effectuée
\emph{simultanément pour toute paire} $(q_1,q_2)$ d'états de $A$ (avec
$q_1\not\in\{q,q_\infty\}$ et $q_2\not\in\{q,q_0\}$) pour laquelle il
existe une transition de $q_1$ vers $q$ et une transition de
$q$ vers $q_2$, \emph{y compris} s'il n'y avait pas déjà une
transition de $q_1$ vers $q_2$, et \emph{y compris} si $q_1=q_2$ :
pour \emph{chaque} telle paire, on remplace l'étiquette de la
transition $r_{12}$ entre eux par $(r_{12}|r_1(s){*}r_2)$.  (On
prendra garde aux cas particuliers suivants : si la transition de $q$
vers lui-même n'existait pas, ce qui revient à prendre $s=\bot$, alors
on remplace la transition $r_{12}$ par $(r_{12}|r_1 r_2)$ en vertu du
fait que $(\bot){*}$ est équivalente à $\underline{\varepsilon}$ ; et
si la transition de $q_1$ vers $q_2$ n'existait pas, ce qui revient à
prendre $r_{12}=\bot$, alors on la crée avec l'étiquette
$r_1(s){*}r_2$ ; et bien sûr, en combinant ces deux cas particuliers,
s'il n'y avait ni transition de $q$ vers lui-même ni de
$q_1$ vers $q_2$, on crée cette dernière avec l'étiquette $r_1 r_2$.)

La transformation qui vient d'être décrite ne change pas le
fonctionnement de l'automate, car tout chemin dans $A$ peut être
remplacé par un chemin dans $A'$ en effaçant simplement les $q$ (si on
considère $q_1$ et $q_2$ les états avant un bloc de $q$ dans le
chemin, on voit que le chemin $q_1 \to q \to q \to \cdots \to q \to
q_2$ peut se transformer en $q_1 \to q_2$ en consommant un mot qui
vérifie l'expression rationnelle $(r_{12}|r_1(s){*}r_2)$).

En éliminant (successivement dans n'importe quel ordre) tous les états
autres que $q_0$ et $q_\infty$, on aboutit ainsi à un automate ayant
une unique transition $(q_0,r,q_\infty)$, qui est donc essentiellement
l'expression rationnelle $r$.
\end{proof}

\thingy La procédure qu'on a décrite dans la démonstration de cette
proposition s'appelle l'algorithme d'\defin{élimination des états} ou
\index{Kleene (algorithme de)|see{élimination des
    états}}\textbf{algorithme de Kleene}\footnote{Peut-être
  abusivement (le théorème \ref{kleenes-theorem} est indubitablement
  dû à Kleene, mais le contenu algorithmique ne l'est peut-être pas) ;
  il est peut-être plus correct d'attribuer l'algorithme à Brzozowski
  et McCluskey.}.

Il va de soi qu'on peut la simplifier un petit peu : s'il n'y a pas de
transition de $q_1$ vers $q$ ou qu'il n'y en a pas de $q$
vers $q_2$ (c'est-à-dire que soit $r_1$ soit $r_2$ doit être considéré
comme valant $\bot$), on ne touche simplement pas à $r_{12}$ (et si la
transition de $q_1$ vers $q_2$ n'existait pas non plus, il n'y a pas
besoin de la créer) ; de même, s'il n'y a pas de transition de $q$
vers lui-même, on ignore la partie $s{*}$.  En revanche, il faut bien
penser à créer une transition de $q_1$ vers $q_2$, même si elle
n'existait pas au départ, lorsqu'on peut arriver de l'un vers l'autre
en passant par $q$.  Et il faut se souvenir que le cas $q_2=q_1$ est à
traiter aussi.

En général, l'élimination des états conduit à une expression
extrêmement compliquée (exponentielle dans le nombre d'états de
l'automate, au moins dans le pire cas, mais aussi dans beaucoup de cas
« typiques »).

\thingy\label{example-of-state-elimination} À titre d'exemple,
considérons le DFA suivant sur l'alphabet $\{0,1\}$, qui reconnaît les
suites binaires qui représentent un nombre multiple de $3$ écrit en
binaire\footnote{Les états $0,1,2$ représentent respectivement les
  états dans lesquels le nombre binaire lu jusqu'à présent par
  l'automate est congru à $0,1,2$ modulo $3$.} (en convenant que le
mot vide est une représentation binaire du nombre $0$, ce qui est
logique) :

\begin{center}
%%% begin example6 %%%

\begin{tikzpicture}[>=latex,line join=bevel,automaton]
%%
\node (q1) at (97bp,20.28bp) [draw,circle,state] {$1$};
  \node (q0) at (18bp,20.28bp) [draw,circle,state,initial,final,accepting below] {$0$};
  \node (q2) at (176bp,20.28bp) [draw,circle,state] {$2$};
  \draw [->] (q1) ..controls (74.757bp,3.6593bp) and (64.084bp,-1.2803bp)  .. (54bp,1.2803bp) .. controls (50.042bp,2.2853bp) and (46.047bp,3.838bp)  .. node[auto] {$1$} (q0);
  \draw [->] (q2) to[loop above] node[auto] {$1$} (q2);
  \draw [->] (q2) ..controls (153.76bp,3.6593bp) and (143.08bp,-1.2803bp)  .. (133bp,1.2803bp) .. controls (129.04bp,2.2853bp) and (125.05bp,3.838bp)  .. node[auto] {$0$} (q1);
  \draw [->] (q0) to[loop above] node[auto] {$0$} (q0);
  \draw [->] (q0) ..controls (45.659bp,20.28bp) and (57.817bp,20.28bp)  .. node[auto] {$1$} (q1);
  \draw [->] (q1) ..controls (124.66bp,20.28bp) and (136.82bp,20.28bp)  .. node[auto] {$0$} (q2);
%
\end{tikzpicture}

%%% end example6 %%%
\end{center}

On commence par ajouter un état initial $q_0$ et un état
final $q_\infty$, avec des ε-transitions $q_0 \to 0$ et $0\to
q_\infty$.  Pour gagner de la place, nous ne figurerons pas ces deux
états sur les dessins qui suivent, mais il faut s'imaginer qu'ils sont
toujours là.

L'élimination de l'état $2$ conduit à l'automate suivant :

\begin{center}
%%% begin example6b %%%

\begin{tikzpicture}[>=latex,line join=bevel,automaton]
%%
\node (q1) at (97bp,20.28bp) [draw,circle,state] {$1$};
  \node (q0) at (18bp,20.28bp) [draw,circle,state,initial,final,accepting below] {$0$};
  \draw [->] (q1) ..controls (74.757bp,3.6593bp) and (64.084bp,-1.2803bp)  .. (54bp,1.2803bp) .. controls (50.042bp,2.2853bp) and (46.047bp,3.838bp)  .. node[auto] {$1$} (q0);
  \draw [->] (q0) ..controls (45.659bp,20.28bp) and (57.817bp,20.28bp)  .. node[auto] {$1$} (q1);
  \draw [->] (q1) to[loop right] node[auto] {$01{*}0$} (q1);
  \draw [->] (q0) to[loop above] node[auto] {$0$} (q0);
%
\end{tikzpicture}

%%% end example6b %%%
\end{center}

\noindent (Répétons qu'on n'a pas figuré l'état initial $q_0$ ni
l'état final $q_\infty$ : les flèches vers et depuis l'état $0$
doivent se comprendre comme des ε-transitions $q_0\to 0$ et $0\to
q_\infty$.)

L'élimination de l'état $1$ conduit alors à l'automate ayant un unique
état $0$, avec une transition vers lui-même étiquetée
$0|1(01{*}0){*}1$.  Enfin, en éliminant l'état $0$, il ne reste qu'une
transition de l'état initial vers l'état final, étiquetée par
$(0|1(01{*}0){*}1){*}$ : le langage reconnu par l'automate de départ
est donc celui dénoté par l'expression
rationnelle $(0|1(01{*}0){*}1){*}$.

On pouvait aussi choisir d'éliminer l'état $1$ en premier, ce qui
conduit à l'automate suivant :

\begin{center}
%%% begin example6c %%%

\begin{tikzpicture}[>=latex,line join=bevel,automaton]
%%
\node (q0) at (18bp,20.114bp) [draw,circle,state,initial,final,accepting below] {$0$};
  \node (q2) at (104bp,20.114bp) [draw,circle,state] {$2$};
  \draw [->] (q2) ..controls (82.598bp,6.6202bp) and (75.237bp,2.9515bp)  .. (68bp,1.1137bp) .. controls (59.228bp,-1.1137bp) and (49.898bp,1.2372bp)  .. node[auto] {$01$} (q0);
  \draw [->] (q0) ..controls (47.743bp,20.114bp) and (62.773bp,20.114bp)  .. node[auto] {$10$} (q2);
  \draw [->] (q0) to[loop above] node[auto] {$0|11$} (q0);
  \draw [->] (q2) to[loop above] node[auto] {$1|00$} (q2);
%
\end{tikzpicture}

%%% end example6c %%%
\end{center}

\noindent et finalement à l'expression rationnelle
$(0|11|10(1|00){*}01){*}$, qui est équivalente à la précédente.

\bigbreak

\thingy Donnons encore l'exemple du DFAi suivant :

\begin{center}
%%% begin example10 %%%

\begin{tikzpicture}[>=latex,line join=bevel,automaton]
%%
\node (q1) at (97bp,20.28bp) [draw,circle,state] {$1$};
  \node (q0) at (18bp,20.28bp) [draw,circle,state,initial,final,accepting below] {$0$};
  \node (q3) at (255bp,20.28bp) [draw,circle,state] {$3$};
  \node (q2) at (176bp,20.28bp) [draw,circle,state] {$2$};
  \node (q4) at (334bp,20.28bp) [draw,circle,state] {$4$};
  \draw [->] (q1) ..controls (74.757bp,3.6593bp) and (64.084bp,-1.2803bp)  .. (54bp,1.2803bp) .. controls (50.042bp,2.2853bp) and (46.047bp,3.838bp)  .. node[auto] {$b$} (q0);
  \draw [->] (q2) ..controls (203.66bp,20.28bp) and (215.82bp,20.28bp)  .. node[auto] {$a$} (q3);
  \draw [->] (q2) ..controls (153.76bp,3.6593bp) and (143.08bp,-1.2803bp)  .. (133bp,1.2803bp) .. controls (129.04bp,2.2853bp) and (125.05bp,3.838bp)  .. node[auto] {$b$} (q1);
  \draw [->] (q3) ..controls (282.66bp,20.28bp) and (294.82bp,20.28bp)  .. node[auto] {$a$} (q4);
  \draw [->] (q3) ..controls (232.76bp,3.6593bp) and (222.08bp,-1.2803bp)  .. (212bp,1.2803bp) .. controls (208.04bp,2.2853bp) and (204.05bp,3.838bp)  .. node[auto] {$b$} (q2);
  \draw [->] (q0) ..controls (45.659bp,20.28bp) and (57.817bp,20.28bp)  .. node[auto] {$a$} (q1);
  \draw [->] (q4) ..controls (311.76bp,3.6593bp) and (301.08bp,-1.2803bp)  .. (291bp,1.2803bp) .. controls (287.04bp,2.2853bp) and (283.05bp,3.838bp)  .. node[auto] {$b$} (q3);
  \draw [->] (q1) ..controls (124.66bp,20.28bp) and (136.82bp,20.28bp)  .. node[auto] {$a$} (q2);
%
\end{tikzpicture}

%%% end example10 %%%
\end{center}

Il s'agit d'un automate « compteur limité », qui ne sait compter que
de $0$ à $4$, incrémentant son compteur quand il reçoit un $a$ et le
décrémentant quand il reçoit un $b$ (et cessant de fonctionner si le
compteur passe au-dessus du maximum ou en-dessous du minimum), et qui
accepte finalement les mots dont le nombre de $b$ égale le nombre
de $a$ sans qu'il y ait jamais eu plus de $b$ que de $a$ ni plus de
quatre $a$ de plus que de $b$.  (On peut dire aussi qu'il s'agit d'une
\emph{approximation} du langage des expressions bien-parenthésées
définies en \ref{example-well-parenthesized-expressions} plus loin, où
$a$ joue le rôle de parenthèse ouvrante et $b$ de parenthèse
fermante ; l'approximation est due au fait qu'on n'accepte que quatre
niveaux d'imbrication des « parenthèses ».)

Si on élimine les états dans l'ordre $4,3,2,1,0$, on montre que le
langage reconnu par cet automate est décrit par l'expression
rationnelle $(a(a(a(ab){*}b){*}b){*}b){*}$.  Si on les élimine dans
l'ordre $0,1,2,3,4$, en revanche, on obtient une expression de taille
considérable\footnote{À savoir : $\underline{\varepsilon} \penalty0 \;
  | \penalty0 \; a(ba){*}b \penalty0 \; | \penalty0 \;
  a(ba){*}a(b(ba){*}a){*}b(ba){*}b \penalty0 \; | \penalty0 \;
  a(ba){*}a(b(ba){*}a){*}a(b(b(ba){*}a){*}a){*}b(b(ba){*}a){*}b(ba){*}b
  \penalty0 \; | \penalty0 \;
  a(ba){*}a(b(ba){*}a){*}a(b(b(ba){*}a){*}a){*}a(b(b(b(ba){*}a){*}a){*}a){*}b(b(b(ba){*}a){*}a){*}b(b(ba){*}a){*}b(ba){*}b$
  (mais si on la regarde d'assez près, on peut comprendre comment elle
  fonctionne).}  et beaucoup moins transparente.

\bigbreak

Récapitulons le contenu essentiel à retenir comme conséquence
immédiate de \ref{rational-languages-are-recognizable}
et \ref{recognizable-languages-are-rational} :

\begin{thm}[Kleene]\label{kleenes-theorem}\index{Kleene (théorème de)}
La classe des langages rationnels et celle des langages
reconnaissables coïncident.  (On pourra donc considérer ces termes
comme synonymes.)
\end{thm}

\smallskip

Il faut cependant retenir que s'il y a, mathématiquement, équivalence
entre ces deux classes de langages, cette équivalence \emph{a un coût
  algorithmique}, c'est-à-dire que la conversion d'une expression
rationnelle en automate (surtout si on souhaite un DFA), ou à plus
forte raison d'un automate en expression rationnelle, a une complexité
exponentielle dans le pire cas.  Il est donc pertinent, en
informatique, de ne pas considérer les descriptions d'un langage par
une expression rationnelle, un DFA, et un NFA, comme interchangeables.

\thingy Une conséquence de \ref{kleenes-theorem} est, par exemple, que
le complémentaire d'un langage rationnel est rationnel
(cf. \ref{dfa-complement}) ou que l'intersection de deux langages
rationnels est rationnelle (cf. \ref{dfa-union-and-intersection}), et
que ces opérations se calculent algorithmiquement.  Il est donc
possible, en principe, à partir d'une expression rationnelle $r$ (sur
un alphabet $\Sigma$), de fabriquer algorithmiquement une expression
rationnelle $r'$ (« négation » de $r$) qui dénote le langage
complémentaire de celui dénoté par $r$ ; et de même, à partir
d'expressions rationnelles $r_1,r_2$, de fabriquer algorithmiquement
une expression rationnelle $r''$ (« conjonction » de $r_1$ et $r_2$)
qui dénote le langage intersection de ceux dénotés par $r_1$ et $r_2$.

\smallskip

Le coût de ces opérations, cependant, est astronomique : doublement
exponentiel, puisqu'il s'agit de convertir l'expression en NFA, de
déterminiser le NFA en DFA (à un premier coût exponentiel),
d'effectuer l'opération sur un DFA, et de reconvertir le DFA en
expression rationnelle (à un deuxième coût exponentiel).  Ce n'est pas
faute d'avoir les bons algorithmes : on peut montrer qu'il existe
effectivement des familles d'expressions rationnelles pour lesquelles
la négation ou la conjonction produisent une augmentation de taille
doublement exponentielle.

%% Gelade & Neven (STACS 2008),
%% "Succinctness of the Complement and Intersection of Regular Expressions"


\subsection{Le lemme de pompage}\label{subsection-pumping-lemma}

\thingy On ne dispose à ce stade-là d'aucun moyen pour montrer qu'un
langage \emph{n'est pas} rationnel.  La
proposition \ref{pumping-lemma} qui va suivre, et qui s'appelle
couramment « lemme de pompage » (une traduction abusive de l'anglais
« pumping lemma ») constitue le moyen le plus fréquent permettant d'y
arriver : il énonce une condition \emph{nécessaire} pour qu'un langage
soit rationnel, si bien qu'on peut arriver à montrer qu'un langage
n'est pas rationnel en invalidant cette condition (généralement en
procédant par l'absurde).

\begin{prop}[lemme de pompage pour les langages rationnels]\label{pumping-lemma}\index{pompage (lemme de)}
Soit $L$ un langage rationnel.  Il existe alors un entier $k$ tel que
tout mot de $t \in L$ de longueur $|t| \geq k$ admette une
factorisation $t = uvw$ en trois facteurs $u,v,w \in \Sigma^*$ où :
\begin{itemize}
\item[(i)] $|v| \geq 1$ (c'est-à-dire $v\neq\varepsilon$),
\item[(ii)] $|uv| \leq k$,
\item[(iii)] pour tout $i\geq 0$ on a $uv^iw \in L$.
\end{itemize}
\end{prop}
\begin{proof}
Soit $A$ un DFA (complet) qui reconnaît $L$, et soit $k$ son nombre
d'états : on va montrer que $k$ vérifie les propriétés énoncées.  Pour
cela, soit $t = x_1 \cdots x_n$ un mot de $L$ de longueur $n \geq k$,
et soient $q_0,\ldots,q_n$ les états traversés par $A$ pendant la
consommation de $t$, autrement dit, $q_0$ est l'état initial, et $q_j
= \delta(q_{j-1}, x_j)$ pour chaque $1\leq j\leq n$ ; l'état $q_n =
\delta^*(q_0, t)$ est final puisque $t \in L$.  Comme $n+1 > k$ et
comme l'automate $A$ a $k$ états, par le principe des tiroirs, il
existe $j_1\neq j_2$ tels que $q_{j_1} = q_{j_2}$ : pour être plus
précis, soit $j_2$ le plus petit possible tel que les états
$q_0,\ldots,q_{j_2}$ ne soient pas tous distincts, autrement dit, le
premier état répété, et soit $j_1$ la précédente occurrence (forcément
unique) de cet état, c'est-à-dire l'indice tel que $j_1<j_2$ et
$q_{j_1} = q_{j_2}$.

\begin{center}
\begin{tikzpicture}[>=latex,line join=bevel,automaton]
\draw[draw=none,fill=green!20!white] (-58.476bp,-58.476bp) rectangle (58.476bp,58.476bp);
\draw[dotted] (29.238bp,50.642bp) arc (60:-100:58.476bp);
\draw[draw=none,fill=red!20!white] (-188.476bp,-25bp) rectangle (-68.476bp,25bp);
\draw[draw=none,fill=blue!20!white] (-48.476bp,-25bp) rectangle (110bp,25bp);
\node[anchor=south west] at (-188.476bp,-25bp) {$u$};
\node[anchor=south east] at (58.476bp,-58.476bp) {$v$};
\node[anchor=south east] at (110bp,-25bp) {$w$};
\node (q0) at (-198.476bp,0bp) [draw,circle,state,fill=white,initial] {\vbox to0pt{\vss\hbox to0pt{$q_0$\hss}}\phantom{$q_{j_0}$}};
\node (q1) at (-138.476bp,0bp) [draw,circle,state,fill=white] {\vbox to0pt{\vss\hbox to0pt{$q_1$\hss}}\phantom{$q_{j_0}$}};
\node (q2) at (-98.476bp,0bp) {$\cdots$};
\node (qj1) at (-58.476bp,0bp) [draw,circle,state,fill=white] {\vbox to0pt{\vss\hbox to0pt{$q_{j_1}$\hss}}\phantom{$q_{j_0}$}};
\node (qj1p1) at (-44.795bp,37.588bp) [draw,circle,state,fill=white] {\phantom{$q_{j_0}$}};
\node (qj1p2) at (-10.154bp,57.588bp) [draw,circle,state,fill=white] {\phantom{$q_{j_0}$}};
\node (qj1p3) at (29.238bp,50.642bp) {};
\node (qj2m2) at (-10.154bp,-57.588bp) {};
\node (qj2m1) at (-44.795bp,-37.588bp) [draw,circle,state,fill=white] {\phantom{$q_{j_0}$}};
\node (qj2p1) at (0bp,0bp) [draw,circle,state,fill=white] {\phantom{$q_{j_0}$}};
\node (qj2p2) at (60bp,0bp) {$\cdots$};
\node (qn) at (120bp,0bp) [draw,circle,state,fill=white,final] {\vbox to0pt{\vss\hbox to0pt{$q_n$\hss}}\phantom{$q_{j_0}$}};
\draw[->] (q0) -- node[auto] {\footnotesize $x_1$} (q1);
\draw[->] (q1) -- (q2);
\draw[->] (q2) -- node[auto] {\footnotesize $x_{j_1}$} (qj1);
\draw[->] (qj1) -- node[auto] {\footnotesize $x_{j_1+1}$} (qj1p1);
\draw[->] (qj1p1) -- node[auto] {\footnotesize $x_{j_1+2}$} (qj1p2);
\draw[->] (qj1p2) -- (qj1p3);
\draw[->] (qj2m2) -- (qj2m1);
\draw[->] (qj2m1) -- node[auto] {\footnotesize $x_{j_2}$} (qj1);
\draw[->] (qj1) -- node[auto] {\footnotesize $x_{j_2+1}$} (qj2p1);
\draw[->] (qj2p1) -- (qj2p2);
\draw[->] (qj2p2) --  node[auto] {\footnotesize $x_n$}(qn);
\end{tikzpicture}
\end{center}

Posons $u = x_1\cdots x_{j_1}$ (le préfixe de $t$ de longueur $j_1$,
qui est le mot vide si $j_1 = 0$) et $v = x_{j_1+1}\cdots x_{j_2}$ (de
longueur $j_2 - j_1$), et enfin $w = x_{j_2+1}\cdots x_n$ (le suffixe
de $t$ de longueur $n-j_2$, avec la convention $w = \varepsilon$ si
$j_2 = n$).  Ceci définit bien une factorisation $t = uvw$.

On a bien (i) $|v| \geq 1$ puisque $j_2 > j_1$.  On a par ailleurs
(ii) $|uv| \leq k$ puisque $|uv| = j_2$ et que les $j_2$ états
$q_0,\ldots,q_{j_2-1}$ sont distincts (c'est la minimalité de $j_2$)
de sorte que $j_2 \leq k$ (toujours par le principe des tiroirs).

Montrons enfin (iii).  On rappelle tout d'abord que
$\delta^*(q,x_1\cdots x_j) =\penalty0 \delta(\cdots\penalty500
\delta(\delta(q,x_1),x_2)\cdots,x_j)$.  Remarquons que
$\delta^*(q_0,u) = \delta^*(q_0,x_1\cdots x_{j_1}) = q_{j_1}$, et que
$\delta^*(q_{j_1},v) = \delta^*(q_{j_1},x_{j_1+1}\cdots x_{j_2}) =
q_{j_2} = q_{j_1}$.  De cette dernière égalité, on tire
$\delta^*(q_{j_1},v^i) = q_{j_1}$ pour tout $i \geq 0$ (par récurrence
sur $i$).  Enfin, $\delta^*(q_{j_1},w) = \delta^*(q_{j_2},w) =
\delta^*(q_{j_2}, x_{j_2+1}\cdots x_n) = q_n$ (qui est un état final).
En mettant ces faits ensemble, on a $\delta^*(q_0, uv^iw) =
\delta^*(q_{j_1}, v^iw) = \delta^*(q_{j_1}, w) = q_n$, et puisque
$q_n$ est final, ceci montre que le mot $uv^iw$ est accepté par $A$,
i.e., $uv^iw \in L$.
\end{proof}

\thingy On attire l'attention sur l'alternation des quantificateurs.
Le lemme de pompage énonce le fait que :
\begin{itemize}
\item\emph{pour tout} langage rationnel $L$,
\item\emph{il existe} un entier $k\geq 0$ tel que
\item\emph{pour tout} mot $t\in L$ de longueur $|t|\geq k$,
\item\emph{il existe} une factorisation $t=uvw$ vérifiant les
  propriétés (i) $|v|\geq 1$, (ii) $|uv|\leq k$ et (iii) qui suit :
\item\emph{pour tout} $i\geq 0$ on a $uv^iw \in L$.
\end{itemize}

La complexité logique d'un énoncé étant justement mesurée par le
nombre d'alernations de quantificateurs (passages entre « pour tout »
et « il existe » ou vice versa), celui-ci mérite une attention
particulière.  Rappelons donc, du point de vue logique, que, quand on
veut \emph{appliquer} un résultat de ce genre, on \emph{choisit
  librement} les objets introduits par un quantificateur universel
(« pour tout »), mais \emph{on ne choisit pas} ceux qui sont
introduits par un quantificateur existentiel (« il existe ») (ces
derniers sont, si on veut, choisis par l'énoncé qu'on applique : on ne
fait que recevoir leur existence) ; les choses sont inversées quand on
doit démontrer un tel énoncé, mais la démonstration a été faite
ci-dessus et il est donc plus fréquent de devoir appliquer le lemme de
pompage.

\thingy Le modèle d'une démonstration par l'absurde pour montrer qu'un
langage $L$ n'est pas rationnel est donc quelque chose comme ceci :
\begin{itemize}
\item on entame un raisonnement par l'absurde en supposant que $L$ est
  rationnel, et on choisit $L$ pour appliquer le lemme de pompage
  (parfois on l'applique à autre chose, comme l'intersection de $L$
  avec un langage connu pour être rationnel, mais en général ce
  sera $L$),
\item le lemme de pompage fournit un $k$ (on \emph{ne choisit donc
  pas} ce $k$, il est donné par le lemme),
\item on choisit alors un mot $t \in L$ de longueur $\geq k$, et c'est
  là que réside la difficulté principale de la démonstration,
\item le lemme de pompage fournit une factorisation $t = uvw$, qu'on
  \emph{ne choisit pas} non plus mais qu'on peut analyser, souvent en
  utilisant (i) et (ii),
\item et on cherche à appliquer la propriété (iii), ce qui implique de
  choisir un $i$, pour arriver à une contradiction (typiquement : le
  mot $uv^i w$ n'est pas dans le langage alors qu'il est censé y
  être).
\end{itemize}

\smallskip

Donnons maintenant un exemple d'utilisation du lemme :

\begin{prop}\label{example-of-pumping-lemma}
Soit $\Sigma = \{a,b\}$.  Le langage $L = \{a^n b^n : n\in\mathbb{N}\}
= \{\varepsilon, ab, aabb, aaabbb,\ldots\}$ constitué des mots formés
d'un certain nombre ($n$) de $a$ suivis du même nombre de $b$ n'est
pas rationnel.
\end{prop}
\begin{proof}
Appliquons la proposition \ref{pumping-lemma} au langage $L$
considéré : appelons $k$ l'entier dont le lemme de pompage garantit
l'existence.  Considérons le mot $t := a^k b^k$ : il doit alors
exister une factorisation $t = uvw$ pour laquelle on a (i) $|v|\geq
1$, (ii) $|uv|\leq k$ et (iii) $uv^iw \in L$ pour tout $i\geq 0$.  La
propriété (ii) assure que $uv$ est formé d'un certain nombre de
répétitions de la lettre $a$ (car tout préfixe de longueur $\leq k$ de
$a^k b^k$ est de cette forme) ; disons $u = a^\ell$ et $v = a^m$, si
bien que $w = a^{k-\ell-m} b^k$.  La propriété (i) donne $m\geq 1$.
Enfin, la propriété (iii) affirme que le mot $uv^iw = a^{k+(i-1)m}
b^k$ appartient à $L$ ; mais dès que $i\neq 1$, ceci est faux : il
suffit donc de prendre $i=0$ pour avoir une contradiction.
\end{proof}

\thingy L'idée intuitive derrière la démonstration qu'on vient de
faire est la suivante : un automate fini ne dispose que d'une quantité
finie (bornée) de mémoire, donc ne peut « compter » que jusqu'à un
nombre borné (au-delà, il va retomber sur un état déjà atteint par un
plus petit nombre de $a$, et sera incapable de vérifier si le nombre
de $b$ est égal).  C'est ce type de raisonnement que le lemme de
pompage permet de formaliser.  Généralement parlant, on doit garder à
l'esprit le fait que toutes sortes de langages ne sont pas rationnels
pour la raison informelle que les identifier demande une quantité de
mémoire qui pourrait être arbitrairement grande, et que pour le
montrer rigoureusement, le lemme de pompage sera souvent utile.


\subsection{L'automate minimal, et la minimisation}

\thingy On sait maintenant convertir une expression rationnelle en un
automate équivalent, et réciproquement convertir un automate en
expression rationnelle équivalente.  Il reste un problème auquel nous
n'avons pas donné de réponse : comment savoir si deux automates
\emph{donnés} ou deux expressions rationnelles données (ou un de
chaque) sont équivalents ?

\smallskip

Pour cela, on va introduire un DFA particulier, \emph{canonique},
associé à un langage rationnel, qu'on pourra calculer
algorithmiquement, et qui sera véritablement associé au langage,
c'est-à-dire que deux descriptions équivalentes du même langage
donneront exactement le \emph{même} automate canonique ; par
conséquent, pour tester l'égalité de deux langages, quelle que soit
leur description, il suffira de calculer leurs automates canoniques et
de les comparer.  En fait, cet automate canonique est simplement le
DFA ayant le plus petit nombre d'états reconnaissant le langage en
question ; ce qui est remarquable, et qui n'est pas du tout évident
\textit{a priori}, c'est qu'il est effectivement canonique
(c'est-à-dire qu'il n'existe qu'un seul DFA ayant un nombre minimal
d'états reconnaissant le langage) et qu'on peut le calculer
algorithmiquement.

\begin{thm}[Myhill-Nerode]\label{myhill-nerode}\index{Myhill-Nerode (théorème de)}
Soit $L$ un langage.  Pour $w\in \Sigma^*$, notons $w^{-1} L := \{t
\in \Sigma^* : wt \in L\}$ (autrement dit, l'ensemble des mots qu'on
peut concaténer à $w$ pour obtenir un mot de $L$).  Considérons la
relation d'équivalence $\equiv$ sur $\Sigma^*$ définie par $u \equiv
v$ si et seulement si $u^{-1}L = v^{-1}L$ ; ce qui signifie, si on
préfère, que $\forall t\in \Sigma^*\,((ut\in L) \liff (vt \in L))$.

Alors :
\begin{itemize}
\item le langage $L$ est rationnel si et seulement si la relation
  d'équivalence $\equiv$ possède un nombre \emph{fini} $k$ de classes
  d'équivalence,
\item lorsque c'est le cas, il existe un DFA (complet) ayant $k$ états
  qui reconnaît $L$, il est unique à renommage des
  états\footnote{C'est-à-dire, si on préfère ce terme, isomorphisme
    d'automates.} près, et il n'existe pas de DFA (complet) ayant $<k$
  états qui reconnaisse $L$.
\end{itemize}
\end{thm}
\begin{proof}
Supposons d'abord que l'ensemble $\Sigma^*/{\equiv}$ des classes
d'équivalence pour $\equiv$ soit fini : appelons-le $Q$, et expliquons
comment on peut construire un DFA reconnaissant $L$ et dont l'ensemble
des états soit $Q$.  Notons $[u] := \{v \in \Sigma^* :\penalty-100
u\equiv v\}$ pour la classe d'équivalence de $u$ pour $\equiv$.
Posons $q_0 := [\varepsilon]$ la classe du mot vide.  Remarquons que
si $u\equiv v$ (pour $u,v\in\Sigma^*$), alors on a $u\in L$ si et
seulement si $v\in L$ (en effet, la définition de $\equiv$ est que
$(ut\in L) \liff (vt \in L)$ pour tout $t$, et on applique ça
à $t=\varepsilon$) : autrement dit, une classe $[u] \in Q$ est soit
entièrement incluse dans $L$ soit disjointe de $L$ ; appelons $F$
l'ensemble $\{[u] : u\in L\}$ des classes incluses dans $L$.  Enfin
remarquons que si $u\equiv v$ (pour $u,v\in\Sigma^*$) et $x\in\Sigma$,
on a encore $ux \equiv vx$ (en effet, $uxt \in L \liff vxt \in L$ pour
tout $t\in L$) : il y a donc un sens à définir $\delta([u], x) =
[ux]$.  On a ainsi fabriqué un DFA $A = (Q,q_0,F,\delta)$.
Vues les définitions de $q_0$ et $\delta$, il est clair que
$\delta^*(q_0,w) = [w]$ pour cet automate, et vue la définition de
$F$, on a $\delta^*(q_0,w) \in F$ si et seulement si $w\in L$.  Ceci
montre bien que $A$ reconnaît le langage $L$.  On a donc prouvé que si
$\Sigma^*/{\equiv}$ est fini, le langage $L$ est rationnel et même il
existe un DFA ayant $k := \#(\Sigma^*/{\equiv})$ états qui le
reconnaît.

Supposons maintenant que $B = (Q_B, q_{0,B}, F_B, \delta_B)$ soit un
DFA reconnaissant $L$.  Définissons une nouvelle relation
d'équivalence $\mathrel{\equiv_B}$ sur $\Sigma^*$ par : $u
\mathrel{\equiv_B} v$ si et seulement si $\delta_B^*(q_{0,B}, u) =
\delta_B^*(q_{0,B}, v)$ (autrement dit, les mots $u$ et $v$ mettent
l'automate $B$ dans le même état).  Si on a $u \mathrel{\equiv_B} v$,
on a aussi $u \equiv v$ : en effet, pour tout $t\in L$ on a
$\delta_B^*(q_{0,B}, ut) = \delta_B^*(\delta_B^*(q_{0,B}, u), t) =
\delta_B^*(\delta_B^*(q_{0,B}, v), t) = \delta_B^*(q_{0,B}, vt)$, et
notamment le membre de gauche appartient à $F_B$ si et seulement si le
membre de droite y appartient, c'est-à-dire que $ut\in L \liff vt\in
L$.  On vient donc de montrer que $u \mathrel{\equiv_B} v$ implique $u
\equiv v$ (la relation $\equiv_B$ est \emph{plus fine} que $\equiv$) :
si on préfère, chaque classe d'équivalence pour $\equiv$ est donc une
réunion de classes d'équivalences pour $\equiv_B$.  Notamment, comme
$\equiv_B$ a un nombre fini de classes d'équivalence (puisque
$\delta_B^*(q_{0,B}, u)$ ne peut prendre qu'un nombre fini de valeurs,
celles dans $Q_B$), il en va de même de $\equiv$, et plus précisément,
comme $\equiv_B$ a au plus $\#Q_B$ classes d'équivalence, il en va de
même de $\equiv$, c'est-à-dire $k \leq \#Q_B$.  Il n'existe donc pas
de DFA ayant $<k$ états reconnaissant $L$.

Enfin, si $B$ a $k$ états et reconnaît $L$, cela signifie que les deux
relations $\equiv$ et $\equiv_B$ coïncident, et on définit une
bijection $\psi$ entre le $Q = Q_A$ du paragraphe précédent et $Q_B$
en associant à une classe $[w] \in Q$ l'état $\psi([w]) :=
\delta_B^*(q_{0,B},w)$ (qui ne dépend que de la classe de $w$
pour $\equiv_B$ et on vient de voir que c'est la classe de $w$
modulo $\equiv$, c'est-à-dire $[w]$ : ceci est donc bien défini).
Cette bijection $\psi$ vérifie $\psi(q_0) = \psi([\varepsilon]) =
\delta_B^*(q_{0,B},\varepsilon) = q_{0,B}$ ; on a $[w] \in F$ si et
seulement si $w\in L$, c'est-à-dire si et seulement si
$\delta_B^*(q_{0,B},w) \in F_B$ autrement dit $\psi([w]) \in F_B$.  Et
enfin, pour $w\in\Sigma^*$ et $x\in \Sigma$, on a $\psi(\delta([w],x))
= \psi([wx]) = \delta_B^*(q_{0,B},wx) = \delta_B(\delta_B^*(q_0,w), x)
= \delta_B(\psi([w]), x)$.  Bref, $\psi$ préserve l'état initial, les
états finaux, et la relation de transition : c'est donc bien un
isomorphisme d'automates (un renommage des états).
\end{proof}

\thingy Ce théorème affirme donc qu'il existe (à renommage des états
près) un unique DFA (complet) ayant un nombre minimal d'états parmi
ceux qui reconnaissent le langage rationnel $L$ : on l'appelle
\defin[minimal (automate)]{automate minimal} ou \index{canonique
  (automate)|see{minimal}}\textbf{automate canonique} du langage $L$.  La
démonstration ci-dessus en donne une construction à partir d'une
relation d'équivalence, mais cette démonstration n'est pas
algorithmique : on va voir comment on peut le construire de façon
algorithmique à partir d'un DFA quelconque qui reconnaît $L$.

\begin{prop}\label{dfa-minimization}
Soit $B = (Q_B, q_{0,B}, F_B, \delta_B)$ un DFA (complet !)
reconnaissant un langage $L$, et dont tous les états sont accessibles
(cf. \ref{definition-dfa-accessible-state}).  Alors l'automate
minimal $A$ de $L$ peut s'obtenir en fusionnant dans $B$ chaque
classe d'équivalence pour la relation d'équivalence $\equiv$ définie
sur $Q_B$ par
\[
q \equiv q' \;\liff\; \forall t\in\Sigma^*\,(\delta^*_B(q,t)\in F_B \liff \delta^*_B(q',t)\in F_B)
\]
(Fusionner chaque classe d'équivalence signifie qu'on construit
l'automate $A$ dont l'ensemble d'états est l'ensemble $Q_A :=
Q_B/{\equiv}$ des classes d'équivalence, l'état initial $q_{0,A}$ est
la classe $[q_{0,B}]$ de $q_{0,B}$ pour $\equiv$, les états finaux
sont les classes contenant au moins un état final et qui d'ailleurs
sont entièrement constituées d'états finaux, et la relation de
transition est donnée par $\delta_A([q],x) = [\delta_B(q,x)]$ en
notant $[q]$ la classe de $q$ pour $\equiv$.)

De plus, cet automate $A$ (ou de façon équivalente, la
relation $\equiv$) peut se déduire algorithmiquement de $B$.
\end{prop}
\begin{proof}
On a vu dans le cours de la démonstration de \ref{myhill-nerode} que
l'automate minimal a pour ensemble d'états $\Sigma^*/{\equiv_L}$ où
$\equiv_L$ désigne la relation d'équivalence (alors notée $\equiv$)
définie par $u \mathrel{\equiv_L} v$ lorsque $\forall t\in \Sigma^*
((ut\in L) \liff (vt \in L))$.  Si $q,q' \in Q_B$, comme $q,q'$ sont
accessibles, il existe $w,w'\in\Sigma^*$ tels que
$\delta_B^*(q_{0,B},w) = q$ et $\delta_B^*(q_{0,B},w') = q'$ ; et on a
$q \equiv q'$ si et seulement si $\delta^*_B(q,t)\in F_B \liff
\delta^*_B(q',t)\in F_B$ pour tout $t\in\Sigma^*$, c'est-à-dire
$\delta^*_B(q_{0,B},wt)\in F_B \liff \delta^*_B(q_{0,B},w't)\in F_B$,
c'est-à-dire $wt\in L \liff w't\in L$, autrement dit, $w
\mathrel{\equiv_L} w'$.  L'application $\varphi$ qui à un état $[w]_L$
de l'automate minimal associe la classe de $\delta_B^*(q_{0,B},w)$
pour $\equiv$ est donc une bijection, et comme dans la démonstration
de \ref{myhill-nerode} on vérifie qu'elle préserve l'état initial, les
états finaux, et la relation de transition.  On obtient donc bien
l'automate minimal en fusionnant chaque classe d'équivalence
pour $\equiv$.

Montrons maintenant comment on peut construire $\equiv$
algorithmiquement.  Pour cela, on va définir des relations
d'équivalence $\equiv_i$ pour $i\in\mathbb{N}$ par
\[
q \mathrel{\equiv_i} q' \;\liff\; \forall t\, (|t|\leq i \limp (\delta^*_B(q,t)\in F_B \liff \delta^*_B(q',t)\in F_B))
\]
C'est-à-dire par récurrence
\[
\begin{aligned}
q \mathrel{\equiv_0} q' \;&\liff\; (q\in F_B \liff q'\in F_B)\\
q \mathrel{\equiv_{i+1}} q' \;&\liff\; (q\mathrel{\equiv_{i}}q' \hbox{~et~} \forall x\in\Sigma\,(\delta_B(q,x) \mathrel{\equiv_i} \delta_B(q',x)))\\
\end{aligned}
\]
(la première équivalence vient de ce que $\delta^*_B(q,\varepsilon) =
q$, et la seconde de ce que $\delta^*_B(q,xt) =
\delta^*_B(\delta_B(q,x), t)$).

Il est trivial que $q \mathrel{\equiv_{i+1}} q'$ implique $q
\mathrel{\equiv_i} q'$, c'est-à-dire que $\equiv_{i+1}$ est plus fine
que $\equiv_i$, mais $\equiv$ est plus fine que toutes, et comme elle
n'a qu'un nombre fini de classes d'équivalence, le nombre de classes
ne peut croître strictement qu'un nombre fini de fois, et il doit donc
stationner : il existe $i$ tel que $({\equiv_{i+1}}) = ({\equiv_i})$,
et à ce moment-là, la seconde équivalence de la récurrence montre que
$({\equiv_j}) = ({\equiv_i})$ pour tout $j\geq i$ et donc $(\equiv) =
({\equiv_i})$.

On peut donc calculer $\equiv$ selon l'algorithme suivant : calculer
$\equiv_0$, et par récurrence calculer les $\equiv_i$ jusqu'à ce que
la relation ne change plus, $({\equiv_{i+1}}) = ({\equiv_i})$, auquel
cas la dernière relation calculée est la relation recherchée $(\equiv)
= ({\equiv_i})$, et l'automate minimal $A$ s'obtient en fusionnant
chaque classe d'équivalence pour $\equiv$.
\end{proof}

\thingy L'algorithme décrit par la proposition \ref{dfa-minimization}
porte le nom d'algorithme \index{Moore (algorithme
  de)|see{minimisation}}\textbf{de Moore} ou \defin[minimisation
  (algorithme de)]{de minimisation} ou \textbf{de réduction}.

Voici comment on peut le mettre en œuvre de façon plus concrète :
\begin{itemize}
\item s'assurer qu'on a affaire à un DFA \underline{\emph{complet sans
    état inaccessible}} (si nécessaire, déterminiser l'automate s'il
  n'est pas déterministe, le compléter s'il est incomplet, et
  supprimer les états inaccessibles s'il y en a) ;
\item appeler $\Pi$ la partition des états en deux classes : les
  états finaux d'un côté, et les non-finaux de l'autre ;
\item répéter l'opération suivante tant que la partition $\Pi$
  change : pour chaque classe $C$ de $\Pi$ et chaque lettre $x \in
  \Sigma$, s'il existe deux états $q,q' \in C$ tels que $\delta(q,x)$
  et $\delta(q',x)$ ne tombent pas dans la même classe de $\Pi$,
  séparer cette classe $C$ selon la classe de $\delta(\tiret,x)$
  (autrement dit, $q,q'$ restent dans la même classe pour la nouvelle
  partition lorsqu'ils sont dans la même classe pour l'ancienne et que
  pour chaque lettre $x$ leurs images $\delta(q,x)$ et $\delta(q',x)$
  sont aussi dans la même classe) ;
\item si $\Pi$ est la (dernière) partition ainsi obtenue, l'ensemble
  des états de l'automate construit est l'ensemble des classes
  de $\Pi$, l'état initial est la classe de l'état initial, les états
  finaux sont les classes qui contiennent un état final (ils sont
  alors forcément tous finaux), et la fonction de transition est
  obtenue en prenant la fonction de transition sur un représentant
  quelconque de la classe (la classe ne doit pas dépendre du
  représentant).
\end{itemize}

\smallskip

La dernière étape (construction de l'automate) permet de vérifier
qu'on a correctement terminé l'étape précédente (raffinement de la
partition) : si deux états dans la même classe ont une transition
sortante d'étiquette $x$ et qui mènent vers des classes différentes,
c'est que ces états auraient dû être séparés.  Il est donc utile de
refaire un contrôle à ce niveau.

\thingy À titre d'exemple d'exécution de l'algorithme de minimisation,
considérons l'automate suivant :

\begin{center}
%%% begin example7 %%%

\begin{tikzpicture}[>=latex,line join=bevel,automaton]
%%
\node (q1) at (98bp,18bp) [draw,circle,state] {$1$};
  \node (q0) at (18bp,67bp) [draw,circle,state,initial] {$0$};
  \node (q3) at (186bp,18bp) [draw,circle,state] {$3$};
  \node (q2) at (98bp,105bp) [draw,circle,state] {$2$};
  \node (q5) at (266bp,67bp) [draw,circle,state,final] {$5$};
  \node (q4) at (186bp,105bp) [draw,circle,state,final] {$4$};
  \draw [->] (q2) ..controls (125.5bp,78.193bp) and (148.96bp,54.459bp)  .. node[auto] {$c$} (q3);
  \draw [->] (q0) ..controls (45.307bp,79.816bp) and (60.14bp,87.042bp)  .. node[auto] {$a$} (q2);
  \draw [->] (q2) to[loop above] node[auto] {$a$} (q2);
  \draw [->] (q5) to[loop below] node[auto] {$a,b,c$} (q5);
  \draw [->] (q4) ..controls (213.31bp,92.184bp) and (228.14bp,84.958bp)  .. node[auto] {$c$} (q5);
  \draw [->] (q1) ..controls (128.25bp,18bp) and (144.18bp,18bp)  .. node[auto] {$a,c$} (q3);
  \draw [->] (q4) to[loop above] node[auto] {$a,b$} (q4);
  \draw [->] (q1) to[loop below] node[auto] {$b$} (q1);
  \draw [->] (q3) ..controls (213bp,34.334bp) and (228.89bp,44.318bp)  .. node[auto] {$b$} (q5);
  \draw [->] (q3) to[loop below] node[auto] {$a,c$} (q3);
  \draw [->] (q0) ..controls (45.002bp,50.666bp) and (60.894bp,40.682bp)  .. node[auto] {$c$} (q1);
  \draw [->] (q0) to[loop above] node[auto] {$b$} (q0);
  \draw [->] (q2) ..controls (128.25bp,105bp) and (144.18bp,105bp)  .. node[auto] {$b$} (q4);
%
\end{tikzpicture}

%%% end example7 %%%
\end{center}

Cet automate est bien déterministe, complet et sans état inaccessible.
Dans un premier temps, on partitionne les états en états non-finaux et
finaux, soit $\{0,1,2,3\}$ d'un côté et $\{4,5\}$ de l'autre.
Ensuite, on doit séparer la classe $\{0,1,2,3\}$ en deux selon que la
transition étiquetée par $b$ tombe dans la classe $\{0,1,2,3\}$
elle-même ou bien dans la classe $\{4,5\}$ : on arrive donc à trois
classes, $\{0,1\}$, $\{2,3\}$ et $\{4,5\}$.  Enfin, on doit séparer la
classe $\{0,1\}$ en deux selon que la transition étiquetée par $c$
tombe dans la classe $\{0,1\}$ elle-même ou bien dans la classe
$\{2,3\}$ : on arrive alors à quatre classes, $\{0\}$, $\{1\}$,
$\{2,3\}$ et $\{4,5\}$, et à l'automate minimal suivant :

\begin{center}
%%% begin example7m %%%

\begin{tikzpicture}[>=latex,line join=bevel,automaton]
%%
\node (q1) at (98bp,18bp) [draw,circle,state] {$1$};
  \node (q0) at (18bp,76bp) [draw,circle,state,initial] {$0$};
  \node (q23) at (190bp,76bp) [draw,circle,state] {$2\equiv 3$};
  \node (q45) at (278bp,76bp) [draw,circle,state,final] {$4\equiv 5$};
  \draw [->] (q1) to[loop below] node[auto] {$b$} (q1);
  \draw [->] (q23) to[loop above] node[auto] {$a,c$} (q23);
  \draw [->] (q1) ..controls (126.79bp,35.911bp) and (146.9bp,48.867bp)  .. node[auto] {$a,c$} (q23);
  \draw [->] (q0) ..controls (48.315bp,77.182bp) and (65.168bp,77.757bp)  .. (80bp,78bp) .. controls (106.51bp,78.434bp) and (136.64bp,77.795bp)  .. node[auto] {$a$} (q23);
  \draw [->] (q23) ..controls (222.17bp,76bp) and (234.88bp,76bp)  .. node[auto] {$b$} (q45);
  \draw [->] (q0) ..controls (44.591bp,56.971bp) and (61.407bp,44.466bp)  .. node[auto] {$c$} (q1);
  \draw [->] (q0) to[loop above] node[auto] {$b$} (q0);
  \draw [->] (q45) to[loop above] node[auto] {$a,b,c$} (q45);
%
\end{tikzpicture}

%%% end example7m %%%
\end{center}

\medskip

Il est intéressant de voir comment des petits changements sur
l'automate initial modifient la minimisation.  Si on fait pointer la
transition étiquetée par $c$ de l'état $1$ vers lui-même (au lieu
d'aller vers l'état $3$), c'est-à-dire :

\begin{center}
%%% begin example7b %%%

\begin{tikzpicture}[>=latex,line join=bevel,automaton]
%%
\node (q1) at (98bp,18bp) [draw,circle,state] {$1$};
  \node (q0) at (18bp,67bp) [draw,circle,state,initial] {$0$};
  \node (q3) at (178bp,18bp) [draw,circle,state] {$3$};
  \node (q2) at (98bp,105bp) [draw,circle,state] {$2$};
  \node (q5) at (258bp,67bp) [draw,circle,state,final] {$5$};
  \node (q4) at (178bp,105bp) [draw,circle,state,final] {$4$};
  \draw [->] (q2) ..controls (123.52bp,77.652bp) and (143.78bp,55.049bp)  .. node[auto] {$c$} (q3);
  \draw [->] (q0) ..controls (45.307bp,79.816bp) and (60.14bp,87.042bp)  .. node[auto] {$a$} (q2);
  \draw [->] (q2) to[loop above] node[auto] {$a$} (q2);
  \draw [->] (q5) to[loop below] node[auto] {$a,b,c$} (q5);
  \draw [->] (q4) ..controls (205.31bp,92.184bp) and (220.14bp,84.958bp)  .. node[auto] {$c$} (q5);
  \draw [->] (q1) ..controls (126.11bp,18bp) and (138.58bp,18bp)  .. node[auto] {$a$} (q3);
  \draw [->] (q4) to[loop above] node[auto] {$a,b$} (q4);
  \draw [->] (q1) to[loop below] node[auto] {$b,c$} (q1);
  \draw [->] (q3) ..controls (205bp,34.334bp) and (220.89bp,44.318bp)  .. node[auto] {$b$} (q5);
  \draw [->] (q3) to[loop below] node[auto] {$a,c$} (q3);
  \draw [->] (q0) ..controls (45.002bp,50.666bp) and (60.894bp,40.682bp)  .. node[auto] {$c$} (q1);
  \draw [->] (q0) to[loop above] node[auto] {$b$} (q0);
  \draw [->] (q2) ..controls (126.11bp,105bp) and (138.58bp,105bp)  .. node[auto] {$b$} (q4);
%
\end{tikzpicture}

%%% end example7b %%%
\end{center}

\noindent alors la minimisation ne sépare pas les états $0$ et $1$, et
on obtient (le même automate qu'en \ref{discussion-example2}, à
savoir) :

\begin{center}
%%% begin example7bm %%%

\begin{tikzpicture}[>=latex,line join=bevel,automaton]
%%
\node (q45) at (198bp,21bp) [draw,circle,state,final] {$4\equiv 5$};
  \node (q01) at (22bp,21bp) [draw,circle,state,initial] {$0\equiv 1$};
  \node (q23) at (110bp,21bp) [draw,circle,state] {$2\equiv 3$};
  \draw [->] (q45) to[loop above] node[auto] {$a,b,c$} (q45);
  \draw [->] (q23) ..controls (142.17bp,21bp) and (154.88bp,21bp)  .. node[auto] {$b$} (q45);
  \draw [->] (q23) to[loop above] node[auto] {$a,c$} (q23);
  \draw [->] (q01) to[loop above] node[auto] {$b,c$} (q01);
  \draw [->] (q01) ..controls (54.17bp,21bp) and (66.885bp,21bp)  .. node[auto] {$a$} (q23);
%
\end{tikzpicture}

%%% end example7bm %%%
\end{center}

En revanche, si on change l'automate pour rendre l'état $4$ non-final
(avec ou sans la modification précédemment évoquée), la minimisation
aboutit sur la partition triviale en six états, c'est-à-dire que
l'automate est, en fait, déjà minimal.

{\footnotesize\thingy\textbf{Cas dégénérés :} Si aucun état d'un DFA
  n'est final (c'est-à-dire s'il reconnaît le langage
  vide $\varnothing$), alors l'algorithme de minimisation termine
  immédiatement avec une seule classe d'équivalence, et donne donc un
  automate minimal ayant un unique état, initial mais non final, et
  des transitions étiquetées par toutes les lettres de cet état vers
  lui-même.  Si \emph{tous} les états d'un DFA sont finaux (il
  reconnaît le langage $\Sigma^*$ de tous les mots), de même,
  l'algorithme de minimisation termine immédiatement avec une seule
  classe d'équivalence, et donne donc un automate minimal ayant un
  unique état, à la fois initial et final (et toujours des transitions
  étiquetées par toutes les lettres de cet état vers lui-même).\par}

\bigbreak

Énonçons ici le fait évoqué plus haut comme application de
l'algorithme de minimisation :

\begin{cor}\label{equivalence-of-regexps-is-decidable}
On peut décider algorithmiquement si deux automates finis (de
n'importe quelle sorte), ou deux expressions rationnelles, ou un
automate et une expression rationnelle, sont équivalents (au sens de
reconnaître/dénoter le même langage).
\end{cor}
\begin{proof}
D'après ce qu'on a déjà vu, et quitte à transformer une expression
rationnelle en εNFA (cf. \ref{rational-languages-are-recognizable}),
et quitte à éliminer les ε-transitions
(cf. \ref{removal-of-epsilon-transitions}) et à déterminiser
(cf. \ref{determinization-of-nfa}), on peut construire
algorithmiquement un DFA reconnaissant le même langage que chacune des
deux données.  La question devient donc de savoir si deux DFA
reconnaissent le même langage.  Or d'après \ref{dfa-minimization}, on
sait transformer un DFA en DFA minimal reconnaissant le même langage,
et d'après \ref{myhill-nerode} on sait que ce DFA est unique à
renumérotation près des états.  On est donc ramené au problème
suivant : donnés deux DFA $A$ et $A'$ (complets et sans états
inaccessibles), trouver s'ils sont le même à renumérotation près
(i.e., s'ils sont isomorphes).

La correspondance entre états de $A$ et de $A'$ peut se construire
état par état : on fait correspondre l'état initial de $A$ à celui
de $A'$, puis pour chaque état $q$ de $A$ mis en correspondance avec
un état $q'$ de $A'$, on fait correspondre chacun des états
$\delta(q,x)$ avec chacun des états $\delta'(q',x)$ où $x$ parcourt
les lettres de l'alphabet.  Si on aboutit ainsi à une contradiction
(deux états de l'un des automates veulent être mis en correspondance
avec le même état de l'autre), on renvoie un échec ; sinon, tous les
états de $A$ et ceux de $A'$ seront mis en correspondance bijective,
et les langages reconnus sont les mêmes.
\end{proof}


\section{Grammaires hors contexte}\label{section-context-free-grammars}

\setcounter{comcnt}{0}

\thingy Alors que les langages et expressions rationnelles servent,
dans le monde informatique, principalement à définir des outils de
recherche de « motifs » (pour la recherche et le remplacement dans un
texte, la validation d'entrées, la syntaxe de bas niveau, etc.), les
grammaires formelles, dont les plus importantes sont les grammaires
\emph{hors contexte} que nous allons maintenant considérer, ont pour
principal intérêt de définir des \emph{syntaxes structurées}, par
exemple la syntaxe d'un langage informatique (typiquement un langage
de programmation).  En fait, l'étude des grammaires formelles en
général, et des grammaires hors contexte en particulier, a été
démarrée\footnote{Certains font remonter leur origine bien plus loin :
  on peut trouver une forme de grammaire hors contexte dans la manière
  dont l'Indien Pāṇini (probablement au IV\textsuperscript{e} siècle
  av. notre ère) décrit la grammaire du sanskrit.} en 1956 par le
linguiste (et activiste politique) Noam Chomsky pour servir dans
l'analyse des langues naturelles ; mais c'est plus en informatique
qu'en linguistique qu'elles ont trouvé leur utilité, à commencer
surtout par la définition de la syntaxe du langage ALGOL 60.

\smallskip

Cette fois-ci, on ne s'intéressera pas simplement au langage défini
(par la grammaire hors contexte, dit langage « algébrique »), mais
aussi à la manière dont ces mots s'obtiennent par la grammaire, et
donc, à la manière d'\emph{analyser} (en anglais, \emph{to parse}) un
mot / programme en une structure de données qui le représente : pour
cela, on va définir la notion d'\emph{arbre d'analyse}.

\thingy Un exemple typique d'usage de grammaire hors contexte pour
définir la syntaxe d'un langage de programmation hypothétique pourrait
ressembler à ceci :
\[
\begin{aligned}
\mathit{Statement} &\rightarrow \mathtt{begin}\ \mathit{Block}\ \mathtt{end}\\
\mathit{Statement} &\rightarrow \mathtt{if}\ \mathit{Expression}\ \mathtt{then}\ \mathit{Block}\ \mathtt{fi}\\
\mathit{Statement} &\rightarrow \mathtt{if}\ \mathit{Expression}\ \mathtt{then}\ \mathit{Block}\ \mathtt{else}\ \mathit{Block}\ \mathtt{fi}\\
\mathit{Block} &\rightarrow \varepsilon\\
\mathit{Block} &\rightarrow \mathit{Statement}\ \mathit{Block}\\
\end{aligned}
\]
Il faut comprendre ces règles de la manière suivante : « pour définir
une instruction ($\mathit{Statement}$), on peut mettre un
$\mathtt{begin}$ suivi d'un bloc ($\mathit{Block}$) suivi d'un
$\mathtt{end}$, ou bien un $\mathtt{if}$ suivi d'une expression
($\mathit{Expression}$) suivi d'un $\mathtt{then}$ suivi d'un bloc,
éventuellement suivie encore d'un $\mathtt{else}$ et d'un nouveau
bloc, et enfin un $\mathtt{fi}$ ; pour définir un bloc
($\mathit{Block}$), on peut soit ne rien mettre du tout, soit mettre
une instruction suivi d'un nouveau bloc ».

\smallskip

Notre but va être d'expliquer quel genre de règles de ce genre on peut
autoriser, comment elles se comportent, quels types de langages elles
définissent, et comment on peut analyser (essentiellement, retrouver
la manière dont on a construit) un texte produit par une application
de telles règles.

\thingy Dans un contexte informatique, l'usage des grammaires
formelles se fait à un niveau différent de celui des chaînes de
caractères où nous nous sommes placés dans les sections précédentes :
cette fois, les lettres de notre alphabet ne seront généralement pas
des caractères mais des \defin[token]{tokens} du langage dont on
définit la grammaire, un « token » pouvant être une unité variable
selon le langage, mais généralement soit un caractère spécial, soit un
mot-clé du langage (\texttt{begin}, \texttt{end}, \texttt{for}, etc.,
constituent donc généralement un unique token dans l'analyse, et ne
sont pas divisés en leurs lettres individuelles), soit encore une
unité (« nombre », « identificateur », éventuellement « commentaire »)
qui aura été découpée dans une première phase d'analyse, appelée
« analyse lexicale ».

Nous ne rentrerons pas ici dans ces considérations, et nous nous en
tiendrons à l'approche mathématique où on continue à avoir affaire à
des lettres d'un alphabet $\Sigma$ abstrait.


\subsection{Définition, notations et premier exemple}\label{subsection-context-free-grammars}

\thingy\label{definition-context-free-grammar} Une \index{hors
  contexte (grammaire)|see{grammaire hors contexte}}\defin{grammaire
  hors contexte} (on dit parfois aussi « sans contexte » ; en anglais,
« context-free grammar » ou \index{CFG|see{grammaire hors
    contexte}}« CFG » ; ou encore grammaire de \textbf{type 2}) sur un
alphabet $\Sigma$ est la donnée
\begin{itemize}
\item d'un second alphabet $N$, disjoint de $\Sigma$, appelé ensemble
  des \defin[nonterminal]{nonterminaux},
\item d'un élément $S \in N$ appelé \defin[axiome (d'une
  grammaire)]{axiome} ou \index{initial
  (symbole)|see{axiome}}\textbf{symbole initial} de la grammaire,
\item d'un ensemble \emph{fini} $R \subseteq N \times (\Sigma\cup
  N)^*$ de couples $(T,\alpha)$ où $T$ est un nonterminal et $\alpha$
  un mot sur l'alphabet $\Sigma\cup N$, les éléments $(T,\alpha)$ de
  $R$ étant appelés les \index{règle (d'une
    grammaire)|see{production}}\textbf{règles} ou
  \defin[production]{productions} de la grammaire.
\end{itemize}

\thingy Une grammaire hors contexte fait donc intervenir deux
alphabets : l'alphabet $\Sigma$ sur lequel sera le langage (encore à
définir), et l'alphabet $N$ (disjoint de $\Sigma$) qui intervient dans
la définition de la grammaire.  Pour fixer la terminologie, on
appellera \defin{symbole} un élément de $\Sigma \cup N$, ces symboles
étant dits \defin[terminal]{terminaux} lorsqu'ils appartiennent
à $\Sigma$ (ou simplement « lettres »), et \textbf{nonterminaux}
lorsqu'ils appartiennent à $N$.  Un mot sur $\Sigma\cup N$
(c'est-à-dire, une suite finie de symboles) sera appelé
\defin{pseudo-mot} (ou parfois une « forme »), tandis que le terme
« \index{mot}mot » sans précision supplémentaire sera utilisé pour
désigner un mot sur $\Sigma$ (autrement dit, un mot est un pseudo-mot
ne comportant que des symboles terminaux).

\smallskip

Pour redire les choses autrement, les symboles terminaux sont les
lettres des mots du langage qu'on cherche à définir ; les symboles
nonterminaux sont des symboles qui servent uniquement à titre
transitoire dans l'utilisation de la grammaire, et qui sont destinés à
disparaître finalement.  Un « pseudo-mot » est un mot pouvant contenir
des nonterminaux, tandis qu'un « mot » sans précision du contraire ne
contient que des terminaux.

\medskip

{\footnotesize Typographiquement, on aura tendance à utiliser des
  lettres minuscules pour désigner les symboles terminaux, et
  majuscules pour désigner les symboles nonterminaux ; et on aura
  tendance à utiliser des lettres grecques minuscules pour désigner
  les pseudo-mots.  Il n'est malheureusement pas possible d'être
  complètement systématique.\par}

\thingy Intuitivement, il faut comprendre la grammaire de la manière
suivante.  Les règles $(T,\alpha)$, où on rappelle que $T$ est un
symbole nonterminal et $\alpha$ un pseudo-mot, doivent se comprendre
comme « le symbole $T$ peut être remplacé par $\alpha$ » (et on notera
$T \rightarrow \alpha$, cf. ci-dessous).  On part de l'axiome $S$, et
on peut effectuer librement des substitutions $T \rightarrow \alpha$
(consistant à remplacer un nonterminal $T$ par un pseudo-mot $\alpha$)
autorisées par les règles : le langage défini par la grammaire est
l'ensemble de tous les mots (ne comportant plus aucun nonterminal)
qu'on peut obtenir de la sorte.

Rendons maintenant cette définition précise :

\thingy\label{derivations-and-contexts} Lorsque $G$ est une grammaire
hors contexte comme en \ref{definition-context-free-grammar}, on note
$T \mathrel{\rightarrow_G} \alpha$, ou simplement $T \rightarrow
\alpha$ (lorsque $G$ est clair) pour signifier que $(T,\alpha)$ est
une règle de $G$.

\smallskip

On définit une relation $\Rightarrow$ en posant $\gamma T \gamma'
\Rightarrow \gamma\alpha\gamma'$ pour toute règle $(T,\alpha)$ et tous
pseudo-mots $\gamma,\gamma'$ : autrement dit, formellement, on définit
$\lambda\Rightarrow\xi$ lorsqu'il existe $\gamma,\gamma' \in
(\Sigma\cup N)^*$ et $(T,\alpha) \in R$ (ensemble des règles de $G$)
tels que $\lambda = \gamma T \gamma'$ et $\xi = \gamma\alpha\gamma'$.
Concrètement, $\lambda\Rightarrow\xi$ signifie donc que $\xi$ est
obtenu en remplaçant un nonterminal $T$ du pseudo-mot $\lambda$ par la
partie droite d'une production $T \rightarrow \alpha$ de la
grammaire $G$ : on dit encore que $\lambda \Rightarrow \xi$ est une
\defin{dérivation immédiate} de $G$.  On pourra dire que $T
\rightarrow \alpha$ est la règle \textbf{appliquée} dans cette
dérivation immédiate, que $T$ est le \textbf{symbole réécrit} (souvent
souligné dans l'écriture de la dérivation immédiate), et que $\gamma$
et $\gamma'$ sont respectivement le \index{contexte}\textbf{contexte
  gauche} et le \textbf{contexte droit} de l'application de la
règle\footnote{Pour être extrêmement rigoureux, une dérivation
  immédiate doit comporter la donnée du symbole réécrit (ou de façon
  équivalente, du contexte gauche et/ou droit), car elle ne peut pas
  se déduire de la seule donnée de $\lambda$ et $\xi$ (par exemple,
  dans $XX \Rightarrow XXX$, même si on sait que la règle appliquée
  était $X \rightarrow XX$, on ne peut pas deviner si c'est le $X$ de
  gauche ou de droite qui a été réécrit : il faut donc écrire
  $\underline{X}X \Rightarrow XXX$ ou bien $X\underline{X} \Rightarrow
  XXX$, en soulignant le symbole réécrit, pour distinguer les deux).
  De même, dans une dérivation, on devrait inclure la donnée du
  symbole réécrit à chaque étape.}.  Bien sûr, si nécessaire, on
précisera la grammaire appliquée en écrivant $\lambda
\mathrel{\Rightarrow_G} \xi$ au lieu de simplement
$\lambda\Rightarrow\xi$.

\smallskip

Une suite de pseudo-mots $(\lambda_0,\ldots,\lambda_n)$ telle que
\[
\lambda_0 \Rightarrow \lambda_1 \Rightarrow \cdots \Rightarrow \lambda_n
\]
autrement dit $\lambda_{i-1} \Rightarrow \lambda_i$ pour chaque $1\leq
i\leq n$, est appelée \defin{dérivation} de $\lambda_n$ à partir de
$\lambda_0$ dans la grammaire $G$, et on note $\lambda_0
\mathrel{\Rightarrow^*} \lambda_n$ pour indiquer son existence
(c'est-à-dire, si on préfère, que la relation $\Rightarrow^*$ est la
clôture réflexive-transitive de $\Rightarrow$, i.e., la plus petite
relation binaire réflexive et transitive contenant $\Rightarrow$).
Remarquons que $n=0$ est permis, autrement dit $\lambda
\mathrel{\Rightarrow^*} \lambda$ pour tout pseudo-mot $\lambda$.  Bien
sûr, si nécessaire, on précisera la grammaire appliquée en écrivant
$\lambda \mathrel{\Rightarrow^*_G} \xi$ au lieu de simplement $\lambda
\mathrel{\Rightarrow^*} \xi$.

Concrètement, $\lambda \mathrel{\Rightarrow^*} \xi$ signifie donc
qu'on peut passer de $\lambda$ à $\xi$ en effectuant une suite
(finie !) de dérivations immédiates, c'est-à-dire en remplaçant à
chaque étape un nonterminal $T$ par la partie droite $\alpha$ d'une
règle $T \rightarrow \alpha$ de la grammaire $G$.

\smallskip

Il va de soi qu'un pseudo-mot qui ne comporte que des terminaux, i.e.,
qui est en fait un mot (sur $\Sigma$), ne peut pas être dérivé plus
loin.  Ceci justifie au moins en partie la terminologie de
« terminal ».

\thingy Le \defin[engendré (langage)]{langage engendré} $L(G)$ par une
grammaire hors contexte $G$ est l'ensemble des mots $w$ (ne comportant
plus de nonterminaux !) qui peuvent s'obtenir par dérivation à partir
de l'axiome $S$ de $G$, autrement dit :
\[
L(G) = \{w \in \Sigma^* : S \mathrel{\Rightarrow^*} w\}
\]

\smallskip

Un langage qui peut s'écrire sous la forme $L(G)$ où $G$ est une
grammaire hors contexte est appelé \index{hors contexte
  (langage)|see{algébrique}}\textbf{langage hors contexte} ou
\defin[algébrique (langage)]{algébrique}.

\smallskip

Deux grammaires hors contexte $G$ et $G'$ sont dites \defin[faiblement
  équivalentes (grammaires)]{faiblement équivalentes} ou
\index{langage-équivalentes (grammaires)|see{faiblement
    équivalentes}}\textbf{langage-équivalentes} lorsqu'elles
engendrent le même langage ($L(G) = L(G')$).

\thingy\label{basic-example-context-free-grammar} À titre d'exemple,
considérons la grammaire sur l'alphabet $\Sigma = \{a,b\}$ donnée par
\[
\begin{aligned}
S &\rightarrow aSb\\
S &\rightarrow \varepsilon\\
\end{aligned}
\]
où implicitement $S$ est l'axiome et le seul nonterminal : autrement
dit, $G$ est donnée par $N = \{S\}$, d'axiome $S$, et de règles de
production $(S,aSb)$ et $(S,\varepsilon)$ (où $\varepsilon$, bien
entendu, désigne le mot vide).  Un pseudo-mot pour cette grammaire est
un mot sur l'alphabet $\Sigma\cup N = \{a,b,S\}$, et une dérivation
immédiate consiste \emph{soit} à ajouter un $a$ et un $b$ à gauche et
à droite d'un $S$ dans un pseudo-mot, \emph{soit} à retirer un $S$.

Dans ces conditions, il est clair qu'une dérivation partant de
l'axiome $S$ est constituée d'un certain nombre d'applications de la
première règle suivi d'au plus une application de la seconde (puisque
le nombre d'occurrences de $S$ va alors tomber de $1$ à $0$ et on ne
pourra plus dériver).  Autrement dit, une telle dérivation prend la
forme
\[
S \Rightarrow aSb \Rightarrow aaSbb \Rightarrow \cdots \Rightarrow a^n
S b^n
\]
ou éventuellement
\[
S \Rightarrow aSb \Rightarrow aaSbb \Rightarrow \cdots \Rightarrow a^n
S b^n \Rightarrow a^n b^n
\]
Le langage $L(G)$ défini par la grammaire est donc $\{a^n b^n :
n\in\mathbb{N}\}$.

\smallskip

On a vu en \ref{example-of-pumping-lemma} que ce langage n'est pas
rationnel : il existe donc des langages algébriques qui ne sont pas
rationnels.  En revanche, pour ce qui est de la réciproque, on verra
dans la section suivante que tout langage rationnel est algébrique.

\thingy\label{more-general-formal-grammars}
Mentionnons brièvement qu'il existe des types de grammaires
plus généraux que les grammaires hors contexte.  Les \defin[grammaire
  contextuelle]{grammaires contextuelles} (ou grammaires de
\textbf{type 1}) sont définies par des règles du type $\gamma T
\gamma' \rightarrow \gamma \alpha \gamma'$ où $T$ est un nonterminal,
et $\alpha,\gamma,\gamma'$ des pseudo-mots (= mots sur l'alphabet de
tous les symboles), c'est-à-dire des règles qui autorisent la
réécriture d'un symbole $T$ en $\alpha$ mais uniquement s'il est
entouré d'un certain contexte ($\gamma$ à gauche, $\gamma'$ à droite).
Les \defin[grammaire syntagmatique]{grammaires syntagmatiques
  générales} (ou grammaires de \textbf{type 0}) sont définies par des
règles de réécriture $\lambda \rightarrow \mu$ où $\lambda,\mu$ sont
des pseudo-mots quelconques, c'est-à-dire qu'elles permettent la
réécriture de multiples symboles à la fois.  Dans tous les cas, le
langage défini par la grammaire est l'ensemble de tous les mots (sans
nonterminaux) qui peuvent s'obtenir par application des règles de
remplacement à partir de l'axiome.

\smallskip

Les grammaires de types 0 et 1, avec celles de type 2 c'est-à-dire
hors contexte, et celles de type 3 (= régulières) qui seront définies
en \ref{regular-grammar} ci-dessous, forment une hiérarchie appelée
\defin[Chomsky (hiérarchie de)]{hiérarchie de Chomsky} : plus le
numéro du type est élevé plus la grammaire est contrainte et plus la
classe de langages définie est petite.


\subsection{Langages rationnels et algébriques, stabilité}

\begin{prop}\label{rational-languages-are-algebraic}
Tout langage rationnel est algébrique.  Mieux, on peut déduire
algo\-ri\-thmi\-quement une grammaire hors contexte $G$ d'un εNFA $A$ de
façon à avoir $L(G) = L(A)$.
\end{prop}
\begin{proof}
Soit $A$ un εNFA : on sait qu'on peut supposer sans perte de
généralité qu'il a un unique état initial $q_0$ (quitte à en créer un
et à ajouter pour tout état $q$ anciennement initial une ε-transition
$q_0 \to q$).  Soit $Q$ son ensemble des états et $\delta \subseteq Q
\times (\Sigma\cup\{\varepsilon\}) \times Q$ sa fonction de
transition.  On construit une grammaire hors contexte $G$ dont
l'ensemble des nonterminaux est $Q$, dont l'axiome est $S := q_0$, et
dont les règles sont les suivantes : (A) pour chaque transition
$(q,t,q') \in \delta$ une règle $q \to tq'$ (on rappelle que $t \in
\Sigma\cup\{\varepsilon\}$), et (B) pour chaque état final $q\in F$
une règle $q\to\varepsilon$.

De même que dans l'exemple \ref{basic-example-context-free-grammar},
une dérivation partant de l'axiome va comporter un certain nombre
d'applications de règles de type (A) suivies éventuellement d'une
unique application d'une règle de type (B) (ce qui est nécessaire pour
faire passer le nombre de nonterminaux de $1$ à $0$).  Autrement dit,
elle prend la forme $q_0 \Rightarrow t_1 q_1 \Rightarrow t_1 t_2 q_2
\Rightarrow \cdots \Rightarrow t_1 t_2\cdots t_n q_n$ où
$(q_{i-1},t_i,q_i) \in \delta$ pour $1\leq i\leq n$, suivie
éventuellement par $t_1\cdots t_n q_n \Rightarrow t_1\cdots t_n$
lorsque $q_n \in F$.  Manifestement, les dérivations de la sorte
(terminant sur un mot sur $\Sigma$) sont en bijection avec les chemins
$q_0 \to \cdots \to q_n$ dans $A$ où $q_n \in F$, et le mot $t_1\cdots
t_n$ dérivé est précisément le mot formé en concaténant les
étiquettes du chemin.  On a donc bien $L(G) = L(A)$.
\end{proof}

\thingy\label{regular-grammar} Une grammaire hors contexte telle que
construite dans la preuve de \ref{rational-languages-are-algebraic}
est dite \index{régulière (grammaire)|see{grammaire
    régulière}}\defin[grammaire régulière]{régulière} : autrement dit,
il s'agit d'une grammaire ayant uniquement des règles de deux formes :
(A) $Q \rightarrow tQ'$ où $t \in \Sigma\cup\{\varepsilon\}$, et
(B) $Q \rightarrow \varepsilon$.  Les langages définis par des
grammaires régulières sont donc exactement les langages rationnels.

{\footnotesize La définition des grammaires régulières varie un peu
  d'auteur en auteur, ces différences n'étant pas très importantes.
  Certains auteurs imposent $t \in \Sigma$ dans les règles de
  type (A), ce qui revient à imposer qu'on part d'un NFA sans
  ε-transition dans la construction ; et certains autorisent (ou
  autorisent uniquement) dans le type (B) les règles de la forme $Q
  \rightarrow x$ avec $x\in\Sigma$, ce qui impose des petits
  changements dans la construction ci-dessus.  Dans tous les cas, il
  reste vrai que les langages définis par des grammaires régulières
  sont exactement les langages rationnels.\par}

Les grammaires régulières sont également appelées grammaires de
\textbf{type 3}.

\begin{prop}\label{cfg-union-concatenation-and-star}
Si $L_1,L_2$ sont des langages algébriques (sur un même
alphabet $\Sigma$), alors la \index{union (de langages)}réunion $L_1
\cup L_2$ est algébrique ; de plus, une grammaire hors contexte
l'engendrant se déduit algorithmiquement de grammaires hors contexte
engendrant $L_1$ et $L_2$.

Si $L_1,L_2$ sont des langages algébriques (sur un même
alphabet $\Sigma$), alors la \index{concaténation (de
  langages)}concaténation $L_1 L_2$ est algébrique ; de plus, une
grammaire hors contexte l'engendrant se déduit algorithmiquement de
grammaires hors contexte engendrant $L_1$ et $L_2$.

Si $L$ est un langage algébrique, alors \index{étoile de
  Kleene}l'étoile de Kleene $L^*$ est algébrique ; de plus, une
grammaire hors contexte l'engendrant se déduit algorithmiquement d'une
grammaire hors contexte engendrant $L$.
\end{prop}
\begin{proof}
Pour la réunion : supposons que $G_1$ et $G_2$ engendrent $L_1$ et
$L_2$ respectivement, et ont des ensembles de nonterminaux $N_1$ et
$N_2$ disjoints, avec axiomes $S_1$ et $S_2$ respectivement.  On
construit une grammaire $G$ dont l'ensemble des nonterminaux est $N_1
\cup N_2 \cup \{S\}$ où $S$ est un nouveau nonterminal, choisi comme
axiome, et les règles de production de $G$ sont celles de $G_1$,
celles de $G_2$, et les deux règles supplémentaires $S \rightarrow
S_1$ et $S \rightarrow S_2$.  Il est évident qu'une dérivation de $G$
partant de $S$ est soit de la forme $S \Rightarrow S_1$ suivie d'une
dérivation de $G_1$, soit de la forme $S \Rightarrow S_2$ suivie d'une
dérivation de $G_2$ : ainsi, $L(G) = L_1 \cup L_2$.

Pour la concaténation : la construction est presque exactement la
même, mais on prend pour règle supplémentaire $S \rightarrow S_1 S_2$.

Pour l'étoile : si $G$ engendre $L$ et a pour ensemble de nonterminaux
$N$ et pour axiome $S$, on défit une nouvelle grammaire $G'$ dont
l'ensemble des nonterminaux est $N' = N \cup \{S'\}$ où $S'$ est un
nouveau nonterminal, choisi comme axiome, et les règles de production
de $G'$ sont celles de $G$ et les deux règles supplémentaires $S'
\rightarrow SS'$ et $S' \rightarrow \varepsilon$.  Il est facile de se
convaincre qu'une dérivation de $G'$ partant de $S'$ et utilisant $n$
fois la règle $S' \rightarrow SS'$ se ramène à $n$ dérivations de $G$
partant de $S$, et seule la règle $S' \rightarrow \varepsilon$ permet
de faire disparaître le symbole $S'$.  (Les détails sont omis et
peuvent être rendus plus clairs en utilisant la notion d'arbre de
dérivation.)
\end{proof}

\thingy Ceci fournit une nouvelle démonstration (partant d'une
expression rationnelle plutôt que d'un automate)
de \ref{rational-languages-are-algebraic}, tant il est évident que les
langages $\varnothing$, $\{\varepsilon\}$ et $\{x\}$
(pour $x\in\Sigma$) sont algébriques.

À cause de cette construction, on se permettra parfois, dans l'énoncé
des règles d'une grammaire, d'écrire $T \rightarrow \alpha_1 | \cdots
| \alpha_n$ pour réunir en une seule ligne plusieurs règles $T
\rightarrow \alpha_1$ jusqu'à $T \rightarrow \alpha_n$.  Par exemple,
la grammaire présentée en \ref{basic-example-context-free-grammar}
peut s'écrire $S \rightarrow aSb | \varepsilon$.  Il s'agit d'un
simple raccourci d'écriture (comparer avec une convention semblable
faite pour les automates en \ref{graphical-representation-of-dfa}).

{\footnotesize
On pourrait même se permettre l'abus de notation consistant à écrire
$T \rightarrow U^*$ pour $T \rightarrow UT|\varepsilon$ (c'est-à-dire
$T \rightarrow UT$ et $T \rightarrow \varepsilon$), mais il vaut mieux
éviter car cela pourrait aussi bien signifier $T \rightarrow
TU|\varepsilon$, qui engendre le même langage (i.e., elle est
faiblement équivalente à la précédente) mais dont l'analyse peut être
plus problématique.  De façon encore plus générale, on pourrait
imaginer d'autoriser des règles de la forme $T \rightarrow r$ où $T$
est un nonterminal et $r$ une expression rationnelle quelconque sur
l'alphabet $\Sigma \cup N$ : de telles règles pourraient se ramener
aux règles qu'on a autorisées, quitte à introduire de nouveaux
nonterminaux pour les expressions intermédiaires (par exemple, $T
\rightarrow UV^*$ se transformerait par l'introduction d'un nouveau
nonterminal $V'$ en $T \rightarrow UV'$ et $V' \rightarrow
VV'|\varepsilon$).  Nous éviterons ces extensions aux grammaires hors
contexte pour ne pas introduire de confusion.
\par}

\thingy\label{words-deriving-from-another-nonterminal} Il peut être
utile, pour raisonner sur les langages algébriques, d'introduire,
lorsque $G$ est une grammaire hors contexte et $T$ un nonterminal
quelconque, le langage $L(G,T) := \{w \in \Sigma^* : T
\mathrel{\Rightarrow^*_G} w\}$ des mots qui dérivent de $T$,
c'est-à-dire le langage engendré par la grammaire $G'$ identique à $G$
à ceci près que son axiome est $T$.

{\footnotesize
On a alors $L(G) = L(G,S)$ où $S$ est l'axiome de $G$, et chaque règle
de production $T \rightarrow \alpha_1 | \cdots | \alpha_n$ se traduit
par une équation portant sur $L(G,T)$, par exemple $T \rightarrow aUbV
| cW$ implique $L(G,T) = a\, L(G,U)\, b\, L(G,V) \; \cup \; c L(G,W)$.
De cette manière, une grammaire hors contexte donne lieu à un système
d'équations portant sur des langages (par exemple, la grammaire de
\ref{basic-example-context-free-grammar} donne lieu à l'équation $L(G)
= (a\,L(G)\,b) \cup \{\varepsilon\}$ : on peut montrer que la
définition des grammaires hors contexte revient à considérer la plus
petite (au sens de l'inclusion) solution de ce système d'équations.
\par}

\bigbreak

\thingy À la différence des langages rationnels, il \emph{n'est pas
  vrai} en général que l'intersection de deux langages algébriques
soit algébrique : un contre-exemple est fourni par les deux langages
$\{a^i b^i c^j : i,j\in\mathbb{N}\}$ et $\{a^i b^j c^j :
i,j\in\mathbb{N}\}$ (chacun des deux est algébrique : par exemple, le
premier est la concaténation du langage $\{a^i b^i : i\in\mathbb{N}\}$
dont on a vu en \ref{basic-example-context-free-grammar} qu'il était
algébrique, et du langage $\{c\}^*$, qui est algébrique car
rationnel) ; leur intersection $\{a^i b^i c^i : i\in\mathbb{N}\}$
n'est pas algébrique, comme on le démontrera
en \ref{example-of-pumping-lemma-for-algebraic-languages}.

En conséquence, il n'est pas non plus vrai en général que le
complémentaire d'un langage algébrique soit algébrique.

\smallskip

En revanche, le fait suivant, que nous admettons sans démonstration,
peut s'avérer utile :

\begin{prop}\label{intersection-of-algebraic-and-rational}
L'intersection d'un langage \emph{algébrique} et d'un langage
\emph{rationnel} est algébrique.
\end{prop}


\subsection{Autres exemples de grammaires hors contexte}

\thingy\label{example-lr-non-ll-grammar} Sur l'alphabet $\Sigma =
\{a,b\}$, considérons la grammaire (d'axiome $S$)
\[
\begin{aligned}
S &\rightarrow AT\\
A &\rightarrow aA \;|\; \varepsilon\\
T &\rightarrow aTb \;|\; \varepsilon\\
\end{aligned}
\]
Il est clair que le langage $L(G,A)$ des mots qui dérivent de $A$ est
simplement $\{a\}^* = \{a^i : i\in\mathbb{N}\}$ ; et le langage
$L(G,T)$ est $\{a^j b^j : j\in\mathbb{N}\}$ comme
en \ref{basic-example-context-free-grammar}.  Le langage $L(G) =
L(G,A)\, L(G,T)$ est donc $\{a^i b^j : i \geq j\}$, l'ensemble des
mots de la forme $a^i b^j$ pour lesquels $i \geq j$.

On peut aussi considérer la grammaire $G'$ définie par
\[
\begin{aligned}
S &\rightarrow T \;|\; aS\\
T &\rightarrow aTb \;|\; \varepsilon\\
\end{aligned}
\]
Il n'est pas difficile de se convaincre qu'elle engendre le même
langage $L(G)$ que la précédente (elle est donc faiblement équivalente
à $G$).

{\footnotesize (Ce langage est intéressant sur le plan théorique, car
  bien qu'il soit algébrique, et ici défini par une grammaire $G'$
  inambiguë (cf. \ref{ambiguous-grammar}), il ne peut pas être analysé
  par un analyseur LL (cf. \ref{handwaving-on-ll-and-lr}).)\par}

\thingy\label{example-unambiguous-but-not-deterministic-grammar} Sur
l'alphabet $\Sigma = \{a,b\}$, considérons la grammaire (d'axiome $S$)
\[
\begin{aligned}
S &\rightarrow U \;|\; V\\
U &\rightarrow aUb \;|\; ab\\
V &\rightarrow aVbb \;|\; abb\\
\end{aligned}
\]
Il est clair que le langage $L(G,U)$ des mots qui dérivent de $U$ est
$\{a^i b^i : i\geq 1\}$ comme
en \ref{basic-example-context-free-grammar}, et de façon analogue, le
langage $L(G,V)$ des mots qui dérivent de $V$ est $\{a^i b^{2i} :
i\geq 1\}$.  Le langage $L(G) = L(G,U) \cup L(G,V)$ est donc
l'ensemble des mots de la forme $a^i b^j$ où $i\geq 1$ et $j=i$
\emph{ou bien} $j=2i$.

{\footnotesize (Ce langage est intéressant sur le plan théorique, car
  bien qu'il soit algébrique et défini par une grammaire inambiguë
  (cf. \ref{ambiguous-grammar}), il n'est pas analysable par un
  automate à pile déterministe
  (cf. \ref{handwaving-on-stack-automata}).)\par}

\thingy\label{example-for-intrinsic-ambiguity} Sur l'alphabet $\Sigma
= \{a,b,c\}$, considérons la grammaire (d'axiome $S$)
\[
\begin{aligned}
S &\rightarrow UC \;|\; AV\\
A &\rightarrow aA \;|\; \varepsilon\\
C &\rightarrow cC \;|\; \varepsilon\\
U &\rightarrow aUb \;|\; \varepsilon\\
V &\rightarrow bVc \;|\; \varepsilon\\
\end{aligned}
\]
Il est clair que le langage $L(G,A)$ des mots qui dérivent de $A$ est
simplement $\{a\}^* = \{a^i : i\in\mathbb{N}\}$ et que le langage
$L(G,C)$ des mots qui dérivent de $C$ est simplement $\{c\}^* = \{c^j
: j\in\mathbb{N}\}$ ; le langage $L(G,U)$ est $\{a^i b^i :
i\in\mathbb{N}\}$ comme en \ref{basic-example-context-free-grammar},
et de même $L(G,V) = \{b^j c^j : j\in\mathbb{N}\}$.  Finalement, $L(G)
= L(G,U)\,L(G,C) \;\cup\; L(G,A)\,L(G,V)$ montre que $L(G) = \{a^i b^i
c^j : i,j\in\mathbb{N}\} \cup \{a^i b^j c^j : i,j\in\mathbb{N}\}$ est
le langage des mots de la forme $a^i b^r c^j$ où $r=i$ \emph{ou}
$r=j$.

{\footnotesize (Ce langage est intéressant sur le plan théorique, car
  il est algébrique mais « intrinsèquement ambigu »
  (cf. \ref{intrinsic-ambiguity}).)\par}

\thingy\label{example-well-parenthesized-expressions} Sur l'alphabet
$\Sigma = \{a,b\}$, considérons la grammaire (d'axiome $S$)
\[
S\; \rightarrow\; aSbS \;|\; \varepsilon
\]
ou, ce qui revient au même
\[
\begin{aligned}
S &\rightarrow TS \;|\; \varepsilon\\
T &\rightarrow aSb\\
\end{aligned}
\]
(la première ligne revient par ailleurs à $S \rightarrow T^*$).  On a
alors $L(G) = L(G,T)^*$ où $L(G,T) = a\, L(G)\, b$ désigne le langage
des mots qui peuvent être dérivés à partir de $T$.

Ce langage $L(G)$ s'appelle langage des \defin[bien-parenthésée
  (expression)]{expressions bien-parenthésées} où $a$ représente une
parenthèse ouvrante et $b$ une parenthèse fermante (et $L(G,T)$ est le
langage des blocs parenthésés : autrement dit, une expression
bien-parenthésée est une suite de blocs parenthésés, et un bloc
parenthésé est une expression bien-parenthésée entourée par une
parenthèse ouvrante et une parenthèse fermante — c'est ce que traduit
la grammaire).

Soit dit en passant, ce langage algébrique des expressions
bien-parenthésées n'est pas rationnel : la façon la plus simple de
s'en convaincre est de remarquer que s'il l'était, son intersection
avec le langage $\{a^i b^j : i,j\in\mathbb{N}\}$ dénoté par
l'expression rationnelle $a{*}b{*}$ le serait aussi
(par \ref{dfa-union-and-intersection}), or cette intersection est
précisément le langage $\{a^n b^n : n\in\mathbb{N}\}$ dont on a montré
en \ref{example-of-pumping-lemma} qu'il n'était pas algébrique.

\thingy\label{example-grammar-equal-a-and-b} Sur l'alphabet $\Sigma =
\{a,b\}$, montrons que la grammaire $G$ (d'axiome $S$) donnée par
\[
S\; \rightarrow\; aSbS \;|\; bSaS \;|\; \varepsilon
\]
(c'est-à-dire $N = \{S\}$ où $S$ est l'axiome, et pour règles $S
\rightarrow aSbS$ et $S \rightarrow bSaS$ et $S \rightarrow
\varepsilon$) engendre le langage $L$ formé des mots ayant un nombre
total de $a$ égal au nombre total de $b$, i.e., $L := \{w \in\Sigma^* :
|w|_a = |w|_b\}$ où $|w|_x$ désigne le nombre total d'occurrences de
la lettre $x$ dans le mot $w$ (cf. \ref{number-of-occurrences-of-letter}).

\begin{proof}
Dans un sens il est évident que tout pseudo-mot (et en particulier,
tout mot sur $\Sigma$) obtenu en dérivant $S$ a un nombre de $a$ égal
au nombre de $b$, puisque cette propriété est conservée par toute
dérivation immédiate.  Autrement dit, $L(G) \subseteq L$.

Pour ce qui est de la réciproque, on va montrer par récurrence sur
$|w|$ que tout mot $w \in \Sigma^*$ ayant autant de $a$ que de $b$ est
dans $L(G)$.  Pour $|w|=0$ c'est trivial.  Sinon, on peut supposer
sans perte de généralité que le mot $w$ considéré commence par $a$,
soit $w = aw'$ où $w'$ vérifie $|w'|_b - |w'|_a = 1$.  Considérons
maintenant la fonction $h$ qui à un entier $k$ entre $0$ et $|w'|$
associe le nombre de $b$ moins le nombre de $a$ dans le préfixe de
longueur $k$ de $w'$ (i.e., parmi les $k$ premières lettres de $w'$) :
cette fonction prend les valeurs $h(0) = 0$ et $h(|w'|) = 1$, et elle
varie de $\pm 1$ à chaque fois (i.e., $h(i) = h(i-1) \pm 1$ selon que
la $i$-ième lettre de $w'$ est $b$ ou $a$).  Soit $k_0$ le plus grand
entier tel que $h(k_0) = 0$ : cet entier est bien défini car $h(0)=0$.
On a forcément $h(k_0+1) = 1$ car l'autre possibilité, $h(k_0+1) = -1$
impliquerait que $h$ devrait repasser par la valeur $0$ pour atteindre
la valeur finale $h(|w'|) = 1$, ce qui contredirait la maximalité de
$k_0$.  Bref, si $u$ est le préfixe de $w'$ de longueur $k_0$, on a
$|u|_b - |u|_a = h(k_0) = 0$ et la lettre suivante de $w'$ est un $b$,
si bien que $w' = ubv$ et finalement $w = aubv$ où $u$ et donc $v$ ont
autant que $a$ que de $b$, i.e., appartiennent à $L$, et par
récurrence (ils sont de longueur strictement plus courte que $w$)
appartiennent à $L(G)$.  On peut donc dériver $u$ et $v$ de $S$, donc
$w$ de $aSbS$, et comme $aSbS$ se dérive immédiatement de $S$, on a
bien $w \in L(G)$, ce qui conclut la récurrence.
\end{proof}

Un raisonnement analogue montre que la grammaire $G'$ donnée par
\[
\begin{aligned}
S &\rightarrow aUbS \;|\; bVaS \;|\; \varepsilon\\
U &\rightarrow aUbU \;|\; \varepsilon\\
V &\rightarrow bVaV \;|\; \varepsilon\\
\end{aligned}
\]
engendre le même langage $L = \{w \in\Sigma^* : |w|_a = |w|_b\}$ que
ci-dessus (elle est donc faiblement équivalente à $G$) : l'idée est
que le langage $L(G,U)$ des mots qui dérivent de $U$ est l'ensemble
des expressions bien-parenthésées où $a$ est la parenthèse ouvrante et
$b$ la fermante (avec les notations de la démonstration ci-dessus, ce
sont les mots pour lesquels la fonction $h$, partant de $0$, revient à
$0$, et reste toujours positive ou nulle), tandis que $L(G,V)$ est
l'ensemble des expressions bien-parenthésées où $b$ est la parenthèse
ouvrante et $a$ la fermante (la fonction $h$ reste toujours négative
ou nulle), et tout mot de $L$ peut s'écrire comme concaténation de
tels mots (en divisant selon le signe de $h$).  Une différence entre
$G$ et $G'$ est que la grammaire $G$ est ambiguë (terme qui sera
défini en \ref{ambiguous-grammar}) tandis que $G'$ ne l'est pas.


\subsection{Arbres d'analyse, dérivations gauche et droite}

\thingy Soit $G$ une grammaire hors contexte sur un alphabet $\Sigma$
et $N$ l'ensemble de ses nonterminaux.  Un \defin{arbre d'analyse} (ou
\index{dérivation (arbre de)|see{arbre d'analyse}}\index{arbre de
  dérivation|see{arbre d'analyse}}\textbf{arbre de dérivation} ; en
anglais \textbf{parse tree}) \textbf{incomplet} pour $G$ est un arbre
(fini, enraciné et ordonné\footnote{C'est-à-dire que l'ensemble des
  fils d'un nœud quelconque de l'arbre est muni d'un ordre total, ou,
  ce qui revient au même, qu'ils sont numérotés (disons de la gauche
  vers la droite).}) dont les nœuds sont étiquetés par des éléments de
$\Sigma \cup N \cup \{\varepsilon\}$, vérifiant les propriétés
suivantes :
\begin{itemize}
\item la racine de l'arbre est étiquetée par l'axiome $S$ de $G$ ;
\item si un nœud de l'arbre n'est pas une feuille (i.e., s'il a des
  fils) et si on appelle $T$ son étiquette, alors
\begin{itemize}
\item soit ce nœud a un unique fils étiqueté $\varepsilon$ et il
  existe une règle $T \rightarrow \varepsilon$ dans $G$,
\item soit ce nœud a des fils étiquetés par des éléments $x_1\cdots
  x_n$ (où $n\geq 1$) de $\Sigma \cup N$ et il existe une règle $T
  \rightarrow x_1\cdots x_n$ dans $G$,
\end{itemize}
(dans ces deux sous-cas, il résulte de l'existence d'une règle
$T\rightarrow\cdots$ que $T$ est un nonterminal : seules les feuilles
peuvent être étiquetées par un terminal ou $\varepsilon$).
\end{itemize}

\smallskip

On dit de plus que l'arbre est \textbf{complet}, ou simplement qu'il
s'agit d'un arbre d'analyse, lorsqu'il vérifie la propriété
supplémentaire suivante :
\begin{itemize}
\item toute feuille de l'arbre est étiquetée soit par un terminal
  (= élément de $\Sigma$) soit par $\varepsilon$.
\end{itemize}

(Les feuilles étiquetées $\varepsilon$ servent uniquement à marquer la
complétude de l'arbre : certains auteurs procèdent différemment ou les
omettent.)

\thingy Si $T$ est un arbre d'analyse incomplet, et $\alpha$ le
pseudo-mot obtenu en lisant les étiquettes des \emph{feuilles} de $T$
en profondeur ordonnée (c'est-à-dire en commençant par toutes les
feuilles qui descendent du premier fils de la racine elles-mêmes dans
le même ordre, puis toutes celles qui descendent du second fils, et
ainsi de suite) et en ignorant tous les $\varepsilon$, alors on dit
que $T$ est un arbre d'analyse \emph{de} $\alpha$, ou que $\alpha$ est
le pseudo-mot analysé par $T$.  La définition d'un arbre d'analyse
complet signifie exactement qu'il est l'arbre d'analyse d'un
\emph{mot} (c'est-à-dire, d'un mot sur $\Sigma$).

\thingy\label{example-of-parse-tree} À titre d'exemple, considérons la
grammaire (considérée en \ref{example-well-parenthesized-expressions})
\[
\begin{aligned}
S &\rightarrow TS \;|\; \varepsilon\\
T &\rightarrow aSb\\
\end{aligned}
\]

L'arbre d'analyse suivant est un arbre d'analyse du mot $aabbab$ (et
on verra en \ref{inambiguity-of-well-parenthesized-expressions} que
c'est le seul possible) :
\begin{center}
\tikzstyle{automaton}=[scale=0.5]
%%% begin parsetree1 %%%

\begin{tikzpicture}[>=latex,line join=bevel,automaton]
%%
\node (e3) at (99bp,18bp) [draw,draw=none] {$\varepsilon$};
  \node (S3) at (549bp,234bp) [draw,draw=none] {$S$};
  \node (S2) at (225bp,234bp) [draw,draw=none] {$S$};
  \node (S1) at (441bp,306bp) [draw,draw=none] {$S$};
  \node (S0) at (369bp,378bp) [draw,draw=none] {$S$};
  \node (T2) at (117bp,162bp) [draw,draw=none] {$T$};
  \node (S6) at (99bp,90bp) [draw,draw=none] {$S$};
  \node (T0) at (261bp,306bp) [draw,draw=none] {$T$};
  \node (T1) at (441bp,234bp) [draw,draw=none] {$T$};
  \node (a1) at (369bp,162bp) [draw,draw=none] {$a$};
  \node (a0) at (153bp,234bp) [draw,draw=none] {$a$};
  \coordinate (spacer) at (297bp,162bp);
  \node (a2) at (27bp,90bp) [draw,draw=none] {$a$};
  \node (b0) at (297bp,234bp) [draw,draw=none] {$b$};
  \node (S5) at (441bp,162bp) [draw,draw=none] {$S$};
  \node (b2) at (171bp,90bp) [draw,draw=none] {$b$};
  \node (S4) at (225bp,162bp) [draw,draw=none] {$S$};
  \node (e1) at (243bp,90bp) [draw,draw=none] {$\varepsilon$};
  \node (e0) at (585bp,162bp) [draw,draw=none] {$\varepsilon$};
  \node (e2) at (441bp,90bp) [draw,draw=none] {$\varepsilon$};
  \node (b1) at (513bp,162bp) [draw,draw=none] {$b$};
  \draw [] (S1) ..controls (484.16bp,277.03bp) and (505.72bp,263.05bp)  .. (S3);
  \draw [] (S0) ..controls (397.96bp,348.85bp) and (412.29bp,334.92bp)  .. (S1);
  \draw [] (S5) ..controls (441bp,132.85bp) and (441bp,118.92bp)  .. (e2);
  \draw [] (S2) ..controls (225bp,204.85bp) and (225bp,190.92bp)  .. (S4);
  \draw [] (T2) ..controls (109.76bp,132.85bp) and (106.18bp,118.92bp)  .. (S6);
  \draw [] (T1) ..controls (469.96bp,204.85bp) and (484.29bp,190.92bp)  .. (b1);
  \draw [] (T0) ..controls (217.84bp,277.03bp) and (196.28bp,263.05bp)  .. (a0);
  \draw [] (T2) ..controls (80.802bp,132.85bp) and (62.893bp,118.92bp)  .. (a2);
  \draw [] (T2) ..controls (138.72bp,132.85bp) and (149.46bp,118.92bp)  .. (b2);
  \draw [] (S1) ..controls (441bp,276.85bp) and (441bp,262.92bp)  .. (T1);
  \draw [] (S4) ..controls (232.24bp,132.85bp) and (235.82bp,118.92bp)  .. (e1);
  \draw [] (S3) ..controls (563.48bp,204.85bp) and (570.64bp,190.92bp)  .. (e0);
  \draw [] (S0) ..controls (325.84bp,349.03bp) and (304.28bp,335.05bp)  .. (T0);
  \draw [] (S6) ..controls (99bp,60.846bp) and (99bp,46.917bp)  .. (e3);
  \draw [] (T0) ..controls (275.48bp,276.85bp) and (282.64bp,262.92bp)  .. (b0);
  \draw [] (T0) ..controls (246.52bp,276.85bp) and (239.36bp,262.92bp)  .. (S2);
  \draw [] (T1) ..controls (412.04bp,204.85bp) and (397.71bp,190.92bp)  .. (a1);
  \draw [] (T1) ..controls (441bp,204.85bp) and (441bp,190.92bp)  .. (S5);
  \draw [] (S2) ..controls (181.84bp,205.03bp) and (160.28bp,191.05bp)  .. (T2);
%
\end{tikzpicture}

%%% end parsetree1 %%%
\end{center}

On obtient des exemples d'arbres d'analyse incomplet en effaçant
arbitrairement tous les descendants d'un ou plusieurs nœuds de cet
arbre.

\thingy À toute dérivation à partir de l'axiome $S$ dans une grammaire
hors contexte $G$ on peut associer un arbre d'analyse incomplet : ce
dernier est construit en partant de la racine et, à chaque fois qu'un
symbole est réécrit (dans la dérivation), en faisant pousser des fils
à la feuille correspondante de l'arbre (qui cesse donc d'être une
feuille) pour indiquer la règle appliquée.

De façon plus formelle : si $S =: \lambda_0 \Rightarrow \lambda_1
\Rightarrow \cdots \Rightarrow \lambda_n$ est une dérivation à partir
de l'axiome $S$ dans une grammaire hors contexte $G$, on va lui
associer, par récurrence sur le nombre d'étapes $n$ de la dérivation,
un arbre d'analyse incomplet du pseudo-mot $\lambda_n$.  Si $n=0$,
l'arbre $\mathscr{T}_0$ en question est simplement l'arbre trivial
(ayant pour seul nœud la racine $S$).  Sinon, on part de l'arbre
$\mathscr{T}_{n-1}$ (qu'on a défini par récurrence) pour la dérivation
$\lambda_0 \Rightarrow \cdots \Rightarrow \lambda_{n-1}$ : supposons
que la dernière dérivation immédiate $\lambda_{n-1} \Rightarrow
\lambda_n$ réécrive le symbole nonterminal $T$ par application de la
règle $T \rightarrow \alpha$ avec pour contextes gauche et droit
$\gamma,\gamma'$ (cf. \ref{derivations-and-contexts}), si bien que
$\lambda_{n-1} = \gamma T \gamma'$ et $\lambda_n = \gamma \alpha
\gamma'$ ; le symbole $T$ réécrit correspond à un certain nœud de
l'arbre $\mathscr{T}_{n-1}$, à savoir la feuille, étiquetée $T$, telle
qu'on lise $\gamma$ à sa gauche et $\gamma'$ à sa droite en suivant
les feuilles dans l'ordre de profondeur ; on ajoute alors au nœud en
question des fils correspondant à la règle $T \to \alpha$,
c'est-à-dire soit un unique fils étiqueté $\varepsilon$ si $\alpha =
\varepsilon$, soit $k$ fils étiquetés $x_1,\ldots,x_k$ si $\alpha =
x_1\cdots x_k$ avec $k\geq 1$.

Notons que l'arbre obtenu est complet exactement lorsque la dérivation
aboutit à un mot (sur $\Sigma$).  Notons aussi que le nombre $n$
d'étapes dans la dérivation est égal au nombre de nœuds de l'arbre qui
\emph{ne sont pas} des feuilles.

Deux dérivations (forcément du même mot ou pseudo-mot) auxquelles sont
associées le même arbre d'analyse sont dites \defin[équivalentes
  (dérivations)]{équivalentes}.

\thingy\label{example-of-derivations} À titre d'exemple, l'arbre
illustré en \ref{example-of-parse-tree} est associé à la dérivation
\[
\begin{aligned}
\underline{S} &\Rightarrow \underline{T}S \Rightarrow a\underline{S}bS \Rightarrow a\underline{T}SbS \Rightarrow
aa\underline{S}bSbS \Rightarrow aab\underline{S}bS \Rightarrow aabb\underline{S}\\
&\Rightarrow aabb\underline{T}S \Rightarrow aabba\underline{S}bS \Rightarrow aabbab\underline{S} \Rightarrow aabbab\\
\end{aligned}
\]
(on a souligné à chaque fois le symbole réécrit), mais aussi à la
dérivation
\[
\begin{aligned}
\underline{S} &\Rightarrow T\underline{S} \Rightarrow TT\underline{S} \Rightarrow T\underline{T} \Rightarrow Ta\underline{S}b \Rightarrow \underline{T}ab\\
&\Rightarrow a\underline{S}bab \Rightarrow aT\underline{S}bab \Rightarrow a\underline{T}bab \Rightarrow aa\underline{S}bbab \Rightarrow aabbab\\
\end{aligned}
\]
Ces deux dérivations sont donc équivalentes.  Elles ont toutes les
deux $10$ étapes puisque l'arbre considéré a $10$ nœuds qui ne sont
pas des feuilles.

\thingy On appelle \defin{dérivation gauche} une dérivation (pour une
grammaire hors contexte donnée) dans laquelle le symbole réécrit est
toujours \emph{le nonterminal le plus à gauche} du pseudo-mot
courant : autrement dit, une dérivation gauche est une dérivation
telle que le contexte gauche de chaque dérivation immédiate la
constituant ne comporte que des symboles terminaux (i.e., appartient
à $\Sigma^*$).  Symétriquement, on appelle \textbf{dérivation droite}
une dérivation dans laquelle le symbole réécrit est toujours \emph{le
  nonterminal le plus à droite}, c'est-à-dire que le contexte droit de
chaque dérivation immédiate la constituant ne comporte que des
symboles terminaux.

\smallskip

À titre d'exemple, les deux dérivations données
en \ref{example-of-derivations} sont respectivement une dérivation
gauche et une dérivation droite.

\smallskip

À chaque arbre d'analyse est associée une et une seule dérivation
gauche : on l'obtient de façon évidente en réécrivant à chaque étape
le symbole nonterminal le plus à gauche en suivant la règle indiquée
par l'arbre d'analyse sous le nœud correspondant ; cela revient à
parcourir l'arbre d'analyse en profondeur de gauche à droite, et à
réécrire le symbole correspondant à chaque nœud que l'on rencontre qui
n'est pas une feuille.  De même, à chaque arbre d'analyse est associée
une et une seule dérivation droite.


\subsection{Ambiguïté}

\thingy\label{ambiguous-grammar} On dit qu'une grammaire hors contexte
est \defin[ambiguë (grammaire)]{ambiguë} lorsqu'il existe un mot
(i.e., un mot sur l'alphabet $\Sigma$ des terminaux) qui admet deux
arbres d'analyse \emph{différents}.  Dans le cas contraire, elle est
dite \defin[inambiguë (grammaire)]{inambiguë} : autrement dit, une
grammaire inambiguë est une grammaire $G$ pour laquelle tout $w \in
L(G)$ a un unique arbre d'analyse ; il revient au même de dire que
tout mot de $L(G)$ a une unique dérivation droite, ou encore que tout
mot de $L(G)$ a une unique dérivation gauche.

\smallskip

Les grammaires des langages informatiques réels sont évidemment
(presque ?) toujours inambiguës : on souhaite qu'un programme
(c'est-à-dire, dans la terminologie mathématique, un mot du langage)
admette une unique interprétation, autrement dit, un unique arbre
d'analyse.

\thingy\label{trivial-example-ambiguity} À titre d'exemple, la
grammaire
\[
S\; \rightarrow\; SS \;|\; a
\]
sur l'alphabet $\{a\}$ est ambiguë.  En effet, le mot $aaa$ admet
l'arbre d'analyse
\begin{center}
\tikzstyle{automaton}=[scale=0.5]
%%% begin parsetree2 %%%

\begin{tikzpicture}[>=latex,line join=bevel,automaton]
%%
\node (S3) at (99bp,90bp) [draw,draw=none] {$S$};
  \node (S2) at (99bp,162bp) [draw,draw=none] {$S$};
  \node (S1) at (27bp,162bp) [draw,draw=none] {$S$};
  \node (S0) at (63bp,234bp) [draw,draw=none] {$S$};
  \node (S4) at (171bp,90bp) [draw,draw=none] {$S$};
  \node (a1) at (99bp,18bp) [draw,draw=none] {$a$};
  \node (a0) at (27bp,90bp) [draw,draw=none] {$a$};
  \node (a2) at (171bp,18bp) [draw,draw=none] {$a$};
  \draw [] (S0) ..controls (48.521bp,204.85bp) and (41.357bp,190.92bp)  .. (S1);
  \draw [] (S1) ..controls (27bp,132.85bp) and (27bp,118.92bp)  .. (a0);
  \draw [] (S2) ..controls (127.96bp,132.85bp) and (142.29bp,118.92bp)  .. (S4);
  \draw [] (S2) ..controls (99bp,132.85bp) and (99bp,118.92bp)  .. (S3);
  \draw [] (S0) ..controls (77.479bp,204.85bp) and (84.643bp,190.92bp)  .. (S2);
  \draw [] (S4) ..controls (171bp,60.846bp) and (171bp,46.917bp)  .. (a2);
  \draw [] (S3) ..controls (99bp,60.846bp) and (99bp,46.917bp)  .. (a1);
%
\end{tikzpicture}

%%% end parsetree2 %%%
\end{center}
mais aussi
\begin{center}
\tikzstyle{automaton}=[scale=0.5]
%%% begin parsetree2b %%%

\begin{tikzpicture}[>=latex,line join=bevel,automaton]
%%
\node (S3) at (27bp,90bp) [draw,draw=none] {$S$};
  \node (S2) at (171bp,162bp) [draw,draw=none] {$S$};
  \node (S1) at (99bp,162bp) [draw,draw=none] {$S$};
  \node (S0) at (135bp,234bp) [draw,draw=none] {$S$};
  \node (S4) at (99bp,90bp) [draw,draw=none] {$S$};
  \node (a1) at (99bp,18bp) [draw,draw=none] {$a$};
  \node (a0) at (27bp,18bp) [draw,draw=none] {$a$};
  \node (a2) at (171bp,90bp) [draw,draw=none] {$a$};
  \draw [] (S1) ..controls (70.042bp,132.85bp) and (55.714bp,118.92bp)  .. (S3);
  \draw [] (S0) ..controls (120.52bp,204.85bp) and (113.36bp,190.92bp)  .. (S1);
  \draw [] (S0) ..controls (149.48bp,204.85bp) and (156.64bp,190.92bp)  .. (S2);
  \draw [] (S1) ..controls (99bp,132.85bp) and (99bp,118.92bp)  .. (S4);
  \draw [] (S3) ..controls (27bp,60.846bp) and (27bp,46.917bp)  .. (a0);
  \draw [] (S2) ..controls (171bp,132.85bp) and (171bp,118.92bp)  .. (a2);
  \draw [] (S4) ..controls (99bp,60.846bp) and (99bp,46.917bp)  .. (a1);
%
\end{tikzpicture}

%%% end parsetree2b %%%
\end{center}

Cette grammaire est donc ambiguë.  Remarquons que le langage
qu'elle engendre est $\{a\}^+ = \{a^n : n\geq 1\}$ (car il est évident
que tout mot engendré par la grammaire est dans $\{a\}^+$, et
réciproquement il est facile de fabriquer une dérivation de $a^n$ pour
tout $n\geq 1$, par exemple $\underline{S} \Rightarrow S\underline{S}
\Rightarrow \cdots \Rightarrow S^{n-1} S \Rightarrow \cdots a^n$).

Le \emph{même} langage $\{a\}^+ = \{a^n : n\geq 1\}$ peut aussi être
engendré par la grammaire
\[
S\; \rightarrow\; aS \;|\; a
\]
(faiblement équivalente à la précédente, donc) qui, elle, \emph{n'est
  pas} ambiguë : en effet, la seule manière de dériver $a^n$ consiste
à appliquer $n-1$ fois la règle $S \rightarrow aS$ et finalement une
fois la règle $S \rightarrow a$ (il y a donc une unique dérivation du
mot, et \textit{a fortiori} un unique arbre d'analyse).

\thingy\label{inambiguity-of-well-parenthesized-expressions} La
grammaire $G$ des expressions bien-parenthésées
\[
\begin{aligned}
S &\rightarrow TS \;|\; \varepsilon\\
T &\rightarrow aSb\\
\end{aligned}
\]
considérée en \ref{example-well-parenthesized-expressions} est
inambiguë.  Le point crucial pour s'en convaincre est que dans
l'application de la règle $S \rightarrow TS$ dans l'analyse d'un mot
$w \neq \varepsilon$ engendré par la grammaire, il n'y a qu'une seule
possibilité sur la limite du préfixe $u$ qui dérivera de $T$ (et donc
du suffixe qui dérivera de $S$) : en s'inspirant
de \ref{example-grammar-equal-a-and-b} on peut se convaincre que $u$
est le préfixe de $w$ de la plus petite longueur $>0$ possible
comportant autant de $b$ que de $a$ (par exemple si $w = aabbab$ alors
$u = aabb$) ; ce point étant acquis, tout mot $w\neq\varepsilon$ de la
grammaire $L(G) = L(G,S)$ s'analyse de façon unique comme
concaténation d'un mot $u \in L(G,T)$ (c'est-à-dire dérivant de $T$)
et d'un mot $v \in L(G,S)$, et chacun de ces morceaux s'analyse à son
tour de façon unique (la règle $T \rightarrow aSb$ ne permet
manifestement qu'une seule analyse d'un mot de $L(G,T)$ : une fois
qu'on enlève le $a$ initial et le $b$ final, il reste un mot de
$L(G,S)$).

Notamment, l'arbre représenté en \ref{example-of-parse-tree} est
l'\emph{unique} arbre d'analyse de $aabbab$ pour la grammaire
présentée ci-dessus.

\thingy\label{intrinsic-ambiguity} Il arrive que le \emph{même}
langage puisse être engendré par une grammaire ambiguë et par une
grammaire inambiguë (on a vu un exemple
en \ref{trivial-example-ambiguity}).  L'ambiguïté est donc une
caractéristique de la \emph{grammaire} hors contexte et non du
\emph{langage} algébrique qu'elle engendre.

\smallskip

Cependant, certains langages algébriques ne sont définis \emph{que}
par des grammaires hors contexte ambiguës.  De tels langages sont dits
\defin[intrinsèquement ambigu (langage algébrique)]{intrinsèquement
  ambigus}.  C'est le cas du langage $\{a^i b^i c^j :
i,j\in\mathbb{N}\} \cup \{a^i b^j c^j : i,j\in\mathbb{N}\}$ dont on a
vu en \ref{example-for-intrinsic-ambiguity} qu'il était algébrique :
il n'est pas évident (cela dépasse le cadre de ce cours) de démontrer
qu'il est intrinsèquement ambigu, mais on peut au moins en donner une
explication intuitive : quelle que soit la manière dont on construit
une grammaire engendrant ce langage, elle devra forcément distinguer
le cas de $a^i b^i c^j$ et celui de $a^i b^j c^j$, or ces cas ne sont
pas disjoints, il existe des mots $a^i b^i c^i$ qui sont à
l'intersection des deux, et ce sont ces mots qui forcent ce langage à
être intrinsèquement ambigu.


\subsection{Le lemme de pompage pour les langages algébriques}

On a vu en §\ref{subsection-pumping-lemma} une condition nécessaire
que doivent vérifier les langages rationnels, et qui est souvent utile
pour montrer qu'un langage \emph{n'est pas} rationnel.  Un lemme tout
à fait analogue existe pour les langages algébriques, est qui s'avère
utile dans des circonstances semblables, même si son emploi est plus
difficile ; on prendra garde au fait que, dans l'énoncé suivant,
$x$ et $y$ désignent des mots et non des lettres comme d'habitude :

\begin{prop}[lemme de pompage pour les langages algébriques]\label{pumping-lemma-for-algebraic-languages}\index{pompage (lemme de)}
Soit $L$ un langage algébrique.  Il existe alors un entier $k$ tel que
tout mot de $t \in L$ de longueur $|t| \geq k$ admette une
factorisation $t = uvwxy$ en cinq facteurs $u,v,w,x,y \in \Sigma^*$
où :
\begin{itemize}
\item[(i)] $|vx| \geq 1$ (c'est-à-dire que l'un au moins de $v$ et $x$
  est $\neq\varepsilon$),
\item[(ii)] $|vwx| \leq k$,
\item[(iii)] pour tout $i\geq 0$ on a $uv^iwx^iy \in L$.
\end{itemize}
\end{prop}

\smallskip

Donnons maintenant un exemple d'utilisation du lemme :

\begin{prop}\label{example-of-pumping-lemma-for-algebraic-languages}
Soit $\Sigma = \{a,b,c\}$.  Le langage $L = \{a^n b^n c^n :
n\in\mathbb{N}\} = \{\varepsilon, abc, aabbcc, aaabbbcc,\ldots\}$
n'est pas algébrique.
\end{prop}
\begin{proof}
Appliquons la proposition \ref{pumping-lemma-for-algebraic-languages}
au langage $L$ considéré : appelons $k$ l'entier dont le lemme de
pompage garantit l'existence.  Considérons le mot $t := a^k b^k c^k$ :
il doit alors exister une factorisation $t = uvwxy$ pour laquelle on a
(i) $|vx|\geq 1$, (ii) $|vwx|\leq k$ et (iii) $uv^iwx^iy \in L$ pour
tout $i\geq 0$.  La propriété (ii) assure que le facteur $vwx$ ne peut
pas contenir simultanément les lettres $a$ et $c$ : en effet, tout
facteur de $t$ comportant un $a$ et un $c$ doit avoir aussi le facteur
$b^k$, et donc être de longueur $\geq k+2$.  Supposons que $vwx$ ne
contienne pas la lettre $c$ (l'autre cas étant complètement
analogue) : en particulier, ni $v$ ni $x$ ne la contient, donc le mot
$uv^iwx^iy$, qui est dans $L$ d'après (iii), a le même nombre de $c$
que le mot $t$ initial ; mais comme son nombre de $a$ ou bien de $b$
est différent (d'après (i)), on a une contradiction.
\end{proof}

\thingy La proposition \ref{intersection-of-algebraic-and-rational}
peut s'avérer utile pour montrer qu'un langage n'est pas algébrique,
en permettant de simplifier le langage auquel on va appliquer le lemme
de pompage.

\smallskip

À titre d'exemple, montrons que le langage $L$ formé des mots sur
$\{a,b,c\}$ ayant le même nombre total de $a$, de $b$ et de $c$
(autrement dit $\{w \in \{a,b,c\}^* : |w|_a = |w|_b = |w|_c\}$ où
$|w|_x$ désigne le nombre d'occurrences de la lettre $x$ dans le
mot $w$, cf. \ref{number-of-occurrences-of-letter}) n'est pas
algébrique ; le plus simple pour le voir est de l'intersecter avec le
langage rationnel $M := \{a^i b^j c^k : i,j,k\in\mathbb{N}\}$ (dénoté
par l'expression rationnelle $a{*}b{*}c{*}$) : si $L$ était algébrique
alors d'après \ref{intersection-of-algebraic-and-rational}, le langage
$L\cap M$ le serait aussi ; mais $L\cap M = \{a^i b^i c^i :
i\in\mathbb{N}\}$, et on vient de voir
en \ref{example-of-pumping-lemma-for-algebraic-languages} qu'il n'est
pas algébrique ; c'est donc que $L$ n'est pas non plus algébrique.


\subsection{Notions sur l'analyse algorithmique des langages hors contexte}

\thingy Il est naturel de se poser la question suivante : existe-t-il
un algorithme qui, donnée une grammaire hors contexte $G$, permet de
déterminer si un mot donné appartient au langage $L(G)$ engendré
par $G$, et, si oui, d'en trouver un arbre d'analyse ?  La réponse à
cette question est positive, mais plus délicate que dans le cadre des
langages rationnels où la notion d'automate fini a permis de donner un
point de vue clair (cf. \ref{rational-languages-are-recognizable}).

\thingy\label{handwaving-on-stack-automata} Il existe bien un modèle
de calcul qu'on peut imaginer comme l'analogue pour les langages
algébriques (et grammaires hors contexte) de ce que les automates
finis sont pour les langages rationnels (et expressions régulières) :
il s'agit des \emph{automates à pile}.  Informellement, un automate à
pile non déterministe fonctionne comme un automate fini non
déterministe à transitions spontanées, mais il a, en plus de son état
courant, accès à une « pile », qui contient des éléments d'un autre
alphabet (l'alphabet de pile), et pour choisir la transition à
effectuer, il peut consulter, en plus du symbole proposé, les symboles
au sommet de la pile (jusqu'à une profondeur bornée), et une fois
cette transition effectuée, décider de rajouter, retirer ou remplacer
des symboles au sommet de la pile.

\medskip

{\footnotesize

De façon plus formelle, un \index{automate à pile}automate à pile non déterministe est la
donnée d'un ensemble fini $Q$ d'états, d'un ensemble $I \subseteq Q$
d'états dits initiaux, d'un ensemble $F \subseteq Q$ d'états dits
finaux, d'un ensemble fini $\Gamma$ appelé alphabet de pile, et d'une
relation de transition $\delta \subseteq (Q \times
(\Sigma\cup\{\varepsilon\}) \times \Gamma^*) \times (Q \times
\Gamma^*)$.  Dire que $((q,t,\lambda), (q',\lambda')) \in \delta$
signifie que l'automate peut transitionner de l'état $q$ avec
$\lambda$ au sommet de la pile vers l'état $q'$ avec $\lambda'$ (à la
place de $\lambda$) au sommet de la pile en consommant la lettre $t$
(ou spontanément si $t=\varepsilon$), et l'automate \emph{accepte} un
mot $w$ lorsqu'il existe une suite de transitions d'un état initial
avec pile vide vers un état final avec pile vide qui consomme les
lettres de $w$.

\smallskip

De façon plus précise, l'automate accepte $w$ lorsqu'il existe
$q_0,\ldots,q_n \in Q$ (les états traversés) et $t_1,\ldots,t_n \in
(\Sigma\cup\{\varepsilon\})$ (les symboles consommés) et
$\gamma_1,\ldots,\gamma_n \in \Gamma^*$ (les états intermédiaires de
la pile) et $\lambda_1,\ldots,\lambda_n \in \Gamma^*$ (les mots
dépilés) et $\lambda'_1,\ldots,\lambda'_n \in \Gamma^*$ (les mots
empilés) tels que : $q_0 \in I$ et $q_n\in F$ et et $\lambda_1\gamma_1
= \varepsilon$ et $\lambda'_n\gamma_n = \varepsilon$ et
$((q_{i-1},t_i,\lambda_i),(q_i,\lambda'_i)) \in \delta$ pour
chaque $1\leq i\leq n$ et $w = t_1\cdots t_n$ et enfin
$\lambda_i\gamma_i = \lambda'_{i-1}\gamma_{i-1}$ pour chaque $1\leq
i\leq n$.

(Il existe différentes variations autour de cette notion, certaines
sans importance : notamment, on peut relaxer l'exigence que l'automate
termine le calcul avec la pile vide, cela ne change rien à la classe
des langages acceptés.)

\par}

\medskip

On peut montrer qu'il y a équivalence entre grammaires hors contexte
et automates à pile non déterministes au sens où tout langage engendré
par une grammaire hors contexte est le langage accepté par un automate
à pile non déterministe et réciproquement.  Il n'est d'ailleurs pas
très difficile de construire algorithmiquement un automate à pile non
déterministe qui accepte le langage engendré par une grammaire hors
contexte donnée.  Mais une différence essentielle avec les automates
finis est que cette fois \emph{le non déterministe est essentiel} :
les automates à pile déterministes (qu'il faut définir soigneusement)
acceptent strictement moins de langages que les automates à pile
non déterministes.

\bigbreak

\thingy Une approche simple, quoique terriblement inefficace, pour
résoudre algorithmiquement, \emph{en théorie}, le problème de décider
si un mot $w$ appartient au langage $L(G)$ engendré par une grammaire
hors contexte $G$ (absolument quelconque) est la suivante :
\begin{itemize}
\item réécrire la grammaire (i.e., la remplacer par une grammaire
  équivalente) de manière à ce que le membre de droite de chaque
  production soit \emph{non vide}, quitte à traiter spécialement la
  question de savoir si $\varepsilon \in L(G)$,
\item on a ensuite affaire à une grammaire \emph{monotone},
  c'est-à-dire que l'application d'une règle ne peut qu'augmenter (au
  sens large) la longueur du pseudo-mot en cours de dérivation, ce qui
  permet d'explorer exhaustivement toutes les possibilités et de
  s'arrêter dès qu'on dépasse la longueur $|w|$ à atteindre.
\end{itemize}

\smallskip

Énonçons précisément le résultat en question :

\begin{thm}\label{algebraic-languages-are-decidable}
Il existe un algorithme qui, donnée une grammaire hors contexte $G$
(sur un alphabet $\Sigma$) et un mot $w \in \Sigma^*$, décide si $w
\in L(G)$.  Autrement dit, les langages algébriques sont
\emph{décidables} au sens
de \ref{definition-computable-function-or-set}.

Plus exactement, on montre :
\begin{itemize}
\item Il existe un algorithme qui, donnée une grammaire
  hors contexte $G$ (sur un alphabet $\Sigma$), calcule une grammaire
  $G'$ (sur le même alphabet $\Sigma$ et ayant le même ensemble $N$ de
  nonterminaux que $G$) dont toutes les productions sont soit de la
  forme $T \rightarrow \alpha$ avec $|\alpha| \geq 1$, et une partie
  $E$ de $\{\varepsilon\}$ (c'est-à-dire soit $\varnothing$ soit
  $\{\varepsilon\}$), telles que $L(G) = L(G') \cup E$.
\item Il existe un algorithme qui, donnée une grammaire $G'$ comme on
  vient de le dire, et un mot $w \in \Sigma^*$, décide si $w \in
  L(G')$.
\end{itemize}
\end{thm}
\begin{proof}
Montrons d'abord l'affirmation du premier point.

Pour cela, on va d'abord calculer l'ensemble $N_0$ des nonterminaux
« \index{évanescent (nonterminal)}évanescents » de $G$, un nonterminal $T$ étant dit « évanescent »
lorsque $T \mathrel{\Rightarrow^*} \varepsilon$.  Mais il est évident
que toute dérivation de $\varepsilon$ dans $G$ ne peut faire
intervenir que des nonterminaux évanescents.  On peut donc calculer
l'ensemble des nonterminaux évanescents de la manière suivante : on
commence avec $N_0 = \varnothing$, et tant qu'il existe dans $G$ une
règle $T \rightarrow \alpha$, pour laquelle $T$ n'est pas encore
dans $N_0$, avec $\alpha \in N_0^*$ (c'est-à-dire, ne faisant
intervenir que des nonterminaux connus pour être évanescents ; y
compris $\alpha = \varepsilon$), on ajoute $T$ à $N_0$ et on
recommence.  Cette boucle va évidemment terminer en temps fini (il n'y
a qu'un nombre fini de nonterminaux) et lorsque c'est le cas $N_0$
sera l'ensemble des nonterminaux évanescents.

Une fois calculé l'ensemble $N_0$ des nonterminaux évanescents, on
peut définir $G'$ de la manière suivante : pour chaque production $T
\rightarrow \alpha$ de $G$, on met dans $G'$ toutes les productions $T
\rightarrow \alpha'$ où $\alpha'$ est un sous-(pseudo-)mot
$\neq\varepsilon$ de $\alpha$ obtenu en effaçant un sous-ensemble
quelconque de ses nonterminaux évanescents.  La grammaire $G'$ est
alors « presque » faiblement équivalente à $G$ : tout arbre de
dérivation dans $G'$ donne un arbre de dérivation du même mot
dans $G$, quitte à ajouter, à chaque fois qu'une règle $T \rightarrow
\alpha'$ est utilisée dans $G'$, les nonterminaux évanescents
manquants, qui portent eux-mêmes un sous-arbre de dérivation du mot
vide (puisqu'ils sont, justement, évanescents) ; et réciproquement,
tout arbre de dérivation dans $G$ en donne un dans $G'$ quitte à
effacer tout sous-arbre qui ne porte que des feuilles $\varepsilon$
(ce qui assure que sa racine est évanescente).  La seule subtilité est
qu'on a éventuellement perdu le mot vide dans $L(G)$ (la procédure
qu'on vient de décrire conduisant à effacer la totalité de l'arbre),
mais il suffit de poser $E = \{\varepsilon\}$ exactement lorsque
l'axiome $S$ de $G$ est évanescent pour corriger ce problème.

Montrons maintenant l'affirmation du second point.  Pour cela, on
considère l'ensemble (fini !) $(\Sigma\cup N)^{\leq|w|}$ de tous les
pseudo-mots de longueur $\leq |w|$.  On construit et on explore
progressivement (par exemple par un algorithme de Dijkstra / parcours
en largeur), à partir de $S$, le graphe sur cet ensemble de sommets
dont les arêtes sont les dérivations immédiates de $G'$ (dont le
membre de droite soit de longueur $\leq |w|$) : comme la grammaire
$G'$ est monotone, toute dérivation de $w$ sera un chemin dans le
graphe qu'on vient de dire (elle ne peut pas passer par des
pseudo-mots de longueur $>|w|$), donc on peut détecter sur ce graphe
fini si une telle dérivation existe.
\end{proof}

{\footnotesize

\thingy\label{handwaving-on-dynamical-programming} Expliquons comment
on peut approcher de façon algorithmiquement plus efficace le problème
de reconnaître si $w \in L(G)$, et fournir une autre démonstration du
théorème \ref{algebraic-languages-are-decidable}, tout en continuant à
ne faire aucune hypothèse sur la grammaire hors contexte $G$.

\smallskip

Il s'agit de travailler en deux étapes :
\begin{itemize}
\item D'abord, trouver algorithmiquement une grammaire $G'$ et un $E
  \subseteq \{\varepsilon\}$ tels que $L(G) = L(G') \cup E$, et que
  $G'$ soit en \defin[Chomsky (forme normale de)]{forme normale de
    Chomsky}, c'est-à-dire que toute production de $G'$ est de la
  forme $T \rightarrow UV$ avec $U,V$ (exactement) deux nonterminaux,
  ou bien $T \rightarrow x$ avec $x$ un terminal.
\item Ensuite, utiliser un algorithme de programmation dynamique pour
  calculer, pour chaque facteur $u$ de $w$ (par ordre croissant de
  taille), l'ensemble de tous les terminaux $T$ tels que $T
  \mathrel{\Rightarrow^*} u$ (ce qui répond notamment à la question de
  savoir si $S \mathrel{\Rightarrow^*} w$).
\end{itemize}

\smallskip

Détaillons un peu plus chacune de ces étapes.

Pour transformer la grammaire $G$ en une grammaire $G'$ sous forme
normale de Chomsky, on effectue les transformations suivantes :
\begin{itemize}
\item Introduire pour chaque terminal $x$ un nouveau nonterminal $X$
  et une règle $X \rightarrow x$, et remplacer chaque occurrence
  de $x$ par $X$ dans le membre de droite de toute règle autre que $X
  \rightarrow x$.  Ceci permet de faire en sorte qu'à part les règles
  $X \rightarrow x$, le membre de droite de toute règle soit
  uniquement constitué de nonterminaux.
\item Pour chaque règle $T \rightarrow U_1 \cdots U_n$ dont le membre
  de droite est de longueur $n\geq 3$, introduire $n-2$ nouveaux
  nonterminaux $Z_1,\ldots, Z_{n-2}$ et remplacer la règle $T
  \rightarrow U_1 \cdots U_n$ par les $n-1$ règles $T \rightarrow U_1
  Z_1$ et $Z_i \rightarrow U_{i+1} Z_{i+1}$ pour $1\leq i\leq n-3$ et
  $Z_{n-2} \rightarrow U_{n-1} U_n$.  Ceci permet de faire en sorte
  que le membre de droite de chaque règle soit de longueur $\leq 2$.
\item Éliminer les règles produisant $\varepsilon$ (et calculer
  l'ensemble $E$) exactement comme dans la démonstration du premier
  point de \ref{algebraic-languages-are-decidable}.  À ce stade-là, le
  membre de droite de chaque règle est de longueur exactement
  $1$ ou $2$ (et dans le second cas, constitué de deux nonterminaux).
\item Pour chaque règle $V \rightarrow \alpha$ où $\alpha$ n'est pas
  un unique nonterminal, et chaque terminal tel qu'il existe une suite
  $T \rightarrow \cdots \rightarrow V$ de règles produisant à chaque
  fois un unique nonterminal, et réécrivant $T$ en $V$, introduire une
  règle $T \rightarrow \alpha$, puis supprimer toutes les règles dont
  le membre de droite est un unique nonterminal.  (Cette
  transformation fonctionne exactement comme l'élimination des
  ε-transitions d'un εNFA, cf. \ref{removal-of-epsilon-transitions},
  si on imagine que les règles de la forme $T \rightarrow U$ sont des
  sortes de transitions spontanées d'un nonterminal vers un autre.)
\end{itemize}

L'ordre de ces transformations peut être légèrement varié, mais celui
proposé ci-dessus est sans doute le meilleur.

\smallskip

Une fois la grammaire $G'$ en forme normale de Chomsky connue,
lorsqu'on a un mot $w$ dont on cherche à tester s'il appartient
à $L(G')$, on calcule, pour chaque facteur $u$ de $w$ (identifié par
son point de départ et sa longueur), dans l'ordre croissant de
longueur, l'ensemble $\Lambda(u)$ des nonterminaux $T$ tels que $u \in
L(G',T)$ (on rappelle,
cf. \ref{words-deriving-from-another-nonterminal}, que $L(G',T) :=
\{v \in \Sigma^* : T \mathrel{\Rightarrow^*} v\}$) :
\begin{itemize}
\item Si $|u|=1$, c'est-à-dire $u \in \Sigma$, il s'agit simplement de
  l'ensemble des $T$ tels que la règle $T \rightarrow u$ soit
  dans $G'$.
\item Si $|u|\geq 2$, considérer chaque factorisation $u = v_1 v_2$ en
  deux facteurs de taille $\geq 1$ (il y en a $|u|-1$ possibles), pour
  chacune d'entre elles, pour chaque $X_1 \in \Lambda(v_1)$ (i.e., tel
  que $v_1 \in L(G',X_1)$) et chaque $X_2 \in \Lambda(v_2)$ (i.e., tel
  que $v_2 \in L(G',X_2)$) (ces deux ensembles étant connus par
  récurrence), et chaque $Y$ tel que la règle $Y \rightarrow X_1 X_2$
  soit dans $G'$, mettre l'élément $Y$ dans $\Lambda(u)$.
\end{itemize}

Il n'est pas difficile de se convaincre que ceci construit bien les
ensembles $\Lambda(u)$ annoncés : la deuxième partie revient à
rechercher toutes les dérivations $Y \Rightarrow X_1 X_2
\mathrel{\Rightarrow^*} v_1 v_2 = u$ possibles dans lesquelles $X_1$ a
donné $v_1$ et $X_2$ a donné $v_2$.  Une fois les $\Lambda(u)$ connus,
tester si $w \in L(G')$ revient à tester si $S \in \Lambda(w)$.

\smallskip

L'algorithme qu'on vient de décrire (pour tester si $w \in L(G')$ une
fois $G'$ sous forme normale de Chomsky) porte le nom
d'\defin[Cocke-Younger-Kasami (algorithme de)]{algorithme de
  Cocke–Younger–Kasami} ou algorithme \index{CYK (algorithme
  de)|see{Cocke-Younger-Kasami}}\textbf{CYK}.  Sa complexité est
cubique en la longueur de $w$.

\par}

\bigbreak

\thingy\label{handwaving-on-ll-and-lr} Du point de vue pratique, on ne
cherche pas simplement à savoir si un mot appartient au langage
engendré par une grammaire, mais aussi à en construire un arbre
d'analyse ; par ailleurs, la complexité algorithmique des approches
décrites en \ref{algebraic-languages-are-decidable} et même
en \ref{handwaving-on-dynamical-programming} est inacceptable.  En
contrepartie de ces exigences, on est prêt à accepter de mettre des
contraintes sur la grammaire qui la rendent plus facile à analyser :
\emph{au minimum}, on supposera que la grammaire est inambiguë
(cf. \ref{ambiguous-grammar}), et en fait, on imposera des contraintes
beaucoup plus fortes (par exemple, la grammaire présentée
en \ref{example-unambiguous-but-not-deterministic-grammar}, bien
qu'inambiguë, ne sera pas acceptable car il n'y a pas de manière
déterministe de l'analyser, ni même d'analyser le langage qu'elle
engendre) ; ces contraintes sont assez techniques et difficiles à
décrire : dans la pratique, elles consistent essentiellement à essayer
de fabriquer l'analyseur et à constater si l'algorithme échoue.

\smallskip

Il existe deux principales approches pour construire un analyseur pour
une grammaire hors contexte (sujette à diverses contraintes
supplémentaires) ; dans les deux cas, on construit une sorte
d'automate à pile (cf. \ref{handwaving-on-stack-automata})
déterministe, mais l'utilisation de la pile est très différente dans
les deux cas.   De façon très simplifiée :
\begin{itemize}
\item Les analyseurs \defin[LL (analyse)]{LL} procèdent de façon
  \emph{descendante} (en anglais « top-down »), parcourent le mot
  depuis la gauche (« L ») et génèrent la dérivation gauche (« L ») de
  l'arbre d'analyse fabriqué, en partant de la racine et en descendant
  jusqu'aux feuilles\footnote{En botanique, un arbre a la racine en
    bas et les feuilles en haut ; en informatique, on les représente
    plutôt racine en haut et feuilles en bas.}.  La pile de
  l'analyseur LL sert, intuitivement, à mémoriser les règles qu'on a
  commencé à reconnaître, en partant de l'axiome de la grammaire, et
  qui restent encore à compléter.
\item Les analyseurs \defin[LR (analyse)]{LR} procèdent de façon
  \emph{ascendante} (en anglais « bottom-up »), parcourent le mot
  depuis la gauche (« L ») et génèrent la dérivation droite (« R ») de
  l'arbre d'analyse fabriqué, en partant des feuilles et en remontant
  jusqu'aux racines.  La pile de l'analyseur LR sert, intuitivement, à
  mémoriser les fragments d'arbre déjà construits, en partant des
  feuilles, et qui restent encore à regrouper.
\end{itemize}

\smallskip

On peut écrire un analyseur LL ou (plus difficilement) LR à la main
dans un cas simple, mais en général ces analyseurs sont fabriqués par
des algorithmes systématiques, implémentés dans des programmes tels
que YACC ou Bison (qui produit des analyseurs LR, même si Bison peut
dépasser ce cadre) ou JavaCC (qui produit des analyseurs LL).

\smallskip

L'idée générale à retenir est que les analyseurs LR sont strictement
plus puissants que les analyseurs LL (ils sont capables d'analyser
strictement plus de grammaires, cf. \ref{example-lr-non-ll-grammar}),
mais leur génération est plus difficile et les messages d'erreur
qu'ils retournent en cas de problème de syntaxe sont plus difficiles à
comprendre pour l'utilisateur.

\thingy\label{example-ll-and-lr-analysis} Pour illustrer le
fonctionnement et les différences des analyseurs LL et LR, considérons
une grammaire très simple à analyser comme
\[
\begin{aligned}
S &\rightarrow TS \;|\; c\\
T &\rightarrow aSb\\
\end{aligned}
\]
(il s'agit d'une variante de celle considérée
en \ref{example-well-parenthesized-expressions}, où on a ajouté un $c$
pour rendre l'analyse plus simple en évitant les
productions $\varepsilon$ ; on peut imaginer, si on veut, qu'il s'agit
d'une forme extrêmement primitive de XML où $a$ représente une balise
ouvrante, $b$ une balise fermante, et $c$ une balise vide).

\medskip

L'approche la plus évidente, si on doit écrire une fonction « analyser
un mot comme dérivant de $S$ dans cette grammaire » consiste à coder
deux fonctions mutuellement récursives, « chercher un préfixe qui
dérive de $S$ » et « chercher un préfixe qui dérive de $T$ ».  En
observant que tout mot qui dérive de $T$ doit commencer par la
lettre $a$, ce qui permet de distinguer les mots dérivant des règles
$S\rightarrow TS$ et $S\rightarrow c$, on va écrire :
\begin{itemize}
\item La fonction « rechercher un préfixe qui dérive de $S$ » (prenant
  en entrée un mot $w\in\{a,b\}^*$, et renvoyant un préfixe de $w$ et
  un arbre de dérivation de $w$ à partir de $S$) est définie comme
  suit :
\begin{itemize}
\item si la première lettre de $w$ est $c$, renvoyer le préfixe $c$ et
  l'arbre trivial $S\to c$, sinon :
\item appeler la fonction « rechercher un préfixe qui dérive
  de $T$ » sur $w$, qui retourne un préfixe $u$ de $w$ et un
  arbre $\mathscr{U}$,
\item appeler la fonction « rechercher un préfixe qui dérive de $S$ »
  sur le suffixe correspondant $t$ de $w$ (c'est-à-dire le $t$ tel que
  $w=ut$), qui retourne un préfixe $v$ de $t$ et un arbre
  $\mathscr{V}$,
\item renvoyer le préfixe $uv$ de $w$ ainsi que l'arbre d'analyse dont
  la racine est donnée par la règle $S\rightarrow TS$ et les
  sous-arbres $\mathscr{U}$ et $\mathscr{V}$ (i.e., une racine
  étiquetée $S$ et deux fils étiquetés $T$ et $S$ qui sont chacun
  racines de sous-arbres donnés par $\mathscr{U}$ et $\mathscr{V}$
  respectivement).
\end{itemize}
\item La fonction « rechercher un préfixe qui dérive de $T$ » (prenant
  en entrée un mot $u\in\{a,b\}^*$, et renvoyant un préfixe de $u$ et
  un arbre de dérivation de $u$ à partir de $T$) est définie comme
  suit :
\begin{itemize}
\item vérifier que la première lettre est un $a$ (sinon, soulever une
  exception indiquant une erreur d'analyse),
\item appeler la fonction « rechercher un préfixe qui dérive de $S$ »
  sur le suffixe correspondant $x$ de $u$ (c'est-à-dire le $x$ tel que
  $u=ax$), qui retourne un préfixe $w$ de $x$ et un
  arbre $\mathscr{W}$,
\item vérifier que la lettre qui suit $w$ dans $x$ est bien $b$,
  c'est-à-dire que $u$ commence par $awb$ (sinon, soulever une
  exception indiquant une erreur d'analyse),
\item renvoyer le préfixe $awb$ de $u$ ainsi que l'arbre d'analyse
  dont la racine est donnée par la règle $T\rightarrow aSb$ et les
  sous-arbres $a$, $\mathscr{W}$ et $b$ (i.e., une racine étiquetée
  $T$ et trois fils étiquetés $a$, $S$ et $b$, celui du milieu étant
  racine d'un sous-arbre donné par $\mathscr{W}$).
\end{itemize}
\end{itemize}

\smallskip

Cette approche réussit sur cette grammaire très simple (où on peut
notamment se convaincre que l'éventuel préfixe dérivant de $S$ ou de
$T$ est toujours défini de façon unique).  L'analyseur qu'on vient de
décrire s'appelle un « analyseur par descente récursive ».  Or, plutôt
qu'utiliser la récursivité du langage de programmation (c'est-à-dire
la pile système), on peut aussi utiliser une pile comme structure de
données, et transformer les fonctions récursives en fonctions
itératives\footnote{Ceci est un fait général : tout ensemble de
  fonctions récursives peut se réécrire pour utiliser une pile comme
  structure de données explicite à la place de la pile système.}.  On
obtient ainsi essentiellement un automate à pile, qui utilise sa pile
pour retenir les règles qu'il a commencé à analyser (à partir de la
racine de l'arbre d'analyse en cours de construction).  On a
essentiellement construit un analyseur LL, ou plus exactement LL($1$)
(le « $1$ » indiquant qu'on se contente de lire une unique lettre du
mot pour décider quelle règle chercher à analyser), pour ce langage.
C'est ici l'approche « descendante » : l'arbre se construit à partir
de la racine et la pile sert à retenir les règles qu'on a commencé à
reconnaître.

\medbreak

L'approche « ascendante » de la même grammaire serait plutôt la
suivante : on parcourt le mot de gauche à droite en gardant de côté
une pile (initialement vide) qui pourra contenir les symboles $a,S,T$,
les deux derniers étant alors associés à des arbres d'analyse ;
\begin{itemize}
\item si on lit un $a$, on se contente de l'empiler,
\item si on lit un $c$, on crée un arbre d'analyse $S\rightarrow c$ et
  on empile le $S$, puis, tant que la pile contient $T$ et $S$ en son
  sommet, on dépile ces deux symboles, on rassemble les deux arbres
  d'analyse associés en les mettant sous un $S\rightarrow TS$ et on
  empile le $S$ correspondant,
\item si on lit un $b$, on vérifie que les deux symboles au sommet de
  la pile sont $a$ et $S$ (sinon on soulève une erreur d'analyse), on
  les dépile et on rassemble l'arbre d'analyse associé au $S$ en le
  mettant sous un $T\rightarrow aSb$, et enfin on empile un $T$
  associé à cet arbre,
\item enfin, si on arrive à la fin du mot, la pile ne doit contenir
  qu'un unique symbole $S$ (sinon on soulève une erreur d'analyse), et
  l'arbre d'analyse final est celui qui lui est associé.
\end{itemize}

Il est un peu difficile d'expliquer en général comment construire un
tel analyseur, mais sur cet exemple précis il est facile de se
convaincre qu'il fonctionne et de comprendre pourquoi : il s'agit
essentiellement là d'un analyseur LR (en fait, LR($0$), le « $0$ »
indiquant qu'on n'a jamais eu besoin de regarder au-delà du symbole
courant pour décider quoi faire).  C'est ici l'approche
« ascendante » : l'arbre se construit à partir des feuilles et la pile
sert à retenir les nonterminaux au sommet des morceaux d'arbre déjà
construits (et éventuellement les arbres eux-mêmes).


%
%
%

\section{Introduction à la calculabilité}\label{section-computability}

\subsection{Présentation générale}

\thingy\textbf{Discussion préalable.} On s'intéresse ici à la question
de savoir ce qu'un \defin[algorithme (en calculabilité)]{algorithme} peut ou ne peut pas faire.
Pour procéder de façon rigoureuse, il faudrait formaliser la notion
d'algorithme (par exemple à travers le concept de machine de Turing) :
on a préféré rester informel sur cette définition — par exemple « un
algorithme est une série d'instruction précises indiquant des calculs
à effectuer étape par étape et qui ne manipulent, à tout moment, que
des données finies » ou « un algorithme est quelque chose qu'on
pourrait, en principe, implémenter sur un ordinateur » — étant entendu
que cette notion est déjà bien connue et comprise, au moins dans la
pratique.  Les démonstrations du fait que tel ou tel problème est
décidable par un algorithme ou que telle ou telle fonction est
calculable par un algorithme deviennent beaucoup moins lisibles quand
on les formalise avec une définition rigoureuse d'algorithme
(notamment, programmer une machine de Turing est encore plus
fastidieux que programmer un ordinateur en assembleur, donc s'il
s'agit d'exhiber un algorithme, c'est probablement une mauvaise idée
de l'écrire sous forme de machine de Turing).

\smallskip

Néanmoins, il est essentiel de savoir que ces formalisations
existent : on peut par exemple évoquer le paradigme du
$\lambda$-calcul de Church (la première formalisation rigoureuse de la
calculabilité), les fonctions générales récursives (=$\mu$-récursives)
à la Herbrand-Gödel-Kleene, les machines de Turing (des machines à
états finis capables de lire, d'écrire et de se déplacer sur un ruban
infini contenant des symboles d'un alphabet fini dont à chaque instant
tous sauf un nombre fini sont des blancs), les machines à registres,
le langage « FlooP » de Hofstadter, etc.  Toutes ces formalisations
sont équivalentes (au sens où, par exemple, elles conduisent à la même
notion de fonction calculable ou calculable partielle, définie
ci-dessous).  La \defin[Church-Turing (thèse de)]{thèse de
  Church-Turing} affirme, au moins informellement, que tout ce qui est
effectivement calculable par un algorithme\footnote{Voire, dans
  certaines variantes de la thèse, tout ce qui est physiquement
  calculable dans notre Univers (y compris par des processus
  quantiques).}  est calculable par n'importe laquelle de ces notions
formelles d'algorithmes, qu'on peut rassembler sous le nom commun de
\defin[calculable (fonction)]{calculabilité au sens de Church-Turing},
ou « calculabilité » tout court.

\smallskip

Notamment, quasiment tous les langages de programmation
informatique\footnote{C, C++, Java, Python, JavaScript, Lisp, OCaml,
  Haskell, Prolog, etc.  Certains langages se sont même révélés
  Turing-complets alors que ce n'était peut-être pas voulu : par
  exemple, HTML+CSS.}, au moins si on ignore les limites des
implémentations et qu'on les suppose capables de manipuler des
entiers, chaînes de caractère, tableaux, etc., de taille arbitraire
(mais toujours finie)\footnote{Autre condition : ne pas utiliser de
  générateur aléatoire matériel.}, sont « Turing-complets »,
c'est-à-dire équivalents dans leur pouvoir de calcul à la
calculabilité de Church-Turing.  Pour imaginer intuitivement la
calculabilité, on peut donc choisir le langage qu'on préfère et
imaginer qu'on programme dedans.  Dans la pratique, pour qu'un langage
soit Turing-complet, il lui suffit d'être capable de manipuler des
entiers de taille arbitraire, de les comparer et de calculer les
opérations arithmétiques dessus, et d'effectuer des tests et des
boucles.

\bigbreak

\thingy Il faut souligner qu'on s'intéresse uniquement à la question
de savoir ce qu'un algorithme peut ou ne peut pas faire
(calculabilité), pas au temps ou aux autres ressources qu'il peut
prendre pour le faire (complexité), et on ne cherche donc pas à rendre
les algorithmes efficaces en quelque sens que ce soit.  Par exemple,
pour arguër qu'il existe un algorithme qui décide si un entier naturel
$n$ est premier ou non, il suffit de dire qu'on peut calculer tous les
produits $pq$ avec $2\leq p,q\leq n-1$ et tester si l'un d'eux est
égal à $n$, peu importe que cet algorithme soit absurdement
inefficace.

\smallskip

De même, nos algorithmes sont capables de manipuler des entiers
arbitrairement grands : ceci permet de dire, par exemple, que toute
chaîne binaire peut être considérée comme un entier, peu importe le
fait que cet entier ait peut-être des milliards de chiffres ; dans les
langages informatiques réels, on a rarement envie de considérer toute
donnée comme un entier, mais en calculabilité on peut se permettre de
le faire.

\thingy\label{computability-all-data-are-integers} Soulignons et
développons le point esquissé au paragraphe précédent : plutôt que de
travailler avec des « mots » (éléments de $\Sigma^*$ avec $\Sigma$ un
alphabet fini), en calculabilité, il est possible, dès que cela est
commode, de remplacer ceux-ci par des entiers naturels.

Il suffit pour cela de choisir un « codage », parfois appelé
\index{Gödel (codage de)|see{codage de Gödel}}\defin{codage de Gödel}
permettant de représenter un élément de
$\Sigma^*$ par un entier naturel : la chose importante est que la
conversion d'un mot en entier ou d'un entier en mot soit calculable
par un algorithme (même très inefficace), si bien que toute
manipulation algorithmique sur des mots puisse être convertie en
manipulation sur des entiers.  Il est commode pour simplifier les
raisonnements (quoique pas indispensable) de supposer que le codage
est bijectif, c'est-à-dire qu'on dispose d'une bijection $\mathbb{N}
\to \Sigma^*$ algorithmiquement calculable et dont la réciproque l'est
aussi.

Il existe toutes sortes de manières d'obtenir un tel codage.  La plus
simple est sans doute la suivante : une fois choisi un ordre total
arbitraire sur l'alphabet (fini !) $\Sigma$, on peut énumérer les mots
sur $\Sigma$ par ordre de taille et, pour une taille donnée, par ordre
lexicographique (par exemple : $\varepsilon, a, b,\penalty0 aa, ab,
ba, bb, aaa, aab...$ sur l'alphabet $\{a,b\}$) ; cette énumération est
visiblement faisable algorithmiquement, et quitte à numéroter au fur
et à mesure qu'on énumère, on obtient un codage comme souhaité (par
exemple $0 \leftrightarrow\varepsilon$, $1 \leftrightarrow a$, $2
\leftrightarrow b$, $3 \leftrightarrow ab$, etc.).  Insistons sur le
fait que ce n'est là qu'une possibilité parmi d'autres et que les
détails du codage sont sans importance tant qu'il est calculable ;
l'important est que ce soit faisable en théorie et donc qu'on peut
considérer qu'au lieu de mots on a affaire à des entiers naturels.

Il va de soi que la concaténation de deux mots, la longueur d'un mot,
le miroir d'un mot, sont tous calculables algorithmiquement.

\smallskip

Grâce au codage de Gödel, on peut considérer que, dans cette partie,
le terme de \index{langage}« langage » désigne non plus une partie de
$\Sigma^*$ mais une partie de $\mathbb{N}$.

{\footnotesize (Le remplacement des mots par des entiers naturels en
  utilisant un codage comme on vient de le dire est assez standard en
  calculabilité.  Il ne doit pas dérouter : on peut imaginer qu'on a
  affaire à $\Sigma^*$ à chaque fois qu'il est question
  de $\mathbb{N}$ ci-dessous, et on a toujours affaire à la théorie
  des langages.  Le point central est justement que cela n'a pas
  d'importance ; comme $\mathbb{N}$ est un objet mathématiquement plus
  simple, c'est surtout pour cela qu'il est utilisé à la place.)\par}

\bigbreak

\thingy\textbf{Terminaison des algorithmes.} Un algorithme qui
effectue un calcul utile doit certainement terminer en temps fini.
Néanmoins, même si on voudrait ne s'intéresser qu'à ceux-ci, il n'est
pas possible d'ignorer le « problème » des algorithmes qui ne
terminent jamais (et ne fournissent donc aucun résultat).  C'est le
point central de la calculabilité (et du théorème de
Turing \ref{undecidability-of-halting-problem} ci-dessous) qu'on ne
peut pas se débarrasser des algorithmes qui ne terminent pas : on ne
peut pas, par exemple, formaliser une notion suffisante\footnote{Tout
  dépend, évidemment, de ce qu'on appelle « suffisant » : il existe
  bien des notions de calculabilité, plus faibles que celle de
  Church-Turing, où tout calcul termine, voir par exemple la notion de
  fonction « primitive récursive » ou le langage « BlooP » de
  Hofstadter ; mais de telles notions ne peuvent pas disposer d'une
  machine universelle comme expliqué plus loin (en raison d'un
  argument diagonal), donc elles sont nécessairement incomplètes en un
  certain sens.} de calculabilité dans laquelle tout algorithme
termine toujours ; ni développer un langage de programmation
suffisamment général dans lequel il est impossible qu'un programme
« plante » indéfiniment ou parte en boucle infinie.

{\footnotesize (Cette subtilité est d'ailleurs sans doute en partie
  responsable de la difficulté historique à dégager la bonne notion
  d'« algorithme » : on a commencé par développer des notions
  d'algorithmes terminant forcément, comme les fonctions primitives
  récursives, et on se rendait bien compte que ces notions étaient
  forcément toujours incomplètes.)\par}

\bigbreak

\begin{defn}\label{definition-computable-function-or-set}
On dit qu'une fonction $f\colon\mathbb{N}\to\mathbb{N}$ est
\defin[calculable (fonction)]{calculable} (ou \index{récursive
  (fonction)|see{calculable}}« récursive ») lorsqu'il existe un
algorithme qui prend en entrée $n\in\mathbb{N}$, termine toujours en
temps fini, et calcule (renvoie) $f(n)$.

\smallskip

On dit qu'un ensemble $A \subseteq \mathbb{N}$ (un « langage »,
cf. \ref{computability-all-data-are-integers}) est \defin[décidable (langage)]{décidable}
(ou \index{calculable (langage)|see{décidable}}« calculable » ou
\index{récursif (langage)|see{décidable}}« récursif ») lorsque sa
fonction indicatrice $\mathbf{1}_A \colon \mathbb{N} \to \mathbb{N}$
(valant $1$ sur $A$ et $0$ sur son complémentaire) est calculable.
Autrement dit : lorsqu'il existe un algorithme qui prend en entrée
$n\in\mathbb{N}$, termine toujours en temps fini, et renvoie
« oui » ($1$) si $n\in A$, et « non » ($0$) si $n\not\in A$ (on dira
que l'algorithme « décide » $A$).

\smallskip

On dit qu'une fonction partielle
$f\colon\mathbb{N}\dasharrow\mathbb{N}$ (c'est-à-dire une fonction
définie sur une partie de $\mathbb{N}$, appelé ensemble de définition
de $f$) est \defin[calculable partielle (fonction)]{calculable
  partielle} (ou « récursive partielle ») lorsqu'il existe un
algorithme qui prend en entrée $n\in\mathbb{N}$, termine en temps fini
si et seulement si $f(n)$ est définie, et dans ce cas calcule
(renvoie) $f(n)$.  (Une fonction calculable est donc simplement une
fonction calculable partielle qui est toujours définie : on dira
parfois « calculable totale » pour souligner ce fait.)

On utilisera la notation $f(n)\downarrow$ pour signifier le fait que
la fonction calculable partielle $f$ est définie en $n$, c'est-à-dire,
que l'algorithme en question termine.

\smallskip

On dit qu'un ensemble $A \subseteq \mathbb{N}$ est
\defin{semi-décidable} (ou
\index{semi-calculable|see{semi-décidable}}« semi-calculable » ou
\index{semi-récursif|see{semi-décidable}}« semi-récursif »)
lorsque la fonction partielle $\mathbb{N}\dasharrow\mathbb{N}$ définie
exactement sur $A$ et y valant $1$, est calculable partielle.
Autrement dit : lorsqu'il existe un algorithme qui prend en entrée
$n\in\mathbb{N}$, termine en temps fini si et seulement si $n \in A$,
et renvoie « oui » ($1$) dans ce cas\footnote{En fait, la valeur
  renvoyée n'a pas d'importance ; on peut aussi définir un ensemble
  semi-décidable comme l'ensemble de définition d'une fonction
  calculable partielle.}  (on dira que l'algorithme
« semi-décide » $A$).
\end{defn}

\smallskip

On s'est limité ici à des fonctions d'une seule variable (entière),
mais il n'y a pas de difficulté à étendre ces notions à plusieurs
variables, et de parler de fonction calculable $\mathbb{N}^k \to
\mathbb{N}$ (voire $\mathbb{N}^* \to \mathbb{N}$ avec $\mathbb{N}^*$
l'ensemble des suites finies d'entiers naturels) ou de fonction
calculable partielle de même type : de toute manière, on peut « coder »
un couple d'entiers naturels comme un seul entier naturel (par exemple
par $(m,n) \mapsto 2^m(2n+1)-1$, qui définit une bijection calculable
$\mathbb{N}^2 \to \mathbb{N}$), ou bien sûr un nombre fini quelconque
(même variable), ce qui permet de faire « comme si » on avait toujours
affaire à un seul paramètre entier.

{\footnotesize\thingy\textbf{Complément :} Comme on n'a pas défini
  formellement la notion d'algorithme, il peut être utile de signaler
  explicitement les faits suivants (qui devraient être évidents sur
  toute notion raisonnable d'algorithme) : les fonctions constantes
  sont calculables ; les opérations arithmétiques usuelles sont
  calculables ; les projections $(n_1,\ldots,n_k) \mapsto n_i$ sont
  calculables, ainsi que la fonction qui à $(m,n,p,q)$ associe $p$ si
  $m=n$ et $q$ sinon ; toute composée de fonctions calculables
  (partielle ou totale) est calculable idem ; si $\underline{m}
  \mapsto g(\underline{m})$ est calculable (partielle ou totale,
  où $\underline{m}$ désigne une liste d'entiers) et
  que $(\underline{m}, n, v) \mapsto h(\underline{m}, n, v)$ l'est,
  alors la fonction $f$ définie par récurrence par $f(\underline{m},0)
  = g(\underline{m})$ et $f(\underline{m},n+1) = h(\underline{m}, n,
  f(\underline{m},n))$ est encore calculable idem (algorithmiquement,
  il s'agit juste de boucler $n$ fois) ; et enfin, si $(\underline{m},
  n) \mapsto g(\underline{m},n)$ est calculable partielle, alors la
  fonction $f$ (aussi notée $\mu_n g$) définie par $f(\underline{m}) =
  \min\{n : g(\underline{m},n) = 0 \land \forall n'<n
  (g(\underline{m},n')\downarrow)\}$ (et non définie si ce $\min$
  n'existe pas) est calculable partielle (algorithmiquement, on teste
  $g(\underline{m},0),g(\underline{m},1),g(\underline{m},2)\ldots$
  jusqu'à tomber sur $0$).  Ces propriétés peuvent d'ailleurs servir à
  \emph{définir} rigoureusement la notion de fonction calculable,
  c'est le modèle des fonctions « générales récursives ».  (Dans ce
  qui précède, la notation $\underline{m}$ signifie
  $m_1,\ldots,m_k$.)\par}

\thingy\textbf{Exemples :} L'ensemble des nombres pairs, des carrés
parfaits, des nombres premiers, sont décidables, c'est-à-dire qu'il
est algorithmique de savoir si un nombre est pair, parfait, ou
premier.  Quitte éventuellement à coder les mots d'un alphabet fini
comme des entiers naturels
(cf. \ref{computability-all-data-are-integers}), tout langage
rationnel, et même tout langage défini par une grammaire hors
contexte, est décidable (cf. \ref{rational-languages-are-recognizable}
et \ref{algebraic-languages-are-decidable}).

On verra plus bas des exemples d'ensembles qui ne le sont pas, et qui
sont ou ne sont pas semi-décidables.

\medbreak

Les deux propositions suivantes, outre leur intérêt intrinsèque,
servent à donner des exemples du genre de manipulation qu'on peut
faire avec la notion de calculabilité et d'algorithme :

\begin{prop}\label{decidable-iff-semidecidable-and-complement}
Un ensemble $A \subseteq \mathbb{N}$ est décidable si et seulement si
$A$ et $\mathbb{N}\setminus A$ sont tous les deux semi-décidables.
\end{prop}
\begin{proof}
Il est évident qu'un ensemble décidable est semi-décidable (si un
algorithme décide $A$, on peut l'exécuter puis effectuer une boucle
infinie si la réponse est « non » pour obtenir un algorithme qui
semi-décide $A$) ; il est également évident que le complémentaire d'un
ensemble décidable est décidable (quitte à échanger les réponses
« oui » et « non » dans un algorithme qui le décide).  Ceci montre
qu'un ensemble décidable est semi-décidable de complémentaire
semi-décidable, i.e., la partie « seulement si ».  Montrons maintenant
le « si » : si on dispose d'algorithmes $T_1$ et $T_2$ qui
semi-décident respectivement $A$ et son complémentaire, on peut lancer
leur exécution en parallèle sur $n \in \mathbb{N}$ (c'est-à-dire
exécuter une étape de $T_1$ puis une étape de $T_2$, puis de $T_1$, et
ainsi de suite jusqu'à ce que l'un des deux termine) : comme il y en a
toujours (exactement) un qui termine, selon lequel c'est, ceci permet
de décider algorithmiquement si $n \in A$ ou $n \not\in A$.
\end{proof}

\begin{prop}
Un ensemble $A \subseteq \mathbb{N}$ non vide est semi-décidable si et
seulement si il existe une fonction calculable $f \colon \mathbb{N}
\to \mathbb{N}$ dont l'image ($f(\mathbb{N})$) vaut $A$ (on dit aussi
que $A$ est « \defin{calculablement énumérable} » ou
« \index{récursivement énumérable|see{calculablement
    énumérable}}récursivement énumérable »).
\end{prop}
\begin{proof}
Montrons qu'un ensemble semi-décidable non vide est calculablement
énumérable.  Fixons $n_0 \in A$ une fois pour toutes.  Soit $T$ un
algorithme qui semi-décide $A$.  On définit une fonction $f \colon
\mathbb{N}^2 \to \mathbb{N}$ de la façon suivante : $f(m,n) = n$
lorsque l'algorithme $T$, exécuté sur l'entrée $n$, termine au
plus $m$ étapes ; sinon, $f(m,n) = n_0$.  On a bien sûr $f(m,n) \in A$
dans tous les cas ; par ailleurs, si $n \in A$, comme l'algorithme $T$
appliqué à $n$ doit terminer, on voit que pour $m$ assez grand on a
$f(m,n) = n$, donc $n$ est bien dans l'image de $f$.  Ceci montre que
$f(\mathbb{N}^2) = A$.  Passer à $f\colon \mathbb{N} \to\mathbb{N}$
est alors facile en composant par une bijection calculable $\mathbb{N}
\to \mathbb{N}^2$ (par exemple la réciproque de $(m,n) \mapsto
2^m(2n+1)-1$).

Réciproquement, si $A$ est calculablement énumérable, disons $A =
f(\mathbb{N})$ avec $f$ calculable, on obtient un algorithme qui
semi-décide $A$ en calculant successivement $f(0)$, $f(1)$, $f(2)$,
etc., jusqu'à trouver un $k$ tel que $f(k)=n$ (où $n$ est l'entrée
proposée), auquel cas l'algorithme renvoie « oui » (et sinon, il ne
termine jamais puisqu'il effectue une boucle infinie à la recherche
d'un tel $k$).
\end{proof}

{\footnotesize\thingy\textbf{Éclaircissement :} Les deux démonstrations
  ci-dessus font appel à la notion intuitive d'« étape » de
  l'exécution d'un algorithme.  Un peu plus précisément, pour chaque
  entier $m$ et chaque algorithme $T$, il est possible d'« exécuter au
  plus $m$ étapes » de l'algorithme $T$, c'est-à-dire commencer
  l'exécution de celui-ci, et si elle n'est pas finie au bout de $m$
  étapes, s'arrêter (on n'aura pas le résultat de l'exécution de $T$,
  juste l'information « ce n'est pas encore fini » et d'éventuels
  résultats intermédiaires, mais on peut décider de faire autre chose,
  y compris reprendre l'exécution plus tard).  La longueur d'une
  « étape » n'est pas spécifiée et n'a pas d'importance, les choses
  qui importent sont que (A) le fait d'exécuter les $m$ premières
  étapes de $T$ termine toujours (c'est bien l'intérêt), et (B) si
  l'algorithme $T$ termine effectivement, alors pour $m$ suffisamment
  grand, exécuter au plus $m$ étapes donne bien le résultat final
  de $T$ comme résultat.\par}

{\footnotesize\thingy\textbf{Complément/exercice :} Un ensemble $A
  \subseteq \mathbb{N}$ infini est décidable si et seulement si il
  existe une fonction calculable $f \colon \mathbb{N} \to \mathbb{N}$
  \underline{strictement croissante} dont l'image vaut $A$.
  (Esquisse : si $A$ est décidable, on peut trouver son $n$-ième
  élément par ordre croissant en testant l'appartenance à $A$ de tous
  les entiers naturels dans l'ordre jusqu'à trouver le $n$-ième qui
  appartienne ; réciproquement, si on a une telle fonction, on peut
  tester l'appartenance à $A$ en calculant les valeurs de la fonction
  jusqu'à tomber sur l'entier à tester ou le dépasser.)  En mettant
  ensemble ce fait et la proposition, on peut en déduire le fait
  suivant : tout ensemble semi-décidable infini a un sous-ensemble
  décidable infini (indication : prendre une fonction qui énumère
  l'ensemble et jeter toute valeur qui n'est pas strictement plus
  grande que toutes les précédentes).
  Cf. exercice \ref{decidable-iff-image-of-computable-increasing}.\par}

\subsection{Machine universelle, et problème de l'arrêt}

\thingy\textbf{Codage et machine universelle.}  Les algorithmes sont
eux-mêmes représentables par des mots sur un alphabet fini donc, si on
préfère (cf. \ref{computability-all-data-are-integers}), par des
entiers naturels : on parle aussi de \defin{codage de Gödel} des
algorithmes/programmes par des entiers.  On obtient donc une
énumération $\varphi_0, \varphi_1, \varphi_2, \varphi_3\ldots$ de
toutes les fonctions calculables partielles (la fonction $\varphi_e$
étant la fonction que calcule l'algorithme [codé par l'entier] $e$,
avec la convention que si cet algorithme est syntaxiquement invalide
ou erroné pour une raison quelconque, la fonction $\varphi_e$ est
simplement non-définie partout).  Les détails de cette énumération
dépendent de la formalisation utilisée pour la calculabilité et ne
sont pas importants\footnote{Par exemple, ceux qui aiment le langage
  Java peuvent considérer que $\varphi_e$ désigne le résultat de
  l'exécution du programme Java représenté par le mot codé par $e$, si
  ce programme est valable (remplacer Java par tout autre langage
  préféré).}.

On aura tendance, dans la suite, à écrire « le $e$-ième algorithme »,
voire « l'algorithme $e$ » (ou « le programme $e$ »), pour désigner
l'algorithme codé par l'entier naturel $e$ (c'est-à-dire qu'on
identifie librement un programme et son codage de Gödel) : on peut
donc écrire que $\varphi_e(n)$ est le résultat de l'exécution du
programme $e$ quand on lui fournit $n$ en entrée (ou non défini si
cette exécution ne termine pas).

\medskip

Un point crucial dans cette numérotation des algorithmes par des
entiers naturels $e$ est l'existence d'une \defin[universelle
  (machine)]{machine universelle}, c'est-à-dire d'un algorithme $U$
qui prend en entrée un entier $e$ (codant un algorithme $T$) et un
entier $n$, et effectue la même chose que $T$ sur l'entrée $n$ (i.e.,
$U$ termine sur les entrées $e$ et $n$ si et seulement si $T$ termine
sur l'entrée $n$, et, dans ce cas, renvoie la même valeur).  Autrement
dit, le calcul de $\varphi_e(n)$ en fonction de $e$ et $n$ est
lui-même algorithmique.

\smallskip

\underline{Informatiquement}, ceci représente le fait que les
programmes informatiques sont eux-mêmes représentables
informatiquement : dans un langage de programmation Turing-complet, on
peut écrire un \emph{interpréteur} $U$ pour le langage lui-même (ou
pour un autre langage Turing-complet), c'est-à-dire un programme qui
prend en entrée la représentation $e$ d'un autre programme et qui
exécute ce programme (sur une entrée $n$).

\underline{Mathématiquement}, on peut le formuler comme le fait que la
fonction (partielle) $(e,n) \mapsto \varphi_e(n)$ (= résultat du
$e$-ième algorithme appliqué sur l'entrée $n$) est elle-même
calculable partielle.

\underline{Philosophiquement}, cela signifie que la notion d'exécution
d'un algorithme est elle-même algorithmique : on peut écrire un
algorithme qui, donnée une description (formelle !) d'un algorithme et
une entrée à laquelle l'appliquer, effectue l'exécution de
l'algorithme fourni sur l'entrée fournie.

\smallskip

On ne peut pas démontrer ce résultat ici faute d'une description
rigoureuse d'un modèle de calcul précis, mais il n'a rien de
conceptuellement difficile (même s'il peut être fastidieux à écrire
dans les détails : écrire un interpréteur d'un langage de
programmation demande un minimum d'efforts).  La machine universelle
dépend, bien sûr, du codage des algorithmes par des entiers naturels.

{\footnotesize\thingy\textbf{Compléments :} Les deux résultats
  classiques suivants sont pertinents en lien avec la numérotation des
  fonctions calculables partielles.  $\bullet$ Le \emph{théorème de la
    forme normale de Kleene} assure qu'il existe un ensemble
  \underline{décidable} $\mathscr{T} \subseteq \mathbb{N}^4$ tel que
  $\varphi_e(n)$ soit défini si et seulement si il existe $m,v$ tels
  que $(e,n,m,v) \in \mathscr{T}$, et dans ce cas $\varphi_e(n) = v$
  (pour s'en convaincre, il suffit de définir $\mathscr{T}$ comme
  l'ensemble des $(e,n,m,v)$ tels que le $e$-ième algorithme exécuté
  sur l'entrée $n$ termine en au plus $m$ étapes et renvoie le
  résultat $v$ : le fait qu'on dispose d'une machine universelle et
  qu'on puisse exécuter $m$ étapes d'un algorithme assure que cet
  ensemble est bien décidable — il est même « primitif
  récursif »). $\bullet$ Le \index{s-m-n (théorème)}\emph{théorème
    s-m-n} assure qu'il existe une fonction calculable $s$ telle que
  $\varphi_{s(e,\underline{m})}(\underline{n}) =
  \varphi_e(\underline{m},\underline{n})$ (intuitivement, donné un
  algorithme qui prend plusieurs entrées et des valeurs
  $\underline{m}$ de certaines de ces entrées, on peut fabriquer un
  nouvel algorithme dans lequel ces valeurs ont été fixées — c'est à
  peu près trivial — mais de plus, cette transformation est
  \emph{elle-même algorithmique}, i.e., on peut algorithmiquement
  substituer des valeurs $\underline{m}$ dans un programme [codé par
    l'entier] $e$ : c'est intuitivement clair, mais cela ne peut pas
  se démontrer avec les seules explications données ci-dessus sur
  l'énumération des fonctions calculables partielles, il faut regarder
  précisément comment le codage standard est fait pour une
  formalisation de la calculabilité).\par}

\smallskip

La machine universelle n'a rien de « magique » : elle se contente de
suivre les instructions de l'algorithme $T$ qu'on lui fournit, et
termine si et seulement si $T$ termine.  Peut-on savoir à l'avance si
$T$ terminera ?  C'est le fameux « problème de l'arrêt ».

\medbreak

Intuitivement, le « problème de l'arrêt » est la question
« l'algorithme suivant termine-t-il sur l'entrée suivante » ?

\begin{defn}
On appelle \index{arrêt (problème de l')|see{problème de
    l'arrêt}}\defin{problème de l'arrêt} (ou « langage de l'arrêt »)
l'ensemble des couples $(e,n)$ tels que le $e$-ième algorithme termine
sur l'entrée $n$, i.e., $\{(e,n) \in \mathbb{N}^2 :
\varphi_e(n)\downarrow\}$ (où la notation « $\varphi_e(n)\downarrow$ »
signifie que $\varphi_e(n)$ est défini, i.e., l'algorithme termine).
Quitte à coder les couples d'entiers naturels par des entiers naturels
(par exemple par $(e,n) \mapsto 2^e(2n+1)-1$), on peut voir le problème
de l'arrêt comme une partie de $\mathbb{N}$.  On peut aussi
préférer\footnote{Même si au final c'est équivalent, c'est \textit{a
    priori} plus fort de dire que $\{e \in \mathbb{N} :
  \varphi_e(e)\downarrow\}$ n'est pas décidable que de dire que
  $\{(e,n) \in \mathbb{N}^2 : \varphi_e(n)\downarrow\}$ ne l'est pas.}
définir le problème de l'arrêt comme $\{e \in \mathbb{N} :
\varphi_e(e)\downarrow\}$, on va voir dans la démonstration ci-dessous
que c'est cet ensemble-là qui la fait fonctionner.
\end{defn}

{\footnotesize (On pourrait aussi définir le problème de l'arrêt comme
  $\{e \in \mathbb{N} : \varphi_e(0)\downarrow\}$ si on voulait, ce
  serait moins pratique pour la démonstration, mais cela ne changerait
  rien au résultat comme on peut le voir en appliquant le théorème
  s-m-n.)\par}

\begin{thm}[Turing]\index{Turing (théorème de)}\label{undecidability-of-halting-problem}
Le problème de l'arrêt est semi-décidable mais non décidable.
\end{thm}
\begin{proof}
Le problème de l'arrêt est semi-décidable en vertu de l'existence
d'une machine universelle : donnés $e$ et $n$, on exécute le $e$-ième
algorithme sur l'entrée $n$ (c'est ce que fait la machine
universelle), et s'il termine on renvoie « oui » (et s'il ne termine
pas, bien sûr, la seule possibilité est de ne pas terminer).

Montrons par l'absurde que le problème de l'arrêt n'est pas décidable.
S'il l'était, on pourrait définir un algorithme qui, donné un entier
$e$, effectue les calculs suivants : (1º) utiliser le problème de
l'arrêt (supposé décidable !) pour savoir, algorithmiquement en temps
fini, si le $e$-ième algorithme termine quand on lui passe son propre
numéro $e$ en entrée, i.e., si $\varphi_e(e)\downarrow$, et ensuite
(2º) si oui, effectuer une boucle infinie, et si non, terminer, en
renvoyant, disons, $42$.  L'algorithme qui vient d'être décrit aurait
un certain numéro, disons, $p$, et la description de l'algorithme fait
que, quel que soit $e$, la valeur $\varphi_p(e)$ est indéfinie si
$\varphi_e(e)$ est définie tandis que $\varphi_p(e)$ est définie (de
valeur $42$) si $\varphi_e(e)$ est indéfinie.  En particulier, en
prenant $e=p$, on voit que $\varphi_p(p)$ devrait être défini si et
seulement si $\varphi_p(p)$ n'est pas défini, ce qui est une
contradiction.
\end{proof}

La démonstration ci-dessus est une instance de l'« argument diagonal »
de Cantor, qui apparaît souvent en mathématiques.  (La « diagonale »
en question étant le fait qu'on considère $\varphi_e(e)$, i.e., on
passe le numéro $e$ d'un algorithme en argument à cet algorithme
lui-même, donc on regarde la diagonale de la fonction de deux
variables $(e,n) \mapsto \varphi_e(n)$ ; en modifiant les valeurs sur
cette diagonale, on produit une fonction qui ne peut pas se trouver
dans une ligne $\varphi_p$.)  Une variante facile du même argument
permet de fabriquer des ensembles non semi-décidables (voir
\ref{example-set-of-total-functions} ci-dessous), ou bien on peut
appliquer ce qui précède :

\begin{cor}
Le complémentaire du problème de l'arrêt n'est pas semi-décidable.
\end{cor}
\begin{proof}
On a vu que le problème de l'arrêt n'est pas décidable, et qu'un
ensemble est décidable si et seulement si il est semi-décidable et que
son complémentaire l'est aussi : comme le problème de l'arrêt est bien
semi-décidable, son complémentaire ne l'est pas.
\end{proof}

{\footnotesize\thingy\textbf{Complément :} L'argument diagonal est
  aussi au cœur du (voire, équivalent au) \index{Kleene (théorème de récursion de)}\emph{théorème de récursion
    de Kleene}, qui affirme que pour toute fonction calculable
  partielle $h\colon\mathbb{N}^2\dasharrow\mathbb{N}$, il existe $p$
  tel que $\varphi_p(n) = h(p,n)$ pour tout $n$ (la signification
  intuitive de ce résultat est qu'on peut supposer qu'un programme a
  accès à son propre code source $p$, i.e., on peut programmer comme
  s'il recevait en entrée un entier $p$ codant ce code source ; ceci
  permet par exemple — de façon anecdotique mais amusante — d'écrire
  des programmes, parfois appelés « quines », qui affichent leur
  propre code source sans aller le chercher sur disque ou autre
  tricherie).  \textit{Démonstration :} donné $e \in \mathbb{N}$, on
  considère $s(e,m)$ tel que $\varphi_{s(e,m)}(n) = \varphi_e(m,n)$ :
  le théorème s-m-n (cf. ci-dessus) assure qu'une telle fonction
  calculable $(e,m) \mapsto s(e,m)$ existe, et $(e,n) \mapsto
  h(s(e,e), n)$ est alors aussi calculable partielle ; il existe donc
  $q$ tel que $\varphi_q(e,n) = h(s(e,e), n)$ : on pose $p = s(q,q)$,
  et on a $\varphi_p(n) = \varphi_q(q,n) = h(s(q,q), n) = h(p, n)$,
  comme annoncé. \smiley\ La non-décidabilité du problème de l'arrêt
  s'obtient en appliquant (de nouveau par l'absurde) ce résultat à
  $h(e, n)$ la fonction qui n'est pas définie si $\varphi_e(n)$ l'est
  et qui vaut $42$ si $\varphi_e(n)$ n'est pas définie.\par}

\medbreak

La non-décidabilité du problème de l'arrêt est un résultat
fondamental, car très souvent les résultats de non-décidabilité soit
sont démontrés sur un modèle semblable, soit s'y ramènent
directement : pour montrer qu'un certain ensemble $A$ (un
« problème ») n'est pas décidable, on cherche souvent à montrer que si
un algorithme décidant $A$ existait, on pourrait s'en servir pour
construire un algorithme résolvant le problème de l'arrêt.

{\footnotesize\thingy\label{example-set-of-total-functions}\textbf{Bonus
    / exemple(s) :} L'ensemble des $e \in \mathbb{N}$ tels que la
  fonction calculable partielle $\varphi_e$ soit \underline{totale}
  (i.e., définie sur tout $\mathbb{N}$) n'est pas semi-décidable.  En
  effet, s'il l'était, d'après ce qu'on a vu, il serait
  « calculablement énumérable », c'est-à-dire qu'il existerait une
  fonction calculable $f\colon\mathbb{N}\to\mathbb{N}$ dont l'image
  soit exactement l'ensemble des $e$ pour lesquels $\varphi_e$ est
  totale, i.e., toute fonction calculable totale s'écrirait sous la
  forme $\varphi_{f(k)}$ pour un certain $k$.  Mais la fonction $n
  \mapsto \varphi_{f(n)}(n) + 1$ est calculable totale, donc il
  devrait exister un $m$ tel que cette fonction s'écrive
  $\varphi_{f(m)}$, c'est-à-dire $\varphi_{f(m)}(n) =
  \varphi_{f(n)}(n) + 1$, et on aurait alors en particulier
  $\varphi_{f(m)}(m) = \varphi_{f(m)}(m) + 1$, une
  contradiction. $\bullet$ Son complémentaire, c'est-à-dire l'ensemble
  des $e \in \mathbb{N}$ tels que la fonction calculable partielle
  $\varphi_e$ \underline{ne soit pas} totale, n'est pas non plus
  semi-décidable.  En effet, supposons qu'il existe un algorithme qui,
  donné $e$, termine si et seulement si $\varphi_e$ n'est pas totale.
  Donnés $e$ et $m$, considérons l'algorithme qui prend une entrée
  $n$, \emph{ignore} celle-ci, et effectue le calcul $\varphi_e(m)$ :
  ceci définit une fonction calculable partielle (soit totale et
  constante, soit définie nulle part !) $\varphi_{s(e,m)}$ où $s$ est
  calculable (on applique ici le théorème s-m-n) — en appliquant à
  $s(e,m)$ l'algorithme supposé semi-décider si une fonction récursive
  partielle est non-totale, on voit qu'ici il semi-décide si
  $\varphi_e(m)$ est non-défini, autrement dit on semi-décide le
  complémentaire du problème de l'arrêt, et on a vu que ce n'était pas
  possible !\par}

{\footnotesize\thingy\textbf{Exercice :} Considérons une fonction $h$
  qui à $e$ associe un nombre au moins égal au nombre d'étapes
  (cf. ci-dessus) du calcul de $\varphi_e(e)$, si celui-ci termine, et
  une valeur quelconque si $\varphi_e(e)$ n'est pas défini.  Alors $h$
  n'est pas calculable.  (Indication : si elle l'était, on pourrait
  décider si $\varphi_e(e)$ est défini en exécutant son calcul pendant
  $h(e)$ étapes.)  On peut même montrer que $H(n) := \max\{h(i) :
  i\leq n\}$ domine asymptotiquement n'importe quelle fonction
  calculable mais c'est un peu plus difficile.\par}

\bigbreak

{\footnotesize\thingy\textbf{Application à la logique :} Sans rentrer
  dans les détails de ce que signifie un « système formel », on peut
  esquisser, au moins informellement, les arguments suivants.
  Imaginons qu'on ait formalisé la notion de démonstration
  mathématique (c'est-à-dire qu'on les écrit comme des mots dans un
  alphabet indiquant quels axiomes et quelles règles logiques sont
  utilisées) : même sans savoir quelle est exactement la logique
  formelle, le fait de \emph{vérifier} qu'une démonstration est
  correcte doit certainement être algorithmique (il s'agit simplement
  de vérifier que chaque règle a été correctement appliquée),
  autrement dit, l'ensemble des démonstrations est décidable.
  L'ensemble des théorèmes, lui, est semi-décidable (on a un
  algorithme qui semi-décide si un certain énoncé est un théorème en
  énumérant toutes les chaînes de caractères possibles et en cherchant
  s'il s'agit d'une démonstration valable dont la conclusion est
  l'énoncé recherché).  Or l'ensemble des théorèmes n'est pas
  décidable : en effet, si on avait un algorithme qui permet de
  décider si un énoncé mathématique est un théorème, on pourrait
  appliquer cet algorithme à l'énoncé formel (*)« le $e$-ième
  algorithme termine sur l'entrée $e$ », en observant qu'un tel
  énoncé, s'il est vrai, est forcément démontrable (i.e., si
  l'algorithme termine, on peut \emph{démontrer} ce fait en écrivant
  étape par étape l'exécution de l'algorithme pour constituer une
  démonstration qu'il a bien été appliqué jusqu'au bout et a terminé),
  et en espérant que s'il est démontrable alors il est vrai : on
  aurait alors une façon de décider le problème de l'arrêt, une
  contradiction.  Mais du coup, l'ensemble des non-théorèmes ne peut
  pas être semi-décidable ; or comme l'ensemble des énoncés $P$ tels
  que $\neg P$ (« non-$P$ », la négation logique de $P$) soit un
  théorème est semi-décidable (puisque l'ensemble des théorèmes
  l'est), ils ne peuvent pas coïncider.  Ceci montre qu'il existe un
  énoncé tel que ni $P$ ni $\neg P$ ne sont des théorèmes : c'est une
  forme du \emph{théorème de Gödel} que Turing cherchait à démontrer ;
  mieux : en appliquant aux énoncés du type (*), on montre ainsi qu'il
  existe un algorithme qui \emph{ne termine pas} mais dont la
  non-terminaison \emph{n'est pas démontrable}.  (Modulo quelques
  hypothèses qui n'ont pas été explicitées sur le système formel dans
  lequel on travaille.)\par}


%
%
%

\appendix
\section{Exercices}\label{section-exercises}

\subsection{Langages rationnels et automates}

\exercice

Soit $\Sigma = \{0,1\}$.  On appelle \emph{mot binaire} un mot sur
l'alphabet $\Sigma$, et mot binaire \emph{normalisé} un mot binaire
qui \emph{soit} commence par $1$, \emph{soit} est exactement égal
à $0$.

(1) Montrer que le langage $L_n = \{0, 1, 10, 11, 100, 101,\ldots\}$
des mots binaires normalisés est rationnel en exhibant directement une
expression rationnelle qui le dénote, et montrer qu'il est
reconnaissable en exhibant directement un automate fini qui le
reconnaît.

(2) On définit la \emph{valeur numérique} d'un mot binaire
$x_{n-1}\cdots x_0$ comme $\sum_{i=0}^{n-1} x_i 2^i$ (où $x_i$ vaut
$0$ ou $1$ et est numéroté de $0$ pour le chiffre le plus à droite
à $n-1$ pour le plus à gauche) ; la valeur numérique du mot
vide $\varepsilon$ est $0$.

Parmi les langages suivants, certains sont rationnels.  Dire lesquels
et justifier brièvement pourquoi ils le sont (on ne demande pas de
justifier pourquoi ceux qui ne sont pas rationnels ne le sont pas) :

(a) le langage $L_a$ des mots binaires dont la valeur numérique est
paire,

(b) le langage $L_b$ des mots binaires \emph{normalisés} dont la
valeur numérique est paire,

(c) le langage $L_c$ des mots binaires dont la valeur numérique est
multiple de $3$ (indication : selon que $n$ est congru à $0$, $1$ ou
$2$ modulo $3$, et selon que $x$ vaut $0$ ou $1$, à quoi est congru
$2n+x$ modulo $3$ ?),

(d) le langage $L_d$ des mots binaires dont la valeur numérique est un
nombre premier,

(e) le langage $L_e$ des mots binaires dont la valeur numérique est
une puissance de $2$, i.e., de la forme $2^i$ pour $i\in\mathbb{N}$,

\begin{corrige}
(1) On peut écrire $L_n = L_{0|1(0|1){*}}$, langage dénoté par
  l'expression rationnelle $0|1(0|1){*}$.  Ce langage est reconnu, par
  exemple, par le DFAI suivant :
\begin{center}
%%% begin ex1p1 %%%

\begin{tikzpicture}[>=latex,line join=bevel,automaton]
%%
\node (qY) at (97bp,105bp) [draw,circle,state,final] {$Y$};
  \node (qX) at (18bp,74bp) [draw,circle,state,initial] {$X$};
  \node (qZ) at (97bp,18bp) [draw,circle,state,final] {$Z$};
  \draw [->] (qX) ..controls (45.279bp,84.58bp) and (58.943bp,90.081bp)  .. node[auto] {$0$} (qY);
  \draw [->] (qX) ..controls (44.52bp,55.44bp) and (60.758bp,43.63bp)  .. node[auto] {$1$} (qZ);
  \draw [->] (qZ) to[loop below] node[auto] {$0,1$} (qZ);
%
\end{tikzpicture}

%%% end ex1p1 %%%
\end{center}

(2) (a) Le langage $L_a$ est rationnel car il s'agit du langage des
mots binaires qui soit sont le mot vide soit finissent par $0$ : il
est dénoté par l'expression rationnelle
$\underline{\varepsilon}|(0|1){*}0$.\spaceout (b) On a $L_b = L_a \cap
L_n$ et on a vu que $L_a$ et $L_n$ sont rationnels, donc $L_b$ l'est
aussi (on peut aussi exhiber une expression rationnelle qui le
dénote : $0|1(0|1){*}0$).

(c) Ajouter un $0$ ou un $1$ à la fin d'un mot binaire de valeur
numérique $n$ le transforme en un mot de valeur numérique $2n+x$
où $x$ est le chiffre affixé.  Considérons les six combinaisons entre
les trois cas possibles de la valeur numérique $n$ modulo $3$ et les
deux cas possibles de la valeur de $x$ :
\begin{center}
\begin{tabular}{c|c|c}
$n\equiv?\pmod{3}$&$x=?$&$2n+x\equiv?\pmod{3}$\\\hline
$0$&$0$&$0$\\
$0$&$1$&$1$\\
$1$&$0$&$2$\\
$1$&$1$&$0$\\
$2$&$0$&$1$\\
$2$&$1$&$2$\\
\end{tabular}
\end{center}
Ceci définit un DFA dont les trois états correspondent aux trois
valeurs possibles de $n$ modulo $3$, la transition $n\to n'$ étiquetée
par $x$ correspond au passage de $n$ à $2n+x$ modulo $3$,
c'est-à-dire :
\begin{center}
%%% begin example6 %%%

\begin{tikzpicture}[>=latex,line join=bevel,automaton]
%%
\node (q1) at (97bp,20.28bp) [draw,circle,state] {$1$};
  \node (q0) at (18bp,20.28bp) [draw,circle,state,initial,final,accepting below] {$0$};
  \node (q2) at (176bp,20.28bp) [draw,circle,state] {$2$};
  \draw [->] (q1) ..controls (74.757bp,3.6593bp) and (64.084bp,-1.2803bp)  .. (54bp,1.2803bp) .. controls (50.042bp,2.2853bp) and (46.047bp,3.838bp)  .. node[auto] {$1$} (q0);
  \draw [->] (q2) to[loop above] node[auto] {$1$} (q2);
  \draw [->] (q2) ..controls (153.76bp,3.6593bp) and (143.08bp,-1.2803bp)  .. (133bp,1.2803bp) .. controls (129.04bp,2.2853bp) and (125.05bp,3.838bp)  .. node[auto] {$0$} (q1);
  \draw [->] (q0) to[loop above] node[auto] {$0$} (q0);
  \draw [->] (q0) ..controls (45.659bp,20.28bp) and (57.817bp,20.28bp)  .. node[auto] {$1$} (q1);
  \draw [->] (q1) ..controls (124.66bp,20.28bp) and (136.82bp,20.28bp)  .. node[auto] {$0$} (q2);
%
\end{tikzpicture}

%%% end example6 %%%
\end{center}
(On a marqué l'état $0$ comme initial car le mot vide a une valeur
numérique congrue à $0$ modulo $3$, et seul $0$ comme final car on
veut reconnaître les multiples de $3$.)

(d) Le langage $L_d$ n'est pas rationnel (on pourrait le démontrer à
l'aide du lemme de pompage, mais ce n'est pas très facile).

(e) Le langage $L_e$ est rationnel car il s'agit du langage dénoté par
l'expression rationnelle $0{*}10{*}$.
\end{corrige}

%

\exercice

Soit $\Sigma = \{a\}$.  Montrer que le langage $L = \{a^2, a^3, a^5,
a^7, a^{11}, a^{13}\ldots\}$ constitué des mots ayant un nombre
\emph{premier} de $a$, n'est pas rationnel.

\begin{corrige}
Supposons par l'absurde que $L$ soit rationnel.  D'après le lemme de
pompage, il existe un certain $k$ tel que tout mot de $L$ de longueur
$\geq k$ se factorise sous la forme $uvw$ avec (i) $|v|\geq 1$,
(ii) $|uv|\leq k$ et (iii) $uv^iw \in L$ pour tout $i\geq 0$.  Soit
$p$ un nombre premier supérieur ou égal à $k$ (qui existe car
l'ensemble des nombres premiers est infini) : le mot $a^p \in L$
admet une factorisation comme on vient de dire.  Posons $|u| =: m$ et
$|v| =: n$, si bien que $|w| = p-m-n$.  On a alors $n\geq 1$
d'après (i), et $|uv^iw| = m + in + (p-m-n) = p+(i-1)n$ est premier
pour tout $i\geq 0$ d'après (iii).  En particulier pour $i=p+1$ on
voit que $p + pn = p(n+1)$ est premier, ce qui contredit le fait qu'il
s'agit d'un multiple non-trivial ($n+1\geq 2$) de $p$.
\end{corrige}

%

\exercice

On considère l'automate suivant :
\begin{center}
%%% begin ex1p2a %%%

\begin{tikzpicture}[>=latex,line join=bevel,automaton]
%%
\node (q1) at (98bp,72bp) [draw,circle,state] {$1$};
  \node (q0) at (18bp,47bp) [draw,circle,state,initial] {$0$};
  \node (q3) at (258bp,72bp) [draw,circle,state] {$3$};
  \node (q2) at (178bp,72bp) [draw,circle,state] {$2$};
  \node (q5) at (178bp,18bp) [draw,circle,state] {$5$};
  \node (q4) at (98bp,18bp) [draw,circle,state] {$4$};
  \node (q7) at (338bp,47bp) [draw,circle,state,final] {$7$};
  \node (q6) at (258bp,18bp) [draw,circle,state] {$6$};
  \draw [->] (q2) ..controls (206.11bp,72bp) and (218.58bp,72bp)  .. node[auto] {$a$} (q3);
  \draw [->] (q3) ..controls (285.85bp,63.395bp) and (299.33bp,59.074bp)  .. node[auto] {$\varepsilon$} (q7);
  \draw [->] (q7) to[loop below] node[auto] {$a,b$} (q7);
  \draw [->] (q6) ..controls (285.62bp,27.897bp) and (299.46bp,33.043bp)  .. node[auto] {$\varepsilon$} (q7);
  \draw [->] (q4) ..controls (126.11bp,18bp) and (138.58bp,18bp)  .. node[auto] {$b$} (q5);
  \draw [->] (q0) to[loop below] node[auto] {$a,b$} (q0);
  \draw [->] (q0) ..controls (45.62bp,37.103bp) and (59.462bp,31.957bp)  .. node[auto] {$\varepsilon$} (q4);
  \draw [->] (q5) ..controls (206.11bp,18bp) and (218.58bp,18bp)  .. node[auto] {$b$} (q6);
  \draw [->] (q0) ..controls (45.849bp,55.605bp) and (59.331bp,59.926bp)  .. node[auto] {$\varepsilon$} (q1);
  \draw [->] (q1) ..controls (126.11bp,72bp) and (138.58bp,72bp)  .. node[auto] {$a$} (q2);
%
\end{tikzpicture}

%%% end ex1p2a %%%
\end{center}

(0) Décrire brièvement le langage accepté par l'automate en question.

(1) Cet automate est-il déterministe ?  Si non, le déterminiser.

(2) Minimiser l'automate déterminisé (on doit trouver un DFA ayant
quatre états).  Décrire brièvement la signification de ces quatre
états, de façon à vérifier qu'il accepte le même langage que décrit
en (0).

(3) Éliminer les états de l'automate d'origine de façon à obtenir une
expression rationnelle dénotant le langage reconnu par le langage
décrit en (0).

\begin{corrige}
(0) L'automate proposé accepte les mots ayant soit deux $a$
  consécutifs (en passant par le chemin $0\to 1\to 2\to 3\to 7$) soit
  deux $b$ consécutifs (en passant par le chemin $0\to 4\to 5\to 6\to
  7$).

(1) L'automate ayant des ε-transitions, il ne peut pas être
  déterministe : on a affaire à un εNFA.  Avant de le déterminiser, on
  élimine ses ε-transitions : la ε-fermeture de $0$ est $\{0,1,4\}$,
  celle de $3$ est $\{3,7\}$ et celle de $6$ est $\{6,7\}$ ; on est
  amené au NFA suivant :
\begin{center}
%%% begin ex1p2b %%%

\begin{tikzpicture}[>=latex,line join=bevel,automaton]
%%
\node (q0) at (18bp,47bp) [draw,circle,state,initial] {$0$};
  \node (q3) at (178bp,72bp) [draw,circle,state,final,accepting above] {$3$};
  \node (q2) at (98bp,72bp) [draw,circle,state] {$2$};
  \node (q5) at (98bp,18bp) [draw,circle,state] {$5$};
  \node (q7) at (268bp,47bp) [draw,circle,state,final] {$7$};
  \node (q6) at (178bp,18bp) [draw,circle,state,final,accepting below] {$6$};
  \draw [->] (q2) ..controls (126.11bp,72bp) and (138.58bp,72bp)  .. node[auto] {$a$} (q3);
  \draw [->] (q3) ..controls (208.21bp,63.701bp) and (225.92bp,58.67bp)  .. node[auto] {$a,b$} (q7);
  \draw [->] (q7) to[loop below] node[auto] {$a,b$} (q7);
  \draw [->] (q6) ..controls (208.34bp,27.667bp) and (226.26bp,33.576bp)  .. node[auto] {$a,b$} (q7);
  \draw [->] (q5) ..controls (126.11bp,18bp) and (138.58bp,18bp)  .. node[auto] {$b$} (q6);
  \draw [->] (q0) ..controls (45.62bp,37.103bp) and (59.462bp,31.957bp)  .. node[auto] {$b$} (q5);
  \draw [->] (q0) ..controls (45.849bp,55.605bp) and (59.331bp,59.926bp)  .. node[auto] {$a$} (q2);
  \draw [->] (q0) to[loop below] node[auto] {$a,b$} (q0);
%
\end{tikzpicture}

%%% end ex1p2b %%%
\end{center}
(Les états $1$ et $4$ étant inaccessibles, ils ont été retirés.  Il ne
faut pas oublier que $3$ et $6$ sont finaux puisqu'ils ont l'état
final $7$ dans leur ε-fermeture.)

On peut maintenant procéder à la déterminisation.  Pour abréger les
noms des états, on note, par exemple, $023$ pour $\{0,2,3\}$.  En
construisant de proche en proche, on obtient le DFA suivant :
\begin{center}
\scalebox{0.8}{%PLEASE!  There HAS to be a better way to do this!
%%% begin ex1p2c %%%

\begin{tikzpicture}[>=latex,line join=bevel,automaton]
%%
\begin{scope}
  \pgfsetstrokecolor{black}
  \definecolor{strokecol}{rgb}{1.0,1.0,1.0};
  \pgfsetstrokecolor{strokecol}
  \definecolor{fillcol}{rgb}{1.0,1.0,1.0};
  \pgfsetfillcolor{fillcol}
\end{scope}
\begin{scope}
  \pgfsetstrokecolor{black}
  \definecolor{strokecol}{rgb}{1.0,1.0,1.0};
  \pgfsetstrokecolor{strokecol}
  \definecolor{fillcol}{rgb}{1.0,1.0,1.0};
  \pgfsetfillcolor{fillcol}
\end{scope}
\begin{scope}
  \pgfsetstrokecolor{black}
  \definecolor{strokecol}{rgb}{1.0,1.0,1.0};
  \pgfsetstrokecolor{strokecol}
  \definecolor{fillcol}{rgb}{1.0,1.0,1.0};
  \pgfsetfillcolor{fillcol}
\end{scope}
\begin{scope}
  \pgfsetstrokecolor{black}
  \definecolor{strokecol}{rgb}{1.0,1.0,1.0};
  \pgfsetstrokecolor{strokecol}
  \definecolor{fillcol}{rgb}{1.0,1.0,1.0};
  \pgfsetfillcolor{fillcol}
\end{scope}
  \node (q0) at (18bp,121.73bp) [draw,circle,state,initial] {$0$};
  \node (q027) at (382bp,39.732bp) [draw,circle,state,final] {$027$};
  \node (q056) at (188bp,58.732bp) [draw,circle,state,final,accepting below] {$056$};
  \node (q023) at (188bp,166.73bp) [draw,circle,state,final,accepting above] {$023$};
  \node (q057) at (382bp,185.73bp) [draw,circle,state,final] {$057$};
  \node (q0567) at (285bp,58.732bp) [draw,circle,state,final,accepting above] {$0567$};
  \node (q02) at (100bp,160.73bp) [draw,circle,state] {$02$};
  \node (q0237) at (285bp,166.73bp) [draw,circle,state,final,accepting below] {$0237$};
  \node (q05) at (100bp,68.732bp) [draw,circle,state] {$05$};
  \draw [->] (q057) ..controls (367.95bp,143.33bp) and (356.68bp,115.58bp)  .. (340bp,95.732bp) .. controls (334.17bp,88.79bp) and (326.73bp,82.582bp)  .. node[auto] {$b$} (q0567);
  \draw [->] (q0) ..controls (45.352bp,134.58bp) and (60.273bp,141.85bp)  .. node[auto] {$a$} (q02);
  \draw [->] (q0237) to[loop above] node[auto] {$a$} (q0237);
  \draw [->] (q023) ..controls (213.27bp,203.89bp) and (232.48bp,226.69bp)  .. (256bp,236.73bp) .. controls (279.71bp,246.86bp) and (289.57bp,244.97bp)  .. (314bp,236.73bp) .. controls (330.09bp,231.3bp) and (345.38bp,220.3bp)  .. node[auto] {$b$} (q057);
  \draw [->] (q023) ..controls (222.58bp,166.73bp) and (234.64bp,166.73bp)  .. node[auto] {$a$} (q0237);
  \draw [->] (q02) ..controls (129.57bp,162.73bp) and (142.05bp,163.6bp)  .. node[auto] {$a$} (q023);
  \draw [->] (q0567) to[loop below] node[auto] {$b$} (q0567);
  \draw [->] (q05) ..controls (129.65bp,65.401bp) and (142.24bp,63.937bp)  .. node[auto] {$b$} (q056);
  \draw [->] (q056) ..controls (222.58bp,58.732bp) and (234.64bp,58.732bp)  .. node[auto] {$b$} (q0567);
  \draw [->] (q0237) ..controls (322.23bp,165.23bp) and (331.57bp,165.88bp)  .. (340bp,167.73bp) .. controls (343.63bp,168.53bp) and (347.33bp,169.66bp)  .. node[auto] {$b$} (q057);
  \draw [->] (q056) ..controls (217.15bp,28.639bp) and (235.83bp,12.746bp)  .. (256bp,5.7322bp) .. controls (280.35bp,-2.7337bp) and (288.93bp,-0.26493bp)  .. (314bp,5.7322bp) .. controls (327.41bp,8.9409bp) and (341.18bp,15.299bp)  .. node[auto,below] {$a$} (q027);
  \draw [->] (q05) ..controls (100bp,100.68bp) and (100bp,117.05bp)  .. node[auto] {$a$} (q02);
  \draw [->] (q057) ..controls (397.04bp,148.37bp) and (402.36bp,127.76bp)  .. (400bp,109.21bp) .. controls (398.44bp,96.972bp) and (395.4bp,83.809bp)  .. node[auto] {$a$} (q027);
  \draw [->] (q027) ..controls (362.81bp,75.635bp) and (351.82bp,95.041bp)  .. (340bp,110.73bp) .. controls (332.19bp,121.1bp) and (322.64bp,131.57bp)  .. node[auto] {$a$} (q0237);
  \draw [->] (q0567) ..controls (324.14bp,51.105bp) and (336.81bp,48.57bp)  .. node[auto,below] {$a$} (q027);
  \draw [->] (q027) ..controls (382bp,87.723bp) and (382bp,124.52bp)  .. node[auto] {$b$} (q057);
  \draw [->] (q02) ..controls (115.85bp,134.94bp) and (120.35bp,122.5bp)  .. (118bp,110.98bp) .. controls (116.93bp,105.75bp) and (115.19bp,100.36bp)  .. node[auto] {$b$} (q05);
  \draw [->] (q0) ..controls (45.218bp,104.36bp) and (61.545bp,93.546bp)  .. node[auto] {$b$} (q05);
%
\end{tikzpicture}

%%% end ex1p2c %%%
}
\end{center}
(À titre d'exemple, la transition étiquetée $b$ partant de l'état
$0237$ conduit à l'état $057$ car les transitions étiquetées $b$ dans
le NFA précédent et partant des états parmi $\{0,2,3,7\}$ sont $0\to
0$, $0\to 5$, $3\to 7$ et $7\to 7$.  Tous les états contenant l'un des
symboles $3,6,7$ sont finaux.)

(2) On a affaire à un DFA (sous-entendu : \emph{complet}) dont tous
les états sont accessibles, on peut donc appliquer directement
l'algorithme de Moore.  Une première partition sépare les états
finaux, soit $023, 0237, 027, 056, 0567, 057$ des non-finaux, soit $0,
02, 05$.  L'étape suivante distingue $02$ parce que sa $a$-transition
conduit à une classe différente que celles de $0$ et $05$, et $05$
parce que sa $b$-transition conduit à une classe différente de $0$ et
$02$.  Les étapes suivantes ne changent rien.  Finalement, on arrive à
un automate à quatre états :
\begin{center}
%%% begin ex1p2d %%%

\begin{tikzpicture}[>=latex,line join=bevel,automaton]
%%
\begin{scope}
  \pgfsetstrokecolor{black}
  \definecolor{strokecol}{rgb}{1.0,1.0,1.0};
  \pgfsetstrokecolor{strokecol}
  \definecolor{fillcol}{rgb}{1.0,1.0,1.0};
  \pgfsetfillcolor{fillcol}
  \filldraw (0bp,0bp) -- (0bp,112bp) -- (200bp,112bp) -- (200bp,0bp) -- cycle;
\end{scope}
  \node (q02) at (100bp,93bp) [draw,circle,state] {$02$};
  \node (q0) at (18bp,56bp) [draw,circle,state,initial] {$0$};
  \node (qA) at (182bp,60bp) [draw,circle,state,final] {$A$};
  \node (q05) at (100bp,19bp) [draw,circle,state] {$05$};
  \draw [->] (q0) ..controls (45.691bp,68.345bp) and (60.407bp,75.151bp)  .. node[auto] {$a$} (q02);
  \draw [->] (qA) to[loop below] node[auto] {$a,b$} (qA);
  \draw [->] (q02) ..controls (129.2bp,81.367bp) and (143.35bp,75.529bp)  .. node[auto] {$a$} (qA);
  \draw [->] (q05) ..controls (128.81bp,33.252bp) and (143.85bp,40.959bp)  .. node[auto,below] {$b$} (qA);
  \draw [->] (q05) ..controls (100bp,46.09bp) and (100bp,55.041bp)  .. node[auto] {$a$} (q02);
  \draw [->] (q02) ..controls (116.71bp,70.13bp) and (120.2bp,60.948bp)  .. (118bp,52.251bp) .. controls (117.33bp,49.598bp) and (116.4bp,46.928bp)  .. node[auto] {$b$} (q05);
  \draw [->] (q0) ..controls (45.691bp,43.655bp) and (60.407bp,36.849bp)  .. node[auto] {$b$} (q05);
%
\end{tikzpicture}

%%% end ex1p2d %%%
\end{center}
où $A$ représente la classe de tous les états finaux de l'automate
précédent.

La signification des quatre états est la suivante : l'état $0$
signifie que l'automate n'a encore rien lu, l'état $02$ signifie que
l'automate vient de lire un $a$, le $05$ signifie qu'il vient de lire
un $b$, le $A$ signifie qu'il a lu deux $a$ consécutifs ou bien deux
$b$ consécutifs.  Sur cette description, il est clair que l'automate
accepte les mots contenant deux $a$ consécutifs ou bien deux $b$
consécutifs.

(3) L'élimination des états $1$ à $6$ peut se faire dans un ordre
quelconque et conduit à l'automate (à transitions étiquetées par des
expressions rationnelles) suivant :
\begin{center}
%%% begin ex1p2a2 %%%

\begin{tikzpicture}[>=latex,line join=bevel,automaton]
%%
\node (q0) at (18bp,20.549bp) [draw,circle,state,initial] {$0$};
  \node (q7) at (104bp,20.549bp) [draw,circle,state,final] {$7$};
  \draw [->] (q0) ..controls (47.743bp,20.549bp) and (62.773bp,20.549bp)  .. node[auto] {$aa$} (q7);
  \draw [->] (q0) ..controls (42.511bp,4.2138bp) and (55.796bp,-1.5495bp)  .. (68bp,1.5495bp) .. controls (71.958bp,2.5545bp) and (75.953bp,4.1071bp)  .. node[auto,below] {$bb$} (q7);
  \draw [->] (q0) to[loop below] node[auto] {$a,b$} (q0);
  \draw [->] (q7) to[loop below] node[auto] {$a,b$} (q7);
%
\end{tikzpicture}

%%% end ex1p2a2 %%%
\end{center}

Il y a plusieurs flèches entre les mêmes états : quitte à les
remplacer par des disjonctions, on obtient :
\begin{center}
%%% begin ex1p2a3 %%%

\begin{tikzpicture}[>=latex,line join=bevel,automaton]
%%
\node (q0) at (18bp,18bp) [draw,circle,state,initial] {$0$};
  \node (q7) at (119bp,18bp) [draw,circle,state,final] {$7$};
  \draw [->] (q0) ..controls (51.292bp,18bp) and (73.384bp,18bp)  .. node[auto] {$aa|bb$} (q7);
  \draw [->] (q0) to[loop below] node[auto] {$a|b$} (q0);
  \draw [->] (q7) to[loop below] node[auto] {$a|b$} (q7);
%
\end{tikzpicture}

%%% end ex1p2a3 %%%
\end{center}

Enfin, on doit éliminer les états $0$ et $7$ eux-mêmes : pour cela, on
ajoute un nouvel état initial qui ne soit la cible d'aucune flèche et
un nouvel état final d'où ne part aucune flèche,
\begin{center}
%%% begin ex1p2a4 %%%

\begin{tikzpicture}[>=latex,line join=bevel,automaton]
%%
\node (qi) at (18bp,18bp) [draw,circle,state,initial] {$\phantom{0}$};
  \node (q0) at (97bp,18bp) [draw,circle,state] {$0$};
  \node (q7) at (198bp,18bp) [draw,circle,state] {$7$};
  \node (qf) at (277bp,18bp) [draw,circle,state,final] {$\phantom{7}$};
  \draw [->] (qi) ..controls (45.659bp,18bp) and (57.817bp,18bp)  .. node[auto] {$\varepsilon$} (q0);
  \draw [->] (q7) ..controls (225.66bp,18bp) and (237.82bp,18bp)  .. node[auto] {$\varepsilon$} (qf);
  \draw [->] (q0) ..controls (130.29bp,18bp) and (152.38bp,18bp)  .. node[auto] {$aa|bb$} (q7);
  \draw [->] (q0) to[loop below] node[auto] {$a|b$} (q0);
  \draw [->] (q7) to[loop below] node[auto] {$a|b$} (q7);
%
\end{tikzpicture}

%%% end ex1p2a4 %%%
\end{center}

Et l'élimination des états $0$ et $7$ dans un ordre quelconque conduit
finalement à l'expression rationnelle $(a|b){*}(aa|bb)(a|b){*}$.
\end{corrige}

%

\exercice

Soit $\Sigma = \{a,b\}$.

(1) Tracer l'automate de Thompson de l'expression rationnelle
$a{*}(baa{*}){*}$.  Cet automate est-il déterministe ?

(2) En éliminer les transitions spontanées.

(3) Déterminiser l'automate obtenu (on demande un automate complet).

(4) Minimiser l'automate obtenu (on demande un automate complet).

(5) Vérifier le résultat en décrivant en français le langage dénoté
par l'expression rationnelle initiale et reconnu par l'automate
finalement obtenu.

\begin{corrige}
(1) L'automate de Thompson de $a{*}(baa{*}){*}$ doit comporter $14$
  états puisque cette expression rationnelle contient $7$ symboles
  parenthèses non comptées.  Il est le suivant (où on a omis les
  $\varepsilon$ sur les transitions spontanées) :
\begin{center}
\scalebox{0.34}{%
%%% begin ex3p1 %%%

\begin{tikzpicture}[>=latex,line join=bevel,automaton]
%%
\node (q1) at (97bp,45bp) [draw,circle,state] {$1$};
  \node (q0) at (18bp,18bp) [draw,circle,state,initial] {$0$};
  \node (q3) at (255bp,18bp) [draw,circle,state] {$3$};
  \node (q2) at (176bp,45bp) [draw,circle,state] {$2$};
  \node (q5) at (413bp,48bp) [draw,circle,state] {$5$};
  \node (q4) at (334bp,18bp) [draw,circle,state] {$4$};
  \node (q7) at (571bp,86bp) [draw,circle,state] {$7$};
  \node (q6) at (492bp,86bp) [draw,circle,state] {$6$};
  \node (q9) at (729bp,86bp) [draw,circle,state] {$9$};
  \node (q8) at (650bp,86bp) [draw,circle,state] {$8$};
  \node (q11) at (891.49bp,125bp) [draw,circle,state] {$11$};
  \node (q10) at (809.5bp,125bp) [draw,circle,state] {$10$};
  \node (q13) at (1055.5bp,25bp) [draw,circle,state,final] {$13$};
  \node (q12) at (973.49bp,63bp) [draw,circle,state] {$12$};
  \draw [->] (q0) ..controls (59.7bp,17.371bp) and (103.03bp,16.838bp)  .. (140bp,17bp) .. controls (169.56bp,17.13bp) and (203.43bp,17.448bp)  .. node[auto] {{}} (q3);
  \draw [->] (q12) ..controls (939.5bp,48.432bp) and (914.89bp,40bp)  .. (892.49bp,40bp) .. controls (491bp,40bp) and (491bp,40bp)  .. (491bp,40bp) .. controls (474.34bp,40bp) and (455.77bp,41.875bp)  .. node[auto] {{}} (q5);
  \draw [->] (q2) ..controls (203.44bp,35.73bp) and (216.64bp,31.102bp)  .. node[auto] {{}} (q3);
  \draw [->] (q10) ..controls (835.62bp,135.26bp) and (845.21bp,137.3bp)  .. (854bp,136bp) .. controls (856.94bp,135.57bp) and (859.98bp,134.94bp)  .. node[auto] {$a$} (q11);
  \draw [->] (q7) ..controls (598.66bp,86bp) and (610.82bp,86bp)  .. node[auto] {$a$} (q8);
  \draw [->] (q9) ..controls (788.12bp,80.488bp) and (892.85bp,70.553bp)  .. node[auto] {{}} (q12);
  \draw [->] (q8) ..controls (677.66bp,86bp) and (689.82bp,86bp)  .. node[auto] {{}} (q9);
  \draw [->] (q11) ..controls (915.87bp,107.43bp) and (926.56bp,99.325bp)  .. (935.99bp,92bp) .. controls (940.47bp,88.526bp) and (945.22bp,84.783bp)  .. node[auto] {{}} (q12);
  \draw [->] (q2) ..controls (152.62bp,39.018bp) and (146.07bp,37.685bp)  .. (140bp,37bp) .. controls (134.92bp,36.427bp) and (129.53bp,36.741bp)  .. node[auto] {{}} (q1);
  \draw [->] (q6) ..controls (519.66bp,86bp) and (531.82bp,86bp)  .. node[auto] {{}} (q7);
  \draw [->] (q12) ..controls (1002.2bp,49.853bp) and (1016.2bp,43.176bp)  .. node[auto] {{}} (q13);
  \draw [->] (q9) ..controls (756.02bp,98.925bp) and (770.16bp,105.95bp)  .. node[auto] {{}} (q10);
  \draw [->] (q3) ..controls (282.66bp,18bp) and (294.82bp,18bp)  .. node[auto] {{}} (q4);
  \draw [->] (q11) ..controls (866.59bp,118.91bp) and (860.06bp,117.66bp)  .. (854bp,117bp) .. controls (848.92bp,116.45bp) and (843.55bp,116.73bp)  .. node[auto] {{}} (q10);
  \draw [->] (q5) ..controls (440.12bp,60.892bp) and (454.27bp,67.873bp)  .. node[auto] {$b$} (q6);
  \draw [->] (q4) ..controls (366.77bp,7.9414bp) and (390.7bp,2bp)  .. (412bp,2bp) .. controls (412bp,2bp) and (412bp,2bp)  .. (974.49bp,2bp) .. controls (992.87bp,2bp) and (1012.7bp,7.6737bp)  .. node[auto] {{}} (q13);
  \draw [->] (q0) ..controls (45.436bp,27.27bp) and (58.635bp,31.898bp)  .. node[auto] {{}} (q1);
  \draw [->] (q1) ..controls (121.23bp,55.953bp) and (131.02bp,58.496bp)  .. (140bp,57bp) .. controls (143.13bp,56.478bp) and (146.36bp,55.711bp)  .. node[auto] {$a$} (q2);
  \draw [->] (q4) ..controls (361.28bp,28.238bp) and (374.94bp,33.562bp)  .. node[auto] {{}} (q5);
%
\end{tikzpicture}

%%% end ex3p1 %%%
}
\end{center}

Cet automate n'est pas déterministe : un automate comportant des
ε-transitions est forcément non-déterministe.

\smallbreak

(2) Tous les états autres que $0$ (car il est initial) et $2,6,8,11$
(car des transitions non spontanées y aboutissent) vont disparaître ;
les ε-fermetures $C(q)$ de ces états sont les suivantes :

\begin{center}
\begin{tabular}{r|l}
$q$&ε-fermeture $C(q)$\\
\hline
$0$&$\{0,1,3,4,5,13\}$\\
$2$&$\{1,2,3,4,5,13\}$\\
$6$&$\{6,7\}$\\
$8$&$\{5,8,9,10,12,13\}$\\
$11$&$\{5,10,11,12,13\}$\\
\end{tabular}
\end{center}

En remplaçant chaque transition $q^\sharp \to q'$ étiquetée
d'un $x\in\Sigma$ dans l'automate par une transition $q\to q'$ pour
chaque état $q$ tel que $q^\sharp \in C(q)$, on obtient le NFA
suivant :

\begin{center}
%%% begin ex3p1b %%%

\begin{tikzpicture}[>=latex,line join=bevel,automaton]
%%
\node (q0) at (18bp,31.498bp) [draw,circle,state,initial,final,accepting below] {$0$};
  \node (q2) at (97bp,61.498bp) [draw,circle,state,final] {$2$};
  \node (q11) at (335.5bp,19.498bp) [draw,circle,state,final,accepting below] {$11$};
  \node (q8) at (255bp,49.498bp) [draw,circle,state,final,accepting above] {$8$};
  \node (q6) at (176bp,31.498bp) [draw,circle,state] {$6$};
  \draw [->] (q0) ..controls (45.279bp,41.737bp) and (58.943bp,47.06bp)  .. node[auto] {$a$} (q2);
  \draw [->] (q2) to[loop above] node[auto] {$a$} (q2);
  \draw [->] (q8) ..controls (282.44bp,39.394bp) and (295.8bp,34.29bp)  .. node[auto] {$a$} (q11);
  \draw [->] (q11) to[loop above] node[auto] {$a$} (q11);
  \draw [->] (q11) ..controls (297.15bp,8.7001bp) and (264.63bp,2.1768bp)  .. (237bp,7.4983bp) .. controls (224.85bp,9.8374bp) and (212.04bp,14.613bp)  .. node[auto,near start] {$b$} (q6);
  \draw [->] (q0) ..controls (47.643bp,24.415bp) and (64.2bp,20.964bp)  .. (79bp,19.498bp) .. controls (94.922bp,17.922bp) and (99.078bp,17.922bp)  .. (115bp,19.498bp) .. controls (125.98bp,20.586bp) and (137.94bp,22.768bp)  .. node[auto,below] {$b$} (q6);
  \draw [->] (q2) ..controls (124.28bp,51.26bp) and (137.94bp,45.936bp)  .. node[auto] {$b$} (q6);
  \draw [->] (q8) ..controls (234.03bp,34.749bp) and (226.51bp,30.579bp)  .. (219bp,28.498bp) .. controls (214.15bp,27.154bp) and (208.87bp,26.751bp)  .. node[auto,near start] {$b$} (q6);
  \draw [->] (q6) ..controls (198.21bp,42.911bp) and (205.25bp,45.859bp)  .. (212bp,47.498bp) .. controls (216.59bp,48.613bp) and (221.55bp,49.292bp)  .. node[auto] {$a$} (q8);
%
\end{tikzpicture}

%%% end ex3p1b %%%
\end{center}

Les états $0,2,8,11$ sont finaux car ce sont eux qui ont $13$ dans
leur ε-fermeture.

\smallbreak

(3) L'automate ainsi obtenu est déjà déterministe incomplet ; pour le
déterminiser-compléter, il n'y a qu'à ajouter un puits $\bot$ avec la
seule transition qui manque, c'est-à-dire une transition étiquetée par
$b$ depuis l'état $6$ (et des transitions de $\bot$ vers lui-même
étiquetées $a$ et $b$).  Nous ne représentons pas l'automate à $6$
états ainsi fabriqué.

\smallbreak

(4) On part de l'algorithme déterministe complet obtenu à la
question (3), et on lui applique l'algorithme de minimisation.  On
sépare d'abord ses états en deux classes, les finaux $\{0,2,8,11\}$ et
les non-finaux $\{6,\bot\}$.  La transition étiquetée $a$ sépare les
états $6$ et $\bot$ car le premier aboutit dans la classe
$\{0,2,8,11\}$ tandis que le second aboutit dans la classe
$\{6,\bot\}$.  On vérifie ensuite qu'aucune transition ne sépare des
états.  L'automate minimal est donc le suivant :
\begin{center}
%%% begin ex3p1c %%%

\begin{tikzpicture}[>=latex,line join=bevel,automaton]
%%
\node (q6) at (97bp,20.28bp) [draw,circle,state] {$6$};
  \node (qbot) at (176bp,20.28bp) [draw,circle,state] {$\bot$};
  \node (qF) at (18bp,20.28bp) [draw,circle,state,initial,final,accepting below] {$F$};
  \draw [->] (q6) ..controls (74.757bp,3.6593bp) and (64.084bp,-1.2803bp)  .. (54bp,1.2803bp) .. controls (50.042bp,2.2853bp) and (46.047bp,3.838bp)  .. node[auto] {$a$} (qF);
  \draw [->] (q6) ..controls (124.66bp,20.28bp) and (136.82bp,20.28bp)  .. node[auto] {$b$} (qbot);
  \draw [->] (qbot) to[loop above] node[auto] {$a,b$} (qbot);
  \draw [->] (qF) to[loop above] node[auto] {$a$} (qF);
  \draw [->] (qF) ..controls (45.659bp,20.28bp) and (57.817bp,20.28bp)  .. node[auto] {$b$} (q6);
%
\end{tikzpicture}

%%% end ex3p1c %%%
\end{center}
où l'état $F$ représente la classe $0\equiv 2\equiv 8\equiv 11$.

\smallbreak

(5) Il s'agit du langage constitué des mots n'ayant jamais deux $b$
consécutifs ni de $b$ final,
c'est-à-dire des mots dans lesquels chaque $b$ est suivi
d'au moins un $a$ : l'expression rationnelle initiale le présente
comme le langage constitués des mots formés d'un nombre quelconque de
$a$ puis d'un nombre quelconque de répétitions d'un $b$ suivi d'au
moins un $a$.  L'automate final interdit les suites de deux $b$
consécutifs comme ceci : l'état $F$ correspond à la situation où on ne
vient pas de rencontrer un $b$ (=la lettre précédente était un $a$ ou
bien on vient de commencer le mot) et on n'en a jamais rencontré deux,
l'état $6$ à la situation où on vient de rencontrer un $b$ et on n'en
a jamais rencontré deux, et l'état $\bot$ à la situation où on a
rencontré au moins une fois deux $b$ consécutifs.  Avec cette
description, il est clair que l'automate fait ce qui était demandé.
\end{corrige}

%

\exercice

Donner plusieurs (au moins trois) expressions rationnelles
équivalentes dénotant le langage reconnu par l'automate suivant sur
l'alphabet $\Sigma = \{a,b\}$ :

\begin{center}
%%% begin ex3p2 %%%

\begin{tikzpicture}[>=latex,line join=bevel,automaton]
%%
\node (q1) at (97bp,20.28bp) [draw,circle,state] {$1$};
  \node (q0) at (18bp,20.28bp) [draw,circle,state,initial,final,initial above,accepting below] {$0$};
  \node (q2) at (176bp,20.28bp) [draw,circle,state] {$2$};
  \draw [->] (q1) ..controls (74.757bp,3.6593bp) and (64.084bp,-1.2803bp)  .. (54bp,1.2803bp) .. controls (50.042bp,2.2853bp) and (46.047bp,3.838bp)  .. node[auto] {$b$} (q0);
  \draw [->] (q2) to[loop right] node[auto] {$b$} (q2);
  \draw [->] (q2) ..controls (153.76bp,3.6593bp) and (143.08bp,-1.2803bp)  .. (133bp,1.2803bp) .. controls (129.04bp,2.2853bp) and (125.05bp,3.838bp)  .. node[auto] {$a$} (q1);
  \draw [->] (q0) to[loop left] node[auto] {$a$} (q0);
  \draw [->] (q0) ..controls (45.659bp,20.28bp) and (57.817bp,20.28bp)  .. node[auto] {$b$} (q1);
  \draw [->] (q1) ..controls (124.66bp,20.28bp) and (136.82bp,20.28bp)  .. node[auto] {$a$} (q2);
%
\end{tikzpicture}


%%% end ex3p2 %%%
\end{center}
(On pourra considérer les ordres suivants d'élimination des états :
(A) $2,1,0$, ensuite (B) $1,2,0$ et enfin (C) $0,2,1$.)

\begin{corrige}
(A) Si on commence par éliminer l'état $2$ (en considérant l'automate
comme un automate à transitions étiquetées par des expressions
rationnelles), l'état $1$ reçoit une transition vers lui-même
étiquetée $ab{*}a$.  Si on élimine l'état $1$, l'état $0$ reçoit à la
place une transition vers lui-même étiquetée par $b(ab{*}a){*}b$,
qu'on peut fusionner avec la transition vers lui-même déjà existante
étiquetée par $a$ pour une seule étiquetée par $a|b(ab{*}a){*}b$.
Quitte éventuellement à ajouter un nouvel état initial (conduisant à
$0$ par une transition spontanée) et un nouvel état final (vers lequel
$0$ conduit par une transition spontanée) et à éliminer n'état $0$, on
obtient finalement l'expression rationnelle
\[
(a|b(ab{*}a){*}b){*}
\]

\smallbreak

(B) Si on commence par éliminer l'état $1$, il apparaît une transition
$0\to 2$ étiquetée $ba$ et une $2\to 0$ étiquetée $ab$ (si on veut
appliquer l'algorithme de façon puremnet mécanique, l'état $1$ n'a pas
de transition vers lui-même, c'est-à-dire qu'on pourrait l'étiqueter
par $\bot$, symbole d'expression rationnelle qui dénote le lagange
vide, et l'expression rationnelle $\bot{*}$ est équivalente
à $\varepsilon$) ; mais il ne faut pas oublier que l'état $2$ reçoit
lui aussi une transition vers lui-même (en passant par $1$) étiquetée
$aa$, qu'on peut fusionner avec la transition vers lui-même déjà
existante étiquetée par $b$ pour obtenir une transition étiquetée
$b|aa$ ; de même, l'état $0$ reçoit une transition étiquetée $bb$,
qu'on peut fusionner avec celle existante pour obtenir $a|bb$.
L'élimination de l'état $2$ fait alors apparaître une transition de
$0$ vers lui-même étiquetée $ba(b|aa){*}ab$, qu'on peut fusionner avec
la transition vers lui-même déjà étiquetée par $a|bb$ pour une seule
étiquetée par $a|bb|ba(b|aa){*}ab$.  On obtient finalement
\[
(a|bb|ba(b|aa){*}ab){*}
\]
(en particulier, cette expression est équivalente à celle obtenue
précédemment).

\smallbreak

(C) Si on préfère commencer par éliminer l'état $0$, il faut au préalable
ajouter un nouvel état initial $I$ (conduisant à $0$ par une
transition spontanée) et un nouvel état final $F$ (vers lequel $0$
conduit par une transition spontanée).  L'élimination de l'état $0$
fait apparaître une transition de $I$ vers $1$ étiquetée $a{*}b$, une
transition $1$ vers $F$ étiquetée $ba{*}$, et une transition de l'état
$1$ vers lui-même étiquetée $ba{*}b$, et enfin une transition de $I$
vers $F$ étiquetée $a{*}$.  L'élimination de l'état $2$ fait
apparaître une transition de $1$ vers lui-même étiquetée $ab{*}a$,
qu'on peut fusionner avec celle déjà existante étiquetée $ba{*}b$ pour
obtenir une transition $(ab{*}a|ba{*}b)$.  Finalement, l'élimination
de l'état $1$ done l'expression rationnelle
\[
a{*}|a{*}b(ab{*}a|ba{*}b){*}ba{*}
\]
(toujours équivalente aux précédentes).
\end{corrige}

%

\exercice

On considère l'automate fini $M$ sur l'alphabet $\Sigma = \{a,b\}$
représenté par la figure suivante :
\begin{center}
\begin{tikzpicture}[>=latex,line join=bevel,automaton]
\node (X) at (-80bp,0bp) [draw,circle,state,initial] {$X$};
\node (Y) at (80bp,0bp) [draw,circle,state,final] {$Y$};
\node (A1) at (-40bp,50bp) [draw,circle,state] {$A\phantom{'}$};
\node (A2) at (40bp,50bp) [draw,circle,state] {$A'$};
\node (B1) at (-40bp,-50bp) [draw,circle,state] {$B\phantom{'}$};
\node (B2) at (40bp,-50bp) [draw,circle,state] {$B'$};
\draw[->] (X) -- node[auto]{$\varepsilon$} (A1);
\draw[->] (X) -- node[auto,below left]{$\varepsilon$} (B1);
\draw[->] (A1) -- node[auto]{$a$} (A2);
\draw[->] (B1) -- node[auto]{$b$} (B2);
\draw[->] (A1) to[loop above] node[auto]{$b$} (A1);
\draw[->] (A2) to[loop above] node[auto]{$b$} (A2);
\draw[->] (B1) to[loop below] node[auto]{$a$} (B1);
\draw[->] (B2) to[loop below] node[auto]{$a$} (B2);
\draw[->] (A2) -- node[auto]{$\varepsilon$} (Y);
\draw[->] (B2) -- node[auto,below right]{$\varepsilon$} (Y);
\end{tikzpicture}
\end{center}

(0) De quelle sorte d'automate s'agit-il ?  (Autrement dit : est-il
déterministe ou non ? avec transitions spontanées ou non ?)

(1a) Décrire brièvement, en français, le langage $L$ reconnu
(=accepté) par l'automate $M$, puis donner une expression
rationnelle qui le dénote.  (On pourra préférer traiter la question
(1b) d'abord.)

(1b) Pour chacun des mots suivants, dire s'ils sont dans $L$ ou non :
$\varepsilon$, $a$, $b$, $ab$, $aa$, $aab$, $aabb$, $abab$, $ababa$.
(Note : il est recommandé de réutiliser ces mots pour vérifier
rapidement les réponses aux questions suivantes et ainsi détecter
d'éventuelles erreurs lors des transformations des automates.)

(2) Éliminer les transitions spontanées de l'automate $M$.  (On
supprimera les états devenus inutiles.)  On appellera $M_2$ l'automate
obtenu.

(3) Déterminiser l'automate $M_2$ obtenu en (2), si nécessaire.  (On
demande un automate déterministe complet.)  On appellera $M_3$
l'automate déterminisé.

Pour simplifier le travail du correcteur, on demande de représenter
$M_3$ de sorte que les transitions étiquetées par $a$ soient, dans la
mesure du possible, horizontales de la gauche vers la droite, et
celles étiquetées par $b$, verticales du haut vers le bas.

(4) Minimiser l'automate $M_3$ obtenu en (3), si nécessaire
(justifier).

(5) Donner un automate (de n'importe quelle sorte) qui reconnaît le
langage $\overline{L} = \Sigma^*\setminus L$ complémentaire de $L$.

(6) Décrire brièvement, en français, ce langage
complémentaire $\overline{L}$.

(7) (Question bonus, plus longue, à ne traiter qu'en dernier.)
Calculer une expression rationnelle qui dénote ce langage
complémentaire $\overline{L}$.  (Ne pas hésiter à introduire des
notations intermédiaires.)

\begin{corrige}
(0) L'automate $M$ est un automate fini non-déterministe à transitions
  spontanées, ou ε-NFA (le concept d'« automate déterministe à
  transitions spontanées » n'aurait tout simplement pas de sens).

\smallbreak

(1a) Le chemin par les états $X,A,A',Y$ accepte les mots exactement
un $a$, c'est-à-dire le langage dénoté par $b{*}ab{*}$.  Le chemin par
les états $X,B,B',Y$ accepte les mots comportant exactement un $b$,
c'est-à-dire le langage dénoté par $a{*}ba{*}$.  L'automate $M$ dans
son ensemble accepte les mots comportant exactement un $a$ ou
(inclusif) exactement un $b$ (i.e. $L = \{w\in\Sigma^* : |w|_a = 1
\penalty0\ \textrm{ou}\penalty0\ |w|_b = 1\}$ si $|w|_x$ désigne le
nombre total d'occurrences de la lettre $x$ dans le mot $w$).  C'est
le langage dénoté par l'expression rationnelle $b{*}ab{*} | a{*}ba{*}$
(nous notons ici et ailleurs $|$ pour la disjonction, qu'on peut aussi
noter $+$).

(1b) Parmi les mots proposés, $a$, $b$, $ab$ et $aab$ appartiennent à
$L$, tandis que $\varepsilon$, $aa$, $aabb$, $abab$ et $ababa$ n'y
appartiennent pas.

\smallbreak

(2) La ε-fermeture (arrière) de l'état $X$ est $\{X,A,B\}$ ; la
ε-fermeture de l'état $A'$ est $\{A',Y\}$ et celle de l'état $B'$ est
$\{B',Y\}$ ; les autres états sont leur propre ε-fermeture (i.e.,
celle-ci est un singleton).  L'élimination des transitions spontanées
conduit donc à l'automate $M_2$ suivant :
\begin{center}
\begin{tikzpicture}[>=latex,line join=bevel,automaton]
\node (X) at (-80bp,0bp) [draw,circle,state,initial] {$X$};
\node (A1) at (-40bp,50bp) [draw,circle,state] {$A\phantom{'}$};
\node (A2) at (40bp,50bp) [draw,circle,state,final] {$A'$};
\node (B1) at (-40bp,-50bp) [draw,circle,state] {$B\phantom{'}$};
\node (B2) at (40bp,-50bp) [draw,circle,state,final] {$B'$};
\draw[->] (X) -- node[auto]{$b$} (A1);
\draw[->] (X) -- node[auto,below left]{$a$} (B1);
\draw[->] (X) -- node[auto,below right]{$a$} (A2);
\draw[->] (X) -- node[auto,above right]{$b$} (B2);
\draw[->] (A1) -- node[auto]{$a$} (A2);
\draw[->] (B1) -- node[auto]{$b$} (B2);
\draw[->] (A1) to[loop above] node[auto]{$b$} (A1);
\draw[->] (A2) to[loop above] node[auto]{$b$} (A2);
\draw[->] (B1) to[loop below] node[auto]{$a$} (B1);
\draw[->] (B2) to[loop below] node[auto]{$a$} (B2);
\end{tikzpicture}
\end{center}
(On a supprimé l'état $Y$ qui est devenu inutile car aucune transition
non-spontanée n'y conduit.)

\smallbreak

(3) L'algorithme de déterminisation conduit à l'automate $M_3$ suivant
où, pour plus de lisibilité, les états finaux ont été marqués en les
entourant deux fois plutôt que par une flèche sortante :
\begin{center}
\begin{tikzpicture}[>=latex,line join=bevel,automaton]
\node (s00) at (0bp,0bp) [draw,rounded corners,state,initial] {$\{X\}$};
\node (s01) at (75bp,0bp) [draw,rounded corners,state,accepting by double] {$\{A',B\}$};
\node (s02) at (150bp,0bp) [draw,rounded corners,state] {$\{B\}$};
\node (s10) at (0bp,-50bp) [draw,rounded corners,state,accepting by double] {$\{A,B'\}$};
\node (s11) at (75bp,-50bp) [draw,rounded corners,state,accepting by double] {$\{A',B'\}$};
\node (s12) at (150bp,-50bp) [draw,rounded corners,state,accepting by double] {$\{B'\}$};
\node (s20) at (0bp,-100bp) [draw,rounded corners,state] {$\{A\}$};
\node (s21) at (75bp,-100bp) [draw,rounded corners,state,accepting by double] {$\{A'\}$};
\node (s22) at (150bp,-100bp) [draw,rounded corners,state] {$\varnothing$};
\draw[->] (s00) -- node[auto]{$a$} (s01);
\draw[->] (s01) -- node[auto]{$a$} (s02);
\draw[->] (s10) -- node[auto]{$a$} (s11);
\draw[->] (s11) -- node[auto]{$a$} (s12);
\draw[->] (s20) -- node[auto]{$a$} (s21);
\draw[->] (s21) -- node[auto]{$a$} (s22);
\draw[->] (s02) to[loop right] node[auto]{$a$} (s02);
\draw[->] (s12) to[loop right] node[auto]{$a$} (s12);
\draw[->] (s22) to[loop right] node[auto]{$a$} (s22);
\draw[->] (s00) -- node[auto]{$b$} (s10);
\draw[->] (s10) -- node[auto]{$b$} (s20);
\draw[->] (s01) -- node[auto]{$b$} (s11);
\draw[->] (s11) -- node[auto]{$b$} (s21);
\draw[->] (s02) -- node[auto]{$b$} (s12);
\draw[->] (s12) -- node[auto]{$b$} (s22);
\draw[->] (s20) to[loop below] node[auto]{$b$} (s20);
\draw[->] (s21) to[loop below] node[auto]{$b$} (s21);
\draw[->] (s22) to[loop below] node[auto]{$b$} (s22);
\end{tikzpicture}
\end{center}

Pour la commodité de la suite de la correction, on renomme les états
de cet automate $M_3$ de la façon suivante :
\begin{center}
\begin{tikzpicture}[>=latex,line join=bevel,automaton]
\node (s00) at (0bp,0bp) [draw,rounded corners,state,initial] {$00$};
\node (s01) at (75bp,0bp) [draw,rounded corners,state,accepting by double] {$01$};
\node (s02) at (150bp,0bp) [draw,rounded corners,state] {$02$};
\node (s10) at (0bp,-50bp) [draw,rounded corners,state,accepting by double] {$10$};
\node (s11) at (75bp,-50bp) [draw,rounded corners,state,accepting by double] {$11$};
\node (s12) at (150bp,-50bp) [draw,rounded corners,state,accepting by double] {$12$};
\node (s20) at (0bp,-100bp) [draw,rounded corners,state] {$20$};
\node (s21) at (75bp,-100bp) [draw,rounded corners,state,accepting by double] {$21$};
\node (s22) at (150bp,-100bp) [draw,rounded corners,state] {$22$};
\draw[->] (s00) -- node[auto]{$a$} (s01);
\draw[->] (s01) -- node[auto]{$a$} (s02);
\draw[->] (s10) -- node[auto]{$a$} (s11);
\draw[->] (s11) -- node[auto]{$a$} (s12);
\draw[->] (s20) -- node[auto]{$a$} (s21);
\draw[->] (s21) -- node[auto]{$a$} (s22);
\draw[->] (s02) to[loop right] node[auto]{$a$} (s02);
\draw[->] (s12) to[loop right] node[auto]{$a$} (s12);
\draw[->] (s22) to[loop right] node[auto]{$a$} (s22);
\draw[->] (s00) -- node[auto]{$b$} (s10);
\draw[->] (s10) -- node[auto]{$b$} (s20);
\draw[->] (s01) -- node[auto]{$b$} (s11);
\draw[->] (s11) -- node[auto]{$b$} (s21);
\draw[->] (s02) -- node[auto]{$b$} (s12);
\draw[->] (s12) -- node[auto]{$b$} (s22);
\draw[->] (s20) to[loop below] node[auto]{$b$} (s20);
\draw[->] (s21) to[loop below] node[auto]{$b$} (s21);
\draw[->] (s22) to[loop below] node[auto]{$b$} (s22);
\end{tikzpicture}
\end{center}
Ici, l'état $0\bullet$ signifie que l'automate n'a pas rencontré
de $a$, l'état $1\bullet$ qu'il en a rencontré exactement un, et
l'état $2\bullet$ qu'il en a rencontré au moins deux ; les états
$\bullet0$, $\bullet1$ et $\bullet2$ ont la même signification pour la
lettre $b$.

\smallbreak

(4) L'automate $M_3$ est déjà minimal.  En effet, l'algorithme de
minimisation commence par séparer les classes $\{01,10,11,12,21\}$
(états finaux) et $\{00,02,20,22\}$ (états non-finaux) ; ensuite, la
transition étiquetée par $a$ sépare la classe $\{01,10,11,12,21\}$ en
$\{01,21\}$ (qui vont vers un état non-final) et $\{10,11,12\}$ (qui
vont vers un état final), et la classe $\{00,02,20,22\}$ en
$\{00,20\}$ (qui vont vers un état final) et $\{02,22\}$ (qui vont
vers un non-final).  La transition étiquetée par $b$ sépare ensuite en
deux chacune des trois classes $\{00,20\}$, $\{01,21\}$ et $\{02,22\}$
(car le premier élément va dans la classe $\{10,11,12\}$ tandis que le
second reste dans la même classe) et sépare en trois la classe
$\{10,11,12\}$.  On a donc séparé chacun des états.

\smallbreak

(5) Pour reconnaître le complémentaire du langage reconnu par un
automate fini déterministe complet, il suffit d'échanger états finaux
et non-finaux : on peut donc prendre l'automate dessiné en (3) avec,
cette fois, la convention que les états simplement entourés sont
finaux (et les doublement entourés sont non-finaux).
Appelons-le $M_5$.

\emph{Attention :} échanger états finaux et non-finaux ne marche pas
pour reconnaître le complémentaire du langage reconnu par un automate
non-déterministe ou incomplet (car la négation de « il existe un
chemin qui va vers un état final » est « aucun chemin ne va vers un
état final » et pas « il existe un chemin qui va vers un état
non-final »).

\smallbreak

(6) Puisque $L$ est le langage formé des mots comportant exactement un
$a$ ou (inclusif) exactement un $b$, son complémentaire $\overline{L}$
est formé des mots ayant un nombre différent de $1$ de $a$ \emph{et}
un nombre différent de $1$ de $b$ ; si on préfère, il s'agit du
langage comportant ($0$ ou au moins $2$ fois la lettre $a$) \emph{et}
($0$ ou au moins $2$ fois la lettre $b$).

\smallbreak

(7) L'élimination des états n'est pas trop complexe car l'automate
$M_5$ a très peu de boucles.  Éliminons simultanément tous les états
non-finaux ($01$, $10$, $11$, $12$ et $21$), et profitons-en pour
créer un nouvel (et unique) état final $F$ :
\begin{center}
\begin{tikzpicture}[>=latex,line join=bevel,automaton]
\node (s00) at (0bp,0bp) [draw,rounded corners,state,initial] {$00$};
\node (s02) at (100bp,0bp) [draw,rounded corners,state] {$02$};
\node (s20) at (0bp,-60bp) [draw,rounded corners,state] {$20$};
\node (s22) at (100bp,-60bp) [draw,rounded corners,state] {$22$};
\node (F) at (160bp,-60bp) [draw,circle,state,final] {$F$};
\draw[->] (s00) -- node[auto]{$aa$} (s02);
\draw[->] (s20) -- node[auto]{$ab{*}a$} (s22);
\draw[->] (s02) to[loop above] node[auto]{$a$} (s02);
\draw[->] (s00) -- node[auto]{$bb$} (s20);
\draw[->] (s02) -- node[auto]{$ba{*}b$} (s22);
\draw[->] (s20) to[loop left] node[auto]{$b$} (s20);
\draw[->] (s22) to[loop below] node[auto]{$a|b$} (s22);
\draw[->] (s22) -- node[auto]{$\varepsilon$} (F);
\draw[->] (s00) -- node[auto]{$r$} (s22);
\draw[->,out=0,in=90] (s02) to node[auto]{$\varepsilon$} (F);
\draw[->,out=270,in=270] (s20) to node[auto,below]{$\varepsilon$} (F);
\draw[->] (s00) to[out=45,in=180] (100bp,40bp) to[out=0,in=45] node[auto,above]{$\varepsilon$} (F);
\end{tikzpicture}
\end{center}
où $r := abaa{*}b \,|\, abbb{*}a \,|\, baaa{*}b \,|\, babb{*}a$
(correspondant aux quatre façons de passer de $00$ à $22$ dans le
graphe ci-dessus).  Éliminons l'état $02$ et l'état $20$ :
\begin{center}
\begin{tikzpicture}[>=latex,line join=bevel,automaton]
\node (s00) at (0bp,0bp) [draw,rounded corners,state,initial] {$00$};
\node (s22) at (100bp,-60bp) [draw,rounded corners,state] {$22$};
\node (F) at (160bp,-60bp) [draw,circle,state,final] {$F$};
\draw[->] (s22) -- node[auto]{$\varepsilon$} (F);
\draw[->] (s22) to[loop above] node[auto]{$a|b$} (s22);
\draw[->] (s00) -- node[auto]{$r'$} (s22);
\draw[->,out=0,in=90] (s00) to node[auto]{$\varepsilon|aaa{*}|bbb{*}$} (F);
\end{tikzpicture}
\end{center}
où $r' := r \penalty0\,|\, aaa{*}ba{*}b \penalty0\,|\, bbb{*}ab{*}a =
abaa{*}b \penalty0\,|\, abbb{*}a \penalty0\,|\, baaa{*}b
\penalty0\,|\, babb{*}a \penalty0\,|\, aaa{*}ba{*}b \penalty0\,|\,
bbb{*}ab{*}a$.  On obtient finalement l'expression rationnelle
suivante pour $\overline{L}$ :
\[
\varepsilon\,|\,aaa{*}\,|\,bbb{*}\,|\,(abaa{*}b \,|\,
abbb{*}a \,|\, baaa{*}b \,|\, babb{*}a
\,|\, aaa{*}ba{*}b \,|\, bbb{*}ab{*}a)(a|b){*}
\]
Pour comprendre cette expression rationnelle, la disjonction de plus
haut niveau correspond aux quatre possibilités : (i) $0$ fois la
lettre $a$ et $0$ fois la lettre $b$, (ii) au moins $2$ fois la lettre
$a$ et $0$ fois la lettre $b$, (iii) $0$ fois la lettre $a$ et au
moins $2$ fois la lettre $b$, et (iv) au moins $2$ fois la lettre $a$
et au moins $2$ fois la lettre $b$.  Pour mieux comprendre
l'expression du cas (iv), on peut remarquer que $abaa{*}b
\penalty0\,|\, baaa{*}b \penalty0\,|\, aaa{*}ba{*}b$ dénote le langage
formé des mots comportant au moins deux $a$ et exactement deux $b$ et
qui finissent par un $b$, et symétriquement $abbb{*}a \penalty0\,|\,
babb{*}a \penalty0\,|\,bbb{*}ab{*}a$ dénote le langage formé des mots
comportant au moins deux $b$ et exactement deux $a$ et qui finissent
par un $a$ : l'expression du cas (iv) correspond donc à écrire un mot
ayant au moins deux $a$ et au moins deux $b$ comme le premier préfixe
qui vérifie cette propriété suivi d'un suffixe quelconque.  (On
pouvait utiliser directement ce raisonnement pour produire
l'expression.)
\end{corrige}

%

\exercice

Dans cet exercice, on pose $\Sigma = \{a,b\}$.

On considère l'expression rationnelle $r$ suivante : $(ab|ba){*}$ (sur
l'alphabet $\Sigma$).  On appelle $L := L(r)$ le langage qu'elle
dénote.

(0) Donner quatre exemples de mots de $L$ et quatre exemples de mots
n'appartenant pas à $L$.

\begin{corrige}
Par exemple, les mots $\varepsilon$, $ab$, $ba$ et $abba$
appartiennent à $L$.  Les mots $a$, $b$, $aa$ et $aba$ n'appartiennent
pas à $L$.
\end{corrige}

(1) Traiter l'une \emph{ou} l'autre des questions suivantes :
(i) construire l'automate de Glushkov $\mathscr{A}_1$ de $r$ ;
(ii) construire l'automate de Thompson de $r$, puis éliminer les
transitions spontanées (= $\varepsilon$-transitions) de ce dernier (on
retirera les états devenus inutiles) : on appellera $\mathscr{A}_1$
l'automate ainsi obtenu.

\begin{corrige}
L'automate $\mathscr{A}_1$ obtenu est le suivant :
\begin{center}
\begin{tikzpicture}[>=latex,line join=bevel,automaton]
\node (S0) at (0bp,0bp) [draw,circle,state,initial,final,accepting below] {$0$};
\node (S1) at (70bp,35bp) [draw,circle,state] {$1$};
\node (S2) at (70bp,-35bp) [draw,circle,state] {$2$};
\node (S3) at (150bp,35bp) [draw,circle,state,final] {$3$};
\node (S4) at (150bp,-35bp) [draw,circle,state,final] {$4$};
\draw [->] (S0) -- node[auto,near end] {$a$} (S1);
\draw [->] (S0) -- node[auto,below] {$b$} (S2);
\draw [->] (S1) -- node[auto,below] {$b$} (S3);
\draw [->] (S2) -- node[auto] {$a$} (S4);
\draw [->] (S3) -- node[auto,above,very near end] {$b$} (S2);
\draw [->] (S4) -- node[auto,very near end] {$a$} (S1);
\draw [->] (S3) to[out=135,in=45] node[auto,above] {$a$} (S1);
\draw [->] (S4) to[out=225,in=315] node[auto] {$b$} (S2);
\end{tikzpicture}
\end{center}
C'est l'automate de Glushkov.  Si on commence par construire
l'automate de Thompson, on obtient le suivant (à $12$ états) :
\begin{center}
\scalebox{0.70}{%
\begin{tikzpicture}[>=latex,line join=bevel,automaton]
\node (S0) at (0bp,0bp) [draw,circle,state,initial] {};
\node (S0u) at (60bp,0bp) [draw,circle,state] {};
\node (S0a) at (120bp,40bp) [draw,circle,state] {};
\node (S1) at (180bp,40bp) [draw,circle,state] {};
\node (S1u) at (240bp,40bp) [draw,circle,state] {};
\node (S3) at (300bp,40bp) [draw,circle,state] {};
\node (S0b) at (120bp,-40bp) [draw,circle,state] {};
\node (S2) at (180bp,-40bp) [draw,circle,state] {};
\node (S2u) at (240bp,-40bp) [draw,circle,state] {};
\node (S4) at (300bp,-40bp) [draw,circle,state] {};
\node (SF) at (360bp,0bp) [draw,circle,state] {};
\node (SFu) at (420bp,0bp) [draw,circle,state,final] {};
\draw [->] (S0) -- (S0u);
\draw [->] (S0u) -- (S0a);
\draw [->] (S0a) -- node[auto] {$a$} (S1);  \draw [->] (S1) -- (S1u);
\draw [->] (S1u) -- node[auto] {$b$} (S3);  \draw [->] (S3) -- (SF);
\draw [->] (S0u) -- (S0b);
\draw [->] (S0b) -- node[auto] {$b$} (S2);  \draw [->] (S2) -- (S2u);
\draw [->] (S2u) -- node[auto] {$a$} (S4);  \draw [->] (S4) -- (SF);
\draw [->] (SF) -- (SFu);
\draw [->] (SF) to[out=90,in=0] (210bp,70bp) to[out=180,in=90] (S0u);
\draw [->] (S0) to[out=270,in=180] (120bp,-70bp) -- (300bp,-70bp) to[out=0,in=270] (SFu);
\end{tikzpicture}
}
\end{center}
(où toutes les flèches non étiquetées sont des transitions spontanées,
c'est-à-dire qu'il faut les imaginer étiquetées par $\varepsilon$) qui
par élimination des transitions spontanées donne celui $\mathscr{A}_1$
qu'on a tracé au-dessus (les seuls états qui subsistent après
élimination des transitions spontanées sont l'état initial et les
quatre états auxquels aboutissent une transition non spontanée).
\end{corrige}

(2) À partir de $\mathscr{A}_1$, construire un automate fini
déterministe complet reconnaissant $L$.  On appellera $\mathscr{A}_2$
l'automate en question.

\begin{corrige}
L'automate $\mathscr{A}_1$ est déterministe incomplet : il s'agit donc
simplement de lui ajouter un état « puits » pour le rendre complet, ce
qui donne l'automate $\mathscr{A}_2$ suivant :
\begin{center}
\begin{tikzpicture}[>=latex,line join=bevel,automaton]
\node (S0) at (0bp,0bp) [draw,circle,state,initial,final,accepting below] {$0$};
\node (S1) at (70bp,35bp) [draw,circle,state] {$1$};
\node (S2) at (70bp,-35bp) [draw,circle,state] {$2$};
\node (S3) at (150bp,35bp) [draw,circle,state,final] {$3$};
\node (S4) at (150bp,-35bp) [draw,circle,state,final] {$4$};
\node (Sink) at (220bp,0bp) [draw,circle,state] {$\bot$};
\draw [->] (S0) -- node[auto,near end] {$a$} (S1);
\draw [->] (S0) -- node[auto,below] {$b$} (S2);
\draw [->] (S1) -- node[auto] {$b$} (S3);
\draw [->] (S2) -- node[auto,below] {$a$} (S4);
\draw [->] (S3) -- node[auto,above,very near end] {$b$} (S2);
\draw [->] (S4) -- node[auto,very near end] {$a$} (S1);
\draw [->] (S3) to[out=135,in=45] node[auto,above] {$a$} (S1);
\draw [->] (S4) to[out=225,in=315] node[auto] {$b$} (S2);
\draw [->] (S1) -- node[auto,very near end] {$a$} (Sink);
\draw [->] (S2) -- node[auto,below,very near end] {$b$} (Sink);
\draw [->] (Sink) to[loop right] node[auto] {$a,b$} (Sink);
\end{tikzpicture}
\end{center}
\vskip-\baselineskip\vskip-1ex\strut
\end{corrige}

(3) Minimiser l'automate $\mathscr{A}_2$.  On appellera
$\mathscr{A}_3$ l'automate canonique ainsi obtenu.

(On obtient un automate ayant $4$ états.)

\begin{corrige}
L'algorithme de minimisation identifie les trois états finaux
($0,3,4$) de l'automate $\mathscr{A}_2$, ce qui donne
l'automate $\mathscr{A}_3$ suivant (on note $0$ pour l'état résultant
de l'identification de $0,3,4$) :
\begin{center}
\begin{tikzpicture}[>=latex,line join=bevel,automaton]
\node (S0) at (0bp,0bp) [draw,circle,state,initial,final,accepting below] {$0$};
\node (S1) at (70bp,35bp) [draw,circle,state] {$1$};
\node (S2) at (70bp,-35bp) [draw,circle,state] {$2$};
\node (Sink) at (140bp,0bp) [draw,circle,state] {$\bot$};
\draw [->] (S0) to[out=75,in=180] node[auto,near end] {$a$} (S1);
\draw [->] (S0) to[out=285,in=180] node[auto,below] {$b$} (S2);
\draw [->] (S1) to[out=270,in=15] node[auto,above] {$b$} (S0);
\draw [->] (S2) to[out=90,in=345] node[auto] {$a$} (S0);
\draw [->] (S1) -- node[auto,near end] {$a$} (Sink);
\draw [->] (S2) -- node[auto,below,near end] {$b$} (Sink);
\draw [->] (Sink) to[loop right] node[auto] {$a,b$} (Sink);
\end{tikzpicture}
\end{center}
\vskip-\baselineskip\vskip-1ex\strut
\end{corrige}

(4) Construire un automate $\mathscr{A}_4$ reconnaissant le langage $M
:= \Sigma^*\setminus L$ complémentaire de $L$.

\begin{corrige}
L'automate $\mathscr{A}_3$ étant un automate fini déterministe complet
reconnaissant $L$, il suffit de marquer finaux les états non-finaux et
vice versa pour obtenir l'automate $\mathscr{A}_4$ suivant :
\begin{center}
\begin{tikzpicture}[>=latex,line join=bevel,automaton]
\node (S0) at (0bp,0bp) [draw,circle,state,initial] {$0$};
\node (S1) at (70bp,35bp) [draw,circle,state,final,accepting above] {$1$};
\node (S2) at (70bp,-35bp) [draw,circle,state,final,accepting below] {$2$};
\node (Sink) at (140bp,0bp) [draw,circle,state,final,accepting below] {$\bot$};
\draw [->] (S0) to[out=75,in=180] node[auto,near end] {$a$} (S1);
\draw [->] (S0) to[out=285,in=180] node[auto,below] {$b$} (S2);
\draw [->] (S1) to[out=270,in=15] node[auto,above] {$b$} (S0);
\draw [->] (S2) to[out=90,in=345] node[auto] {$a$} (S0);
\draw [->] (S1) -- node[auto,near end] {$a$} (Sink);
\draw [->] (S2) -- node[auto,below,near end] {$b$} (Sink);
\draw [->] (Sink) to[loop right] node[auto] {$a,b$} (Sink);
\end{tikzpicture}
\end{center}
\vskip-\baselineskip\vskip-1ex\strut
\end{corrige}

(5) En appliquant à $\mathscr{A}_4$ la méthode d'élimination des
états, déterminer une expression rationnelle dénotant le langage $M$.

\begin{corrige}
Pour appliquer la méthode d'élimination des états, on commence par
modifier l'automate pour avoir un unique état initial $I$, qui ne soit
pas final, et qui n'ait aucune transition qui y aboutisse, et un
unique état final $F$, qui ne soit pas initial, et sans aucune
transition qui en part :
\begin{center}
\begin{tikzpicture}[>=latex,line join=bevel,automaton]
\node (SI) at (-70bp,0bp) [draw,circle,state,initial] {$I$};
\node (S0) at (0bp,0bp) [draw,circle,state] {$0$};
\node (S1) at (70bp,35bp) [draw,circle,state] {$1$};
\node (S2) at (70bp,-35bp) [draw,circle,state] {$2$};
\node (Sink) at (140bp,0bp) [draw,circle,state] {$\bot$};
\node (SF) at (210bp,0bp) [draw,circle,state,final] {$F$};
\draw [->] (SI) -- node[auto] {$\underline{\varepsilon}$} (S0);
\draw [->] (S0) to[out=75,in=180] node[auto,near end] {$a$} (S1);
\draw [->] (S0) to[out=285,in=180] node[auto,below] {$b$} (S2);
\draw [->] (S1) to[out=270,in=15] node[auto,above] {$b$} (S0);
\draw [->] (S2) to[out=90,in=345] node[auto] {$a$} (S0);
\draw [->] (S1) -- node[auto,near end] {$a$} (Sink);
\draw [->] (S2) -- node[auto,below,near end] {$b$} (Sink);
\draw [->] (Sink) to[loop above] node[auto] {$a|b$} (Sink);
\draw [->] (S1) to[out=15,in=135] node[auto,near end] {$\underline{\varepsilon}$} (SF);
\draw [->] (S2) to[out=345,in=225] node[auto,near end] {$\underline{\varepsilon}$} (SF);
\draw [->] (Sink) -- node[auto] {$\underline{\varepsilon}$} (SF);
\end{tikzpicture}
\end{center}
On élimine ensuite l'état $\bot$ :
\begin{center}
\begin{tikzpicture}[>=latex,line join=bevel,automaton]
\node (SI) at (-70bp,0bp) [draw,circle,state,initial] {$I$};
\node (S0) at (0bp,0bp) [draw,circle,state] {$0$};
\node (S1) at (70bp,35bp) [draw,circle,state] {$1$};
\node (S2) at (70bp,-35bp) [draw,circle,state] {$2$};
\node (SF) at (210bp,0bp) [draw,circle,state,final] {$F$};
\draw [->] (SI) -- node[auto] {$\underline{\varepsilon}$} (S0);
\draw [->] (S0) to[out=75,in=180] node[auto,near end] {$a$} (S1);
\draw [->] (S0) to[out=285,in=180] node[auto,below] {$b$} (S2);
\draw [->] (S1) to[out=270,in=15] node[auto,above] {$b$} (S0);
\draw [->] (S2) to[out=90,in=345] node[auto] {$a$} (S0);
\draw [->] (S1) to[out=15,in=135] node[auto,near end] {$\underline{\varepsilon}|a(a|b){*}$} (SF);
\draw [->] (S2) to[out=345,in=225] node[auto,near end] {$\underline{\varepsilon}|b(a|b){*}$} (SF);
\end{tikzpicture}
\end{center}
Puis les états $1$ et $2$ peuvent être éliminés simultanément :
\begin{center}
\begin{tikzpicture}[>=latex,line join=bevel,automaton]
\node (SI) at (-70bp,0bp) [draw,circle,state,initial] {$I$};
\node (S0) at (0bp,0bp) [draw,circle,state] {$0$};
\node (SF) at (210bp,0bp) [draw,circle,state,final] {$F$};
\draw [->] (SI) -- node[auto] {$\underline{\varepsilon}$} (S0);
\draw [->] (S0) -- node[auto] {$a(\underline{\varepsilon}|a(a|b){*})|b(\underline{\varepsilon}|b(a|b){*})$} (SF);
\draw [->] (S0) to[loop above] node[auto] {$ab|ba$} (S0);
\end{tikzpicture}
\end{center}
Et l'expression rationnelle finalement obtenue est la suivante :
\[
(ab|ba){*}(a(\underline{\varepsilon}|a(a|b){*})|b(\underline{\varepsilon}|b(a|b){*}))
\]
On pouvait aussi donner, entre autres choses :
\[
(ab|ba){*}(a|aa(a|b){*}|b|bb(a|b){*})
\]
(obtenue en éliminant les états $1$ et $2$ avant $\bot$).
\end{corrige}

%
%
%

\subsection{Langages algébriques et grammaires hors contexte}

\exercice

Considérons le fragment simplifié suivant de la grammaire d'un langage
de programmation hypothétique :
\[
\begin{aligned}
\mathit{Instruction} &\rightarrow \mathtt{foo} \;|\; \mathtt{bar} \;|\; \mathtt{qux} \;|\; \mathit{Conditional}\\
|&\phantom{\rightarrow} \mathtt{begin}\ \mathit{InstrList}\ \mathtt{end}\\
\mathit{Conditional} &\rightarrow \mathtt{if}\ \mathit{Expression}\ \mathtt{then}\ \mathit{Instruction}\ \mathtt{else}\ \mathit{Instruction}\\
|&\phantom{\rightarrow} \mathtt{if}\ \mathit{Expression}\ \mathtt{then}\ \mathit{Instruction}\\
\mathit{InstrList} &\rightarrow \mathit{Instruction} \;|\; \mathit{Instruction}\ \mathit{InstrList}\\
\mathit{Expression} &\rightarrow \mathtt{true} \;|\; \mathtt{false} \;|\; \mathtt{happy} \;|\; \mathtt{trippy}\\
\end{aligned}
\]
(Ici, les « lettres » ou tokens ont été écrits comme des mots, par
exemple $\mathtt{foo}$ est une « lettre » : les terminaux sont écrits
en police à espacement fixe tandis que les nonterminaux sont en
italique et commencent par une majuscule.  On prendra
$\mathit{Instruction}$ pour axiome.)

(1) Donner l'arbre d'analyse de :
$\mathtt{if}\penalty0\ \mathtt{happy}\penalty0\ \mathtt{then}\penalty0\ \mathtt{if}\penalty0\ \mathtt{trippy}\penalty0\ \mathtt{then}\penalty0\ \mathtt{foo}\penalty0\ \mathtt{else}\penalty0\ \mathtt{bar}\penalty0\ \mathtt{else}\penalty0\ \mathtt{qux}$ ;
expliquer brièvement pourquoi il n'y en a qu'un.

(2) Donner deux arbres d'analyse distincts de :
$\mathtt{if}\penalty0\ \mathtt{happy}\penalty0\ \mathtt{then}\penalty0\ \mathtt{if}\penalty0\ \mathtt{trippy}\penalty0\ \mathtt{then}\penalty0\ \mathtt{foo}\penalty0\ \mathtt{else}\penalty0\ \mathtt{bar}$.
Que peut-on dire de la grammaire présentée ?

(3) En supposant que, dans ce langage,
$\mathtt{begin}\penalty0\ I\penalty0\ \penalty0\mathtt{end}$ (où $I$
est une instruction) a le même effet que $I$ seul, comment un
programmeur peut-il réécrire l'instruction considérée en (2) pour
obtenir un comportant équivalent à l'une ou l'autre des deux
interprétations ?

(4) Modifier légèrement la grammaire proposée de manière à obtenir une
grammaire faiblement équivalente dans laquelle seul l'un des arbres
d'analyse obtenus en (2) est possible (i.e., une grammaire qui force
cette interprétation-là par défaut) ; on pourra être amené à
introduire des nouveaux nonterminaux pour des variantes de
$\mathit{Instruction}$ et $\mathit{Conditional}$ qui interdisent
récursivement les conditionnelles sans $\mathtt{else}$.

\begin{corrige}
(1) L'arbre d'analyse de
  $\mathtt{if}\penalty0\ \mathtt{happy}\penalty0\ \mathtt{then}\penalty0\ \mathtt{if}\penalty0\ \mathtt{trippy}\penalty0\ \mathtt{then}\penalty0\ \mathtt{foo}\penalty0\ \mathtt{else}\penalty0\ \mathtt{bar}\penalty0\ \mathtt{else}\penalty0\ \mathtt{qux}$
  est le suivant (en notant $I$, $C$ et $E$ pour
  $\mathit{Instruction}$, $\mathit{Condition}$ et
  $\mathit{Expression}$ respectivement) :
\begin{center}
\tikzstyle{automaton}=[scale=0.4]
%%% begin ex2p1 %%%

\begin{tikzpicture}[>=latex,line join=bevel,automaton]
%%
\node (foo) at (279bp,18bp) [draw,draw=none] {$\mathtt{foo}$};
  \node (then1) at (207bp,90bp) [draw,draw=none] {$\mathtt{then}$};
  \node (if0) at (27bp,234bp) [draw,draw=none] {$\mathtt{if}$};
  \node (if1) at (63bp,90bp) [draw,draw=none] {$\mathtt{if}$};
  \node (bar) at (423bp,18bp) [draw,draw=none] {$\mathtt{bar}$};
  \node (I1) at (243bp,234bp) [draw,draw=none] {$I$};
  \node (I0) at (207bp,378bp) [draw,draw=none] {$I$};
  \node (I3) at (423bp,90bp) [draw,draw=none] {$I$};
  \node (I2) at (279bp,90bp) [draw,draw=none] {$I$};
  \node (I4) at (387bp,234bp) [draw,draw=none] {$I$};
  \node (trippy) at (135bp,18bp) [draw,draw=none] {$\mathtt{trippy}$};
  \node (else0) at (315bp,234bp) [draw,draw=none] {$\mathtt{else}$};
  \node (else1) at (351bp,90bp) [draw,draw=none] {$\mathtt{else}$};
  \node (qux) at (387bp,162bp) [draw,draw=none] {$\mathtt{qux}$};
  \node (C1) at (243bp,162bp) [draw,draw=none] {$C$};
  \node (then0) at (171bp,234bp) [draw,draw=none] {$\mathtt{then}$};
  \node (C0) at (207bp,306bp) [draw,draw=none] {$C$};
  \node (E1) at (135bp,90bp) [draw,draw=none] {$E$};
  \node (E0) at (99bp,234bp) [draw,draw=none] {$E$};
  \node (happy) at (99bp,162bp) [draw,draw=none] {$\mathtt{happy}$};
  \draw [] (I2) ..controls (279bp,60.846bp) and (279bp,46.917bp)  .. (foo);
  \draw [] (C0) ..controls (192.52bp,276.85bp) and (185.36bp,262.92bp)  .. (then0);
  \draw [] (I0) ..controls (207bp,348.85bp) and (207bp,334.92bp)  .. (C0);
  \draw [] (C0) ..controls (250.16bp,277.03bp) and (271.72bp,263.05bp)  .. (else0);
  \draw [] (C1) ..controls (286.16bp,133.03bp) and (307.72bp,119.05bp)  .. (else1);
  \draw [] (C1) ..controls (257.48bp,132.85bp) and (264.64bp,118.92bp)  .. (I2);
  \draw [] (E0) ..controls (99bp,204.85bp) and (99bp,190.92bp)  .. (happy);
  \draw [] (C1) ..controls (199.84bp,133.03bp) and (178.28bp,119.05bp)  .. (E1);
  \draw [] (C1) ..controls (228.52bp,132.85bp) and (221.36bp,118.92bp)  .. (then1);
  \draw [] (C0) ..controls (263.23bp,285.79bp) and (310.83bp,268.84bp)  .. (351bp,252bp) .. controls (353.93bp,250.77bp) and (356.97bp,249.43bp)  .. (I4);
  \draw [] (C1) ..controls (186.77bp,141.79bp) and (139.17bp,124.84bp)  .. (99bp,108bp) .. controls (96.068bp,106.77bp) and (93.027bp,105.43bp)  .. (if1);
  \draw [] (C0) ..controls (150.77bp,285.79bp) and (103.17bp,268.84bp)  .. (63bp,252bp) .. controls (60.068bp,250.77bp) and (57.027bp,249.43bp)  .. (if0);
  \draw [] (E1) ..controls (135bp,60.846bp) and (135bp,46.917bp)  .. (trippy);
  \draw [] (C1) ..controls (299.23bp,141.79bp) and (346.83bp,124.84bp)  .. (387bp,108bp) .. controls (389.93bp,106.77bp) and (392.97bp,105.43bp)  .. (I3);
  \draw [] (C0) ..controls (221.48bp,276.85bp) and (228.64bp,262.92bp)  .. (I1);
  \draw [] (I4) ..controls (387bp,204.85bp) and (387bp,190.92bp)  .. (qux);
  \draw [] (I1) ..controls (243bp,204.85bp) and (243bp,190.92bp)  .. (C1);
  \draw [] (I3) ..controls (423bp,60.846bp) and (423bp,46.917bp)  .. (bar);
  \draw [] (C0) ..controls (163.84bp,277.03bp) and (142.28bp,263.05bp)  .. (E0);
%
\end{tikzpicture}

%%% end ex2p1 %%%
\end{center}

Il est le seul possible car une fois acquis que les deux $\mathtt{if}$
comportent chacun un $\mathtt{else}$, il se construit ensuite en
descendant de façon unique (l'instruction est forcément une condition,
qui s'analyse en $\mathtt{if}\ E\ \mathtt{then}\ I\ \mathtt{else}\ I$
de façon unique, et chacun des morceaux s'analyse de nouveau de façon
unique).

\vskip .5explus.1fil

(2) Un arbre d'analyse possible consiste à associer le
$\mathtt{else}\penalty0\ \mathtt{bar}$ avec
$\mathtt{if}\penalty0\ \mathtt{trippy}\penalty0\ \mathtt{then}\penalty0\ \mathtt{foo}$ :
\begin{center}
\tikzstyle{automaton}=[scale=0.4]
%%% begin ex2p1a %%%

\begin{tikzpicture}[>=latex,line join=bevel,automaton]
%%
\node (bar) at (423bp,18bp) [draw,draw=none] {$\mathtt{bar}$};
  \node (then1) at (207bp,90bp) [draw,draw=none] {$\mathtt{then}$};
  \node (if0) at (27bp,234bp) [draw,draw=none] {$\mathtt{if}$};
  \node (if1) at (63bp,90bp) [draw,draw=none] {$\mathtt{if}$};
  \node (I1) at (243bp,234bp) [draw,draw=none] {$I$};
  \node (I0) at (135bp,378bp) [draw,draw=none] {$I$};
  \node (I3) at (423bp,90bp) [draw,draw=none] {$I$};
  \node (I2) at (279bp,90bp) [draw,draw=none] {$I$};
  \node (trippy) at (135bp,18bp) [draw,draw=none] {$\mathtt{trippy}$};
  \node (else1) at (351bp,90bp) [draw,draw=none] {$\mathtt{else}$};
  \node (foo) at (279bp,18bp) [draw,draw=none] {$\mathtt{foo}$};
  \node (C1) at (243bp,162bp) [draw,draw=none] {$C$};
  \node (then0) at (171bp,234bp) [draw,draw=none] {$\mathtt{then}$};
  \node (C0) at (135bp,306bp) [draw,draw=none] {$C$};
  \node (E1) at (135bp,90bp) [draw,draw=none] {$E$};
  \node (E0) at (99bp,234bp) [draw,draw=none] {$E$};
  \node (happy) at (99bp,162bp) [draw,draw=none] {$\mathtt{happy}$};
  \draw [] (I2) ..controls (279bp,60.846bp) and (279bp,46.917bp)  .. (foo);
  \draw [] (C0) ..controls (149.48bp,276.85bp) and (156.64bp,262.92bp)  .. (then0);
  \draw [] (I0) ..controls (135bp,348.85bp) and (135bp,334.92bp)  .. (C0);
  \draw [] (C1) ..controls (228.52bp,132.85bp) and (221.36bp,118.92bp)  .. (then1);
  \draw [] (C1) ..controls (286.16bp,133.03bp) and (307.72bp,119.05bp)  .. (else1);
  \draw [] (I3) ..controls (423bp,60.846bp) and (423bp,46.917bp)  .. (bar);
  \draw [] (E0) ..controls (99bp,204.85bp) and (99bp,190.92bp)  .. (happy);
  \draw [] (C1) ..controls (199.84bp,133.03bp) and (178.28bp,119.05bp)  .. (E1);
  \draw [] (C1) ..controls (186.77bp,141.79bp) and (139.17bp,124.84bp)  .. (99bp,108bp) .. controls (96.068bp,106.77bp) and (93.027bp,105.43bp)  .. (if1);
  \draw [] (C0) ..controls (91.844bp,277.03bp) and (70.277bp,263.05bp)  .. (if0);
  \draw [] (E1) ..controls (135bp,60.846bp) and (135bp,46.917bp)  .. (trippy);
  \draw [] (C1) ..controls (299.23bp,141.79bp) and (346.83bp,124.84bp)  .. (387bp,108bp) .. controls (389.93bp,106.77bp) and (392.97bp,105.43bp)  .. (I3);
  \draw [] (C0) ..controls (178.16bp,277.03bp) and (199.72bp,263.05bp)  .. (I1);
  \draw [] (C1) ..controls (257.48bp,132.85bp) and (264.64bp,118.92bp)  .. (I2);
  \draw [] (I1) ..controls (243bp,204.85bp) and (243bp,190.92bp)  .. (C1);
  \draw [] (C0) ..controls (120.52bp,276.85bp) and (113.36bp,262.92bp)  .. (E0);
%
\end{tikzpicture}

%%% end ex2p1a %%%
\end{center}
un autre consiste à associer le $\mathtt{else}\penalty0\ \mathtt{bar}$
avec
$\mathtt{if}\penalty0\ \mathtt{happy}\penalty0\ \mathtt{then}\penalty0\ ...$ :
\begin{center}
\tikzstyle{automaton}=[scale=0.4]
%%% begin ex2p1b %%%

\begin{tikzpicture}[>=latex,line join=bevel,automaton]
%%
\node (bar) at (387bp,162bp) [draw,draw=none] {$\mathtt{bar}$};
  \node (then1) at (279bp,90bp) [draw,draw=none] {$\mathtt{then}$};
  \node (if0) at (27bp,234bp) [draw,draw=none] {$\mathtt{if}$};
  \node (if1) at (135bp,90bp) [draw,draw=none] {$\mathtt{if}$};
  \node (I1) at (243bp,234bp) [draw,draw=none] {$I$};
  \node (I0) at (207bp,378bp) [draw,draw=none] {$I$};
  \node (I2) at (351bp,90bp) [draw,draw=none] {$I$};
  \node (I4) at (387bp,234bp) [draw,draw=none] {$I$};
  \node (trippy) at (207bp,18bp) [draw,draw=none] {$\mathtt{trippy}$};
  \node (else0) at (315bp,234bp) [draw,draw=none] {$\mathtt{else}$};
  \node (foo) at (351bp,18bp) [draw,draw=none] {$\mathtt{foo}$};
  \node (C1) at (243bp,162bp) [draw,draw=none] {$C$};
  \node (then0) at (171bp,234bp) [draw,draw=none] {$\mathtt{then}$};
  \node (C0) at (207bp,306bp) [draw,draw=none] {$C$};
  \node (E1) at (207bp,90bp) [draw,draw=none] {$E$};
  \node (E0) at (99bp,234bp) [draw,draw=none] {$E$};
  \node (happy) at (99bp,162bp) [draw,draw=none] {$\mathtt{happy}$};
  \draw [] (I2) ..controls (351bp,60.846bp) and (351bp,46.917bp)  .. (foo);
  \draw [] (C0) ..controls (192.52bp,276.85bp) and (185.36bp,262.92bp)  .. (then0);
  \draw [] (I0) ..controls (207bp,348.85bp) and (207bp,334.92bp)  .. (C0);
  \draw [] (C0) ..controls (250.16bp,277.03bp) and (271.72bp,263.05bp)  .. (else0);
  \draw [] (C1) ..controls (286.16bp,133.03bp) and (307.72bp,119.05bp)  .. (I2);
  \draw [] (I4) ..controls (387bp,204.85bp) and (387bp,190.92bp)  .. (bar);
  \draw [] (E0) ..controls (99bp,204.85bp) and (99bp,190.92bp)  .. (happy);
  \draw [] (C1) ..controls (228.52bp,132.85bp) and (221.36bp,118.92bp)  .. (E1);
  \draw [] (C0) ..controls (263.23bp,285.79bp) and (310.83bp,268.84bp)  .. (351bp,252bp) .. controls (353.93bp,250.77bp) and (356.97bp,249.43bp)  .. (I4);
  \draw [] (C1) ..controls (199.84bp,133.03bp) and (178.28bp,119.05bp)  .. (if1);
  \draw [] (C0) ..controls (150.77bp,285.79bp) and (103.17bp,268.84bp)  .. (63bp,252bp) .. controls (60.068bp,250.77bp) and (57.027bp,249.43bp)  .. (if0);
  \draw [] (E1) ..controls (207bp,60.846bp) and (207bp,46.917bp)  .. (trippy);
  \draw [] (C0) ..controls (221.48bp,276.85bp) and (228.64bp,262.92bp)  .. (I1);
  \draw [] (C1) ..controls (257.48bp,132.85bp) and (264.64bp,118.92bp)  .. (then1);
  \draw [] (I1) ..controls (243bp,204.85bp) and (243bp,190.92bp)  .. (C1);
  \draw [] (C0) ..controls (163.84bp,277.03bp) and (142.28bp,263.05bp)  .. (E0);
%
\end{tikzpicture}

%%% end ex2p1b %%%
\end{center}
La grammaire présentée est donc ambiguë.

\vskip .5explus.1fil

(3) Pour forcer la première interprétation (le
$\mathtt{else}\penalty0\ \mathtt{bar}$ se rapporte au
$\mathtt{if}\penalty0\ \mathtt{trippy}$), on peut écrire :
$\mathtt{if}\penalty0\ \mathtt{happy}\penalty0\ \mathtt{then}\penalty0\ \mathtt{begin}\penalty0\ \mathtt{if}\penalty0\ \mathtt{trippy}\penalty0\ \mathtt{then}\penalty0\ \mathtt{foo}\penalty0\ \mathtt{else}\penalty0\ \mathtt{bar}\penalty0\ \mathtt{end}$.

Pour forcer la seconde interprétation (le
$\mathtt{else}\penalty0\ \mathtt{bar}$ se rapporte au
$\mathtt{if}\penalty0\ \mathtt{happy}$), on peut écrire :
$\mathtt{if}\penalty0\ \mathtt{happy}\penalty0\ \mathtt{then}\penalty0\ \mathtt{begin}\penalty0\ \mathtt{if}\penalty0\ \mathtt{trippy}\penalty0\ \mathtt{then}\penalty0\ \mathtt{foo}\penalty0\ \mathtt{end}\penalty0\ \mathtt{else}\penalty0\ \mathtt{bar}$.

\vskip .5explus.1fil

(4) Pour forcer la première interprétation (le $\mathtt{else}$ se
rapporte au $\mathtt{if}$ le plus proche possible), on peut modifier
la grammaire comme suit :
\[
\begin{aligned}
\mathit{Instruction} &\rightarrow \mathtt{foo} \;|\; \mathtt{bar} \;|\; \mathtt{qux} \;|\; \mathit{Conditional}\\
|&\phantom{\rightarrow} \mathtt{begin}\ \mathit{InstrList}\ \mathtt{end}\\
\mathit{InstrNoSC} &\rightarrow \mathtt{foo} \;|\; \mathtt{bar} \;|\; \mathtt{qux} \;|\; \mathit{CondNoSC}\\
|&\phantom{\rightarrow} \mathtt{begin}\ \mathit{InstrList}\ \mathtt{end}\\
\mathit{Conditional} &\rightarrow \mathtt{if}\ \mathit{Expression}\ \mathtt{then}\ \mathit{InstrNoSC}\ \mathtt{else}\ \mathit{Instruction}\\
|&\phantom{\rightarrow} \mathtt{if}\ \mathit{Expression}\ \mathtt{then}\ \mathit{Instruction}\\
\mathit{CondNoSC} &\rightarrow \mathtt{if}\ \mathit{Expression}\ \mathtt{then}\ \mathit{InstrNoSC}\ \mathtt{else}\ \mathit{InstrNoSC}\\
\mathit{InstrList} &\rightarrow \mathit{Instruction} \;|\; \mathit{Instruction}\ \mathit{InstrList}\\
\mathit{Expression} &\rightarrow \mathtt{true} \;|\; \mathtt{false} \;|\; \mathtt{happy} \;|\; \mathtt{trippy}\\
\end{aligned}
\]
L'idée est d'obliger une instruction conditionnelle qui apparaîtrait
après le $\mathtt{then}$ d'une conditionnelle complète à être
elle-même complète (elle ne peut pas être courte, car alors le
$\mathtt{else}$ devrait se rattacher à elle), et ce, récursivement.
On peut montrer que la grammaire ci-dessus est inambiguë et faiblement
équivalente à celle de départ.

On peut aussi fabriquer une grammaire inambiguë, faiblement
équivalente à celle de départ, qui force l'autre interprétation (le
$\mathtt{else}$ se rapporte au $\mathtt{if}$ le plus lointain
possible), mais c'est nettement plus complexe (l'idée générale pour
apparier un $\mathtt{else}$ avec un $\mathtt{if}...\mathtt{else}$ dans
cette logique est de demander que \emph{soit} le $\mathtt{else}$ n'est
suivi d'aucun autre $\mathtt{else}$, \emph{soit} toute instruction
conditionnelle entre le $\mathtt{then}$ et le $\mathtt{else}$ est
elle-même complète).  Contrairement à la grammaire précédente, cette
grammaire, bien qu'inambiguë, est probablement impossible à analyser
avec un analyseur LR (ou même, déterministe).
\end{corrige}


%

\exercice\label{non-square-words-is-algebraic}

Soit $\Sigma = \{a,b\}$.  On considère le langage $M$ des mots qui
\emph{ne s'écrivent pas} sous la forme $ww$ avec $w\in\Sigma^*$
(c'est-à-dire sous la forme d'un carré ; autrement dit, le langage $M$
est le \emph{complémentaire} du langage $Q$ des carrés considéré dans
l'exercice \ref{square-words-not-algebraic}) : par exemple, $M$
contient les mots $a$, $b$, $ab$, $aab$ et $aabb$ mais pas
$\varepsilon$, $aa$, $abab$ ni $abaaba$.

(0) Expliquer pourquoi tout mot sur $\Sigma$ de longueur impaire est
dans $M$, et pourquoi un mot $x_1\cdots x_{2n}$ de longueur paire $2n$
est dans $M$ si et seulement si il existe $i$ tel que $x_i \neq
x_{n+i}$.

On considère par ailleurs la grammaire hors contexte $G$
(d'axiome $S$)
\[
\begin{aligned}
S &\rightarrow A \;|\; B \;|\; AB \;|\; BA\\
A &\rightarrow a \;|\; aAa \;|\; aAb \;|\; bAa \;|\; bAb\\
B &\rightarrow b \;|\; aBa \;|\; aBb \;|\; bBa \;|\; bBb\\
\end{aligned}
\]

(1) Décrire le langage $L(G,A)$ des mots dérivant de $A$ dans la
grammaire $G$ (autrement dit, le langage engendré par la grammaire
identique à $G$ mais ayant pour axiome $A$).  Décrire de même
$L(G,B)$.

(2) Montrer que tout mot de longueur impaire est dans le
langage $L(G)$ engendré par $G$.

(3) Montrer que tout mot $t \in M$ de longueur paire est dans $L(G)$.
(Indication : si $t = x_1\cdots x_{2n}$ est de longueur paire $2n$ et
que $x_i \neq x_{n+i}$, on peut considérer la factorisation de $t$ en
$x_1\cdots x_{2i-1}$ et $x_{2i}\cdots x_{2n}$.)

(4) Montrer que tout mot de $L(G)$ de longueur paire est dans $M$.

(5) En déduire que $M$ est algébrique.

\begin{corrige}
(0) En remarquant que si $n = |w|$ alors $|ww| = 2n$, on constate que
  tout mot de la forme $ww$ est de longueur paire, et de plus, que
  pour un mot de longueur $2n$, être de la forme $ww$ signifie que son
  préfixe de longueur $n$ soit égal à son suffixe de longueur $n$ ;
  c'est-à-dire, si $t = x_1\cdots x_{2n}$, que $x_1\cdots x_n =
  x_{n+1}\cdots x_{2n}$, ce qui signifie exactement $x_i = x_{n+i}$
  pour tout $1\leq i\leq n$.

(1) La règle $A \rightarrow a \,|\, aAa \,|\, aAb \,|\, bAa \,|\, bAb$
  permet de faire à partir de $A$ une dérivation qui ajoute un nombre
  quelconque de fois une lettre ($a$ ou $b$) de chaque part de $A$, et
  finalement remplace ce $A$ par $a$.  On obtient donc ainsi
  exactement les mots de longueur impaire ayant un $a$ comme lettre
  centrale : $L(G,A) = \{w_1aw_2 : |w_1|=|w_2|\}$.  De même, $L(G,B) =
  \{w_1bw_2 : |w_1|=|w_2|\}$.

(2) Tout mot de longueur impaire est soit dans $L(G,A)$ soit dans
  $L(G,B)$ selon que sa lettre centrale est un $a$ ou un $b$.  Il est
  donc dans $L(G)$ en vertu des règles $S\rightarrow A$ et
  $S\rightarrow B$.

(3) Soit $t = x_1\cdots x_{2n}$ un mot de $M$ de longueur paire $2n$ :
  d'après (0), il existe $i$ tel que $x_i \neq x_{n+i}$.  Posons alors
  $u = x_1\cdots x_{2i-1}$ et $v = x_{2i}\cdots x_{2n}$.  Chacun de
  $u$ et de $v$ est de longueur impaire.  De plus, leurs lettres
  centrales sont respectivement $x_i$ et $x_{n+i}$, et elles sont
  différentes : l'une est donc un $a$ et l'autre un $b$ ; mettons sans
  perte de généralité que $x_i = a$ et $x_{n+i} = b$.  Alors $u \in
  L(G,A)$ d'après (1) et $v \in L(G,B)$ : le mot $t = uv$ s'obtient
  donc par la règle $S \rightarrow AB$ (suivie de dérivations de $u$ à
  partir de $A$ et de $v$ a partir de $B$) : ceci montre bien $t \in
  L(G)$.

(4) On a vu en (1) que tout mot dérivant de $A$ ou de $B$ est de
  longueur impaire.  Un mot $t$ de $L(G)$ de longueur paire $2n$
  dérive donc forcément de $AB$ ou de $BA$.  Sans perte de généralité,
  supposons qu'il dérive de $AB$, et on veut montrer qu'il appartient
  à $M$.  Appelons $u$ le facteur de $t$ qui dérive de $A$ et $v$ le
  facteur de $t$ qui dérive de $B$ : on sait alors (toujours
  d'après (1)) que $u$ s'écrit sous la forme $x_1\cdots x_{2i-1}$ où
  la lettre centrale $x_i$ vaut $a$, et que $v$ s'écrit sous la forme
  (quitte à continuer la numérotation des indices) $x_{2i}\cdots
  x_{2n}$ où la lettre centrale $x_{n+i}$ vaut $b$.  Alors $x_{n+i}
  \neq x_i$ donc le mot $t$ est dans $M$ d'après (0).

(5) On a $M = L(G)$ car d'après les questions précédentes, tout mot de
  longueur impaire est dans les deux et qu'un mot de longueur paire
  est dans l'un si et seulement si il est dans l'autre.  On a donc
  montré que $M$ est algébrique.
\end{corrige}


%

\exercice\label{square-words-not-algebraic}

Soit $\Sigma = \{a,b\}$.  Montrer que le langage $Q := \{ww :
w\in\Sigma^*\}$ constitué des mots de la forme $ww$ (autrement dit,
des carrés ; par exemple, $\varepsilon$, $aa$, $abab$, $abaaba$ ou
encore $aabbaabb$ sont dans $Q$) n'est pas algébrique.  On pourra pour
cela considérer son intersection avec le langage $L_0$ dénoté par
l'expression rationnelle $a{*}b{*}a{*}b{*}$ et appliquer le lemme de
pompage.

\begin{corrige}
Supposons par l'absurde que $Q$ soit algébrique : alors son
intersection avec le langage rationnel $L_0 = \{a^m b^n a^{m'} b^{n'} :
m,n,m',n' \in \mathbb{N}\}$ est encore algébrique.  Or $Q \cap L_0 =
\{a^m b^n a^m b^n : m,n \in \mathbb{N}\}$.  On va maintenant utiliser
le lemme de pompage pour arriver à une contradiction.

Appliquons le lemme de pompage pour les langages algébriques au
langage $Q \cap L_0 = \{a^m b^n a^m b^n : m,n \in \mathbb{N}\}$
considéré : appelons $k$ l'entier dont le lemme de pompage garantit
l'existence.  Considérons le mot $t := a^k b^k a^k b^k$ : dans la
suite de cette démonstration, on appellera « bloc » de $t$ un des
quatre facteurs $a^k$, $b^k$, $a^k$ et $b^k$.  D'après la propriété de
$k$ garantie par le lemme de pompage, il doit exister une
factorisation $t = uvwxy$ pour laquelle on a (i) $|vx|\geq 1$,
(ii) $|vwx|\leq k$ et (iii) $uv^iwx^iy \in Q \cap L_0$ pour
tout $i\geq 0$.

Chacun de $v$ et de $x$ doit être contenu dans un seul bloc, i.e.,
doit être de la forme $a^\ell$ ou $b^\ell$, sinon sa répétition ($v^i$
ou $x^i$ pour $i\geq 2$, qui appartient à $L_0$ d'après (iii)) ferait
apparaître plus d'alternations entre $a$ et $b$ que le langage $L_0$
ne le permet.  Par ailleurs, la propriété (ii) assure que le facteur
$vwx$ ne peut rencontrer qu'un ou deux blocs de $t$ (pas plus).
Autrement dit, $v$ et $x$ sont contenus dans deux blocs de $t$ qui
sont identiques ou bien consécutifs\footnote{La formulation est
  choisie pour avoir un sens même si $v$ ou $x$ est le mot vide (ce
  qui est possible \textit{a priori}).}.

D'après la propriété (i), au moins l'un de $v$ et de $x$ n'est pas le
mot vide.  Si ce facteur non trivial est dans le premier bloc $a^k$,
l'autre ne peut pas être dans l'autre bloc $a^k$ d'après ce qui vient
d'être dit : donc $uv^iwx^iy$ est de la forme $a^{k'} b^n a^k b^k$
avec $k'>k$ si $i>1$, qui n'appartient pas à $Q\cap L_0$, une
contradiction.  De même, le facteur non trivial est dans le premier
bloc $b^k$, l'autre ne peut pas être dans l'autre bloc $b^k$ : donc
$uv^iwx^iy$ est de la forme $a^m b^{k'} a^{m'} b^k$ avec $k'>k$ si
$i>1$, qui n'appartient pas à $Q\cap L_0$, de nouveau une
contradiction.  Les deux autres cas sont analogues.
\end{corrige}

\textbf{Remarque :} Les exercices \ref{non-square-words-is-algebraic}
et \ref{square-words-not-algebraic} mis ensemble donnent un exemple
explicite d'un langage $M$ algébrique dont le complémentaire $Q$ n'est
pas algébrique.

%

\exercice

On considère la grammaire hors contexte $G$ suivante (d'axiome $S$)
sur l'alphabet $\Sigma = \{a,b\}$ :
\[
S \rightarrow aSS \;|\; b
\]
Soit $L = L(G)$ le langage algébrique qu'elle engendre.

(0) Donner quelques exemples de mots dans $L$.

(1) Expliquer pourquoi un $w\in \Sigma^*$ appartient à $L$ si et
seulement si : soit $w = b$, soit il existe $u,v\in L$ tels que $w =
auv$.

(2) En déduire par récurrence sur la longueur $|w|$ de $w$ que si $w
\in L$ alors on a $wz \not\in L$ pour tout $z \in \Sigma^+$
(c'est-à-dire $z \in \Sigma^*$ et $z \neq \varepsilon$).  Autrement
dit : en ajoutant des lettres à la fin d'un mot de $L$ on obtient un
mot qui n'appartient jamais à $L$.  (Indication : si on a $auvz =
au'v'$ on pourra considérer les cas (i) $|u'|>|u|$, (ii) $|u'|<|u|$ et
(iii) $|u'|=|u|$.)

(3) En déduire que si $auv = au'v'$ avec $u,v,u',v'\in L$ alors $u=u'$
et $v=v'$.  (Indication analogue.)

(4) En déduire que $G$ est inambiguë, c'est-à-dire que chaque mot
$w\in L$ a un unique arbre d'analyse pour $G$ (on pourra reprendre
l'analyse de la question (1) et procéder de nouveau par récurrence
sur $|w|$).

(5) En s'inspirant des questions précédentes, décrire un algorithme
simple (en une seule fonction récursive) qui, donné un mot
$w\in\Sigma^*$ renvoie la longueur du préfixe de $w$ appartenant à $L$
s'il existe (il est alors unique d'après la question (2)) ou bien
« échec » s'il n'existe pas ; expliquer comment s'en servir pour
décider si $w\in L$ (i.e., écrire une fonction qui répond vrai ou faux
selon que $w\in L$ ou $w\not\in L$).

\begin{corrige}
(0) Quelques exemples de mots dans $L$ sont $b$, $abb$, $aabbb$,
  $ababb$ ou encore $aabbabb$.

(1) Si $w \in L$, considérons un arbre d'analyse $\mathscr{W}$ de $w$
  pour $G$ : sa racine, étiquetée $S$, doit avoir des fils étiquetés
  selon l'une des deux règles de la grammaire, soit $S\rightarrow b$
  ou bien $S\rightarrow aSS$, autrement dit, soit elle a un unique
  fils étiqueté $b$, soit elle a trois fils étiquetés respectivement
  $a,S,S$.  Dans le premier cas, le mot $w$ est simplement $b$ ; dans
  le second, les sous-arbres $\mathscr{U}, \mathscr{V}$ ayant pour
  racine les deux fils étiquetés $S$ sont encore des arbres d'analyse
  pour $G$, et si on appelle $u$ et $v$ les mots dont ils sont des
  arbres d'analyse (c'est-à-dire, ceux obtenus en lisant les feuilles
  de $\mathscr{U}$ et $\mathscr{V}$ respectivement par ordre de
  profondeur), alors on a $w = auv$ et $u,v\in L$ (puisqu'ils ont des
  arbres d'analyse pour $G$).

La réciproque est analogue : le mot $b$ appartient trivialement à $L$,
et si $u,v\in L$, ils ont des arbres d'analyse $\mathscr{U},
\mathscr{V}$ pour $G$, et on peut fabriquer un arbre d'analyse pour $w
:= auv$ qui a une racine étiquetée $S$ ayant trois fils étiquetés
$a,S,S$, ces deux derniers ayant pour descendants des sous-arbres
donnés par $\mathscr{U}$ et $\mathscr{V}$.

\smallbreak

(2) Montrons par récurrence sur $|w|$ que si $w \in L$ alors on a $wz
\not\in L$ pour tout $z \in \Sigma^+$.  La récurrence permet de
supposer la conclusion déjà connue pour tout mot de longueur $<|w|$.
D'après la question (1), le mot $w$ est soit $b$ soit de la forme
$auv$ avec $u,v\in L$ et trivialement $|u|<|w|$ et $|v|<|w|$.  Si
$w=b$, il est évident qu'aucun mot de la forme $bz$ ne peut appartenir
à $L$ (la question (1) montre que les seuls mots de $L$ sont le mot
$b$ et des mots commençant par $a$).  Il reste le cas $w = auv$ : on
veut montrer que $wz$, c'est-à-dire $auvz$, n'appartient pas à $L$.
Mais s'il y a appartenait, toujours d'après la question (1), il serait
de la forme $au'v'$ (le cas $b$ étant trivialement exclu), où $u',v'
\in L$ ; notamment, $uvz = u'v'$.  Distinguons trois cas : (i) soit
$|u'|>|u|$, mais alors $u$ est un préfixe strict de $u'$, c'est-à-dire
que $u'$ peut s'écrire sous la forme $u' = ut$ pour un $t\in
\Sigma^+$, et par l'hypothèse de récurrence, on a $u'\not\in L$, une
contradiction ; (ii) soit $|u'|<|u|$, mais alors $u'$ est un préfixe
strict de $u$, c'est-à-dire que $u$ peut s'écrire sous la forme $u =
u't$ pour un $t\in\Sigma^+$, et par l'hypothèse de récurrence (puisque
$|u'|<|u|<|w|$), on a $u\not\in L$, de nouveau une contradiction ;
(iii) soit $|u'|=|u|$, donc $u'=u$ (puisqu'ils sont préfixes de même
longueur du même mot $uvz = u'v'$), et on a alors $v' = vz$, mais
comme $v\in L$, l'hypothèse de récurrence entraîne $v'\not\in L$,
encore une contradiction.

\smallbreak

(3) Montrons que si $auv = au'v'$ avec $u,v,u',v'\in L$ alors $u=u'$
et $v=v'$.  On a notamment $uv = u'v'$.  Distinguons trois cas :
(i) soit $|u'|>|u|$, mais alors $u$ est un préfixe strict de $u'$,
c'est-à-dire que $u'$ peut s'écrire sous la forme $u' = ut$ pour un
$t\in \Sigma^+$, et par la question (2), on a $u'\not\in L$, une
contradiction ; (ii) soit $|u'|<|u|$, mais alors $u'$ est un préfixe
strict de $u$, c'est-à-dire que $u$ peut s'écrire sous la forme $u =
u't$ pour un $t\in\Sigma^+$, et par la question (2), on a $u\not\in
L$, de nouveau une contradiction ; (iii) soit $|u'|=|u|$, donc $u'=u$
(puisqu'ils sont préfixes de même longueur du même mot $uv = u'v'$),
et on a alors $v' = v$, la conclusion annoncéee.

\smallbreak

(4) Soit $w \in L$ : on veut montrer qu'il a un unique arbre d'analyse
pour $G$.  On procède par récurrence sur $|w|$, ce qui permet de
supposer la conclusion connue pour tout mot de longueur $<|w|$.  Comme
on l'a expliqué en (1), il y a deux possibilités pour un arbre
d'analyse de $w$ : soit la racine a un unique fils étiqueté $b$ et le
mot analysé est $w=b$, soit la racine a trois fils étiquetés $a,S,S$,
et des deux derniers fils partent des arbres d'analyse
$\mathscr{U},\mathscr{V}$ de mots $u,v\in L$ tels que $w=auv$.  Ces
deux cas sont évidemment incompatibles : il reste donc simplement à
expliquer que dans le dernier, $\mathscr{U}$ et $\mathscr{V}$ sont
uniquement déterminés.  Or la question (3) assure que $u,v$ (tels que
$w=auv$) sont uniquement déterminés, et l'hypothèse de récurrence
permet de conclure (comme $|u|<|w|$ et $|v|<|w|$) que les arbres
d'analyse $\mathscr{U}$ et $\mathscr{V}$ de $u$ et $v$ sont uniquement
déterminés, comme on le voulait.

\smallbreak

(5) Donné un mot $w\in\Sigma^*$, la fonction « rechercher préfixe
  dans $L$ » suivante renvoie la longueur du préfixe de $w$
appartenant à $L$, ou bien « échec » si ce préfixe n'existe pas :
\begin{itemize}
\item si $w=\varepsilon$, renvoyer échec,
\item si la première lettre de $w$ est $b$, renvoyer $1$,
\item sinon (la première lettre de $w$ est $a$), soit $x$ le suffixe
  de $w$ correspondant (c'est-à-dire $w = ax$, ou si on préfère, $x$
  enlève la première lettre de $w$),
\item appeler la fonction elle-même (rechercher préfixe dans $L$)
  sur $x$ :
\item si elle échoue, renvoyer échec,
\item si elle retourne $k$, soit $u$ le préfixe de $x$ de longueur
  $k$, et $y$ le suffixe correspondant (c'est-à-dire $x = uy$, ou si
  on préfère, $y$ enlève les $k$ premières lettres de $x$),
\item appeler la fonction elle-même (rechercher préfixe dans $L$)
  sur $y$ :
\item si elle échoue, renvoyer échec,
\item si elle retourne $\ell$, retourner $1+k+\ell$ (en effet, on a $w
  = auvz$ où $u$ est de longueur $k$ et $v$ est de longueur $\ell$).
\end{itemize}

Pour savoir si un mot appartient à $w$, il s'agit simplement de
vérifier que la valeur retournée (=la longueur du préfixe appartenant
à $L$) n'est pas un échec et est égale à la longueur $|w|$.

La terminaison de cet algorithme est claire par récurrence sur la
longueur (chaque appel récursif est fait sur un mot de longueur
strictement plus courte), et sa correction est garantie par les
questions précédentes : les cas $b$ et $auv$ sont disjoints et dans le
dernier cas, $u$ et $v$ sont uniquement déterminés (c'est ce
qu'affirme la non-ambiguïté de la grammaire).

(Il s'agit ici du cas le plus simple d'un analyseur LL, et
l'algorithme présenté ci-dessus est essentiellement un analyseur LL(1)
camouflé sous forme d'analyseur par descente récursive.)
\end{corrige}

%

\exercice

On considère la grammaire hors-contexte $G$ d'axiome $S$ et de
nonterminaux $N = \{S, T, U, V\}$ sur l'alphabet $\Sigma =
\{{\#},{@},{(},{)}, x, y, z\}$ (soit $7$ lettres) donnée par
\[
\begin{aligned}
S &\rightarrow T \;|\; S\mathbin{\#}T\\
T &\rightarrow U \;|\; U\mathbin{@}T\\
U &\rightarrow V \;|\; (\,S\,)\\
V &\rightarrow x \;|\; y \;|\; z\\
\end{aligned}
\]
(On pourra imaginer que $\#$ et $@$ sont deux opérations binaires, les
parenthèses servent comme on s'y attend, et $x,y,z$ sont trois opérandes
auxquelles on peut vouloir appliquer ces opérations.)

La grammaire en question est inambiguë (on ne demande pas de le
démontrer).

(1) Donner les arbres d'analyse (=de dérivation) des mots suivants :
(a) $x\mathbin{\#}y\mathbin{\#}z$, (b) $x\mathbin{\#}y\mathbin{@}z$,
(c) $x\mathbin{@}y\mathbin{\#}z$, (d) $x\mathbin{@}y\mathbin{@}z$,
(e) $(x\mathbin{\#}y)\mathbin{@}z$.

\begin{corrige}
On obtient les arbres d'analyse suivants :
(a) % S<S<S<T<U<V<x.>>>>#.T<U<V<y.>>>>#.T<U<V<z.>>>>
\begin{tikzpicture}[line join=bevel,baseline=(S0.base)]
\node (S0) at (53.333bp,0.000bp) [draw=none] {$S$};
\node (S1) at (20.000bp,-30.000bp) [draw=none] {$S$};  \draw (S0) -- (S1);
\node (S2) at (0.000bp,-60.000bp) [draw=none] {$S$};  \draw (S1) -- (S2);
\node (T0) at (0.000bp,-80.000bp) [draw=none] {$T$};  \draw (S2) -- (T0);
\node (U0) at (0.000bp,-100.000bp) [draw=none] {$U$};  \draw (T0) -- (U0);
\node (V0) at (0.000bp,-120.000bp) [draw=none] {$V$};  \draw (U0) -- (V0);
\node (x0) at (0.000bp,-140.000bp) [draw=none] {$x$};  \draw (V0) -- (x0);
\node (o0) at (20.000bp,-60.000bp) [draw=none] {$\#$};  \draw (S1) -- (o0);
\node (T1) at (40.000bp,-60.000bp) [draw=none] {$T$};  \draw (S1) -- (T1);
\node (U1) at (40.000bp,-80.000bp) [draw=none] {$U$};  \draw (T1) -- (U1);
\node (V1) at (40.000bp,-100.000bp) [draw=none] {$V$};  \draw (U1) -- (V1);
\node (y0) at (40.000bp,-120.000bp) [draw=none] {$y$};  \draw (V1) -- (y0);
\node (o1) at (60.000bp,-30.000bp) [draw=none] {$\#$};  \draw (S0) -- (o1);
\node (T2) at (80.000bp,-30.000bp) [draw=none] {$T$};  \draw (S0) -- (T2);
\node (U2) at (80.000bp,-50.000bp) [draw=none] {$U$};  \draw (T2) -- (U2);
\node (V2) at (80.000bp,-70.000bp) [draw=none] {$V$};  \draw (U2) -- (V2);
\node (z0) at (80.000bp,-90.000bp) [draw=none] {$z$};  \draw (V2) -- (z0);
\end{tikzpicture}
\hskip1cmplus5cmminus.5cm
(b) % S<S<T<U<V<x.>>>>#.T<U<V<y.>>@.T<U<V<z.>>>>>
\begin{tikzpicture}[line join=bevel,baseline=(S0.base)]
\node (S0) at (26.667bp,0.000bp) [draw=none] {$S$};
\node (S1) at (0.000bp,-30.000bp) [draw=none] {$S$};  \draw (S0) -- (S1);
\node (T0) at (0.000bp,-50.000bp) [draw=none] {$T$};  \draw (S1) -- (T0);
\node (U0) at (0.000bp,-70.000bp) [draw=none] {$U$};  \draw (T0) -- (U0);
\node (V0) at (0.000bp,-90.000bp) [draw=none] {$V$};  \draw (U0) -- (V0);
\node (x0) at (0.000bp,-110.000bp) [draw=none] {$x$};  \draw (V0) -- (x0);
\node (o0) at (20.000bp,-30.000bp) [draw=none] {$\#$};  \draw (S0) -- (o0);
\node (T1) at (60.000bp,-30.000bp) [draw=none] {$T$};  \draw (S0) -- (T1);
\node (U1) at (40.000bp,-60.000bp) [draw=none] {$U$};  \draw (T1) -- (U1);
\node (V1) at (40.000bp,-80.000bp) [draw=none] {$V$};  \draw (U1) -- (V1);
\node (y0) at (40.000bp,-100.000bp) [draw=none] {$y$};  \draw (V1) -- (y0);
\node (o1) at (60.000bp,-60.000bp) [draw=none] {$@$};  \draw (T1) -- (o1);
\node (T2) at (80.000bp,-60.000bp) [draw=none] {$T$};  \draw (T1) -- (T2);
\node (U2) at (80.000bp,-80.000bp) [draw=none] {$U$};  \draw (T2) -- (U2);
\node (V2) at (80.000bp,-100.000bp) [draw=none] {$V$};  \draw (U2) -- (V2);
\node (z0) at (80.000bp,-120.000bp) [draw=none] {$z$};  \draw (V2) -- (z0);
\end{tikzpicture}
\hskip1cmplus5cmminus.5cm
(c) % S<S<T<U<V<x.>>@.T<U<V<y.>>>>>#.T<U<V<z.>>>>
\begin{tikzpicture}[line join=bevel,baseline=(S0.base)]
\node (S0) at (53.333bp,0.000bp) [draw=none] {$S$};
\node (S1) at (20.000bp,-30.000bp) [draw=none] {$S$};  \draw (S0) -- (S1);
\node (T0) at (20.000bp,-50.000bp) [draw=none] {$T$};  \draw (S1) -- (T0);
\node (U0) at (0.000bp,-80.000bp) [draw=none] {$U$};  \draw (T0) -- (U0);
\node (V0) at (0.000bp,-100.000bp) [draw=none] {$V$};  \draw (U0) -- (V0);
\node (x0) at (0.000bp,-120.000bp) [draw=none] {$x$};  \draw (V0) -- (x0);
\node (o0) at (20.000bp,-80.000bp) [draw=none] {$@$};  \draw (T0) -- (o0);
\node (T1) at (40.000bp,-80.000bp) [draw=none] {$T$};  \draw (T0) -- (T1);
\node (U1) at (40.000bp,-100.000bp) [draw=none] {$U$};  \draw (T1) -- (U1);
\node (V1) at (40.000bp,-120.000bp) [draw=none] {$V$};  \draw (U1) -- (V1);
\node (y0) at (40.000bp,-140.000bp) [draw=none] {$y$};  \draw (V1) -- (y0);
\node (o1) at (60.000bp,-30.000bp) [draw=none] {$\#$};  \draw (S0) -- (o1);
\node (T2) at (80.000bp,-30.000bp) [draw=none] {$T$};  \draw (S0) -- (T2);
\node (U2) at (80.000bp,-50.000bp) [draw=none] {$U$};  \draw (T2) -- (U2);
\node (V2) at (80.000bp,-70.000bp) [draw=none] {$V$};  \draw (U2) -- (V2);
\node (z0) at (80.000bp,-90.000bp) [draw=none] {$z$};  \draw (V2) -- (z0);
\end{tikzpicture}
\hskip1cmplus5cmminus.5cm
(d) % S<T<U<V<x.>>@.T<U<V<y.>>@.T<U<V<z.>>>>>>
\begin{tikzpicture}[line join=bevel,baseline=(S0.base)]
\node (S0) at (26.667bp,0.000bp) [draw=none] {$S$};
\node (T0) at (26.667bp,-20.000bp) [draw=none] {$T$};  \draw (S0) -- (T0);
\node (U0) at (0.000bp,-50.000bp) [draw=none] {$U$};  \draw (T0) -- (U0);
\node (V0) at (0.000bp,-70.000bp) [draw=none] {$V$};  \draw (U0) -- (V0);
\node (x0) at (0.000bp,-90.000bp) [draw=none] {$x$};  \draw (V0) -- (x0);
\node (o0) at (20.000bp,-50.000bp) [draw=none] {$@$};  \draw (T0) -- (o0);
\node (T1) at (60.000bp,-50.000bp) [draw=none] {$T$};  \draw (T0) -- (T1);
\node (U1) at (40.000bp,-80.000bp) [draw=none] {$U$};  \draw (T1) -- (U1);
\node (V1) at (40.000bp,-100.000bp) [draw=none] {$V$};  \draw (U1) -- (V1);
\node (y0) at (40.000bp,-120.000bp) [draw=none] {$y$};  \draw (V1) -- (y0);
\node (o1) at (60.000bp,-80.000bp) [draw=none] {$@$};  \draw (T1) -- (o1);
\node (T2) at (80.000bp,-80.000bp) [draw=none] {$T$};  \draw (T1) -- (T2);
\node (U2) at (80.000bp,-100.000bp) [draw=none] {$U$};  \draw (T2) -- (U2);
\node (V2) at (80.000bp,-120.000bp) [draw=none] {$V$};  \draw (U2) -- (V2);
\node (z0) at (80.000bp,-140.000bp) [draw=none] {$z$};  \draw (V2) -- (z0);
\end{tikzpicture}
\hskip1cmplus5cmminus.5cm
(e) % S<T<U<(.S<S<T<U<V<x.>>>>#.T<U<V<y.>>>>).>@.T<U<V<z.>>>>>
\begin{tikzpicture}[line join=bevel,baseline=(S0.base)]
\node (S0) at (86.667bp,0.000bp) [draw=none] {$S$};
\node (T0) at (86.667bp,-20.000bp) [draw=none] {$T$};  \draw (S0) -- (T0);
\node (U0) at (40.000bp,-50.000bp) [draw=none] {$U$};  \draw (T0) -- (U0);
\node (p0) at (0.000bp,-80.000bp) [draw=none] {$($};  \draw (U0) -- (p0);
\node (S1) at (40.000bp,-80.000bp) [draw=none] {$S$};  \draw (U0) -- (S1);
\node (S2) at (20.000bp,-110.000bp) [draw=none] {$S$};  \draw (S1) -- (S2);
\node (T1) at (20.000bp,-130.000bp) [draw=none] {$T$};  \draw (S2) -- (T1);
\node (U1) at (20.000bp,-150.000bp) [draw=none] {$U$};  \draw (T1) -- (U1);
\node (V0) at (20.000bp,-170.000bp) [draw=none] {$V$};  \draw (U1) -- (V0);
\node (x0) at (20.000bp,-190.000bp) [draw=none] {$x$};  \draw (V0) -- (x0);
\node (o0) at (40.000bp,-110.000bp) [draw=none] {$\#$};  \draw (S1) -- (o0);
\node (T2) at (60.000bp,-110.000bp) [draw=none] {$T$};  \draw (S1) -- (T2);
\node (U2) at (60.000bp,-130.000bp) [draw=none] {$U$};  \draw (T2) -- (U2);
\node (V1) at (60.000bp,-150.000bp) [draw=none] {$V$};  \draw (U2) -- (V1);
\node (y0) at (60.000bp,-170.000bp) [draw=none] {$y$};  \draw (V1) -- (y0);
\node (p1) at (80.000bp,-80.000bp) [draw=none] {$)$};  \draw (U0) -- (p1);
\node (o1) at (100.000bp,-50.000bp) [draw=none] {$@$};  \draw (T0) -- (o1);
\node (T3) at (120.000bp,-50.000bp) [draw=none] {$T$};  \draw (T0) -- (T3);
\node (U3) at (120.000bp,-70.000bp) [draw=none] {$U$};  \draw (T3) -- (U3);
\node (V2) at (120.000bp,-90.000bp) [draw=none] {$V$};  \draw (U3) -- (V2);
\node (z0) at (120.000bp,-110.000bp) [draw=none] {$z$};  \draw (V2) -- (z0);
\end{tikzpicture}
\end{corrige}

(2) En considérant que $(w)$ a le même effet ou la même valeur
que $w$ : (a) L'une des deux opérations $\#$ et $@$ est-elle
prioritaire\footnote{On dit qu'une opération binaire $\boxtimes$ est
  \emph{prioritaire} sur $\boxplus$ lorsque « $u\boxtimes v\boxplus
  w$ » se comprend comme « $(u\boxtimes v)\boxplus w$ » et
  « $u\boxplus v\boxtimes w$ » comme « $u\boxplus (v\boxtimes w)$ ».}
sur l'autre ?  Laquelle ?\quad (b) L'opération $\#$
s'associe-t-elle\footnote{On dit qu'une opération binaire $\boxdot$
  \emph{s'associe à gauche} lorsque « $u\boxdot v\boxdot w$ » se
  comprend comme « $(u\boxdot v)\boxdot w$ », et \emph{s'associe à
    droite} lorsque « $u\boxdot v\boxdot w$ » se comprend comme
  « $u\boxdot (v\boxdot w)$ ».} à gauche ou à droite ?\quad
(c) L'opération $@$ s'associe-t-elle à gauche ou à droite ?

\begin{corrige}
(a) Les arbres d'analyse (b) et (c) de la question (1) montrent que
  l'opération $@$ est prioritaire sur $\#$ (elle est toujours plus
  profonde dans l'arbre, c'est-à-dire appliquée en premier aux
  opérandes).\quad (b) L'arbre d'analyse (a) de la question (1) montre
  que $\#$ s'associe à gauche (le $\#$ entre les deux opérandes gauche
  est plus profond, c'est-à-dire appliqué en premier).\quad
  (c) Symétriquement, l'arbre d'analyse (d) de la question (1) montre
  que $@$ s'associe à droite.
\end{corrige}

\smallskip

(La question (3) est indépendante de (1) et (2).  Par ailleurs, on
pourra, si on souhaite s'épargner des confusions, choisir d'appeler
$a$ la lettre « parenthèse ouvrante » de $\Sigma$ et $b$ la lettre
« parenthèse fermante » afin de les distinguer des parenthèses
mathématiques.)

(3) Soit $L = L(G)$ le langage engendré par $G$.  Soit $M$ le langage
des mots sur $\Sigma$ formés d'un certain nombre de parenthèses
ouvrantes, puis de la lettre $x$, puis d'un certain nombre
(possiblement différent) de parenthèses fermantes : autrement dit, des
mots de la forme « ${(}^i\, x\, {)}^j$ », où on a noté « ${(}^i$ » une
succession de $i$ parenthèses ouvrantes et « ${)}^j$ » une succession
de $j$ parenthèses fermantes (on pourra noter ce mot $a^i x b^j$ si on
préfère).\quad (a) Décrire le langage $L\cap M$ (des mots de $L$ qui
appartiennent aussi à $M$).\quad (b) Ce langage $L\cap M$ est-il
rationnel ?\quad (c) Le langage $L$ est-il rationnel ?

\begin{corrige}
(a) Cherchons pour quelles valeurs $(i,j) \in \mathbb{N}^2$ le mot
  « ${(}^i\, x\, {)}^j$ » de $M$ appartient à $L$.  La dérivation $S
  \Rightarrow T \Rightarrow U \Rightarrow (S) \Rightarrow (T)
  \Rightarrow (U) \Rightarrow ((S)) \Rightarrow \cdots
  \Rightarrow{(}^i\, S\, {)}^i \Rightarrow{(}^i\, T\, {)}^i
  \Rightarrow{(}^i\, U\, {)}^i \Rightarrow{(}^i\, V\, {)}^i
  \Rightarrow{(}^i\, x\, {)}^i$ montre que c'est le cas si $j=i$,
  autrement dit s'il y a autant de parenthèses fermantes qu'ouvrantes.
  Mais réciproquement, la propriété d'avoir autant de parenthèses
  fermantes qu'ouvrantes est préservée par toutes les productions
  de $G$ (et est satisfaite par l'axiome), donc est vérifiée par tout
  mot de $L=L(G)$.  On a donc prouvé que le $L\cap M$ est l'ensemble
  des mots de la forme « ${(}^i\, x\, {)}^i$ » où $i\in \mathbb{N}$.

(b) Le langage $L\cap M$ constitué des mots de la forme « ${(}^n\,
  x\, {)}^n$ » où $i\in \mathbb{N}$ n'est pas rationnel.  Cela résulte
  du lemme de pompage (presque exactement comme l'exemple du langage
  $\{a^n b^n : i\in\mathbb{N}\}$ donné en cours) : donnons l'argument
  en remplaçant la parenthèse ouvrante par $a$ et la parenthèse
  fermante par $b$ pour plus de clarté notationnelle.  Appliquons le
  lemme de pompage pour les langages rationnels au langage $L \cap M =
  \{a^n x b^n : i\in\mathbb{N}\}$ supposé par l'absurde être
  rationnel : appelons $k$ l'entier dont le lemme de pompage garantit
  l'existence.  Considérons le mot $t := a^k x b^k$ : il doit alors
  exister une factorisation $t = uvw$ pour laquelle on a (i) $|v|\geq
  1$, (ii) $|uv|\leq k$ et (iii) $uv^iw \in L\cap M$ pour tout $i\geq
  0$.  La propriété (ii) assure que $uv$ est formé d'un certain nombre
  de répétitions de la lettre $a$ (car tout préfixe de longueur $\leq
  k$ de $a^k x b^k$ est de cette forme) ; disons $u = a^\ell$ et $v =
  a^m$, si bien que $w = a^{k-\ell-m} x b^k$.  La propriété (i)
  donne $m\geq 1$.  Enfin, la propriété (iii) affirme que le mot
  $uv^iw = a^{k+(i-1)m} x b^k$ appartient à $L\cap M$ ; mais dès que
  $i\neq 1$, ceci est faux : il suffit donc de prendre $i=0$ pour
  avoir une contradiction.

(c) Le langage $M$ est rationnel puisqu'il est dénoté par l'expression
  rationnelle $a{*}xb{*}$ (toujours en notant $a$ la parenthèse
  ouvrante et $b$ la parenthèse fermante).  Si le langage $L$ était
  rationnel, le langage $L\cap M$ le serait aussi (on sait que
  l'intersection de deux langages rationnels est rationnelle), ce qui
  n'est pas le cas d'après la question (b).  Le langage $L$ n'est donc
  pas rationnel.
\end{corrige}


%

\exercice

On considère la grammaire hors contexte $G$ suivante (d'axiome $S$)
sur l'alphabet $\Sigma = \{a,b\}$ :
\[
S \rightarrow SaS \;|\; b
\]
Soit $L = L(G)$ le langage algébrique qu'elle engendre.

(0) Donner quelques exemples de mots dans $L$.

(1) Cette grammaire $G$ est-elle ambiguë ?

(2) Décrire simplement le langage $L$.

(3) Donner une grammaire inambiguë engendrant le même langage $L$
que $G$.

\begin{corrige}
(0) Quelques exemples de mots dans $L$ sont $b$, $bab$, $babab$, et
  ainsi de suite (cf. (2)).

(1) La grammaire $G$ est ambiguë : le mot $babab$ a deux arbres
  d'analyse différents, à savoir
\begin{center}
\tikzstyle{automaton}=[scale=0.5]
%%% begin ex3pt1 %%%

\begin{tikzpicture}[>=latex,line join=bevel,automaton]
%%
\node (S3) at (27bp,90bp) [draw,draw=none] {$S$};
  \node (S2) at (243bp,162bp) [draw,draw=none] {$S$};
  \node (S1) at (99bp,162bp) [draw,draw=none] {$S$};
  \node (S0) at (171bp,234bp) [draw,draw=none] {$S$};
  \node (S4) at (171bp,90bp) [draw,draw=none] {$S$};
  \node (a1) at (99bp,90bp) [draw,draw=none] {$a$};
  \node (a0) at (171bp,162bp) [draw,draw=none] {$a$};
  \node (b0) at (243bp,90bp) [draw,draw=none] {$b$};
  \node (b1) at (27bp,18bp) [draw,draw=none] {$b$};
  \node (b2) at (171bp,18bp) [draw,draw=none] {$b$};
  \draw [] (S1) ..controls (70.042bp,132.85bp) and (55.714bp,118.92bp)  .. (S3);
  \draw [] (S0) ..controls (142.04bp,204.85bp) and (127.71bp,190.92bp)  .. (S1);
  \draw [] (S0) ..controls (171bp,204.85bp) and (171bp,190.92bp)  .. (a0);
  \draw [] (S2) ..controls (243bp,132.85bp) and (243bp,118.92bp)  .. (b0);
  \draw [] (S1) ..controls (99bp,132.85bp) and (99bp,118.92bp)  .. (a1);
  \draw [] (S0) ..controls (199.96bp,204.85bp) and (214.29bp,190.92bp)  .. (S2);
  \draw [] (S3) ..controls (27bp,60.846bp) and (27bp,46.917bp)  .. (b1);
  \draw [] (S4) ..controls (171bp,60.846bp) and (171bp,46.917bp)  .. (b2);
  \draw [] (S1) ..controls (127.96bp,132.85bp) and (142.29bp,118.92bp)  .. (S4);
%
\end{tikzpicture}

%%% end ex3pt1 %%%
\end{center}
et son symétrique par rapport à l'axe vertical.

\smallbreak

(2) Montrons que $L$ est égal au langage $L_{(ba){*}b}$ dénoté par
l'expression rationnelle $(ba){*}b$ comme la question (0) le laisse
entrevoir.  Pour cela, il suffit de voir que l'ensemble des
pseudo-mots (:= mots sur l'ensemble des terminaux et des nonterminaux)
qu'on peut dériver de $S$ dans $G$ est le langage $((S|b)a){*}(S|b)$
constitué des pseudo-mots de la forme $xaxax\cdots ax$ où chaque $x$
est soit un $S$ soit un $b$, indépendamment les uns des autres.  Or il
est clair que remplacer un $x$ (et notamment, remplacer un $S$) par
$SaS$ (qui est de la forme $xax$) dans un pseudo-mot de cette forme
donne encore un pseudo-mot de cette forme : ceci montre qu'une
dérivation de $G$, quelle que soit l'application faite de la règle
$S\rightarrow SaS$, donne des pseudo-mots de cette forme, donc tout
pseudo-mot obtenu par dérivation de $S$ dans la grammaire $G$ est de
la forme qu'on vient de dire ; et réciproquement, il est facile de
produire n'importe quel nombre $n$ de répéritions du motif $aS$ en
appliquant $n$ fois la règle $S\rightarrow SaS$ à partir de $S$, et on
peut ensuite librement transformer certains $S$ en $b$.

\smallbreak

(3) La grammaire $G'$ donnée par $S \rightarrow baS\,|\,b$ engendre le
même langage que $G$ d'après la question précédente ; et il est
évident que cette grammaire est inambiguë : toutes les dérivations
sont de la forme $S \Rightarrow baS \Rightarrow babaS \Rightarrow
(ba)^n S$ suivie éventuellement de $\Rightarrow (ba)^n b$.
\end{corrige}

%
%
%

\subsection{Introduction à la calculabilité}

\exercice

On rappelle la définition du problème de l'arrêt : c'est l'ensemble
$H$ des couples\footnote{Pour être rigoureux, on a fixé un codage
  permettant de représenter les programmes $e$, les entrées $x$ à ces
  programmes, et les couples $(e,x)$, comme des éléments
  de $\mathbb{N}$ (ou bien de $\Sigma^*$ sur un alphabet
  $\Sigma\neq\varnothing$ arbitraire).  Il n'est pas nécessaire de
  faire apparaître ce codage dans la description des algorithmes
  proposés, qui peut rester informelle.} $(e,x)$ formés d'un programme
(=algorithme) $e$ et d'une entrée $x$, tels que l'exécution du
programme $e$ sur l'entrée $x$ termine en temps fini.  On a vu en
cours que $H$ était semi-décidable mais non décidable\footnote{Une
  variante consiste à considérer l'ensemble des $e$ tels que $e$
  termine sur l'entrée $e$ : l'ensemble $H$ qu'on vient de définir s'y
  ramène facilement.}.

On considère l'ensemble $H'$ des triplets $(e_1,e_2,x)$ tels que
$(e_1,x) \in H$ \emph{ou bien} $(e_2,x) \in H$ (ou les deux).
Autrement dit, au moins l'un des programmes $e_i$ termine en temps
fini quand on l'exécute sur l'entrée $x$.

(1) L'ensemble $H'$ est-il semi-décidable ?

\begin{corrige}
Nous allons montrer que $H'$ est semi-décidable.

Considérons l'algorithme suivant : donné en entrée un triplet
$(e_1,e_2,x)$, on exécute en parallèle, en fournissant à chacun
l'entrée $x$, les programmes (codés par) $e_1$ et $e_2$ (au moyen
d'une machine universelle), et en s'arrêtant immédiatement si l'un des
deux s'arrête, et on renvoie alors « vrai ».  Manifestement, cet
algorithme termine (en renvoyant « vrai ») si et seulement si
$(e_1,e_2,x) \in H'$, ce qui montre que $H'$ est semi-décidable.
\end{corrige}

(2) L'ensemble $H'$ est-il décidable ?

\begin{corrige}
Nous allons montrer que $H'$ n'est pas décidable.

Supposons par l'absurde qu'il existe un algorithme $T$ qui
décide $H'$.  Soit $e'$ (le code d')un programme quelconque qui
effectue une boucle infinie (quelle que soit son entrée) : comme $e'$
ne termine jamais, un triplet $(e,e',x)$ appartient à $H'$ si et
seulement si $(e,x)$ appartient à $H$.  Considérons maintenant
l'algorithme suivant : donné en entrée un couple $(e,x)$, on applique
l'algorithme $T$ supposé exister pour décider si $(e,e',x) \in H'$ ;
comme on vient de l'expliquer, ceci équivaut à $(e,x) \in H$.  On a
donc fabriqué un algorithme qui décide $H$, ce qui est impossible,
d'où la contradiction annoncée.
\end{corrige}

%

\exercice

Soit $\Sigma$ un alphabet (fini, non vide) fixé.  Les questions
suivantes sont indépendantes (mais on remarquera leur parallélisme).
Ne pas hésiter à décrire les algorithmes de façon succincte et
informelle.

(1) Expliquer, au moyen des résultats vus en cours, pourquoi il existe
un algorithme $A_1$ qui, étant donnée une expression rationnelle $r$
sur $\Sigma$, décide si le langage $L_r$ dénoté par $r$ est différent
du langage $\Sigma^*$ de tous les mots sur $\Sigma$.  (Autrement dit,
l'algorithme $A_1$ doit prendre en entrée une expression rationnelle
$r$, terminer en temps fini, et répondre « vrai » s'il existe un mot
$w\in\Sigma^*$ tel que $w\not\in L_r$ et « faux » si $L_r = \Sigma^*$.
On ne demande pas que l'algorithme soit efficace.)

(2) Expliquer pourquoi il existe un algorithme $A_2$ qui, étant donnée
une grammaire hors contexte $G$ sur $\Sigma$, « semi-décide » si le
langage $L_G$ engendré par $G$ est différent du langage $\Sigma^*$ de
tous les mots.  (« Semi-décider » signifie que l'algorithme $A_2$ doit
prendre en entrée une grammaire hors contexte $G$, terminer en temps
fini en répondant « vrai » s'il existe un mot $w\in\Sigma^*$ tel que
$w\not\in L_G$, et ne pas terminer\footnote{On peut admettre qu'il
  termine parfois en répondant « faux », mais ce ne sera pas utile.}
si $L_G = \Sigma^*$.)  Indication : on peut tester tous les mots
possibles.

(3) Expliquer pourquoi il existe un algorithme $A_3$ comme suit : on
lui fournit en entrée un algorithme $T$ qui décide un langage $L_T
\subseteq \Sigma^*$ (c'est-à-dire que $T$ termine toujours en temps
fini quand on lui présente un mot sur $\Sigma$, et répond « vrai » ou
« faux », et $L_T$ est le langage des mots sur lesquels il répond
« vrai »), et l'algorithme $A_3$ doit semi-décider si $L_T$ est
différent de $\Sigma^*$.  (C'est-à-dire que $A_3$ doit terminer en
répondant « vrai » s'il existe un mot $w\in\Sigma^*$ tel que $w\not\in
L_T$, et ne pas terminer si $L_T = \Sigma^*$.)  Indication : la même
approche permet de traiter les questions (2) et (3).

(4) Expliquer pourquoi il \emph{n'existe pas} d'algorithme $A_4$ qui,
dans les mêmes conditions que $A_3$, décide (au lieu de seulement
semi-décider) si $L_T$ est différent de $\Sigma^*$.  (C'est-à-dire que
$A_4$ est censé terminer toujours, et répondre « vrai » s'il existe un
mot $w\in\Sigma^*$ tel que $w\not\in L_T$, et « faux » si $L_T =
\Sigma^*$.)  Indication : expliquer comment on pourrait utiliser un
tel $A_4$ pour résoudre le problème de l'arrêt, en cherchant à
fabriquer un $T$ qui rejette un mot précisément si un programme donné
s'arrête.

\begin{corrige}
(1) Donnée une expression rationnelle $r$, on sait qu'on peut
  algorithmiquement fabriquer un automate fini non-déterministe à
  transitions spontanées qui reconnaît exactement le langage $L_r$
  dénoté par $r$, et ensuite éliminer les transitions spontanées et
  déterminiser l'automate pour obtenir un automate fini déterministe
  complet reconnaissant $L_r$.  Sur un tel automate, savoir si $L_r
  \neq \Sigma^*$ est trivial : dès lors qu'il existe un état $q$
  non-final accessible, il existe un mot rejeté par l'automate (i.e.,
  n'appartenant pas à $L_r$), à savoir le mot lu en suivant les
  étiquettes d'un chemin quelconque de l'état initial $q_0$
  jusqu'à $q$, et inversement, si tous les états sont finaux, il est
  trivial que l'automate accepte tous les mots.  (On pouvait aussi
  minimiser l'automate et le comparer à l'automate minimal trivial qui
  reconnaît le langage $\Sigma^*$.)

\smallbreak

(2) On sait qu'il existe un algorithme qui, donnée une grammaire
hors-contexte $G$ et un mot $w$, décide si $w \in L_G$.  Pour
semi-décider s'il existe un $w$ tel que $w \not\in L_G$, il suffit de
tester tous les mots possibles : plus exactement, on construit un
algorithme $A_2$ qui effectue une boucle infinie sur tous les
$w\in\Sigma^*$ (il est évidemment algorithmiquement faisable
d'énumérer tous les mots sur $\Sigma$) et, pour chacun, teste si $w
\in L_G$, et si ce n'est pas le cas, termine immédiatement en
répondant « vrai » (on a trouvé un $w$ n'appartenant pas à $L_G$),
tandis que si c'est le cas, l'algorithme $A_2$ ne terminera jamais.

\smallbreak

(3) On procède exactement comme en (2) : par hypothèse on dispose d'un
algorithme $T$ qui, donné un mot $w$, décide si $w \in L_T$.  Pour
semi-décider s'il existe un $w$ tel que $w \not\in L_T$, il suffit de
tester tous les mots possibles : plus exactement, on construit un
algorithme $A_3$ qui effectue une boucle infinie sur tous les
$w\in\Sigma^*$ et, pour chacun, teste si $w \in L_T$ (en lançant
l'algorithme $T$ qui, par hypothèse, termine toujours), et si ce n'est
pas le cas, termine immédiatement en répondant « vrai » (on a trouvé
un $w$ n'appartenant pas à $L_T$), tandis que si c'est le cas,
l'algorithme $A_3$ ne terminera jamais.

\smallbreak

(4) Supposons par l'absurde qu'on dispose d'un algorithme $A_4$ comme
on vient de dire, et montrons pour arriver à une contradiction qu'on
peut s'en servir pour résoudre le problème de l'arrêt.  On se donne
donc un algorithme $S$ et une entrée $x$ de $S$ et on cherche à savoir
(en utilisant $A_4$) si $S$ termine sur l'entrée $x$.  Pour cela, on
va construire un $T$ auquel appliquer $A_4$.

Voici une solution possible : donné un mot $w \in \Sigma^*$, le
programme $T$ ne considère que la longueur $|w|$ de $w$, et lance
(=simule) l'exécution de $S$ sur l'entrée $x$ pour au plus $|w|$
étapes : si l'exécution termine dans le temps imparti, alors $T$
rejette le mot $w$, sinon, il l'accepte (dans tous les cas, $T$
termine et répond « vrai » ou « faux », donc il est une entrée
légitime à $A_4$).  Cette construction fait que $L_T$ rejette au moins
un mot précisément lorsque $S$ termine sur $x$ : si au contraire $S$
ne termine pas sur $x$, alors $L_T = \Sigma^*$.  L'utilisation de
$A_4$ sur $T$ permet donc de savoir algorithmiquement si $S$ termine
sur $x$, ce qui contredit l'indécidabilité du problème de l'arrêt.

\textit{Variante de la même idée :} on appelle « trace d'exécution » de $S$ sur
$x$ un mot $w$ qui code le calcul complet de l'exécution de $S$ sur
$x$ (par exemple, si on voit $S$ comme une machine de Turing, l'état
courant et le contenu du ruban à chaque étape), du début à l'arrêt.
Une telle trace d'exécution existe donc précisément si $S$ termine
sur $x$.  Or il est visiblement décidable de savoir si un mot $w$
donné est une trace d'exécution (il suffit de vérifier qu'à chaque
étape la machine a bien fait ce qu'elle devait faire).  On peut donc
écrire un algorithme $T$ qui termine toujours et accepte précisément
les mots qui \emph{ne sont pas} une trace d'exécution de $S$ sur $x$.
Le fait que $L_T$ soit différent de $\Sigma^*$ signifie alors
exactement qu'une trace d'exécution existe, donc que $S$ termine
sur $x$.  Ainsi l'utilisation de $A_4$ permet de savoir
algorithmiquement si $S$ termine sur $x$, ce qui contredit
l'indécidabilité du problème de l'arrêt.
\end{corrige}

%

\exercice

(1) Soit $\Sigma$ un alphabet (i.e., un ensemble fini).  L'ensemble $L
= \{u^k : u\in\Sigma^*, k\geq 2\}$ des mots qui sont une puissance
$k$-ième pour un $k\geq 2$ est-il décidable ?  Semi-décidable ?

(2) L'ensemble des $e \in \mathbb{N}$ tels que l'exécution du $e$-ième
programme (ou, si on préfère, de la $e$-ième machine de Turing),
exécuté sur l'entrée $42$, termine en au plus $10^{1000+e}$ étapes
est-il décidable ?  Semi-décidable ?

(3) L'ensemble des $e \in \mathbb{N}$ tels que l'exécution du $e$-ième
programme (ou, si on préfère, de la $e$-ième machine de Turing),
exécuté sur l'entrée $42$, termine en temps fini est-il décidable ?
Semi-décidable ?  (On pourra montrer qu'on peut y ramener le problème
de l'arrêt.)

\begin{corrige}
(1) Si $w = u^k$ pour un certain $u\neq\varepsilon$, alors
  nécessairement $k\leq |w|$ puisque $|w|=k\cdot|u|$.  On dispose donc
  de l'algorithme suivant pour décider si $w\in L$ : si
  $w=\varepsilon$, retourner vrai immédiatement ; sinon, pour $k$
  allant de $2$ à $|w|$ et qui divise $|w|$, considérer les $k$
  facteurs successifs de $w$ de longueur $|w|/k$ (c'est-à-dire, pour
  $0\leq i<k$, le facteur de $w$ de longueur $\ell := |w|/k$
  commençant à la position $i \ell$) : s'ils sont tous égaux, renvoyer
  vrai ; si la boucle termine sans avoir trouvé de $k$ qui convienne,
  renvoyer faux.  Le langage proposé est donc décidable (et \textit{a
    fortiori} semi-décidable).

\smallbreak

(2) Donné un $e \in \mathbb{N}$, la fonction $10^{1000+e}$ est
  évidemment calculable.  On peut ensuite lancer l'exécution du
  $e$-ième programme, sur l'entrée $42$, pour au plus ce nombre
  d'étapes (en utilisant la machine universelle, c'est-à-dire, par
  exemple, en simulant la $e$-ième machine de Turing sur une machine
  de Turing).  Si l'exécution termine en le temps imparti, on renvoie
  vrai, sinon, on renvoie faux : ceci montre que l'ensemble proposé
  est bien décidable (et \textit{a fortiori} semi-décidable).

\smallbreak

(3) L'ensemble $A$ proposé est « presque » le problème de l'arrêt.  La
  différence est que le problème de l'arrêt est l'ensemble des couples
  $(e,n)$ tels que le $e$-ième programme termine sur l'entrée $n$
  alors qu'ici on a fixé l'entrée à $42$.  Il s'agit donc de montrer
  que cette limitation ne rend pas pour autant calculable l'ensemble
  considéré.  Or donnés deux entiers $(e,n)$, on peut fabriquer un
  programme $e'$ qui prend en entrée une valeur, \emph{ignore} cette
  valeur, et exécute le $e$-ième programme sur l'entrée $n$ ; de plus
  un tel $e'$ se calcule algorithmiquement\footnote{Techniquement, on
    invoque ici le théorème s-m-n ; mais dans les faits, il s'agit
    essentiellement d'ajouter une ligne « \texttt{let
      n=}(représentation décimale de $n$) » au début d'un programme,
    ce qui est certainement faisable.} à partir de $e$ et $n$.
  L'exécution du programme $e'$ sur l'entrée $42$ (ou n'importe quelle
  autre entrée) se comporte donc comme l'exécution du programme $e$
  sur l'entrée $n$, et notamment, termine si et seulement si elle
  termine.  Autrement dit, $e'$ appartient à l'ensemble $A$ considéré
  dans cette question si et seulement si $(e,n)$ appartient au
  problème de l'arrêt.  Comme on vient de dire qu'on peut calculer
  $e'$ algorithmiquement à partir de $(e,n)$, si l'ensemble $A$ était
  décidable, le problème de l'arrêt le serait, ce qui n'est pas le
  cas.  Donc $A$ n'est pas décidable.  En revanche, $A$ est
  semi-décidable : il suffit de lancer l'exécution du programme $e$
  sur l'entrée $42$ et renvoyer vrai si elle termine (si elle ne
  termine pas, on ne termine pas).
\end{corrige}

%

\exercice

Soit $\Sigma$ un alphabet (i.e., un ensemble fini).  On s'intéresse à
des langages sur $\Sigma$.

(A) Montrer que si deux langages $L_1$ et $L_2$ sont décidables, alors
$L_1\cup L_2$ et $L_1\cap L_2$ et $L_1 L_2$ sont décidables ; montrer
que si un langage $L$ est décidable alors $L^*$ est décidable (pour ce
dernier, on pourra commencer par chercher, si $w \in \Sigma^*$ est un
mot de longueur $n$, comment énumérer toutes les façons de le
factoriser en mots de longueur non nulle).

(B) Montrer que si deux langages $L_1$ et $L_2$ sont semi-décidables,
alors $L_1\cup L_2$ et $L_1\cap L_2$ et $L_1 L_2$ sont
semi-décidables ; montrer que si un langage $L$ est semi-décidable
alors $L^*$ est semi-décidable.

\begin{corrige}
(A) Supposons qu'on dispose d'algorithmes $T_1$ et $T_2$ qui décident
  $L_1$ et $L_2$ respectivement (i.e., donné $w \in \Sigma^*$,
  l'algorithme $T_i$ termine toujours en temps fini, en répondant oui
  si $w\in L_i$ et non si $w\not\in L_i$).

Pour faire un algorithme qui décide $L_1\cup L_2$, donné un mot
$w\in\Sigma^*$, il suffit de lancer successivement $T_1$ et $T_2$ : si
l'un des deux répond oui, on répond oui, sinon on répond non
(autrement dit, on calcule les valeurs de vérité de $w\in L_1$ et
$w\in L_2$ au moyen de $T_1$ et $T_2$, et on calcule ensuite leur
« ou » logique).  De même pour décider $L_1\cap L_2$, il suffit de
lancer successivement $T_1$ et $T_2$, si les deux répondent oui on
répond oui, sinon on répond non (i.e., on calcule les valeurs de
vérité de $w\in L_1$ et $w\in L_2$ au moyen de $T_1$ et $T_2$, et on
calcule ensuite leur « et » logique).

Pour décider $L_1 L_2$, on effectue une boucle sur toutes les
factorisation $w = uv$ de $w$, c'est-à-dire, une boucle sur toutes les
longueurs $0\leq i\leq |w|$ en appelant à chaque fois $u$ le préfixe
de $w$ de longueur $i$ et $v$ le suffixe de $w$ de longueur $|w|-i$,
et pour chaque paire $(u,v)$ ainsi trouvée, on utilise $T_1$ et $T_2$
pour tester si $u\in L_1$ et $v\in L_2$ : si c'est le cas, on termine
l'algorithme en répondant oui (on a $w = uv \in L_1 L_2$) ; si aucune
paire ne convient, on répond non.

L'algorithme pour décider $L^*$ est semblable : il s'agit de tester
toutes les manières de factoriser un mot $w \in \Sigma^*$ en facteurs
de longueur non nulle.  (On peut d'ores et déjà exclure
$w=\varepsilon$ car le mot vide appartient de toute façon à $L^*$.)
Si $n=|w| > 0$, on peut effectuer une boucle pour un nombre de
facteurs $k$ allant de $1$ à $n$, et, pour chaque $k$, effectuer $k$
boucles emboîtées pour déterminer les limites des facteurs
$u_1,\ldots,u_k \in \Sigma^+$ tels que $w = u_1\cdots u_k$ (il suffit
par exemple de faire boucler $i_1,\ldots,i_k$ chacun de $1$ à $n$, et
lorsque $i_1+\cdots+i_k = n$, appeler $u_j$ le facteur de $w$ de
longueur $i_j$ commençant à la position $i_1+\cdots+i_{j-1}$).  Pour
chaque factorisation comme on vient de le dire, on teste si tous les
$u_i$ appartiennent à $L$, et si c'est le cas on renvoie vrai (le mot
$w$ appartient à $L^*$) ; si aucune factorisation ne convient, on
renvoie faux.

(Dans l'algorithme qui précède, on a écarté les factorisations faisant
intervenir le mot vide, car si $w$ est factorisable en mots de $L$ en
faisant intervenir le mot vide, quitte à retirer celui-ci, il est
encore factorisable en mots non vides de $L$.)

\smallbreak

(B) Les algorithmes sont très semblables à ceux de la partie (A) si ce
n'est qu'il faut tenir compte de la possibilité qu'ils puissent ne pas
terminer.  À part pour l'intersection, on doit donc les lancer « en
parallèle » et pas « en série » : lorsqu'on dira qu'on lance deux
algorithmes $T$ et $T'$ « en parallèle », cela signifie qu'on exécute
une étape du calcul de $T$, puis une étape de $T'$, puis de nouveau
une de $T$, et ainsi de suite en alternant entre les deux, jusqu'à ce
que l'un termine et renvoie vrai.

Si $L_1$ et $L_2$ sont semi-dédicables et si $T_1$ et $T_2$ sont des
algorithmes qui les « semi-décident » (i.e., $T_i$ termine en temps
fini et répond oui si $w\in L_i$, et ne termine pas sinon), pour
semi-décider $L_1\cup L_2$, on lance les deux algorithmes $T_1$ et
$T_2$ en parallèle sur le même mot $w$ : si l'un d'eux termine, on
termine en renvoyant vrai (sinon, bien sûr, on ne termine pas).

Pour semi-décider $L_1\cap L_2$, en revanche, il n'y a pas de raison
de procéder en parallèle : on lance d'abord $T_1$ sur le mot $w$ à
tester : si $T_1$ termine, on lance ensuite $T_2$ sur le même mot : si
$T_2$ termine et renvoie vrai, on renvoie vrai ; si l'un des deux
algorithmes $T_i$ lancés séquentiellement ne termine pas, bien sûr, le
calcul dans son ensemble ne terminera pas.

Pour semi-décider $L_1 L_2$ ou $L^*$, on procède comme dans le cas (A)
en lançant en parallèle les algorithmes pour tester toutes les
différentes factorisations possibles $w = uv$ ou bien $w = u_1\cdots
u_k$ (en mots non vides) du mot $w$.
\end{corrige}

%

\exercice

On rappelle qu'une fonction $f\colon \mathbb{N} \to \mathbb{N}$ est
dite \emph{calculable} lorsqu'il existe un algorithme (par exemple, un
programme pour une machine de Turing) prenant en entrée un
$n\in\mathbb{N}$ qui termine toujours en temps fini et renvoie la
valeur $f(n)$.  On rappelle qu'une partie $E$ de $\mathbb{N}$ ou de
$\mathbb{N}^2$ est dite \emph{décidable} lorsque sa fonction
indicatrice est calculable, ou, ce qui revient au même, lorsqu'il
existe un algorithme prenant en entrée un élément de $\mathbb{N}$ ou
de $\mathbb{N}^2$ qui termine toujours en temps fini et renvoie
vrai ($1$) ou faux ($0$) selon que l'élément fourni appartient ou non
à $E$.  On rappelle enfin qu'une partie $E$ de $\mathbb{N}$ ou de
$\mathbb{N}^2$ est dite \emph{semi-décidable} lorsqu'il existe un
algorithme prenant en entrée un élément de $\mathbb{N}$ ou de
$\mathbb{N}^2$ qui termine toujours en temps fini et renvoie
vrai ($1$) si l'élément fourni appartient à $E$, et sinon ne termine
pas (on peut aussi accepter qu'il termine en renvoyant faux, cela ne
change rien).

Soit $f\colon \mathbb{N} \to \mathbb{N}$ : montrer qu'il y a
équivalence entre les affirmations suivantes :
\begin{enumerate}
\item la fonction $f$ est calculable,
\item le graphe $\Gamma_f := \{(n,f(n)) : n\in\mathbb{N}^2\} =
  \{(n,p)\in\mathbb{N}^2 : p=f(n)\}$ de $f$ est décidable,
\item le graphe $\Gamma_f$ de $f$ est semi-décidable.
\end{enumerate}

(Montrer que (3) implique (1) est le plus difficile : on pourra
commencer par s'entraîner en montrant que (2) implique (1).  Pour
montrer que (3) implique (2), on pourra chercher une façon de tester
en parallèle un nombre croissant de valeurs de $p$ de manière à
s'arrêter si l'une quelconque convient.)

\begin{corrige}
Montrons que (1) implique (2) : si on dispose d'un algorithme capable
de calculer $f(n)$ en fonction de $n$, alors il est facile d'écrire un
algorithme capable de décider si $p=f(n)$ (il suffit de calculer
$f(n)$ avec l'algorithme supposé exister, de comparer avec la valeur
de $p$ fournie, et de renvoyer vrai/$1$ si elles sont égales, et
faux/$0$ sinon).

Le fait que (2) implique (3) est évident car tout ensemble décidable
est semi-décidable.

Montrons que (2) implique (1) même si ce ne sera au final pas utile :
supposons qu'on ait un algorithme $T$ qui décide $\Gamma_f$ (i.e.,
donnés $(n,p)$, termine toujours en temps fini, en répondant oui si
$p=f(n)$ et non si $p\neq f(n)$), et on cherche à écrire un algorithme
qui calcule $f(n)$.  Pour cela, donné un $n$, il suffit de lancer
l'algorithme $T$ successivement sur les valeurs $(n,0)$ puis $(n,1)$
puis $(n,2)$ et ainsi de suite (c'est-à-dire faire une boucle infinie
sur $p$ et lancer $T$ sur chaque couple $(n,p)$) jusqu'à trouver un
$p$ pour lequel $T$ réponde vrai : on termine alors en renvoyant la
valeur $p$ qu'on a trouvée, qui vérifie $p=f(n)$ par définition
de $T$.

\smallbreak

Reste à montrer que (3) implique (1) : supposons qu'on ait un
algorithme $T$ qui « semi-décide » $\Gamma_f$ (i.e., donnés $(n,p)$,
termine en temps fini et répond oui si $p=f(n)$, et ne termine pas
sinon), et on cherche à écrire un algorithme qui calcule $f(n)$.  Pour
cela, on va tester les valeurs $0\leq p\leq M$ chacune pour $M$ étapes
et faire tendre $M$ vers l'infini : plus exactement, on utilise
l'algorithme $U$ suivant :
\begin{itemize}
\item pour $M$ allant de $0$ à l'infini,
\begin{itemize}
\item pour $p$ allant de $0$ à $M$,
\begin{itemize}
\item exécuter l'algorithme $T$ sur l'entrée $(n,p)$ pendant au
  plus $M$ étapes,
\item s'il termine en renvoyant vrai ($1$), terminer et renvoyer $p$
  (sinon, continuer les boucles).
\end{itemize}
\end{itemize}
\end{itemize}

(Intuitivement, $U$ essaie de lancer l'algorithme $T$ sur un nombre de
valeurs de $p$ de plus en plus grand et en attendant de plus en plus
longtemps pour voir si l'une d'elles termine.)

Si l'algorithme $U$ défini ci-dessus termine, il renvoie forcément
$f(n)$ (puisque l'algorithme $T$ a répondu vrai, c'est que $p=f(n)$,
et on renvoie la valeur en question) ; il reste à expliquer pourquoi
$U$ termine toujours.  Mais la valeur $f(n)$ existe (même si on ne la
connaît pas) car la fonction $f$ était supposée définie partout, et
lorsque l'algorithme $T$ est lancé sur $(n,f(n))$ il est donc censé
terminer en un certain nombre (fini !) d'étapes : si $M$ est supérieur
à la fois à $f(n)$ et à ce nombre d'étapes, la valeur $f(n)$ va être
prise par $p$ dans la boucle intérieure, et pour cette valeur,
l'algorithme $T$ va terminer sur l'entrée $(n,p)$ en au plus $M$
étapes, ce qui assure que $U$ termine effectivement.

L'algorithme $U$ calcule donc bien la fonction $f$ demandée, ce qui
prouve (1).
\end{corrige}

%

\exercice\label{decidable-iff-image-of-computable-increasing}

Soit $A \subseteq \mathbb{N}$ un ensemble infini.  Montrer qu'il y a
équivalence entre :
\begin{itemize}
\item l'ensemble $A$ est décidable,
\item il existe une fonction calculable \emph{strictement croissante}
  $f\colon\mathbb{N}\to\mathbb{N}$ telle que $f(\mathbb{N}) = A$.
\end{itemize}

\begin{corrige}
Supposons $A$ décidable : on va construire $f$ comme indiqué.  Plus
exactement, on va appeler $f(n)$ le $n$-ième élément de $A$ par ordre
croissant (c'est-à-dire que $f(0)$ est le plus petit élément de $A$,
et $f(1)$ le suivant par ordre de taille, et ainsi de suite ; noter
que $A$ est infini donc cette fonction est bien définie).  Montrons
que $f$ est calculable : donné un entier $n$, on teste successivement
si $0\in A$ puis $1\in A$ puis $2\in A$ et ainsi de suite, à chaque
fois en utilisant un algorithme décidant $A$ (qui est censé exister
par hypothèse) jusqu'à obtenir $n$ fois la réponse « oui » ; plus
exactement :
\begin{itemize}
\item initialiser $m \leftarrow 0$,
\item pour $k$ allant de $0$ à l'infini,
\begin{itemize}
\item interroger l'algorithme qui décide si $k\in A$,
\item s'il répond « oui » :
\begin{itemize}
\item si $m=n$, terminer et renvoyer $k$,
\item sinon, incrémenter $m$ (c'est-à-dire faire $m \leftarrow m+1$).
\end{itemize}
\end{itemize}
\end{itemize}
La boucle termine car $A$ est infini.

Réciproquement, supposons $f$ strictement croissante calculable et
posons $A = f(\mathbb{N})$ : on veut montrer que $A$ est décidable.
Or pour décider si $k \in A$, il suffit de calculer successivement
$f(0)$, $f(1)$, $f(2)$ et ainsi de suite, et de terminer si $f(n)$
atteint ou dépasse le $k$ fixé : s'il l'atteint, on renvoie vrai (on a
trouvé $n$ tel que $f(n)=k$), sinon, on renvoie faux (la valeur $k$ a
été sautée par la fonction $f$ et ne sera donc jamais atteinte).
L'algorithme est donc explicitement :
\begin{itemize}
\item pour $n$ allant de $0$ à l'infini,
\begin{itemize}
\item calculer $f(n)$,
\item si $f(n) = k$, renvoyer vrai,
\item si $f(n) > k$, renvoyer faux.
\end{itemize}
\end{itemize}
La boucle termine car toute fonction strictement croissante
$\mathbb{N}\to\mathbb{N}$ est de limite $+\infty$ en l'infini (donc
$f(n)$ finit forcément par atteindre ou dépasser $k$).
\end{corrige}

%

\exercice

Soit $S(e,n)$ le nombre d'étapes de l'exécution du $e$-ième programme
(ou, si on préfère, de la $e$-ième machine de Turing) quand on lui
fournit le nombre $n$ en entrée, à supposer que cette exécution
termine ; sinon, $S(e,n)$ n'est pas défini.

Soit par ailleurs $M(k)$ le maximum des $S(e,n)$ pour $0\leq e\leq k$
et $0\leq n\leq k$ qui soient définis (et $0$ si aucun d'eux n'est
défini).  Autrement dit, il s'agit du plus petit entier supérieur ou
égal au nombre d'étapes de l'exécution de l'un des programmes $0\leq
e\leq k$ sur l'un des entiers $0\leq n\leq k$ en entrée, lorsqu'ils
terminent.

Montrer que la fonction $M$ n'est pas calculable (i.e., n'est pas
calculable par un algorithme) : on pourra pour cela montrer que la
connaissance de $M$ permet de résoudre le problème de l'arrêt.
Montrer même qu'\emph{aucune} fonction $M'$ telle que $M'(k) \geq
M(k)$ pour tout $k$ n'est calculable.  Montrer que même si $M'$
vérifie simplement $M'(k)\geq M(k)$ pour $k\geq k_0$, alors $M'$ n'est
pas calculable.

\emph{Remarque :} La fonction $M$, ou différentes variantes de
celle-ci, s'appelle fonction du « castor affairé ».  On peut montrer
encore plus fort : si $F$ est une fonction calculable quelconque,
alors il existe $k_0$ tel que $M(k) \geq F(k)$ pour $k\geq k_0$
(autrement dit, la fonction $M$ finit par dépasser n'importe quelle
fonction calculable : Radó, 1962, \textit{On Non-Computable
  Functions}).

\begin{corrige}
Supposons que $M$ soit calculable.  On peut alors résoudre le problème
de l'arrêt de la manière suivante : donné un algorithme $T$, de
numéro $e$, et une entrée $n$ à fournir à cet algorithme, pour savoir
si $T$ s'arrête, on calcule $M(k)$ où $k = \max(e,n)$, on exécute
ensuite l'algorithme $T$ pendant au plus $M(k)$ étapes : s'il termine
dans le temps imparti, on répond vrai (il a terminé), sinon, on répond
faux (il ne terminera jamais).  Cette résolution du problème de
l'arrêt est correcte, car si $T$ termine sur l'entrée $n$, il prendra
par définition $S(e,n)$ étapes, avec $0\leq e\leq k$ et $0\leq n\leq
k$ par définition de $k$, donc $S(e,n) \leq M(k)$ par définition de
$M(k)$ : ceci signifie précisément que si $T$ n'a pas terminé en
$M(k)$ étapes, il ne terminera jamais.

Exactement le même argument montre que $M'$ n'est pas calculable sous
l'hypothèse que $M'(k) \geq M(k)$ pour tout $k$ : s'il l'était, on
pourrait exécuter l'algorithme $T$ pendant au plus $M'(k)$ étapes, et
comme on a $S(e,n) \leq M(k) \leq M'(k)$, la même démonstration
convient.

Enfin, si on suppose seulement $M'(k)\geq M(k)$ pour $k\geq k_0$, la
fonction $M'$ n'est toujours pas calculable : en effet, si on suppose
par l'absurde qu'elle l'est, la fonction $M''$ qui à $k$ associe
$M'(k)$ si $k\geq k_0$ et $M(k)$ sinon, serait encore calculable
puisqu'elle ne diffère de $M'$ qu'en un nombre fini de valeurs, or
changer la valeur en un point d'une fonction calculable donne toujours
une fonction calculable (même si on « ne connaît pas » la valeur à
changer, elle existe, donc l'algorithme modifié existe).  Mais d'après
le paragraphe précédent, $M''$ n'est pas calculable puisqu'elle est
partout supérieure ou égale à $M$.
\end{corrige}

%

\exercice

Dans cet exercice, on s'intéresse au langage $L$ formé des
programmes $e$ (codés, dans une formalisation quelconque de la
calculabilité\footnote{Par exemple, un langage de programmation
  (Turing-complet) quelconque.}, comme des entiers naturels ou comme
des mots sur un alphabet fixé sans importance) qui, quel que soit le
paramètre $n$ qu'on leur fournit en entrée, terminent en temps fini et
retournent la valeur $0$ : soit $L = \{e : (\forall n)
{\varphi_e(n)\downarrow} = 0\}$.  Si l'on préfère, $L$ est l'ensemble
de toutes les façons de coder la fonction constante égale à $0$.  On
appellera aussi $M$ le complémentaire de $L$.  On se demande si $L$ ou
$M$ sont semi-décidables.

\smallskip

(1) Thésée et Hippolyte se disputent pour savoir si $L$ est
semi-décidable.  Thésée pense qu'il l'est, et il tient le raisonnement
suivant pour l'expliquer :

{\narrower

Pour savoir si un programme $e$ est dans l'ensemble $L$, il suffit de
l'examiner pour vérifier que toutes les instructions qui mettent fin
au programme renvoient la valeur $0$ : si c'est le cas, on termine en
répondait « oui » (c'est-à-dire $e\in L$) ; sinon, on rentre dans une
boucle infinie.  Ceci fournit un algorithme qui semi-décide $L$.

\par}

\smallskip

\noindent Hippolyte, elle, pense que $L$ n'est pas semi-décidable.  Son
argument est le suivant :

{\narrower

Si $L$ était semi-décidable, je pourrais m'en servir pour résoudre le
problème de l'arrêt.  En effet, donné un programme $e'$ et une
entrée $m$ sur laquelle je cherche à tester l'arrêt de $e'$, je peux
fabriquer le programme $e$ qui prend une entrée $n$, lance l'exécution
de $e'$ sur l'entrée $m$ pendant au plus $n$ étapes et à la fin
renvoie $1$ si ces $n$ étapes ont suffi à terminer l'exécution
de $e'$, et $0$ sinon.  Dire que $e \in L$ signifie que $e'$ ne
termine jamais sur $m$, donc pouvoir semi-décider $L$ permet de
résoudre algorithmiquement le problème de l'arrêt, ce qui est
impossible.

\par}

\smallskip

\noindent Qui a raison ?  Expliquer précisément quelle est l'erreur
(ou les erreurs) commise(s) par le raisonnement incorrect, et
détailler les éventuels passages incomplets dans le raisonnement
correct.

\begin{corrige}
C'est Hippolyte qui a raison ($L$ n'est pas semi-décidable).

L'argument de Thésée est stupide : une instruction mettant fin au
programme avec une valeur autre que $0$ pourrait ne jamais être
atteinte, et réciproquement, même si toutes les instructions mettant
fin au programme renvoyaient $0$, il pourrait aussi ne jamais
terminer.  (Au passage, l'argument de Thésée semblerait même montrer
que $L$ est décidable, pas juste semi-décidable.)

L'argument d'Hippolyte, lui, est correct : il montre que si $L$ était
semi-décidable, le complémentaire du problème de l'arrêt (l'ensemble
des $e'$ qui ne terminent jamais) serait semi-décidable, ce qui n'est
pas le cas.  (Le complémentaire du problème de l'arrêt n'est pas
semi-décidable, car s'il l'était, le problème de l'arrêt serait
décidable, vu qu'il est déjà semi-décidable.)  Parmi les points qui
méritent éventuellement d'être précisés, on peut mentionner : le fait
que $e$ se déduise algorithmiquement de $e'$ et $m$ ; et le fait qu'il
est algorithmiquement possible de lancer l'exécution d'un programme
sur $n$ étapes (peu importe la définition exacte de « étape » tant que
chacune termine à coup sûr en temps fini) et tester si elle est bien
finie au bout de ce temps.
\end{corrige}

\medskip

(2) Achille et Patrocle se disputent pour savoir si $M$ (le
complémentaire de $L$) est semi-décidable.  Achille pense qu'il l'est,
et il tient le raisonnement suivant pour l'expliquer :

{\narrower

Pour savoir si un programme $e$ est dans l'ensemble $M$, il suffit de
tester successivement les valeurs $\varphi_e(n)$ pour tous les $n$
possibles : si l'on rencontre un $n$ tel que $\varphi_e(n)$ n'est
pas $0$, alors on termine en répondant « oui » (c'est-à-dire $e\in
M$) ; sinon, on ne va jamais terminer, et cela signifie que $e\not\in
M$.  Ceci fournit un algorithme qui semi-décide $M$.

\par}

\smallskip

\noindent Patrocle, lui, pense que $M$ n'est pas semi-décidable.  Son
argument est le suivant :

{\narrower

Si $M$ était semi-décidable, je pourrais m'en servir pour résoudre le
problème de l'arrêt.  En effet, donné un programme $e'$ et une
entrée $m$ sur laquelle je cherche à tester l'arrêt de $e'$, je peux
fabriquer le programme $e$ qui prend une entrée $n$, l'ignore purement
et simplement, exécute $e'$ sur l'entrée $m$ et à la fin renvoie $0$.
Dire que $e\in M$ signifie que $e'$ ne termine pas sur $m$, donc
pouvoir semi-décider $M$ permet de résoudre algorithmiquement le
problème de l'arrêt, ce qui est impossible.

\par}

\smallskip

\noindent Qui a raison ?  Expliquer précisément quelle est l'erreur
(ou les erreurs) commise(s) par le raisonnement incorrect, et
détailler les éventuels passages incomplets dans le raisonnement
correct.

\begin{corrige}
C'est Patrocle qui a raison ($M$ n'est pas semi-décidable).

Le raisonnement d'Achille se fonde sur l'idée que le complémentaire de
« l'ensemble des programmes qui renvoient toujours la valeur $0$ » est
« l'ensemble des programmes qui renvoient parfois une valeur autre
que $0$ », ce qui oublie la possibilité que le programme ne termine
pas.  C'est là la principale erreur (le reste est globalement
correct).

L'argument de Patrocle, lui, est correct : il montre que si $M$ était
semi-décidable, le complémentaire du problème de l'arrêt (l'ensemble
des $e'$ qui ne terminent jamais) serait semi-décidable, ce qui n'est
pas le cas.  (Le complémentaire du problème de l'arrêt n'est pas
semi-décidable, car s'il l'était, le problème de l'arrêt serait
décidable, vu qu'il est déjà semi-décidable.)  On fait appel à des
fonctions constantes, et qui ne peuvent renvoyer que $0$ ou ne pas
terminer, mais ce n'est pas spécialement problématique.  Parmi les
points qui méritent éventuellement d'être précisés, on peut mentionner
le fait que $e$ se déduise algorithmiquement de $e'$ et $m$.
\end{corrige}



%
%
%

\printindex







%
%
%
\end{document}
